\documentclass[book]{magnoliaold}

\magtex{tex_driver={pdftex},
        tex_packages={float,caption,titling,epigraph,minitoc,slashbox,tabularx,cancel,pgfplots,xypic,nicefrac},
        tex_pstricks={pstricks,pst-plot}}
\magfiche{document_nom={Cours de Sup},
          auteur_nom={François Fayard},
          auteur_mail={fayard.prof@gmail.com}}
\magcours{cours_matiere={maths},
          cours_niveau={mpsi},
          cours_chapitre_numero={11},
          cours_chapitre={Dérivation}}
% \magmisenpage{misenpage_presentation={tikzvelvia},
%           misenpage_format={presentation},
%           misenpage_nbcolonnes={1},
%           misenpage_preuve={non},
%           misenpage_sol={non}}
\magmisenpage{}
\maglieudiff{}
\magprocess

\usepackage[scale=3]{ccicons}
\title{{\Huge\bf Cours de Mathématiques}\\\vspace{1cm}
%       \textbf{\Huge 2022--2023}\\
       \textbf{\Huge Maths Sup -- Tome I}\\\vspace{1cm}
      %  \textsc{L. Bouyge-Bono, F. Fayard, V. Lambert}\\\vspace{1cm}
       \textsc{F. Fayard}\\\vspace{1cm}
       \includegraphics[width=8cm]{../../Commun/Images/lazos-bde-2024.png}\\%\vspace{1cm}
       \includegraphics[width=8cm]{../../Commun/Images/lazos.png}}
% \author{{\sc François Fayard}, {\sc Victor Lambert}}
% \date{3 Août 2020 -- 1a20aa}

% \pretitle{%
%   \begin{center}
%   \includegraphics[width=12cm]{/Users/fayard/Desktop/lazos.png}\\[\bigskipamount]
% }
% \posttitle{\end{center}}

\mtcsettitle{minitoc}{}
\mtcsettitle{secttoc}{Table des matières}
\mtcsettitle{parttoc}{Table des matières}

\usetikzlibrary{positioning}
% \lstset{%
%   frame            = tb,    % draw frame at top and bottom of code block
%   tabsize          = 1,     % tab space width
%   numbers          = none,  % display line numbers on the left
%   framesep         = 3pt,   % expand outward
%   framerule        = 0.4pt, % expand outward 
%   commentstyle     = \color{green},      % comment color
%   keywordstyle     = \color{blue},       % keyword color
%   stringstyle      = \color{blue},    % string color
%   backgroundcolor  = \color{colorLazoBlue1Light}, % backgroundcolor color
%   showstringspaces = false,              % do not mark spaces in strings
%   basicstyle       = \ttfamily,
%   breaklines       = true
% }

\input{vc.tex}

\begin{document}

\maketitle

La version de ce document est la \textsc{\GITAbrHash}.\\

Merci à tous les élèves des lycées Janson de Sailly, du Parc et des Lazaristes pour leurs remarques et corrections. Je remercie
particulièrement \nom{Samuel Auroy}, \nom{Antonin Barbier}, \nom{Amin Belfkira}, \nom{Martin Bot}, \nom{Alexandre Brousse}, \nom{Élodie Brun}, \nom{Damien Callendrier}, \nom{Lauren Calvosa}, \nom{Sylvain Crosnier}, \nom{Enguerrand De Jaegere}, \nom{Thibaud De Valicourt}, \nom{Victor Déru}, \nom{Raphaël Des Boscs}, \nom{Grégoire Dhimoïla}, \nom{Léo Duhamel-Callot}, \nom{Sacha Evrard}, \nom{Axel Faou}, \nom{Titouan Francheteau}, \nom{Hélène Ghaleb}, \nom{Cédric Holocher},
\nom{Maxime Joubert}, \nom{Maxime Lombard}, \nom{Mira Maamri}, \nom{Gauthier Malandrin}, \nom{Pierre-Antoine Nguyen},
\nom{Hilaire Oudinot}, \nom{Eliott Pradeleix}, \nom{Yann-Ellie Ravon}, \nom{Sixtine Reynaud}, \nom{Vivien Thienot},
\nom{Carole Vacherand}, \nom{Camille Vialet}, \nom{Paul Vilars} et \nom{Antonin Villepontoux}.\\

Je tiens enfin à remercier mes anciens professeurs et
collègues qui ont eu une influence sur la rédaction de ce document~: \nom{Walter Appel}, \nom{Bruno Arsac}, \nom{Jean-Pierre Barani}, \nom{Vincent Bayle}, \nom{Christophe Bertault}, \nom{Laurence Bouyge}, \nom{Gilles Chaffard}, \nom{Alain Chillès}, \nom{Denis Choimet}, \nom{Vincent Clapiès}, \nom{Gérard Esposito}, \nom{Stéphane Gonnord}, \nom{Victor Lambert}, %
%\nom{Pierre-Louis Lions},
\nom{Frédéric Morlot}, \nom{Franz Ridde}, \nom{Emmanuel Roblet} et \nom{Alain Troesch}.\\
%, \nom{Cédric Villani}.\\
\vfill


\begin{center}
  \ccbysa\\
  \vspace{2ex}
  This work is licensed under a Creative Commons\\
  Attribution-ShareAlike 4.0 International License.\\
  \url{https://creativecommons.org/licenses/by-sa/4.0/legalcode.fr}\\
  \vspace{2ex}
  La dernière version de ce document ainsi que\\
  les sources \LaTeX{} sont disponibles à l'adresse\\
  \url{https://github.com/FayardProf/Maths-MPSI-MP2I}
  \end{center}
  \vspace{2ex}
  \begin{center}
  \textbf{Vous êtes autorisés à~:}
  \end{center}
  \vspace{2ex}
  \begin{itemize}
  \item \textbf{Partager}~: copier, distribuer et communiquer le matériel par tous les moyens et sous tous formats.
  \item \textbf{Adapter}~: remixer, transformer et créer à partir du matériel
  pour toute utilisation, y compris commerciale.
  \end{itemize}
  \vspace{2ex}
  \begin{center}
  \textbf{Selon les conditions suivantes~:}
  \end{center}
  \vspace{2ex}
  \begin{itemize}
    \item \textbf{Attribution}~: Vous devez créditer l'œuvre, intégrer un lien vers la licence et indiquer si des modifications ont été effectuées à l'œuvre. Vous devez indiquer ces informations par tous les moyens raisonnables, sans toutefois suggérer que l'offrant vous soutient ou soutient la façon dont vous avez utilisé son œuvre.
    \item \textbf{Partage dans les mêmes conditions}~: Dans le cas où vous effectuez un remix, que vous transformez, ou créez à partir du matériel composant l'œuvre originale, vous devez diffuser l'œuvre modifiée dans les même conditions, c'est à dire avec la même licence avec laquelle l'œuvre originale a été diffusée.
    \item \textbf{Pas de restrictions complémentaires}~: Vous n'êtes pas autorisé à appliquer des conditions légales ou des mesures techniques qui restreindraient légalement autrui à utiliser l'œuvre dans les conditions décrites par la licence.
    \end{itemize}

\tableofcontents

% \part{Sup}

\chapter{Nombres complexes}
\setcounter{numeroexercicecours}{1}
\documentclass{magnoliaold}

\magtex{tex_driver={pdftex},
        tex_packages={epigraph,pgfplots,caption,float,xypic}}
\magfiche{document_nom={Cours sur les nombres complexes},
          auteur_nom={François Fayard},
          auteur_mail={fayard.prof@gmail.com}}
\magcours{cours_matiere={maths},
          cours_niveau={mpsi},
          cours_chapitre_numero={1},
          cours_chapitre={Nombres Complexes}}
\magmisenpage{misenpage_presentation={tikzvelvia},
          misenpage_format={a4},
          misenpage_nbcolonnes={1},
          misenpage_preuve={non},
          misenpage_sol={non}}
\maglieudiff{}
\magprocess

\begin{document}

%BEGIN_BOOK
\setlength\epigraphwidth{.6\textwidth}
\epigraph{\og La voie la plus courte et la meilleure entre deux vérités du domaine réel passe souvent par le domaine imaginaire. \fg}{--- \textsc{Jacques Hadamard (1865--1963)}}
\hidemesometimes{
\bigskip
\hfill\includegraphics[width=0.2\textwidth]{../../Commun/Images/maths-cours-imaginary_numbers.png}}

\magtoc

\section{Le corps des nombres complexes}
\subsection{Définition, conjugaison, module}
Le carré de tout nombre réel étant positif, l'équation
$$x^2=-1$$
n'admet aucune solution réelle. Nous admettrons qu'il existe un ensemble de
nombres $A$ ayant les propriétés suivantes.
\begin{itemize}
  \item $\R\subset A$.
  \item On peut additionner, soustraire et multiplier les éléments de $A$ en
    utilisant les règles usuelles de l'algèbre.
  \item L'équation $z^2=-1$ admet au moins une solution sur $A$.
\end{itemize}
On note $\ii$ une solution de cette équation.

\begin{definition}[utile=-3]
  On appelle corps des nombres complexes et on note $\C$
  l'ensemble des nombres $x+\ii y$ où $x$ et $y$ sont réels.
\end{definition}

\begin{remarques}
\remarque $\R$ est inclus dans $\C$. Autrement dit, tout nombre réel est un nombre
  complexe.
\remarque $\C$ est stable par les opérations d'addition, de soustraction et
  de multiplication.
\begin{preuve}
Soit $z$ et $z'$ deux nombres complexes. Il existe donc
$x,y,x',y'\in\R$ tels que $z=x+\ii y$ et $z'=x'+\ii y'$. Alors
\begin{eqnarray*}
z+z' &=& x+\ii y+x'+\ii y'\\
     &=& \underbrace{\p{x+x'}}_{\in\R}
         +\ii \underbrace{\p{y+y'}}_{\in\R}\in\C,\\
z-z' &=& \underbrace{\p{x-x'}}_{\in\R}
         +\ii \underbrace{\p{y-y'}}_{\in\R}\in\C,\\
z z' &=& \p{x+\ii y}\p{x'+\ii y'}\\
     &=& x x'+\ii x y'+\ii x' y+\underbrace{\ii^2}_{=-1} y y'\\
     &=& \underbrace{\p{x x'-y y'}}_{\in\R} +
         \ii\underbrace{\p{x y'+x' y}}_{\in\R} \in\C.
\end{eqnarray*}
\end{preuve}
\end{remarques}

\begin{definition}[utile=-3]
  Pour tout nombre complexe $z$, il existe un unique couple de réels $(x,y)$ 
  tel que $z=x+\ii y$. Les réels $x$ et $y$ sont respectivement appelés \emph{partie réelle} et \emph{partie imaginaire} de $z$. On note
  $$\Re(z)\defeq x \quad \text{et} \quad \Im(z)\defeq y$$
et on a donc $z=\Re(z)+\ii\Im(z)$.
\end{definition}
\begin{preuve}
  Soit $z$ un nombre complexe.
  \begin{itemize}
  \item \emph{Existence}~:
    Par définition de $\C$, il existe $x,y\in\R$ tels que $z=x+\ii y$.
  \item \emph{Unicité}~:
    Soit $x_1,y_1,x_2,y_2\in\R$ tels que $z=x_1+\ii y_1$ et $z=x_2+\ii y_2$. Alors
    \begin{equation*}\begin{split}
    & x_1+\ii y_1=x_2+\ii y_2\\
    \text{donc}\quad  & x_1-x_2=\ii\p{y_2-y_1}\\
    \text{donc}\quad  & \p{x_1-x_2}^2=\ii^2\p{y_2-y_1}^2\\   
    \text{donc}\quad  & \p{x_1-x_2}^2=-\p{y_2-y_1}^2\\   
    \text{donc}\quad  & \underbrace{\p{x_1-x_2}^2}_{\geq 0}+
                        \underbrace{\p{y_2-y_1}^2}_{\geq 0}=0
    \end{split}\end{equation*}
    La somme de deux réels positifs n'étant nulle que si ces deux réels sont
    nuls, on en déduit que $x_1=x_2$ et $y_1=y_2$.
  \end{itemize}
\end{preuve}

\begin{remarques}
\remarque On appelle \emph{forme cartésienne} de $z\in\C$, toute écriture de la forme $z=x+\ii y$ où $x,y\in\R$. La proposition précédente affirme que, quel que soit $z\in\C$, une telle écriture existe et est unique.
\remarque Si $x_1,y_1,x_2,y_2\in\R$ sont tels que
  \[x_1+\ii y_1=x_2+\ii y_2,\]
  alors $x_1=x_2$ et $y_1=y_2$. On dit souvent qu'on procède par \emph{identification}, mais cette terminologie est abusive. On utilise en fait l'unicité provenant de la proposition précédente.
\remarque Soit $\mathcal{R}=\p{O,\ve{e_1},\ve{e_2}}$ un repère orthonormé
  direct du plan. À tout nombre complexe $z=x+\ii y$, on associe le point $M$ dont les coordonnées dans le repère
  $\mathcal{R}$ sont $\p{x,y}$. On a donc
  \[\ve{OM}=x\ve{e_1}+y\ve{e_2}\]
  et on dit que $M$ a pour affixe $z$. Pout tout point $M$ du plan, il existe un unique $z\in\C$ tel que
  $M$ a pour affixe $z$; on dit alors que $z$ est \emph{l'affixe} du point $M$. On a ainsi identifié
  $\C$ avec l'ensemble des points du plan.
\remarque Un nombre complexe est réel si et seulement si sa partie imaginaire est nulle. On dit qu'un nombre complexe
  est \emph{imaginaire pur} lorsque sa partie réelle est nulle. L'ensemble des nombres imaginaires purs est
  donc
  \[\ii\R\defeq\ensim{\ii y}{y\in\R}.\]
\remarque De même qu'on ne peut pas écrire d'inégalités entre les points du plan, les inégalités entre nombres
  complexes n'ont aucun sens.
\end{remarques}

\begin{proposition}[utile=-3]
Un nombre complexe est nul si et seulement si sa partie réelle et sa partie imaginaire le sont.
\end{proposition}

\begin{preuve}
Soit $z$ un nombre complexe. Alors $z=\Re(z)+\ii \Im(z)$. 
\begin{itemize}
\item Si $\Re(z)=0$ et $\Im(z)=0$, alors $z=0$.
\item Réciproquement, si $z=0$, alors $\Re(z)+\ii \Im(z)=0+\ii 0$. Par unicité de la forme cartésienne, on en déduit que
  \[\Re(z)=0 \et \Im(z)=0.\]
\end{itemize}
\end{preuve}

\begin{proposition}[utile=-3]
Soit $z$ et $z'$ deux nombres complexes, $\lambda$ et $\mu$ deux réels. Alors
\begin{itemize}
\item $\Re(\lambda z+\mu z')=\lambda\Re(z)+\mu\Re(z'),$
\item $\Im(\lambda z+\mu z')=\lambda\Im(z)+\mu\Im(z').$
\end{itemize}
\end{proposition}

\begin{preuve}
Soit $z$ et $z'$ deux nombres complexes, $\lambda$ et $\mu$ deux réels.
Il existe donc $x,y,x',y'\in\R$ tels que $z=x+\ii y$ et $z'=x'+\ii y'$. Alors
\begin{eqnarray*}
\lambda z+\mu z'
&=& \lambda\p{x+\ii y}+\mu\p{x'+\ii y'}\\
&=& \underbrace{\p{\lambda x+\mu x'}}_{\in\R}+\ii
    \underbrace{\p{\lambda y+\mu y'}}_{\in\R}
\end{eqnarray*}
Puisque $x=\Re(z)$, $y=\Im(z)$, $x'=\Re(z')$, $y'=\Im(z')$, on en déduit que
\[\Re(\lambda z+\mu z')=\lambda\Re(z)+\mu\Re(z') \et
  \Im\p{\lambda z+\mu z'}=\lambda\Im(z)+\mu\Im(z').\]
\end{preuve}

\begin{remarqueUnique}
\remarque Attention, l'identité $\Re(z z')=\Re(z)\Re(z')$ est fausse.
  Par exemple $\Re(\ii\cdot \ii)=-1$ et $\Re(\ii)\Re(\ii)=0$.
\end{remarqueUnique}

\begin{definition}[utile=-3]
  Soit $z$ un nombre complexe. On appelle \emph{conjugué} de $z$ et on note
  $\conj{z}$ le nombre complexe
  $$\conj{z}\defeq x-\ii y$$
  où $x$ et $y$ sont respectivement la partie réelle et imaginaire de $z$.
\end{definition}

\begin{remarqueUnique}
\remarque Si $M$ est le point d'affixe $z\in\C$, le point $M'$ d'affixe $\conj{z}$ est le symétrique de $M$ par rapport à l'axe $(Ox).$
\end{remarqueUnique}

\begin{proposition}[utile=-3]
  Soit $z,z'\in\C$. Alors
 \[\conj{z+z'}=\conj{z}+\conj{z}', \quad \conj{z z'}=\conj{z}\ \conj{z}', \et \conj{\conj{z}}=z.\]
\end{proposition}
\begin{preuve} Soit $z,z'\in\C$. Il existe donc $x,y,x',y'\in\R$ tels que
    $z=x+\ii y$ et $z'=x'+\ii y'$
  \begin{itemize}
  \item On a $z+z'=\p{x+x'}+\ii \p{y+y'}$ donc
    \[\conj{z+z'}=\p{x+x'}-\ii \p{y+y'}.\] Or
    \begin{eqnarray*}
    \conj{z}+\conj{z}'
    &=& x-\ii y+x'-\ii y'\\
    &=& \p{x+x'}-\ii \p{y+y'}\\
    &=& \conj{z+z'}
    \end{eqnarray*}
  \item
    De même $zz'=\p{xx'-yy'}+\ii \p{xy'+x'y}$, donc
    \[\conj{zz'}=\p{xx'-yy'}-\ii \p{xy'+x'y}.\] Or
    \begin{eqnarray*}
    \conj{z}\ \conj{z}'
    &=& \p{x-\ii y}\p{x'-\ii y'}\\
    &=& xx'-\ii xy'-\ii x'y-yy'\\
    &=& \p{xx'-yy'}-\ii \p{xy'+x'y}\\
    &=& \conj{z z'}
    \end{eqnarray*}
  \item Enfin $\conj{z}=x-\ii y=x+\ii(-y)$, donc
    $\conj{\conj{z}}=x-\ii(-y)=x+\ii y=z$.
  \end{itemize}
\end{preuve}

\begin{definition}[utile=-3]
Soit $z$ un nombre complexe. On appelle \emph{module} de $z$ et on note $\abs{z}$ le nombre réel positif
\[\abs{z}\defeq\sqrt{x^2+y^2}\]
où $x$ et $y$ sont respectivement la partie réelle et imaginaire de $z$.
\end{definition}

\begin{remarques}
% \item On voit directement que $\overline{z\conj{z}}=\conj{z}\conj{\conj{z}}=\conj{z}z$.
%   En particulier $z\conj{z}\in\R$. Par contre cette méthode ne permet pas de
%   montrer que $z\conj{z}$ est positif.
\remarque Si $x\in\R$ est considéré comme un nombre complexe, son module est égal à sa valeur absolue.
\remarque Si $M$ est le point d'affixe $z$, le module de $z$ est la distance $OM$. Si $A$ et $B$ sont deux points d'affixes respectives $a$ et $b$, alors $AB=\abs{b-a}$.
\remarque Si $a\in\C$ et $r>0$
  \begin{itemize}
  \item $\enstq{z\in\C}{\abs{z-a}=r}$ est le \emph{cercle} de centre $a$ et de rayon $r$.
  \item $\enstq{z\in\C}{\abs{z-a}\leq r}$ est le \emph{disque fermé} de centre $a$ et de rayon $r$.
  \item $\enstq{z\in\C}{\abs{z-a}<r}$ est le \emph{disque ouvert} de centre $a$ et de rayon $r$.
  \end{itemize}
  \remarque Soit $a,b,c\in\R$ tels qu'au moins l'un des deux réels $a,b$ est non nul.
  Alors, l'ensemble d'équation
  \[ax+by+c=0\]
  est une droite orthogonale au vecteur de coordonnées $(a,b)$.
\remarque Soit $a,b,c\in\R$. Alors l'ensemble d'équation
  \[x^2+y^2+ax+by+c=0\]
  est soit un cercle, soit un point, soit l'ensemble vide.
\end{remarques}

\begin{definition}[utile=-3]
On note $\U$ l'ensemble des nombres complexes de module $1$.
\end{definition}

\begin{remarqueUnique}
\remarque L'identification entre $\C$ et le plan complexe nous amène à identifier $\U$
  avec le cercle de centre $O$ et de rayon 1, appelé \emph{cercle trigonométrique}.
\end{remarqueUnique}

\begin{proposition}
Soit $z$ un nombre complexe. Alors
\[\abs{z}^2=z \conj{z}.\]
De plus, $\abs{z}=0$ si et seulement si $z=0$.
\end{proposition}
\begin{preuve}
Soit $z\in\C$. Il existe $x,y\in\R$ tels que $z=x+\ii y$. On a donc
\begin{eqnarray*}
\abs{z}=0
&\ssi& \abs{z}^2=0\\
&\ssi& \underbrace{x^2}_{\geq 0}+\underbrace{y^2}_{\geq 0}=0\\
&\ssi& x^2=0 \et y^2=0\\
&\ssi& x=0 \et y=0\\
&\ssi& z=0.
\end{eqnarray*}
\end{preuve}

\begin{remarqueUnique}
\remarque Afin d'exploiter cette identité, on cherchera souvent à travailler avec le carré des modules.
% \remarque Si $\mathcal{C}$ est le cercle de centre $a\in\C$ et de rayon $r>0$, alors
%   \begin{eqnarray*}
%   \forall z\in\C\qsep z\in\mathcal{C}
%   &\ssi& \abs{z-a}=r\\
%   &\ssi& \abs{z-a}^2=r^2\\
%   &\ssi& (z-a)\conj{(z-a)}=r^2\\
%   &\ssi& z\conj{z}-\conj{a}z-a\conj{z}+\p{a\conj{a}-r^2}=0.
%   \end{eqnarray*}
%   Réciproquement, si $a\in\C$ et $b\in\R$, l'ensemble d'équation $z\conj{z}-\conj{a}z-a\conj{z}+b=0$ est soit un cercle de centre $a$, soit réduit au point $a$, soit l'ensemble vide.
%   \begin{sol}
% Soit $a\in\C$ et $b\in\R$.
% \begin{eqnarray*}
% \forall z\in\C\qsep z\conj{z}-\conj{a}z-a\conj{z}+b=0
% &\ssi& z\conj{z}-\conj{a}z-a\conj{z}+a\conj{a}=a\conj{a}-b\\
% &\ssi& (z-a)\conj{(z-a)}=a\conj{a}-b\\
% &\ssi& \abs{z-a}^2=\underbrace{a\conj{a}-b}_{\eqdef c\in\R}
% \end{eqnarray*}
% \begin{itemize}
% \item Si $c>0$, on pose $r\defeq\sqrt{c}$. Alors
%   \[\forall z\in\C\qsep z\conj{z}-\conj{a}z-a\conj{z}+b=0 \quad\ssi\quad \abs{z-a}^2=r^2 \quad\ssi\quad \abs{z-a}=r\]
%   donc l'ensemble cherché est le cercle de centre $a$ et de rayon $r$.
% \item Si $c=0$, alors
%   \[\forall z\in\C\qsep z\conj{z}-\conj{a}z-a\conj{z}+b=0 \quad\ssi\quad \abs{z-a}=0 \quad\ssi\quad z=a\]
%   donc l'ensemble cherché est réduit au point $a$.
% \item Si $c<0$, l'ensemble est vide.
% \end{itemize}
%   \end{sol}
\end{remarqueUnique}

% \begin{exoUnique}
% \exo Donner une condition nécessaire et suffisante sur $z\in\C\setminus\ens{\ii}$ pour que
%   \[\frac{z+2}{1+\ii z}\]
%   soit réel.
%   \begin{sol}
%   Soit $z\in\C$. Alors
%   \begin{eqnarray*}
% \abs{z+\ii}=\abs{z-\ii}
% &\ssi& \abs{z+\ii}^2=\abs{z-\ii}^2\\
% &\ssi& (z+\ii)(\conj{z}-\ii)=(z-\ii)(\conj{z}+\ii)\\
% &\ssi& z=\conj{z} \ssi z\in\R
%   \end{eqnarray*}
%   Donc $z$ est solution de l'équation $\abs{z+\ii}=\abs{z-\ii}$ si et seulement si $z$ est réel.
%   \end{sol}
% \end{exoUnique}

\begin{proposition}[utile=-3]
  Soit $z,z'\in\C$. Alors
  \[\abs{z z'}=\abs{z}\abs{z'} \et \abs{\conj{z}}=\abs{z}.\]
\end{proposition}

\begin{preuve}
Soit $z,z'\in\C$. Alors
\[|z z'|^2=z z'\conj{z z'}=z z'\conj{z}\ \conj{z}'=z\conj{z}z'\conj{z}'=|z|^2|z'|^2=\p{|z||z'|}^2.\]
Ainsi, $\abs{z z'}$ et $\abs{z}\abs{z'}$ sont deux réels positifs qui ont même carré, ils sont donc égaux. De plus $|\conj{z}|^2=\conj{z}\ \conj{\conj{z}}=\conj{z}z=|z|^2$. Puisque $\abs{\conj{z}}$ et $\abs{z}$ sont deux réels positifs ayant même carrés, ils sont égaux. Donc $\abs{\conj{z}}=\abs{z}$.
\end{preuve}

% \subsection{Inverse}

\begin{proposition}
  Si $z$ et $z'$ sont deux nombres complexes tels que $z z'=0$, alors
  $z=0$ ou $z'=0$. On dit que $\C$ est \emph{intègre}.
\end{proposition}

\begin{preuve}
Soit $z,z'\in\C$ tels que $zz'=0$. Alors $|zz'|=0$, donc $|z||z'|=0$. Puisque $\R$ est intègre, l'un des deux modules est nul. D'après la proposition précédente, on en déduit que $z=0$ ou $z'=0$.

\end{preuve}
\begin{proposition}
  Soit $z$ un nombre complexe non nul. Alors il existe un unique nombre
  complexe $z'$ tel que $zz'=1$. On note ce nombre $z^{-1}$ ou $1/z$. De plus
  \[\frac{1}{z}=\frac{\conj{z}}{\abs{z}^2}.\]
\end{proposition}

\begin{preuve}
Soit $z$ un nombre complexe non nul.
\begin{itemize}
\item \emph{Unicité}~: Soit $z_1,z_2\in\C$ tels que $z z_1=1$ et $z z_2=1$. Alors $z z_1=z z_2$, donc $z(z_1-z_2)=0$. Puisque $z$ est non nul et que $\C$ est intègre, on en déduit que $z_1-z_2=0$, donc $z_1=z_2$.
\item \emph{Existence}~: On pose
\[z'= \conj{z} \cdot \frac{1}{\abs{z}^2}.\]
Alors
\[z z'=z\conj{z}\cdot\frac{1}{\abs{z}^2} =1.\]
\end{itemize}
\end{preuve}

\begin{remarques}
\remarque Pour obtenir $1/z$ sous forme cartésienne, il suffit de multiplier son
  numérateur et son dénominateur par $\conj{z}$.
\remarque Si $z\in\U$, alors $1/z=\conj{z}$. Pour inverser un nombre complexe de module 1, il suffit donc de le conjuguer.
\end{remarques}

\begin{exoUnique}
\exo Calculer l'inverse de $1+\ii$.
  \begin{sol}
  On a
  \[\frac{1}{1+\ii}=\frac{1-\ii}{(1+\ii)(1-\ii)}=\frac{1-\ii}{\abs{1+\ii}^2}=\frac{1-\ii}{2}=\frac{1}{2}-\frac{1}{2}\ii.\]
  \end{sol}
\end{exoUnique}

\begin{proposition}
  Soit $z$ un nombre complexe non nul. Alors
  $$\conj{\p{\frac{1}{z}}}=\frac{1}{\conj{z}} \quad \text{et} \quad
    \abs{\frac{1}{z}}=\frac{1}{\abs{z}}.$$
\end{proposition}

\begin{preuve}
Soit $z$ un nombre complexe non nul. D'une part
\[\conj{z}\conj{\p{\frac{1}{z}}}=\conj{z\cdot\frac{1}{z}}=\conj{1}=1\]
d'où le premier résultat. D'autre part
\[\abs{\frac{1}{z}}\abs{z}=\abs{\frac{1}{z}\cdot z}=\abs{1}=1\]
d'où le deuxième résultat.
\end{preuve}

\begin{exoUnique}
\exo Soit $a,b\in\C$ tels que $\abs{a}<1$ et $\abs{b}<1$. Montrer que
  \[\abs{\frac{a-b}{1-\conj{a}b}}<1.\]
\begin{sol}
Montrons d'abord que le dénominateur est non nul. On a
\[\abs{\conj{a}b}=\abs{\conj{a}}\abs{b}=\abs{a}\abs{b}<1\]
donc $\conj{a}b\neq 1$. De plus
 \begin{eqnarray*}
 \abs{\frac{a-b}{1-\conj{a}b}}<1&\Longleftrightarrow &|a-b|<|1-\conj{a}b|\\
 &\Longleftrightarrow & |a-b|^2<|1-\conj{a}b|^2\\
 &\Longleftrightarrow & (a-b)\conj{(a-b)}<(1-\conj{a}b)\conj{(1-\conj{a}b)}\\
 &\Longleftrightarrow & \ldots \\
 &\Longleftrightarrow & \abs{a}^2+\abs{b}^2<1+\abs{a}^2\abs{b}^2 \\
 &\Longleftrightarrow & (\abs{a}^2-1)(\abs{b}^2-1)>0
 \end{eqnarray*}
  Le résultat à droite de l'équivalence étant vrai grâce aux hypothèses, le résultat au départ l'est.
  \end{sol}
\end{exoUnique}

\begin{proposition}[utile=-3]
Soit $z$ un nombre complexe. Alors
\[\Re(z)=\frac{z+\conj{z}}{2} \et
  \Im(z)=\frac{z-\conj{z}}{2\ii}.\]
En particulier
\begin{itemize}
\item $z$ est réel si et seulement si $\conj{z}=z$.
\item $z$ est imaginaire pur si et seulement si $\conj{z}=-z$.
\end{itemize}
\end{proposition}
\begin{preuve}
Soit $z$ un nombre complexe. Il existe donc $x,y\in\R$ tels que $z=x+\ii y$. Alors $\conj{z}=x-\ii y$, donc
\[z+\conj{z}=2x \et z-\conj{z}=2\ii y.\]
\end{preuve}

\begin{remarqueUnique}
\remarque En pratique, pour montrer qu'un nombre complexe est réel, une
  bonne méthode est de montrer qu'il est égal à son conjugué. La méthode
  consistant à montrer que sa partie imaginaire est nulle est à proscrire.
\end{remarqueUnique}

\begin{exos}
\exo Soit $a$ et $b$ deux nombres complexes de module 1 tels que $ab\neq -1$. Montrer que
  \[\frac{a+b}{1+ab}\]
  est un nombre réel.
\begin{sol}
Soit $a$ et $b$ deux nombres complexes de module 1 tels que $ab\neq -1$. Alors
\begin{eqnarray*}
\conj{\p{\frac{a+b}{1+ab}}}
&=& \frac{\conj{a}+\conj{b}}{1+\conj{a}\conj{b}}\\
&=& \frac{\frac{1}{a}+\frac{1}{b}}{1+\frac{1}{ab}} \quad \text{car $\abs{a}=\abs{b}=1$.}\\
&=& \frac{b+a}{ab+1} = \frac{a+b}{1+ab}.
\end{eqnarray*}
Donc $(a+b)/(1+ab)$ est réel.
\end{sol}
\exo Donner une condition nécessaire et suffisante sur $z\in\C\setminus\ens{\ii}$ pour que
  \[\frac{z+2}{1+\ii z}\]
  soit réel.
  \begin{sol}
Cette expression a un sens pour tout $z\in\C\setminus\ens{\ii}$. Soit donc un tel $z$.
\begin{eqnarray*}
 \frac{z+2}{1+\ii z} \in \R&\Longleftrightarrow &\conj{\frac{z+2}{1+\ii z}}=\frac{z+2}{1+\ii z}\\
 &\Longleftrightarrow & (\conj{z}+2)(1+iz)=(z+2)(1-i\conj{z})\\
 &\Longleftrightarrow & \ldots \\
 &\Longleftrightarrow & z\conj{z}-\p{-1-\frac{1}{2}i}z-\p{-1+\frac{1}{2}i}\conj{z}=0 
 \end{eqnarray*}
Avec la remarque de la page précédente, on peut conclure que le lieu cherché est le cercle de centre $-1+\ii/2$ et de rayon $\sqrt{5/4}$ privé de $\ii$. Il passe par l'origine et par $z=-2$. 
  \end{sol}
\end{exos}



\subsection{Inégalité triangulaire}


\begin{proposition}[utile=-3]
Soit $a\in\C$. Alors
\[\Re(a)\leq\abs{\Re(a)}\leq\abs{a} \et
  \Im(a)\leq\abs{\Im(a)}\leq\abs{a}.\]
De plus, $\Re(a)=\abs{a}$ si et seulement si $a$ est réel positif.
\end{proposition}

% \begin{exoUnique}
% \exo Résoudre sur $\C$ le système
%   \[\begin{cases}
%     \abs{1+z}\leq 1 &\\
%     \abs{1-z}\leq 1. &
%     \end{cases}\]
%   \begin{sol}
% On raisonne par analyse-synthèse. Commençons par l'analyse et donnons-nous une solution $z\in\C$ du système. Alors $\abs{1+z}\leq 1$ et $\abs{1-z}\leq 1$ donc, en élevant au carré, $1+\abs{z}^{2}+2\Re{z}\leq 1$ et $1+\abs{z}^{2}-2\Re{z}\leq 1$. En ajoutant ces deux inégalités, on obtient $ \abs{z}^{2}\leq 0 $ donc $z=0$. La synthèse est immédiate car $z=0$ est bien solution de ce système. En conclusion l'ensemble des solutions est
% \[\mathcal{S}=\ens{0}.\]
% \end{sol}
% \end{exoUnique}

\begin{proposition}[utile=3, nom=Inégalité triangulaire]
Soit $a$ et $b$ deux nombres complexes. Alors
\[\abs{a+b}\leq\abs{a}+\abs{b}.\]
De plus, l'égalité a lieu si et seulement si  $a$ et $b$ sont positivement
liés, c'est-à-dire lorsque $a=0$ ou lorsqu'il existe $\lambda\in\RP$ tel que
$b=\lambda a$.
\end{proposition}

\begin{preuve}
Soit $a$ et $b$ deux nombres complexes. Alors
\begin{eqnarray*}
\abs{a+b}\leq\abs{a}+\abs{b}
&\ssi& \abs{a+b}^2\leq\p{\abs{a}+\abs{b}}^2, \quad\text{car $\abs{a+b}$ et $\abs{a}+\abs{b}$ sont positifs}\\
&\ssi& \p{a+b}\conj{\p{a+b}}\leq\p{\abs{a}+\abs{b}}^2\\
&\ssi& \p{a+b}\p{\conj{a}+\conj{b}}\leq\p{\abs{a}+\abs{b}}^2\\
&\ssi& a\conj{a}+a\conj{b}+\conj{a}b+b\conj{b}\leq\abs{a}^2+2\abs{a}\abs{b}+\abs{b}^2\\
&\ssi& \abs{a}^2+2\Re\p{\conj{a}b}+\abs{b}^2\leq\abs{a}^2+2\abs{a}\abs{b}+\abs{b}^2\\
&\ssi& \Re\p{\conj{a}b}\leq \abs{\conj{a}b}.
\end{eqnarray*}
Or, d'après la proposition précédente, ce dernier point est vrai. Donc $\abs{a+b}\leq\abs{a}+\abs{b}$.\\

En reprenant les équivalences précédentes, on a
\begin{eqnarray*}
\abs{a+b}=\abs{a}+\abs{b}
&\ssi& \Re\p{\conj{a}b}= \abs{\conj{a}b}\\
&\ssi& \conj{a}b\in\RP.
\end{eqnarray*}
Montrons que ce dernier point est vrai si et seulement si $a=0$ ou qu'il existe $\lambda\in\RP$ tel que $b=\lambda a$.
\begin{itemize}
\item Supposons que $\conj{a}b\in\RP$. On pose $\mu\defeq\conj{a}b\in\RP$. Si $a\neq 0$, alors
\[b=\frac{\mu}{\conj{a}}=\frac{\mu a}{a\conj{a}}=\underbrace{\frac{\mu}{\abs{a}^2}}_{\eqdef\lambda\in\RP}a.\]
\item Réciproquement, supposons que $a=0$ ou qu'il existe $\lambda\in\RP$ tel que $b=\lambda a$. Si $a=0$, alors $\conj{a}b=0\in\RP$. S'il existe $\lambda\in\RP$ tel que $b=\lambda a$, alors $\conj{a}b=\conj{a}\lambda a=\abs{a}^2 \lambda\in\RP$.
\end{itemize}
En conclusion $\abs{a+b}=\abs{a}+\abs{b}$ si et seulement si $a=0$ ou s'il existe $\lambda\in\RP$ tel que $b=\lambda a$.
\end{preuve}

\begin{remarques}
\remarque Si $(ABC)$ est un triangle, alors $AC\leq AB+BC$. En effet, si on note $a$, $b$, $c$ les affixes respectives de $A$, $B$, $C$
\[AC=\abs{c-a}=\abs{c-b+b-a}\leq\abs{c-b}+\abs{b-a}=BC+AB.\]
Cette inégalité explique le nom d'inégalité triangulaire donné à la proposition précédente.
\remarque Attention, il est possible que $a$ et $b$ soient positivement liés sans
  qu'il existe $\lambda\in\RP$ tel que $b=\lambda a$.
\end{remarques}


%% Soit $a$ et $b$ deux nombres complexes.
%% \begin{itemize}
%% \item Montrons que~:
%% $$\abs{a+b}\leq\abs{a}+\abs{b}$$
%% Remarquons, que le résultat est immédiat lorsque $b=0$. On suppose donc
%% désormais que $b\not=0$. Soit $t\in\R$.
%% \begin{eqnarray*}
%%   \abs{a+tb}^2 &=& \p{a+tb}\conj{\p{a+tb}}\\
%%                    &=& \p{a+tb}\p{\conj{a}+t\conj{b}}\\
%%                    &=& a\conj{a}+ta\conj{b}+tb\conj{a}+
%%                        t^2b\conj{b}\\
%%                    &=& b\conj{b} t^2 + \p{a\conj{b}+
%%                        \conj{a\conj{b}}}t+a\conj{a}\\
%%                    &=& \abs{b}^2 t^2 + 2\Re\p{a\conj{b}} t
%%                        +\abs{a}^2
%% \end{eqnarray*}
%% Or $\abs{a+tb}^2\geq 0$ quelque soit $t\in\R$. On en déduit donc que le
%% discriminant du trinôme en $t$ est négatif. Puisque $\abs{b}^2\not= 0$,
%% on en déduit que le discriminant de ce trinôme est négatif. En effet, si tel
%% n'était pas le cas, le trinôme aurait deux racines distinctes et changerait
%% de signe. On a donc~:
%% $$\Delta=\p{2\Re\p{a\conj{b}}}^2-4\abs{b}^2 \abs{a}^2 \leq 0$$
%% Donc~:
%% $$\abs{\Re\p{a\conj{b}}}^2 \leq \abs{a}^2\abs{b}^2$$
%% Puisque $\abs{\Re\p{a\conj{b}}}$ et $\abs{a}\abs{b}$ sont positifs,
%% on en déduit que~:
%% $$\abs{\Re{\p{a\conj{b}}}} \leq \abs{a}\abs{b}$$
%% Calculons maintenant~:
%% \begin{eqnarray*}
%%   \abs{a+b}^2-\p{\abs{a}+\abs{b}}^2 &=&
%%       \p{a+b}\conj{\p{a+b}}-\p{\abs{a}+\abs{b}}^2\\
%%   &=& \abs{a}^2+\abs{b}^2+2\Re\p{a\conj{b}}-
%%       \p{\abs{a}^2+\abs{b}^2+2\abs{a}\abs{b}}\\
%%   &=& 2\p{\Re\p{a\conj{b}}-\abs{a}\abs{b}}
%% \end{eqnarray*}
%% Or $\abs{\Re{\p{a\conj{b}}}} \leq \abs{a}\abs{b}$, donc puisque
%% $\Re{\p{a\conj{b}}} \leq \abs{\Re{\p{a\conj{b}}}}$, on en déduit
%% que~:
%% $$\abs{a+b}^2-\p{\abs{a}+\abs{b}}^2 \leq 0$$
%% et donc~:
%% $$\abs{a+b}^2\leq\p{\abs{a}+\abs{b}}^2$$
%% Or $\abs{a+b}$ et $\abs{a}+\abs{b}$ sont tous deux positifs, donc~:
%% $$\abs{a+b}\leq\abs{a}+\abs{b}$$
%% \item Montrons désormais que cette inégalité devient une égalité si et
%% seulement si $a$ et $b$ sont positivements liés.\\
%% \begin{itemize}
%% \item Si $a$ et $b$ sont positivement liés,il existe $\lambda\in\RP$ tel
%% que $a=\lambda b$ ou $b=\lambda a$. Quitte à les échanger, on peut
%% supposer que $b=\lambda a$. Alors~:
%% \begin{eqnarray*}
%%   \abs{a+b} &=& \abs{a+\lambda a}\\
%%                 &=& \abs{(1+\lambda)a}\\
%%                 &=& \abs{1+\lambda}\abs{a}\\
%%                 &=&\p{1+\lambda}\abs{a}\quad\text{car $1+\lambda\geq 0$}\\
%%                 &=& \abs{a}+\lambda\abs{a}\\
%%                 &=& \abs{a}+\abs{\lambda}\abs{a}
%%                     \quad\text{car $\lambda\geq 0$}\\
%%                 &=& \abs{a}+\abs{\lambda a}\\
%%                 &=& \abs{a}+\abs{b}
%% \end{eqnarray*}
%% \item Réciproquement, supposons que $\abs{a+b}=\abs{a}+\abs{b}$ et
%% montrons que $a$ et $b$ sont positivement liés.\\
%% Si $b$ est nul, $a$ et $b$ sont positivement liés, car $b=0 a$. On
%% suppose donc que $b$ est non nul. En remontant les calculs qui ont permis
%% la démonstration de l'inégalité triangulaire, de l'égalité
%% $\abs{a+b}=\abs{a}+\abs{b}$, il découle que~:
%% $$\Re\p{a\conj{b}}=\abs{a}\abs{b}$$
%% Donc le discriminant du trinôme considéré plus haut est nul.Ce trinôme admet
%% donc une unique racine $\mu$. Par définition $\abs{a+\mu b}=0$, donc
%% $a=-\mu b$. On pose $\lambda=-\mu$. Il reste alors à montrer que
%% $\lambda \geq 0$. Or~:
%% $$\Re\p{a\conj{b}}=\abs{a}\abs{b}$$
%% Donc, un réinjectant dans cette égalité la relation $a=\lambda b$, on
%% trouve~:
%% $$\lambda \abs{b}^2=\abs{\lambda}\abs{b}^2$$
%% Puisque $b\not=0$, on peut simplifier par $\abs{b}^2$, et il vient
%% $\lambda=\abs{\lambda}\geq 0$.
%% \end{itemize}
%% \end{itemize}
%% \end{preuve}

\begin{proposition}[utile=-3,nom={Seconde inégalité triangulaire}]
Soit $a$ et $b$ deux nombres complexes. Alors
\[\abs{\abs{a}-\abs{b}}\leq\abs{a+b}.\]
%\leq\abs{a}+\abs{b}.\]  
\end{proposition}

\begin{preuve}
$|a|=|a+b-b|\leq |a+b|+|-b|$ ... + changement des rôles.

\end{preuve}

\begin{remarqueUnique}
% \remarque Si $ABC$ est un triangle,alors
%   \[\abs{AB-BC}\leq AC\leq AB+BC\]
\remarque La seconde inégalité triangulaire admet plusieurs variantes.
  \begin{itemize}
  \item Si on remplace $b$ par $-b$, on obtient l'inégalité
    \[\abs{\abs{a}-\abs{b}}\leq\abs{a-b}\]
    qui affirme que deux nombres complexes proches ont des modules proches.
  \item En remarquant que si $x$ est réel, $\abs{x}\geq x$, on obtient
    \[\abs{a+b}\geq\abs{a}-\abs{b}.\]
    Cette inégalité affirme que si $b$ a un module petit par rapport à celui de $a$, alors $a+b$ est éloigné de~0.
  \end{itemize}
\end{remarqueUnique}

\begin{exos}
\exo Soit $a$ et $b$ deux nombres complexes distincts. On pose $\delta\defeq\abs{a-b}$. Montrer que les disques ouverts de centre $a$ et $b$ et de rayon $\delta/2$ sont disjoints.
\exo Que peut-on dire de $\abs{z}$ si $\abs{1-z}\leq 1/4$~? Faire un dessin puis une preuve.
\begin{sol}
Par l'absurde, si $z$ est dans l'intersection des deux disques, $|a-b|\leq|a-z|+|z-b|<\delta$.

$|1|\leq |1-z|+|z|$ et $|z|\leq |z-1|+|1|$ donnent $5/4\geq \abs{z}\geq 3/4$.
\end{sol}
% \exo Soit $z_0,a,\ldots,z_n$ des nombres complexes de module 1. Montrer
%   que
%   \[\sum_{k=0}^n \frac{z_k}{2^k}\neq 0\]
%   \begin{sol}
%   Il suffit de remarquer que
%   \[\abs{\sum_{k=0}^n \frac{z_k}{2^k}}\geq
%     1-\abs{\sum_{k=1}^n \frac{z_k}{2^k}}\geq \frac{1}{2^n}\]
%   \end{sol}
\end{exos}

\begin{proposition}
Soit $a_1,\ldots,a_n\in\C$. Alors
\[\abs{\sum_{k=1}^n a_k}\leq \sum_{k=1}^n \abs{a_k}.\]
\end{proposition}


\subsection{Puissance entière, binôme de \nom{Newton}}

\begin{definition}[utile=-3]
  Soit $a\in\C$. On définit $a^n$ pour tout entier naturel $n\in\N$ en posant
  \begin{itemize}
    \item $a^0\defeq 1$,
    \item $\forall n\in\N \qsep a^{n+1}\defeq a^n a$.
  \end{itemize}
\end{definition}

\begin{remarqueUnique}
\remarque En particulier, $0^0=1$.
\end{remarqueUnique}

\begin{proposition}[utile=-3]
  Soit $a,b$ deux nombres complexes, $n$ et $m$ deux entiers
  naturels. Alors
  \[\p{ab}^n=a^n b^n, \quad a^{n+m}=a^n a^m \et 
    \p{a^n}^m=a^{nm}.\]
\end{proposition}

\begin{preuve}
Soit $a\in\C$ et $n\in\N$. Pour tout $m\in\N$, on pose
\[\mathcal{H}_m\defeq\mog a^{n+m}=a^n a^m\mfg.\]
Pour tout $m\in\N$, on pose
\[\mathcal{H}_m\defeq\mog \p{a^n}^m=a^{nm}\mfg.\]
\end{preuve}

\begin{proposition}[utile=-3]
  Soit $a\in\C$ et $n\in\N$. Alors
  \[\conj{a^n}=\conj{a}^n \et \abs{a^n}=\abs{a}^n.\]
\end{proposition}

\begin{exoUnique}
\exo Montrer que si $P(z)\defeq a_n z^n+\cdots+a_1 z+a_0$ est un polynôme à coefficients réels, l'ensemble de ses racines est stable par conjugaison.
\end{exoUnique}

\begin{sol}
Soit $\alpha$ une racine. On a $P(\alpha)=0$, il suffit alors d'écrire $0=\conj{P(\alpha)}=\ldots=P(\conj{\alpha})$ grâce aux règles de calcul de la conjugaison.
\end{sol}

\begin{definition}
  Soit $a$ un nombre complexe non nul. On étend la définition de
  $a^n$ à $n\in\Z$ en posant
  $$a^n\defeq \frac{1}{a^{-n}}$$
  lorsque $n<0$.
\end{definition}

\begin{proposition}
Soit $a,b$ deux nombres complexes non nuls, $n$ et $m$ deux entiers
relatifs. Alors
\[\p{ab}^n=a^n b^n,\quad
  a^{n+m}=a^n a^m \et
  \p{a^n}^m=a^{nm}.\]
\end{proposition}

\begin{proposition}
Soit $a$ un nombre complexe non nul et $n\in\Z$. Alors
\[\conj{a^n}=\conj{a}^n \et \abs{a^n}=\abs{a}^n.\]
\end{proposition}

\begin{definition}[nom={Division euclidienne}]
Soit $a\in\Z$ et $b\in\Ns$. Alors il existe un unique couple $\p{q,r}\in\Z^2$
tel que
\[a=qb+r \et 0\leq r<b.\]
$q$ est appelé \emph{quotient} de la division euclidienne de $a$ par $b$, $r$ son
\emph{reste}.
\end{definition}

\begin{preuve}
\begin{itemize}
\item[$\bullet$] \textbf{Unicité :} Supposons trouvés deux couples $(q_1,r_1)$ et $(q_2,r_2)$. On obtient alors $$b|q_1-q_2|=|r_2-r_1|.$$ Supposons un instant $q_1\neq q_2$, alors $$b\leq b|q_1-q_2|=|r_2-r_1|\leq b-1$$ d'où la contradiction.
\item[$\bullet$] \textbf{Existence :} Commençons par le cas $a\geq 0$ et considérons l'ensemble~:
\[A=\enstq{k\in\N}{kb\leq a}\]
Cet ensemble est non vide ($0\in A$), majoré (par $a$) donc admet un plus grand
élément $q$. Comme $q\in A$ et $q+1 \notin A$, on a $bq\leq a <b(q+1)$ d'où $0\leq a-bq<b$. Il reste à poser $r=a-bq$.\\
Si $a\leq 0$, il suffit d'appliquer ce qui précède à $-a$. Il existe donc $q\in \N$ et $r\in \N$ vérifiant \[-a=qb+r \et 0\leq r<b.\]
Donc $a=(-q)b+(-r)$ ce qui convient si $r=0$ et si $0<r<b$, alors $-b<-r<0$ donc $0<-r+b<b$ et dans ce cas $a=(-q-1)b+(-r+b)$ est de la forme recherchée.
\end{itemize}
\end{preuve}

\begin{exoUnique}
\exo On pose $\jj\defeq-\frac{1}{2}+\ii\frac{\sqrt{3}}{2}$. Calculer $\jj^3$, puis en déduire
  $\jj^{2023}$.
\end{exoUnique}
  

\begin{definition}[utile=-3]
Pour tout entier naturel $n$, on définit la \emph{factorielle} de $n$ que l'on note
$n!$ par
\begin{itemize}
\item $0!\defeq 1$,
\item $\forall n\in\N \qsep \p{n+1}!\defeq\p{n+1}n!$.
\end{itemize}
\end{definition}

\begin{remarqueUnique}
\remarque Si $n\in\Ns$
  \[n!=n\times(n-1)\times\cdots\times 2\times 1.\]
\end{remarqueUnique}

\begin{definition}[utile=-3]
Pour tout couple $\p{k,n}$ d'entiers naturels, on définit $\binom{n}{k}$ que l'on prononce \og $k$ parmi $n$ \fg, comme étant le nombre de parties à $k$ éléments d'un ensemble
à $n$ éléments.
\end{definition}

\begin{remarques}
\remarque Si $k>n$, alors $\binom{n}{k}=0$.
\remarque Si $n\in\N$, alors
\[\binom{n}{0}=1 \et \binom{n}{n}=1.\]
\end{remarques}

\begin{proposition}[utile=-3]
Soit $n\in\N$ et $k\in\intere{0}{n}$. Alors
\[\binom{n}{k}=\binom{n}{n-k}.\]
\end{proposition}

\begin{preuve}
On remarque que choisir $k$ éléments parmi $n$ revient à sélectionner les $n-k$ éléments qu'on ne choisira pas.\end{preuve}

\begin{proposition}[utile=2, nom=Relation de \nom{Pascal}]
Soit $k$ et $n$ deux entiers naturels. Alors
\[\binom{n}{k}+\binom{n}{k+1}=\binom{n+1}{k+1}.\]
\end{proposition}

\begin{preuve}
\'Evacuons tout d'abord certains cas en fonction des valeurs de $k$ : $k>n$, $k=n$.

Désormais on considère $k<n$. Démontrons ce résultat de manière combinatoire. Considérons l'ensemble $\ens{1, 2, . . . ,n+1}$ et dénombrons ses sous-ensembles à $k+1$ éléments. On distingue les ensembles qui contiennent l'élément $n+1$ et les ensembles qui ne le contiennent pas. Il y a$\binom{n}{k+1}$ ensembles qui ne contiennent pas $n+1$ (on choisit $k+1$ éléments dans $\ens{1, 2, . . . ,n}$ et il y en a $\binom{n}{k}$ qui contiennent $n+1$, ce qui fait bien le total souhaité.

\end{preuve}

\begin{remarqueUnique}
\remarque Cette formule est appelée relation de \nom{Pascal}. Elle permet de calculer
  efficacement les $\binom{n}{k}$ en construisant le triangle de \nom{Pascal}. Dans ce
  tableau contenant les $\binom{n}{k}$, où $n$ désigne la ligne et $k$ désigne la colonne,
  on commence par placer une colonne de 1 indiquant le fait que $\binom{n}{0}=1$, puis une diagonale
  de 1 indiquant le fait que $\binom{n}{n}=1$. Les coefficients au-dessus de la diagonale sont nuls
  et ne sont généralement pas représentés. Ceux en dessous de la diagonale sont complétés, ligne
  après ligne en utilisant la relation de Pascal qui affirme que chaque coefficient est la somme
  du coefficient se situant au-dessus de lui et de celui au-dessus à gauche.
  \[\begin{matrix}
    1 &   &    &    &   &  \\
    1 & 1 &    &    &   &  \\
    1 & 2 &  1 &    &   &  \\
    1 & 3 &  3 &  1 &   &  \\
    1 & 4 &  6 &  4 & 1 &  \\
    1 & 5 & 10 & 10 & 5 & 1
    \end{matrix}\]
  Ce triangle permet par exemple de lire sur la dernière ligne que $\binom{5}{1}=5$ et $\binom{5}{2}=10$.
\end{remarqueUnique}

\begin{proposition}[utile=-3]
Soit $n$ un entier naturel et $k\in\intere{0}{n}$. Alors
\[\binom{n}{k}=\frac{n!}{k!\p{n-k}!}.\]
\end{proposition}

\begin{preuve}
Utilisons  le  principe  de double-comptage, très  important  en  combinatoire.  Il consiste à établir une égalité en comptant de deux manières différentes une certaines quantité. Ici, comptons le nombre de suites à $k$ éléments différents qu'on peut créer en utilisant les $n$ éléments de $\ens{1, 2, . . . ,n}$. D'une part, nous avons $n$ choix pour le premier terme de la suite, $n-1$ choix pour le deuxième, et ainsi de suite jusqu'au $k$-ième élément pour lequel nous avons $n-k+1$ choix. Finalement, il y a en tout $n(n-1)\ldots(n-k+1) =\dfrac{n!}{(n-k)!}$ telles suites. D'autre part, pour créer une suite à $k$ éléments, on peut commencer par choisir les $k$ éléments qui vont constituer la suite ($\binom{n}{k}$ possibilités), puis les ordonner ($k!$ manières possibles de les ordonner). Il y a donc en tout $k!\binom{n}{k}$ telles suites.

Ainsi, $\dfrac{n!}{(n-k)!}=k!\binom{n}{k}$, d'où le résultat souhaité.

\end{preuve}

\begin{preuve}
On peut maintenant redémontrer les deux propositions précédentes à l'aide de la "formule" du coefficient binomial que l'on vient d'obtenir :
Laissé en exercice aux élèves.

\end{preuve}
\begin{remarques}
% \remarque Cette formule est très mauvaise pour calculer effectivement
%   $\binom{n}{k}$. Par exemple, si l'on souhaite calculer $\binom{13}{2}$ à
%   l'aide de cette formule, on est amené à calculer $13!$. Si on dispose d'un
%   ordinateur codant les entiers sur 32 bits, le calcul de $13!$ donnera un
%   résultat erroné car $13!\geq 2^{32}$. L'explosion d'\nom{Ariane} 5 lors de
%   son décollage le 4 juin 1996 est due à une erreur de ce type.
\remarque On peut simplifier l'écriture de $\binom{n}{k}$ en
  \[\binom{n}{k}=\frac{n!}{k!\p{n-k}!}=\frac{\frac{n!}{\p{n-k}!}}{k!}=
    \frac{\overbrace{n\p{n-1}\cdots\p{n-(k-1)}}^{\text{$k$ termes}}}{k!}.\]
  En particulier
  \[\binom{n}{1}=n \et \binom{n}{2}=\frac{n\p{n-1}}{2}.\]
\remarque Si $k,n\in\Ns$, on a la formule dite \og du capitaine \fg
  \[\binom{n}{k}=\frac{n}{k}\cdot\binom{n-1}{k-1}.\]
\end{remarques}

\begin{exoUnique}
\exo Soit $k,n\in\N$. Montrer que
  \[\binom{n}{k}\leq\frac{n^k}{k!}.\]
  \begin{sol}
  Soit $k,n\in\N$. Si $k\leq n$, alors
  \[\binom{n}{k}=
    \frac{\overbrace{n\p{n-1}\cdots\p{n-k+1}}^{\text{$k$ termes inférieurs ou égaux à $n$}}}{k!}\leq \frac{n^k}{k!}.\]
  Sinon, l'inégalité est triviale.
  \end{sol}
\end{exoUnique}


\begin{proposition}[utile=3, nom=Binôme de \nom{Newton}]
Soit $a$ et $b$ deux nombres complexes et $n$ un entier naturel. Alors
\[\p{a+b}^n=\sum_{k=0}^n \binom{n}{k}a^{n-k} b^k.\]
\end{proposition}
\begin{preuve}
Soit $a$ et $b$ deux nombres complexes. Montrons ce résultat par récurrence
sur $n$~:
\begin{recurrence}
test $$\p{a+b}^n=\sum_{k=0}^n a^k b^{n-k}$$
\recinit En effet~:
  $$\p{a+b}^0=1 \et a^0 b^0=1$$
\rechere On suppose que $\mathcal{H}_n$ est vraie. Montrons que
  $\mathcal{H}_{n+1}$ est vraie.
  On écrit $(a+b)^{n+1}=(a+b)(a+b)^n$ puis on applique l'H.R, on distribue, on réindexe la première somme...
\end{recurrence}
\end{preuve}

\begin{exoUnique}
\exo Soit $n\in\Ns$. Montrer que
  \[\sum_{k=0}^n  \binom{n}{k}=2^n \quad\et\quad  \sum_{k=0}^n  \p{-1}^k\binom{n}{k}=0.\]
\end{exoUnique}


\begin{proposition}[utile=-3]
  Soit $a$ et $b$ deux nombres complexes. Alors
  \begin{itemize}
  \item $a^2-b^2=\p{a-b}\p{a+b}.$
  \item Plus généralement, pour tout $n\in\Ns$
    \begin{eqnarray*}
  a^n-b^n
  &=& \p{a-b}\p{a^{n-1}+a^{n-2}b+\dots+
      ab^{n-2}+b^{n-1}}\\
  &=& \p{a-b}\p{\sum_{k=0}^{n-1} a^{n-1-k}b^k}.
  \end{eqnarray*}
  \end{itemize}
\end{proposition}

\begin{preuve}
On distribue dans le membre de droite puis on réindexe la deuxième somme.
\end{preuve}

\begin{proposition}[utile=-3]
Soit $a$ un nombre complexe et $n$ un entier naturel. Alors
\[\sum_{k=0}^n a^k=1+a+a^2+\cdots+a^n=
  \begin{cases}
  \dsp\frac{1-a^{n+1}}{1-a} & \text{si $a\neq 1$}\\
  n+1                  & \text{si $a=1$.}
  \end{cases}\]
\end{proposition}

\section{Forme trigonométrique}

\subsection{Exponentielle $\ii\theta$}

\begin{definition}[utile=-3]
Pour tout réel $\theta$, on définit l'exponentielle de $\ii\theta$ par
\[\e^{\ii\theta}\defeq\cos\theta+\ii\sin\theta.\]
\end{definition}

\begin{proposition}[utile=-3]
Soit $\theta_1$ et $\theta_2$ deux réels. Alors
\[\e^{\ii 0}=1 \qquad\text{et}\qquad \e^{\ii\p{\theta_1+\theta_2}}=\e^{\ii\theta_1}\e^{\ii\theta_2}.\]
\end{proposition}

\begin{preuve}
Immédiat avec $\cos(a+b)=\cos a\cos b-\sin a\sin b$ et $\sin(a+b)=\sin a\cos b+\cos a\sin b$.
\end{preuve}

\begin{proposition}[utile=-3]
Soit $\theta\in\R$. Alors
\[\overline{\e^{\ii\theta}}=\e^{-\ii\theta}.\]
De plus, $\e^{\ii\theta}$ est non nul et si $n\in\Z$, alors
\[\frac{1}{\e^{\ii\theta}}=\e^{-\ii\theta} \qquad\text{et}\qquad \e^{\ii n\theta}=
  \p{\e^{\ii \theta}}^n.\]  
\end{proposition}


\begin{proposition}[utile=-3, nom={Formules d'\nom{Euler} et \nom{Moivre}}]
Soit $\theta$ un réel. Alors les formules d'\nom{Euler} s'écrivent
\[\cos\theta=\frac{\e^{\ii \theta}+\e^{-\ii\theta}}{2} \et
  \sin\theta=\frac{\e^{\ii \theta}-\e^{-\ii\theta}}{2\ii}.\]
Pour $n\in\Z$, la formule de \nom{Moivre} nous donne
\[\cos \p{n\theta}+\ii\sin\p{n\theta}=\p{\cos \theta+\ii\sin \theta}^n.\]
\end{proposition}

% \exo Montrer que tout $z\in\U\setminus\ens{-1}$ il existe $a\in\R$ tel que
%   \[z=\frac{1+\ii a}{1-\ii a}.\]
%   \begin{sol}
% Soit $z\in\U\setminus\ens{-1}$. Alors
% \begin{eqnarray*}
% \forall a \in \R\qsep \frac{1+\ii a}{1-\ii a}=z
% &\ssi& a=\frac{z-1}{i(1+z)}
% \end{eqnarray*}
% Montrons que
% \[a\defeq \frac{z-1}{i(1+z)}\]
% est réel. On a
% \begin{eqnarray*}
% \conj{a}
% &=& \conj{\frac{z-1}{i(1+z)}}\\
% &=& -\frac{\conj{z}-1}{i(1+\conj{z})}\\
% &=& -\frac{\frac{1}{z}-1}{i\p{1+\frac{1}{z}}} \quad\text{car $z\in\U$}\\
% &=& \frac{1-z}{\ii\p{z+1}} = a
% \end{eqnarray*}
% Donc $a\in\R$.
%   \end{sol}
% \end{exoUnique}

\begin{proposition}[utile=-3]
\begin{itemize}
\item Soit $\theta\in\R$. Alors $\e^{\ii\theta}=1$ si et seulement si
  $\theta\equiv 0\ \cro{2\pi}$.
\item Plus précisément, étant donnés $\theta_1$ et $\theta_2\in\R$,
  $\e^{\ii\theta_1}=\e^{\ii\theta_2}$ si et seulement si
  $\theta_1\equiv\theta_2\ \cro{2\pi}$.
\end{itemize}
\end{proposition}

\begin{exoUnique}
\exo Déterminer la partie réelle de
  \[\frac{1}{1-\cos\theta-\ii\sin\theta}.\]
  \begin{sol}
  On commence par chercher le domaine de définition de l'expression. On a, pour tout $\theta\in\R$
  \begin{eqnarray*}
1-\cos\theta-\ii\sin\theta=0
&\ssi& 1-\e^{\ii \theta}=0\\
&\ssi& \e^{\ii \theta}=1\\
&\ssi& \theta\equiv 0\ \cro{2\pi}.
  \end{eqnarray*}
  Pour $\theta\in\R\setminus 2\pi\Z$, on a
  \begin{eqnarray*}
\frac{1}{1-\cos\theta-\ii\sin\theta}
&=& \frac{1}{1-\e^{\ii\theta}}\\
&=& \frac{1-\e^{-\ii\theta}}{(1-\e^{\ii\theta})(1-\e^{-\ii\theta})}\\
&=& \frac{1-\cos\theta+\ii\sin\theta}{2(1-\cos\theta)}
  \end{eqnarray*}
  donc
  \[\Re\p{\frac{1}{1-\cos\theta-\ii\sin\theta}}=\frac{1}{2}.\] 
  \end{sol}
\end{exoUnique}


\begin{proposition}[utile=-3,nom={Paramétrisation de $\U$ par \og l'exponentielle $\ii\theta$ \fg}]
L'application qui à $\theta$ associe $\e^{\ii\theta}$ est une surjection de $\R$
dans $\U$. Autrement dit~:
\begin{itemize}
\item Si $\theta\in\R$, $\e^{\ii\theta}\in\U$.
\item Pour tout $u\in\U$, il existe $\theta\in\R$ tel que $u=\e^{\ii\theta}$.
\end{itemize}
\end{proposition}

\begin{preuve}
La première partie vient du fait que $\cos^2+\sin^2=1$.

Soit $u\in \U$, $u=a+ib$. On a $|a|=|\Re(u)|\leq |u|=1$ donc $a\in [-1;1]$. On peut donc fixer $\theta\in \R$ tel que $a=\cos \theta$. Mais alors $b^2=1-\cos^2(\theta)$ donc $b=\pm \sin \theta=\sin(\pm \theta)$. Ce $\theta'$ convient alors.

On voit par la preuve que ce $\theta$ est unique à $2\pi$ près. \underline{Insister là-dessus en vue de la preuve d'après}
\end{preuve}



\subsection{Application à la trigonométrie}
\begin{applications}
\application {\it Factorisation par l'arc moitié}\\
  Étant donné un réel $\theta$
  \begin{eqnarray*}
  1+\e^{\ii\theta} &=& \e^{\ii\frac{\theta}{2}}\p{ \e^{-\ii\frac{\theta}{2}}
                                            +\e^{ \ii\frac{\theta}{2}}}\\
                &=& 2\cos\p{\frac{\theta}{2}}\e^{\ii\frac{\theta}{2}}.
  \end{eqnarray*}
  De même
  \begin{eqnarray*}
  1-\e^{\ii\theta} &=& \e^{\ii\frac{\theta}{2}}\p{ \e^{-\ii\frac{\theta}{2}}
                                            -\e^{ \ii\frac{\theta}{2}}}\\
                &=& -2\ii\sin\p{\frac{\theta}{2}}\e^{\ii\frac{\theta}{2}}.
  \end{eqnarray*}
  Plus généralement, étant donnés $\theta_1,\theta_2\in\R$
  \[\e^{\ii\theta_1}+\e^{\ii\theta_2}=\e^{\ii\frac{\theta_1+\theta_2}{2}}\p{\e^{\ii\frac{\theta_1-\theta_2}{2}}+\e^{-\ii\frac{\theta_1-\theta_2}{2}}}=
    2\cos\p{\frac{\theta_1-\theta_2}{2}}\e^{\ii\frac{\theta_1+\theta_2}{2}},\]
   \[\e^{\ii\theta_1}-\e^{\ii\theta_2}=\e^{\ii\frac{\theta_1+\theta_2}{2}}\p{\e^{\ii\frac{\theta_1-\theta_2}{2}}-\e^{-\ii\frac{\theta_1-\theta_2}{2}}}=
    2\ii\sin\p{\frac{\theta_1-\theta_2}{2}}\e^{\ii\frac{\theta_1+\theta_2}{2}}.\]   
\application {\it Calcul de sommes trigonométriques}\\
  Soit $\theta\in\R$ et $n\in\N$. Calculer
   \[C_n\defeq\sum_{k=0}^n \cos(k\theta) \quad \text{et} \quad
    S_n\defeq\sum_{k=0}^n \sin(k\theta).\]
  \begin{sol}
  On a
  \begin{eqnarray*}
  C_n+\ii S_n &=& \sum_{k=0}^n \cos k\theta+\ii\sin k\theta\\
           &=& \sum_{k=0}^n \e^{\ii k\theta}\\
           &=& \sum_{k=0}^n \p{\e^{\ii\theta}}^k
  \end{eqnarray*}
  \begin{itemize}
  \item Lorsque $\e^{\ii\theta}\neq 1$, c'est-à-dire lorsque
    $\theta\not\equiv 0\ \cro{2\pi}$, on a
    \begin{eqnarray*}
    C_n+\ii S_n &=& \frac{1-\p{\e^{\ii\theta}}^{n+1}}{1-\e^{\ii\theta}}\\
             &=& \frac{1-\e^{\ii \p{n+1}\theta}}{1-\e^{\ii\theta}}\\
             &=& \frac{\e^{\ii\frac{n+1}{2}\theta}}{\e^{\ii\frac{\theta}{2}}}
                 \frac{\e^{-\ii\frac{n+1}{2}\theta}-\e^{\ii\frac{n+1}{2}\theta}}
                      {\e^{-\ii\frac{\theta}{2}}-\e^{\ii\frac{\theta}{2}}}\\
             &=& \e^{\ii\frac{n\theta}{2}}\frac{-2\ii\sin\p{\frac{n+1}{2}\theta}}
                 {-2\ii\sin\p{\frac{\theta}{2}}}\\
             &=& \e^{\ii\frac{n\theta}{2}}\frac{\sin\p{\frac{n+1}{2}\theta}}
                 {\sin\p{\frac{\theta}{2}}}.
    \end{eqnarray*}
    En identifiant parties réelles et imaginaires, on obtient
    $$C_n=\frac{\sin\p{\frac{n+1}{2}\theta}\cos\p{\frac{n}{2}\theta}}
          {\sin\p{\frac{\theta}{2}}} \quad \text{et} \quad
      S_n=\frac{\sin\p{\frac{n+1}{2}\theta}\sin\p{\frac{n}{2}\theta}}
          {\sin\p{\frac{\theta}{2}}}.$$
  \item Dans le cas où $\theta\equiv 0\ \cro{2\pi}$, $\cos (k\theta)=1$ et
    $\sin(k\theta)=0$, donc
    $$C_n=n+1 \quad \text{et} \quad S_n=0.$$
  \end{itemize}      
  \end{sol}
\application {\it Linéarisation de $\cos^n \theta \sin^m \theta$}\\
  Étant donné deux entiers naturels $n$ et $m$, on cherche à exprimer
  $\cos^n \theta \sin^m \theta$ comme combinaison linéaire des $\cos(k\theta)$
  et $\sin(k\theta)$ pour $k\in\N$. Pour cela, on peut utiliser les formules
  d'\nom{Euler} avant de développer l'expression par la formule
  du binôme de \nom{Newton} et de regrouper les termes en utilisant à nouveau
  les formules d'\nom{Euler}. Cette opération sera utile lors
  du calcul de primitives.\\
  {\bf Exemple~:} Linéariser $\sin^6 \theta$ et $\sin\theta \cos^4 \theta$.
    \begin{sol}
    On a
    \begin{eqnarray*}
    \sin^6 \theta &=& \p{\frac{\e^{\ii\theta}-\e^{-\ii\theta}}{2\ii}}^6\\
    &=& -\frac{1}{64}\p{\e^{\ii 6\theta}-6\e^{\ii 4\theta}+15\e^{\ii 2\theta}-20
                        +15\e^{-i2\theta}-6\e^{-i4\theta}+\e^{-i6\theta}}\\
    &=& -\frac{1}{64}\p{\p{\e^{\ii 6\theta}+\e^{-i6\theta}}
                        -6\p{\e^{\ii 4\theta}+\e^{-i4\theta}}
                        +15\p{\e^{\ii 2\theta}+\e^{-i2\theta}}-20}\\
    &=& -\frac{1}{32}\p{\cos 6\theta-6\cos 4\theta+15\cos 2\theta-10}
    \end{eqnarray*}
    \begin{eqnarray*}
    \sin \theta \cos^4 \theta
    &=& \p{\frac{\e^{\ii\theta}-\e^{-\ii\theta}}{2\ii}}
        \p{\frac{\e^{\ii\theta}+\e^{-\ii\theta}}{2}}^4\\
    &=& -\frac{i}{32}\p{\e^{\ii\theta}-\e^{-\ii\theta}}
                     \p{\e^{\ii\theta}+\e^{-\ii\theta}}^4\\
    &=& -\frac{i}{32}\p{\e^{\ii\theta}-\e^{-\ii\theta}}
        \p{\e^{\ii 4\theta}+4\e^{\ii 2\theta}+6+4\e^{-i2\theta}+\e^{-i4\theta}}\\
    &=& -\frac{i}{32}\p{\e^{\ii 5\theta}+3\e^{\ii 3\theta}+2\e^{\ii\theta}-2\e^{-\ii\theta}
                        -3\e^{-i3\theta}-\e^{-i5\theta}}\\
    &=& -\frac{i}{32}\p{2\ii\sin 5\theta+6\ii\sin 3\theta+4\ii\sin \theta}\\
    &=& \frac{1}{16}\p{\sin 5\theta+3\sin 3\theta+2\sin \theta}
    \end{eqnarray*}      
    \end{sol}
\application {\it Expression  de $\cos(n\theta)$ et $\sin(n\theta)$ comme
  polynôme en $\cos \theta$ et $\sin \theta$}\\
  Pour cette opération, une méthode consiste à utiliser la formule de
  \nom{Moivre} avant de développer l'expression obtenue à l'aide du binôme de
  \nom{Newton}.\\
  {\bf Exemple~:} Exprimer $\cos(5\theta)$ comme un polynôme en $\cos\theta$. 
\end{applications}

\begin{exos}
% \exo Exprimer $\tan\p{5\theta}$ en fonction de $\tan\theta$.
% \begin{sol}
% On a
% \begin{eqnarray*}
% \cos\p{5\theta}+\ii\sin\p{5\theta}
% &=& \p{\cos\theta+\ii\sin\theta}^5\\
% &=& (\cos^{5}\theta-10\cos^{3}\theta\sin^{2}\theta+5\cos\theta\sin^{4}\theta)+\\
% & & \ii(5\cos^{4}\theta\sin\theta-10\cos^{2}\theta\sin^{3}\theta+\sin^{5}).
% \end{eqnarray*}
% Donc, en exploitant le fait que
% \[\tan (5\theta)=\frac{\sin (5\theta)}{\cos (5\theta)}\]
% et en divisant le numérateur et le dénominateur par $ \cos^{5}\theta $, on trouve
%   \[\tan\p{5\theta}=
%     \frac{5\tan\theta-10\tan^3\theta+\tan^5\theta}%
%     {1-10\tan^2\theta+5\tan^4\theta}.\]
% \end{sol}
\exo Soit $n\in\N$ et $\theta\in\R$. Calculer
  \[\sum_{k=0}^n \binom{n}{k}\sin\p{k\theta}.\]
  \begin{sol}
  On a
  \begin{eqnarray*}
 \sum_{k=0}^{n}\binom{n}{k}\e^{\ii k\theta}
 &=& \sum_{k=0}^{n}\binom{n}{k}(\e^{\ii\theta})^k\\
 &=& \p{1+\e^{\ii\theta}}^n\\
 &=& \cro{\e^{\frac{\ii\theta}{2}}\p{\e^{\frac{\ii\theta}{2}}+\e^{\frac{-\ii\theta}{2}}}}^n\\
 &=& \e^{\ii\frac{n\theta}{2}}\p{ \e^{\ii\frac{-n\theta}{2}}+\e^{\ii\frac{n\theta}{2}}}^{n}\\
 &=& \e^{\ii\frac{n\theta}{2}}\left( 2\cos\left( \frac{\theta}{2}\right) \right) ^{n}\\
 &=& 2^n \cos^n\p{\frac{\theta}{2}}\e^{\ii\frac{n\theta}{2}}.
  \end{eqnarray*}
  En prenant la partie imaginaire, on en déduit que
  \[\sum_{k=0}^n \binom{n}{k}\cos\p{k\theta}=2^n \cos^n\p{\frac{\theta}{2}}\sin\p{\frac{n\theta}{2}}.\]
  \end{sol}
\exo En exprimant $\cos(5\theta)$ comme un polynôme en $\cos\theta$,
  montrer que $\cos\p{\pi/10}$ est racine d'un polynôme à
  coefficients entiers. En déduire une expression de $\cos\p{\pi/10}$ à l'aide de radicaux.
\begin{sol}
On commence par exprimer $\cos(5\theta)$ comme polynôme en $\cos\theta$. On a
\[\cos (5\theta)=\cos \theta(16\cos^{4}\theta -20\cos^{2}\theta +5).\]
Pour $\theta=\pi/10$, en posant $x\defeq\cos(\pi/10)$, on obtient $x(16x^4-20x^2+5)=0$. Or $\theta\in\interfo{0}{\pi/2}$ donc $x>0$, donc $16x^4-20x^2+5=0$. Or
\begin{eqnarray*}
\forall u\in\R\qsep 16u^4-20u^2+5=0
&\ssi& u^2=\frac{5-\sqrt{5}}{8} \quad\text{ou}\quad  u^2=\frac{5+\sqrt{5}}{8}\\
&\ssi& u=\pm\sqrt{\frac{5-\sqrt{5}}{8}} \quad\text{ou}\quad  u=\pm\sqrt{\frac{5+\sqrt{5}}{8}}
\end{eqnarray*}
Or $x>0$ donc
\[x=\sqrt{\frac{5-\sqrt{5}}{8}} \quad\text{ou}\quad x=\sqrt{\frac{5+\sqrt{5}}{8}}.\]
Puisque $\pi/10\in[0,\pi/6]$, on en déduit que $x>\sqrt{3/4}$. Or, si on évalue $16u^4-20u^2+5$ en $3/4$, on obtient $-19/16<0$, donc $3/4$ est entre les deux racines donc
\[\frac{5-\sqrt{5}}{8}<\frac{3}{4}.\]
En conclusion
\[\cos\p{\frac{\pi}{10}}=\sqrt{\frac{5+\sqrt{5}}{8}}.\]
\end{sol}
\end{exos}



\subsection{Forme trigonométrique}

\begin{definition}
Soit $z\in\Cs$. On appelle \emph{forme trigonométrique} de $z$ toute écriture
\[z=r \e^{\ii\theta}\]
où $r\in\RPs$ et $\theta\in\R$.
\end{definition}

\begin{remarques}
\remarque Si $z=r\e^{\ii\theta}$ est une forme trigonométrique, alors $r=\abs{z}$.
\remarque Tout nombre complexe non nul admet une forme trigonométrique. En pratique,
pour la déterminer, on force la factorisation de $z$ par $\abs{z}$ et on écrit le second terme sous la forme $\e^{\ii\theta}$.
\remarque Il existe deux moyens de représenter un même nombre
  complexe~: la forme cartésienne et la forme trigonométrique. La première
  est particulièrement adaptée aux calculs de sommes, tandis que la seconde
  est particulièrement adaptée aux calculs de produits.
  % Nous verrons cependant, que lorsqu'ils ont même module, il existe un moyen
  % simple pour mettre sous forme trigonométrique la somme de deux nombres
  % complexes écrits sous forme trigonométrique.
\end{remarques}

\begin{exos}
  \exo Mettre $-2-2\sqrt{3}\ii$ sous forme trigonométrique. %Celle de $1+2\ii$~?
  \begin{sol}
  On a
  \[-2-2\sqrt{3}\ii=4\p{-\frac{1}{2}-\ii\frac{\sqrt{3}}{2}}=4\e^{-\ii\frac{2\pi}{3}}.\]
  De même
  \[1+2\ii=\sqrt{5}\underbrace{\p{\frac{1}{\sqrt{5}}+\ii\frac{2}{\sqrt{5}}}}_{\eqdef u}\]
  Puisque $(1/\sqrt{5})^2+(2/\sqrt{5})^2=1$, il existe $\theta\in\interf{0}{\pi/2}$ tel que $u=\e^{\ii\theta}$. Alors
  \[1+2\ii=\sqrt{5}\e^{\ii\theta}.\]
  \end{sol}
  \exo Mettre sous forme trigonométrique
    \[\p{\frac{1+\ii \sqrt{3}}{1-\ii}}^{20}.\]
  \begin{sol}
  \begin{eqnarray*}
  \p{\frac{1+\ii \sqrt{3}}{1-\ii}}^{20}
  &=& \p{\frac{2\e^{\ii\frac{\pi}{3}}}{\sqrt{2}\e^{-\ii\frac{\pi}{4}}}}^{20}\\
  &=& 1024\e^{-\ii\frac{\pi}{3}}\\
  &=& 512(1-\ii\sqrt{3}).
  \end{eqnarray*}
  \end{sol}
  \end{exos}



\begin{definition}[utile=-3]
On appelle \emph{argument} de $z\in\Cs$ tout réel
$\theta$ tel que
\[z=\abs{z}\e^{\ii\theta}.\]
\end{definition}

\begin{proposition}[utile=-3]
Tout nombre complexe non nul $z\in\Cs$ admet au moins un argument. Si $\theta$ est l'un de ses arguments, l'ensemble de ses arguments est
$\theta+2\pi\Z\defeq\ensim{\theta+k2\pi}{k\in\Z}$. On écrit
\[\arg(z)\equiv\theta\ \cro{2\pi}.\]
\end{proposition}

\begin{preuve}
Si $z\neq 0$, alors $z/|z| \in \U$ et on utilise la proposition précédente.
\end{preuve}

\begin{remarques}
\remarque Soit $r_1,r_2>0$ et $\theta_1,\theta_2\in\R$. Alors
  \[r_1\e^{\ii\theta_1}=r_2\e^{\ii\theta_2}\]
  si et seulement si $r_1=r_2$ et $\theta_1\equiv\theta_2\ [2\pi]$. Contrairement à la forme cartésienne, il n'y a donc pas unicité de la forme trigonométrique.
% \remarque Si $z=a+\ii b$ est non nul, alors~:
%   \begin{itemize}
%   \item si $a=0$, $\arg(z)\equiv \frac{\pi}{2}\ \cro{\pi}$. L'argument est donné par le
%     signe de $b$.
%   \item sinon, $\tan\theta=b/a$ ce qui permet de connaitre l'argument de
%     $z$ modulo $\pi$. Pour le connaitre modulo $2\pi$, il suffit de regarder
%     le signe de $a$.
%   \end{itemize}
\remarque Si $z$ est un nombre complexe non nul, il existe un unique
  $\theta\in\interof{-\pi}{\pi}$ tel que $\arg z\equiv \theta\ [2\pi]$.
  On dit que $\theta$ est l'\emph{argument principal} de $z$ et on le note
  $\Arg(z)$.
\remarque Étant donné un nombre complexe $z$, on appelle
   {\it forme trigonométrique généralisée} de $z$ toute écriture du type
    $z=\rho \e^{\ii\theta}$ où $\rho\in\R$ et $\theta\in\R$. Attention, même si $z\neq 0$, on n'a pas nécessairement
    $\arg{z}\equiv\theta\ \cro{2\pi}$. En effet~:
    \begin{itemize}
    \item {Si $\rho>0$}, alors $\rho=\abs{z}$ et
      $\arg{z}\equiv\theta\ \cro{2\pi}$.
    \item {Si $\rho<0$}, alors $\rho=-\abs{z}$ et
      $\arg{z}\equiv\theta+\pi\ \cro{2\pi}$.
    \end{itemize}
  Lorsque l'énoncé demandera explicitement
  de mettre un nombre complexe sous forme trigonométrique, c'est bien
  sous la forme $z=r \e^{\ii\theta}$ avec $r>0$ qu'il faudra le mettre.
  Cependant, lorsqu'on demandera de mettre $z$ sous forme trigonométrique pour
  conduire des calculs, une forme trigonométrique généralisée suffira le plus
  souvent.
\end{remarques}

\begin{exos}
% \exo Résoudre l'équation $z^2=\conj{z}$.
\exo Résoudre l'équation
  \[z^2=\conj{z}\]
  en utilisant la forme cartésienne, puis la forme trigonométrique de $z$.
  \begin{sol}
Avec la forme cartésienne, on a
\begin{eqnarray*}
z^2=\conj{z}
&\ssi& (x+\ii y)^2=x-\ii y\\
&\ssi& (x^2-y^2)+2\ii x y=x-\ii y\\
&\ssi& x^2-y^2=x \et 2xy=-y\\
&\ssi& x^2-y^2=x \et y(2x+1)=0\\
&\ssi& x^2-y^2=x \et \cro{y=0 \ou x=-\frac{1}{2}}\\
&\ssi& \cro{x^2-y^2=x \et y=0} \ou \cro{x^2-y^2=x \et x=-\frac{1}{2}}\\
&\ssi& z=1 \ou z=0 \ou z=-\frac{1}{2}+\ii\frac{\sqrt{3}}{2} \ou z=-\frac{1}{2}-\ii\frac{\sqrt{3}}{2}.
\end{eqnarray*}
  Avec la forme trigonométrique. On sait que 0 est solution. On cherche donc les solutions non nulles. Soit $z\in\Cs$, alors il existe $r\in\RPs$ et $\theta\in\R$ tels que $z=r\e^{\ii\theta}$. On a
\begin{eqnarray*}
z^2=\conj{z}
&\ssi& r^2 \e^{2\ii\theta} = r \e^{-\ii\theta}\\
&\ssi& r^2=r \et 2\theta\equiv -\theta\ [2\pi]\\
&\ssi& r=1 \et \theta\equiv 0\ \cro{\frac{2\pi}{3}}\\
\end{eqnarray*}
Les solutions sont donc $0$, $1$ et $\e^{\pm\ii\frac{2\pi}{3}}$.
  \end{sol}
\exo Résoudre l'équation $z^5=1/\conj{z}$.
\begin{sol}
Cette équation a pour domaine de définition $\Cs$. On cherche donc les solutions non nulles. Soit $z\in\Cs$, alors il existe $r\in\RPs$ et $\theta\in\R$ tels que $z=r\e^{\ii\theta}$. On a
\begin{eqnarray*}
z^5=\frac{1}{\conj{z}}
&\ssi& r^5 \e^{5\ii\theta} = \frac{1}{r} \e^{\ii\theta}\\
&\ssi& r^6=1 \et 5\theta\equiv \theta\ [2\pi]\\
&\ssi& r=1 \et \theta\equiv 0\ \cro{\frac{\pi}{2}}\\
&\ssi& z=1 \ou z=\ii \ou z=-1 \ou z=-\ii.
\end{eqnarray*}
  \end{sol}
\end{exos}

\begin{proposition}[utile=-3]
  Soit $z,z'\in\Cs$ et $n\in\Z$. Alors
  \[\arg(z z')\equiv\arg(z)+\arg(z')\ \cro{2\pi}, \qquad
    \arg\p{\frac{z}{z'}}\equiv\arg(z)-\arg(z')\ \cro{2\pi},\]
  \[\arg(z^n)\equiv n\arg(z)\ \cro{2\pi}, \qquad
    \arg(\conj{z})\equiv -\arg(z)\ [2\pi].\]
\end{proposition}

\begin{preuve}
Il suffit d'écrire la forme trigonométrique.
\end{preuve}

\begin{remarqueUnique}
\remarque En général,
  $\Arg(z z')\neq\Arg(z)+\Arg(z')$. Par exemple $\Arg((-1)(-1))=0$ et $\Arg(-1)+\Arg(-1)=2\pi$.
\end{remarqueUnique}

\subsection{Exponentielle complexe}

\begin{definition}[utile=-3]
Soit $z=x+\ii y$ un nombre complexe où $x$ et $y$ sont respectivement la partie
réelle et imaginaire de $z$. On appelle exponentielle de $z$ et on note $\e^z$
le nombre complexe défini par
\[\e^z\defeq\e^{x}\p{\cos y +\ii \sin y}.\]
\end{definition}

\begin{proposition}[utile=-3]
Soit $z$ et $z'$ deux nombres complexes. Alors
\[\e^{0}=1 \qquad \text{et} \qquad \e^{z+z'}=\e^{z} \e^{z'}.\]
\end{proposition}

\begin{preuve}
On écrit $a$ et $b$ sous forme cartésienne et on déroule.
\end{preuve}

\begin{proposition}[utile=-3]
Soit $z$ un nombre complexe, et $n\in\Z$ . Alors $\e^z$ est non nul,
\[\frac{1}{\e^z}=\e^{-z} \qquad \text{et} \qquad \e^{nz}=\p{\e^z}^n.\]
\end{proposition}

\begin{preuve}
Voir la première partie comme un corollaire de la proposition précédente. 
\end{preuve}

\begin{proposition}[utile=-3]
Soit $z$ un nombre complexe. Alors
\[\conj{\e^z}=\e^{\conj{z}}, \qquad  \abs{\e^z}=\e^{\Re(z)}
\qquad \text{et} \qquad \arg(\e^z)\equiv \Im(z)\ [2\pi].\]
\end{proposition}

\begin{proposition}[utile=-3]
\begin{itemize}
\item Soit $z\in\C$. Alors $\e^z=1$ si et seulement si il existe
  $k\in\Z$ tel que $z=\ii k2\pi$.
\item  Plus précisément, étant donnés $z_1$ et $z_2$ deux nombres
  complexes, $\e^{z_1}=\e^{z_2}$ si et seulement si il existe
  $k\in\Z$ tel que $z_1=z_2+\ii k2\pi$.
\end{itemize}
\end{proposition}

\begin{proposition}[utile=-3]
L'exponentielle est une surjection de $\C$ dans $\Cs$. Autrement dit~:
\begin{itemize}
\item Pour tout $z\in\C$, $\e^{z}\in\Cs$.
\item Pour tout $z\in\Cs$, il existe $z'\in\C$ tel que $\e^{z'}=z$.
\end{itemize}
\end{proposition}

\begin{preuve}
Soit $a\in \Cs$ que l'on écrit $a=re^{i\theta}$ avec $r>0$ et $\theta \in \R$.
\begin{eqnarray*}
\forall z=x+iy \in \C, \quad e^z=a & \Longleftrightarrow &e^z=re^{i\theta} \\
& \Longleftrightarrow & \begin{cases}
    e^x=r\\
    y-\theta \in 2\pi\Z
    \end{cases}\\
& \Longleftrightarrow & \begin{cases}
    x=\ln(r)\\
    \exists k \in \Z, y=\theta+2k\pi
    \end{cases}
\end{eqnarray*}
$a$ admet donc un seul $x$ mais une infinité de $y$ donc une infinité d'antécédents qui sont exactements les $\set{\ln(r)+i(\theta+2k\pi), k\in \Z}$.


\end{preuve}

\begin{remarqueUnique}
\remarque Si $z\in\Cs$, nous venons de voir qu'il existe $z'\in\C$ tel que
  $e^{z'}=z$. Cependant, $z'$ n'est pas unique, ce qui nous empêche de définir
  le logarithme de $z$. Par contre, on peut montrer qu'il existe un unique $z'\in\C$ tel que $\e^{z'}=z$ et
  $\Im(z')\in\interof{-\pi}{\pi}$. Ce nombre est appelé logarithme
  principal de $z$ et noté $\Ln(z)$. De plus $\Ln(z)=\ln\abs{z}+\ii\Arg(z)$. C'est le logarithme calculé par les logiciels de calcul formel ainsi que vos calculatrices. Malheureusement, l'identité $\Ln\p{z z'}=\Ln(z)+\Ln(z')$ est fausse;
  elle n'est vraie que modulo $\ii2\pi$. C'est pourquoi, nous n'emploierons
  jamais de logarithme avec les nombres complexes.
\end{remarqueUnique}

\begin{exoUnique}
\exo Résoudre sur $\C$ l'équation $\e^z=\sqrt{3}+3\ii$.
  \begin{sol}
  $\sqrt{3}+3i=2\sqrt{3}\p{\dfrac{1}{2}+\dfrac{\sqrt 3}{2}i}=2\sqrt{3}e^{i\frac{\pi}{3}}$ donc on trouve $\ln\p{2\sqrt{3}}+\ii\frac{\pi}{3}+\ii k2\pi$.
  \end{sol}
\end{exoUnique}

\section{Racines d'un nombre complexe}
\subsection{L'équation du second degré}

\begin{definition}[utile=-3]
Soit $a$ un nombre complexe. On appelle \emph{racine} de $a$ tout nombre complexe $z$
tel que
\[z^2=a.\]
\end{definition}

\begin{remarques}
\remarque Si $a$ est un réel positif, les racines de $a$ sont $\sqrt{a}$ et $-\sqrt{a}$.
  Si $a$ est un réel négatif, ses racines sont $\ii\sqrt{\abs{a}}$ et
  $-\ii\sqrt{\abs{a}}$.
\remarque Si $a\in\C$, on parlera de racine de $a$, mais on n'écrira jamais $\sqrt{a}$. Cette notation est réservée aux réels positifs.
\end{remarques}

\begin{proposition}[utile=-3]
Soit $a$ un nombre complexe non nul. Alors $a$ admet exactement deux
racines distinctes opposées l'une à l'autre.
\end{proposition}

\begin{remarqueUnique}
\remarque En pratique, pour trouver les racines d'un nombre complexe $a$, on procède ainsi
  \begin{itemize}
    \item {\it Si $a$ s'exprime facilement sous forme trigonométrique.}
      On connait donc $r>0$ et $\theta$ tels que $a=r\e^{\ii\theta}$. Alors
      les racines de $a$ sont $\sqrt{r}\e^{\ii\frac{\theta}{2}}$
      et $-\sqrt{r}\e^{\ii\frac{\theta}{2}}$.
    \item {\it Si $a$ est sous forme cartésienne et qu'il n'est pas
      simple de le mettre sous forme trigonométrique.}
      On recherche les racines de $a$ sous la forme $z\defeq x+\ii y$ en
      effectuant une analyse~: on suppose que $z$ est une racine de $a$
      et on exploite le fait que $\abs{z}^2=\abs{a}$ et que $z^2$ et $a$
      ont même partie réelle. On obtient donc
      \begin{align*}
        x^2+y^2&=\abs{a},\\
        x^2-y^2&=\Re(a).
      \end{align*}
      En résolvant ce système linéaire en $x^2$ et $y^2$, on obtient $4$
      couples $(x,y)$ de solutions parmi lesquelles se trouvent les racines de 
      $a$. Un argument de signe sur les parties imaginaires de $z^2$ et
      $a$ permet d'éliminer deux candidats. La proposition précédente
      nous assure que les deux candidats restants sont bien des racines de
      $a$.
  \end{itemize}
\end{remarqueUnique}

\begin{exos}
\exo Calculer les racines carrées de $1+\ii\sqrt{3}$.
  \begin{sol}
  On trouve après mise sous forme trigonométrique
  \[\pm\sqrt{2}e^{i\frac{\pi}{6}}=\pm\frac{\sqrt{6}+\ii\sqrt{2}}{2}.\]
  \end{sol}
\exo Calculer les racines de $1+\ii$ de deux manières différentes.
  En déduire une expression avec des radicaux emboîtés de
  $\cos(\pi/8)$, $\sin(\pi/8)$ et $\tan(\pi/8)$.
  \begin{sol}
  Avec la forme cartésienne, on doit résoudre $x^2+y^2=\sqrt{2}$ et $x^2-y^2=1$. On trouve 4 couples mais on doit avoir $xy\geq 0$ ce qui fait qu'on garde 
  \[\sqrt{\frac{1+\sqrt{2}}{2}}+i\sqrt{\frac{-1+\sqrt{2}}{2}} \et
    -\sqrt{\frac{1+\sqrt{2}}{2}}-i\sqrt{\frac{-1+\sqrt{2}}{2}}\]
  et par la forme trigonométrique, on trouve $\pm 2^{\frac{1}{4}}e^{i\frac{\pi}{8}}$. Donc
  \[\cos\frac{\pi}{8}=\frac{1}{2^{\frac{1}{4}}}\sqrt{\frac{1+\sqrt{2}}{2}}=
    \frac{\sqrt{2+\sqrt{2}}}{2} \qquad
    \sin\frac{\pi}{8}=\frac{1}{2^{\frac{1}{4}}}\sqrt{\frac{-1+\sqrt{2}}{2}}=
    \frac{\sqrt{2-\sqrt{2}}}{2}\]
  \[\tan\frac{\pi}{8}=\sqrt{2}-1\]
  \end{sol}
\end{exos}

\begin{proposition}[utile=-3]
Soit $a,b,c\in\C$ avec $a\neq 0$.
% \begin{itemize}
%\item 
On considère l'équation
  \[(E) \qquad az^2+bz+c=0.\]
  On appelle discriminant de $(E)$ le nombre complexe $\Delta=b^2-4ac$.
  \begin{itemize}
  \item Si $\Delta\not=0$, le trinôme admet deux racines distinctes
    \[z_1\defeq\frac{-b+\delta}{2a} \qquad \text{et} \qquad
      z_2\defeq\frac{-b-\delta}{2a}.\]
    où $\delta$ est une racine carrée de $\Delta$.
  \item Si $\Delta=0$, le trinôme admet une seule racine, appelée racine
    double
    \[z_0\defeq-\frac{b}{2a}.\]
  \end{itemize}
% \item Lorsque l'équation s'écrit sous la forme~:
%   \[az^2+2bz+c=0\]
%   on utilise le discriminant réduit $\Delta'=b^2-ac$. Dans ce cas~:
%   \begin{itemize}
%   \item Si $\Delta'=0$, le trinôme admet une et une seule racine appelée racine
%     double~:
%     \[z_0=-\frac{b}{a}\]
%   \item Si $\Delta'\not=0$, le trinôme admet exactement deux racines distinctes~:
%     \[a=\frac{-b+\delta'}{a} \qquad \text{et} \qquad
%       b=\frac{-b-\delta'}{a}\]
%     où $\delta'$ est une racine carrée de $\Delta'$.
%   \end{itemize}
% \end{itemize}
\end{proposition}

\begin{preuve}
$\displaystyle z^2+\frac{b}{a}z+\frac{c}{a}=\p{z+\frac{b}{2a}}^2-\p{\frac{\delta}{2a}}^2=\p{z+\frac{b+\delta}{2a}}\p{z+\frac{b-\delta}{2a}}=0$.
\end{preuve}

\begin{remarqueUnique}
\remarque Soit $a,b,c\in\R$ avec $a\neq 0$. On considère l'équation
  \[(E) \qquad az^2+bz+c=0.\]
  \begin{itemize}
  \item Si $\Delta>0$, le trinôme admet deux racines réelles distinctes
    \[z_1\defeq\frac{-b+\sqrt{\Delta}}{2a} \qquad \text{et} \qquad
      z_2\defeq\frac{-b-\sqrt{\Delta}}{2a}.\]
  \item Si $\Delta=0$, le trinôme admet une seule racine réelle, appelée racine
    double
    \[z_0\defeq-\frac{b}{2a}.\]
  \item Si $\Delta<0$, le trinôme admet deux racines complexes conjuguées
    \[z_1\defeq\frac{-b+\ii\sqrt{\abs{\Delta}}}{2a} \qquad \text{et} \qquad
      z_2\defeq\frac{-b-\ii\sqrt{\abs{\Delta}}}{2a}.\]
  \end{itemize}     
\end{remarqueUnique}

\begin{exoUnique}
\exo Soit $\theta\in\R$. Résoudre sur $\C$ l'équation
  $z^2-2\cos(\theta)z+1=0$.
  \begin{sol}
  Les racines sont $\e^{\ii\theta}$ et $\e^{-\ii\theta}$.
  \end{sol}
\end{exoUnique}

\begin{proposition}
Soit $a,b,c\in\C$ avec $a\not=0$ et $z_1,z_2$ deux nombres complexes. Alors $z_1$
et $z_2$ sont les deux racines, éventuellement égales, de l'équation
$az^2+bz+c=0$ si et seulement si
\[z_1+z_2=-\frac{b}{a} \qquad \text{et} \qquad z_1 z_2=\frac{c}{a}.\]
\end{proposition}

\begin{preuve}
$\displaystyle z^2+\frac{b}{a}z+\frac{c}{a}=(z-z_1)(z-z_2)$.
\end{preuve}
\begin{remarqueUnique}
\remarque Si $P,S\in\C$, les solutions du système
  \[\syslin{z_1+z_2&=&S\hfill\cr
             z_1 z_2&=&P,\hfill}\]
  sont $\p{\omega_1,\omega_2}$ et $\p{\omega_2,\omega_1}$ où $\omega_1$
  et $\omega_2$ sont les racines du trinôme $z^2-Sz+P=0$.
\end{remarqueUnique}
\begin{exos}
\exo Déterminer les solutions de l'équation $3z^2-5z+2=0$.
\exo Résoudre sur $\C$ le système
  \[\syslin{u^2+v^2&=&3-2\ii\hfill\cr
                 uv&=&3+\ii.\hfill}\]
  \begin{sol}
  On résout le système                 
  \[\syslin{a+b&=&3-2\ii\hfill\cr
             ab&=&8+6\ii\hfill}\]
  On obtient l'équation du second degré $z^2-\p{3-2\ii}z+\p{8+6\ii}=0$ dont le
  discriminant est $-27-36\ii$. $3-6\ii$ en est une racine donc les racines
  du trinôme sont $3-4\ii$ et $2\ii$. Donc $a=3-4\ii$ et $b=2\ii$, ou le contraire.
  Les racines de $3-4\ii$ sont $\pm\p{2-\ii}$ et les racines de $2\ii$ sont
  $\pm\p{1+\ii}$. Ce qui nous donne au final 4 solutions dans le premier cas.
  On en élimine 2 en regardant le signe de la partie réelle de $uv$. Il
  reste $\p{u,v}=\p{2-\ii,1+\ii}$ et $\p{u,v}=\p{-2+\ii,-1-\ii}$. On rajoute bien sûr
  les deux autres solutions.
  \end{sol}
\end{exos}

\subsection{Racines $n$-ièmes}

\begin{definition}[utile=-3]
Étant donné $n\in\Ns$ et $a\in\C$, on appelle \emph{racine $n$-ième} de $a$ tout
nombre complexe $z$ tel que $z^n=a$. Les racines $n$-ièmes de 1 sont appelées
\emph{racines $n$-ièmes de l'unité} et l'ensemble de ces racines est noté $\U[n]$.
\end{definition}

\begin{remarqueUnique}
\remarque Les racines $n$-ièmes de 1 sont de module 1. Autrement dit, $\U[n]\subset\U$.
\end{remarqueUnique}


% \subsection{Division euclidienne}

\begin{proposition}[utile=-3]
Soit $n\in\Ns$. Il existe exactement $n$ racines $n$-ièmes de l'unité. En posant $\omega\defeq\e^{\ii\frac{2\pi}{n}}$, ce sont
\[1,\omega,\ldots,\omega^{n-1}.\]
\end{proposition}

\begin{preuve}
Soit $z\in\R$. On peut fixer $\theta\in\R$ tel que $z=e^{i\theta}$.
\begin{eqnarray*}
z^n=1 &\Longleftrightarrow& e^{in\theta}=1 \text{ d'après Moivre}\\
&\Longleftrightarrow& \exists k\in \Z, n\theta=2k\pi \\
&\Longleftrightarrow& \exists k\in \Z, \theta=\dfrac{2k\pi}{n}
\end{eqnarray*}
On a montré ainsi que $\U[n]=\ens{e^{\frac{2ik\pi}{n}}, k\in \Z}$.

Soit $k\in \Z$. \'Ecrivons la division euclidienne de $k$ par $n$. $\exists ! (q,r) \in \Z^2 \text{ tel que } k=qn+r \text{ et } 0\leq r<n$. On a alors :
\[e^{\frac{2ik\pi}{n}}=e^{\frac{2i(qn+r)\pi}{n}}=e^{2iq\pi}e^{\frac{2ir\pi}{n}}=e^{\frac{2ir\pi}{n}},\]
ce qui montre que $\U[n]\subset\ens{e^{\frac{2ik\pi}{n}}, 0\leq k\leq n-1}$. L'autre inclusion étant claire, on a bien l'égalité.

Il reste à voir qu'elles sont bien toutes distinctes (et qu'on en a donc exactement $n$). Pour cela, supposons qu'il existe $k$ et $\ell$ dans $\llbracket0,n-1\rrbracket$ tels que $e^{\frac{2ik\pi}{n}} =e^{\frac{2i\ell\pi}{n}}$. Ainsi, $e^{\frac{2i(k-\ell)\pi}{n}}=1$ et donc il existe $m\in\Z$ tel que $\dfrac{2(k-\ell)\pi}{n}=2m\pi$. Mais alors, $n$ divise $k-\ell$. Or $-(n-1)\leq k-\ell \leq n-1$ donc $k-l=0$.
\end{preuve}

\begin{remarques}
\remarque Lorsque l’on doit calculer sur des racines $n$-ièmes, il est souvent plus efficace de les manipuler via leur propriété ($z^n = 1$) plutôt que par leur description ($z=\omega^k$). On ne se rabat sur la description que lorsque la relation $z^n = 1$ ne suffit pas, ou en toute fin de calcul.
\remarque Dans le cas où $n=3$, $\omega$ est noté $\jj$. Les racines
  3-ièmes de l'unité sont donc $1,\jj,\jj^2$. Lorsqu'on travaille avec le nombre
  complexe $\jj$, on exploite les relations
  \[\jj^3=1, \qquad 1+\jj+\jj^2=0 \et \frac{1}{\jj}=\conj{\jj}=\jj^2.\]
% \remarque Attention, si $n\geq 2$, l'égalité $a^n=b^n$ n'implique pas $a=b$.
\remarque Les racines $n$-ièmes de l'unité forment un polygone régulier à $n$ côtés.
\pgfplotsset{
    standard/.style={
        axis x line=middle,
        axis y line=middle,
        enlarge x limits=0.15,
        enlarge y limits=0.15,
        every axis x label/.style={at={(current axis.right of origin)},anchor=north west},
        every axis y label/.style={at={(current axis.above origin)},anchor=north east}
    }
}
\begin{figure}[H]
\centering
\begin{tikzpicture}
\begin{axis}[standard,axis equal,mark=none,
  ymax=1,ymin=-1,
  domain=0:3,samples=4,disabledatascaling,
  xtick=0,xticklabels=0,
  ytick=0,yticklabels=0,colormap={slategraywhite}{rgb255=(0,0,0) rgb255=(0,0,0)}]
\addplot[scatter,color=black]({cos(2*pi*deg(x)/3)},{sin(2*pi*deg(x)/3)});
\draw[help lines] (0,0) circle (1);
\node[above right] at (1,0) {$1$};
\node[above left] at (-0.5,0.8660) {$\jj$};
\node[below left] at (-0.5,-0.8660) {$\jj^2$};
\end{axis}
\end{tikzpicture}
$\qquad$
\begin{tikzpicture}
\begin{axis}[standard,axis equal,mark=none,
  ymax=1,ymin=-1,
  domain=0:7,samples=8,disabledatascaling,
  xtick=0,xticklabels=0,
  ytick=0,yticklabels=0,colormap={slategraywhite}{rgb255=(0,0,0) rgb255=(0,0,0)}]
\addplot[scatter,color=black]({cos(2*pi*deg(x)/7)},{sin(2*pi*deg(x)/7)});
\draw[help lines] (0,0) circle (1);
\node[above right] at (1,0) {$1$};
\node[above right] at (0.62349,0.781831) {$\omega$};
\node[above] at (-0.2225,0.9749) {$\omega^2$};
\node[below right] at (0.62349,-0.781831) {$\omega^{n-1}$};
\node[below] at (-0.2225,-0.9749) {$\omega^{n-2}$};
\end{axis}
\end{tikzpicture}
\end{figure}
% \remarque Si $\omega\defeq\e^{\ii 2\pi/n}$, alors~:
%   $\forall k\in\Z\qsep \omega^k=1 \quad\ssi\quad k\equiv 0\ [n].$
\end{remarques}

\begin{exos}
% \exo Soit $f$ la fonction d'expression
%   \[f(z)\defeq\frac{z^{12}+3z^7-4z^4+6z+1}{z^6+z^4+z^2-1}.\]
%   Calculer $f(\ii)$, $f(\jj)$ et $f(\jj^2)$.
\exo Que dire de deux nombres complexes $a,b\in\C$ tels que $a^3=b^3$~?

\begin{sol}
Ou bien ils sont nuls, ou bien $a/b \in \ens{1,j,j^2}$.
\end{sol}
\exo Soit $n\in\Ns$. Résoudre sur $\C$ l'équation
  \[\p{z+1}^n=\p{z-1}^n.\]
  \begin{sol}
  Déjà, $z\neq 1$.
  \begin{eqnarray*}
  \p{\frac{z+1}{z-1}}^n=1 &\Longleftrightarrow & \exists k\in \llbracket 0,n-1\rrbracket\text{ tel que } \frac{z+1}{z-1}= e^{\frac{2ik\pi}{n}}\\
  &\Longleftrightarrow & z(1-e^{\frac{2ik\pi}{n}})=-1-e^{\frac{2ik\pi}{n}}\\
  &\Longleftrightarrow & z=\frac{1+e^{\frac{2ik\pi}{n}}}{e^{\frac{2ik\pi}{n}}-1} \text{ en reprenant les équivalences avec }k\neq 0
  \end{eqnarray*}
  On trouve $z=-\ii\cotan\p{\frac{k\pi}{n}}$ pour $k\in\intere{1}{n-1}$ en passant finalement à l'angle moitié.
  \end{sol}

\exo Calculer
  \[\sum_{z\in\U[n]} \abs{z-1}.\]
  \begin{sol}
\begin{eqnarray*}
  \sum_{z\in\U[n]} \abs{z-1}&=& \sum_{k=0}^{n-1} \abs{e^{\frac{2ik\pi}{n}}-1}\\
  &=& \sum_{k=0}^{n-1} \abs{e^{\frac{ik\pi}{n}}\p{e^{\frac{ik\pi}{n}}-e^{\frac{-ik\pi}{n}}}}\\
  &=& \sum_{k=0}^{n-1}2\abs{\sin\p{\frac{k\pi}{n}}}\\
  &=& 2\sum_{k=0}^{n-1}\sin\p{\frac{k\pi}{n}}\\
  &=& 2\sum_{k=0}^{n-1}\Im\p{e^{\frac{ik\pi}{n}}}\\
  &=& 2\Im\p{\sum_{k=0}^{n-1}\p{e^{\frac{i\pi}{n}}}^k}\\
  &=& 2\Im\p{\frac{1-e^{i\pi}}{1-e^{i\pi/n}}}\\
  &=& 2\frac{2}{-e^{\frac{i\pi}{2n}}2i\sin(\frac{\pi}{2n})}\\
  &=& 2\frac{ie^{\frac{-i\pi}{2n}}}{\sin(\frac{\pi}{2n})}\\
  &=& 2\cotan\p{\frac{\pi}{2n}}
\end{eqnarray*}
  On peut vérifier avec $n=4$ et le résultat obtenu dans l'exercice 12 :  $2\cotan\p{\frac{\pi}{8}}=\dfrac{2}{\sqrt{2}-1}$.
  \end{sol}
% \exo Résoudre sur $\C$ l'équation $z^2=27\conj{z}$.
%   \begin{sol}
%   On trouve $z=0,27,27j,27j^2$.
%   \end{sol}
\end{exos}

\begin{proposition}[utile=-3]
Soit $n\geq 2$ et $\zeta$ une racine $n$-ième de l'unité, différente de $1$.
Alors
\[1+\zeta+\dots+\zeta^{n-1}=0.\]
En particulier, la somme des racines $n$-ièmes de l'unité est nulle.
\end{proposition}

\begin{preuve}
Somme géométrique et/ou développement du polynôme factorisé.
\end{preuve}

% \begin{exos}
% \exo Calculer
%   \[\cos\frac{\pi}{11}+\cos\frac{3\pi}{11}+\cos\frac{5\pi}{11}+
%     \cos\frac{7\pi}{11}+\cos\frac{9\pi}{11}\]
% \end{exos}

\begin{proposition}[utile=-3]
Soit $a\in\Cs$ et $n\in\Ns$. Alors $a$ admet exactement $n$ racines $n$-ièmes. En posant $\omega\defeq\e^{\ii\frac{2\pi}{n}}$, si $z_0$ est une racine $n$-ième de $a$, les racines $n$-ièmes de $a$ sont
\[z_0,\omega z_0,\ldots,\omega^{n-1}z_0.\]
\end{proposition}

\begin{preuve}
$z_0\neq 0$ et $z/z_0 \in \U[n]$.

\end{preuve}

\begin{remarques}
\remarque Si $a=r\e^{\ii\theta}$ est sous forme trigonométrique, alors 
  \[z_0=\sqrt[n]{r}\e^{\ii\frac{\theta}{n}}\]
  est une racine $n$-ième de $a$.
\remarque Si $n\geq 2$, la somme des racines $n$-ièmes d'un nombre complexe est
  nulle.
\end{remarques}

\begin{exos}
\exo Résoudre sur $\C$ l'équation
  \[27\p{z-1}^6+\p{z+1}^6=0.\]
  \begin{sol}
  Comme $z\neq 1$, cela revient à $\displaystyle \p{\frac{z+1}{z-1}}^6=-27$. $z_0=\sqrt{3}e^{i\pi/6}$ en est une solution. D'où $\frac{z+1}{z-1}=\sqrt{3}e^{i\pi/6}e^{2ik\pi/6}$.
  On trouve
  \[z=\frac{\sqrt{3}\e^{\ii\frac{(1+2k)\pi}{6}}+1}{\sqrt{3}\e^{\ii\frac{(1+2k)\pi}{6}}-1}\]
   pour $k\in\intere{0}{5}$
  % On trouve
  % \[z=-\frac{1+i\sqrt{3}\omega^k}{1-i\sqrt{3}\omega^k}(faute)=
  %   \frac{1-i\sqrt{3}\cos\p{\frac{k\pi}{3}}}{2+\sqrt{3}\sin\p{\frac{k\pi}{3}}}
  %   (juste)\]
  % avec $k\in\intere{0}{5}$ et $\omega=e^{i\frac{\pi}{3}}$.
  \end{sol}
\exo En considérant les racines 7-ièmes de $-1$, montrer que
  \[\cos\frac{\pi}{7}-\cos\frac{2\pi}{7}+\cos\frac{3\pi}{7}=\frac{1}{2}.\]
  \begin{sol}
  On résout $z^7=-1$ en prenant $z_0=e^{i\pi/7}$. On exploite le fait que la somme des racines est nulle en en prenant la partie réelle. On regroupe les cosinus identiques et on conclut en remarquant que $\cos(5\pi/7)=-\cos(2\pi/7)$.
  \end{sol}
\end{exos}


\section{Nombres complexes et géométrie plane}

\subsection{Le plan complexe}
\begin{definition}
Soit $\mathcal{P}$ le plan euclidien orienté et $\mathcal{R}=\p{O,\ve{e_1},\ve{e_2}}$ un repère orthonormé direct.
\begin{itemize}
\item Si $M$ est un point du plan de coordonnées $(x,y)\in\R^2$
  \[\ve{OM}=x \ve{e_1}+y \ve{e_2}\]
  on appelle \emph{affixe} de $M$ le nombre complexe $x+\ii y$.
\item Si $\ve{u}$ est un vecteur de coordonnées $(x,y)\in\R^2$
  \[\ve{u}=x \ve{e_1}+y \ve{e_2}\]
  on appelle \emph{affixe} de $\ve{u}$ le nombre complexe $x+\ii y$.
\end{itemize}
\end{definition}

\begin{remarqueUnique}
\remarque Si $M$ est un point du plan, son affixe est l'affixe du vecteur $\ve{OM}$.
\end{remarqueUnique}

\begin{proposition}
\begin{itemize}
\item Soit $A$ et $B$ deux points du plan d'affixes respectives $a$ et $b\in\C$. Alors $\ve{AB}$ a pour affixe $b-a$.
\item Soit $\ve{u}$ et $\ve{v}$ deux vecteurs d'affixes respectives $u$ et $v\in\C$ et $\lambda,\mu\in\R$. Alors $\lambda\ve{u}+\mu\ve{v}$ a pour affixe $\lambda u+\mu v$.
\end{itemize}
\end{proposition}

% \begin{definition}
% Soit $A_1,\ldots,A_n$ des points du plan d'affixes respectifs $a_1,\ldots,a_n\in\C$ et $\lambda_1,\ldots,\lambda_n\in\R$ tels que
% $\lambda_1+\cdots+\lambda_n\neq 0$. On appelle \emph{barycentre} des points pondérés $(A_1,\lambda_1),\ldots,(A_n,\lambda_n)$ le point d'affixe
% \[z\defeq\frac{\lambda_1 a_1+\cdots+\lambda_n a_n}{\lambda_1+\cdots+\lambda_n}.\]
% \end{definition}

% \begin{remarques}
% \remarque Le milieu de $[AB]$ est le barycentre du système $(A,1), (B,1)$. En particulier, si $a$ et $b$ sont les affixes respectifs de $A$ et $B$, l'affixe du milieu de $[AB]$ est $(a+b)/2$.
% \remarque Si $A$ et $B$ sont deux points du plan, le segment $[AB]$ est l'ensemble des barycentres du système $(A,\lambda)$, $(B,1-\lambda)$ pour $\lambda\in\interf{0}{1}$. La droite $(AB)$ est l'ensemble des barycentres du système $(A,\lambda)$, $(B,1-\lambda)$ pour $\lambda\in\R$.
% \remarque Si $A_1,\ldots,A_n$ sont des points du plan, $\lambda_1,\ldots,\lambda_n\in\R$ sont tels que
% $\lambda_1+\cdots+\lambda_n\neq 0$ et $\lambda\in\Rs$, le barycentre du système $(A_1,\lambda \lambda_1),\ldots,(A_n,\lambda \lambda_n)$ est égal au barycentre du système $(A_1,\lambda_1),\ldots,(A_n,\lambda_n)$.
% \remarque On dit qu'un point $M$ est dans l'\emph{enveloppe convexe} de $A_1,\ldots,A_n$ lorsqu'il existe $\lambda_1,\ldots,\lambda_n\in\RP$ tels que $M$ est le barycentre du système $(A_1,\lambda_1),\ldots,(A_n,\lambda_n)$.
% \end{remarques}


\begin{proposition}[utile=-3]
\begin{itemize}
\item Soit $a$ et $b$ deux nombres complexes. Alors $\abs{a-b}$ est la distance entre les points d'affixes $a$ et $b$.
\item Soit $u$ un nombre complexe. Alors $\abs{u}$ est la norme du vecteur
$\ve{u}$.
\end{itemize}
\end{proposition}

\begin{proposition}[utile=-3]
Soit $A,B,C$ trois points deux à deux distincts d'affixes respectives $a,b,c$. Alors
\[\abs{\frac{c-a}{b-a}} = \frac{AC}{AB} \quad\text{et}\quad   \arg\p{\frac{c-a}{b-a}}\equiv\p{\overrightarrow{AB}, \overrightarrow{AC}}\  [2 \pi].\]
\end{proposition}

\begin{preuve}
Le premier point est immédiat Le second découle de l'interprétation géométrique de l'argument. En notant $\overrightarrow{i}$ le vecteur d'affixe $1$ et en posant $\theta=\arg(z)$, on a $z=|z|(\cos(\theta)+i\sin(\theta))$ donc $\p{\overrightarrow{i}, \overrightarrow{u}})\equiv \arg(z)[2\pi]$. De même, $\p{\overrightarrow{i}, \overrightarrow{u'}})\equiv \arg(z')[2\pi]$. D'où :


\begin{eqnarray*}
\p{\overrightarrow{u}, \overrightarrow{u'}}=& \p{\overrightarrow{i}, \overrightarrow{u'}}-\p{\overrightarrow{i}, \overrightarrow{u}}&[2\pi]\\
=&  \arg(z')-\arg(z)   &[2\pi]\\
=& \arg\p{\frac{z'}{z}} &[2\pi]
\end{eqnarray*}
\end{preuve}

\begin{proposition}[utile=-3]
Soit $A,B,C$ trois points deux à deux distincts d'affixes respectives $a,b,c$. Alors
\begin{itemize}
\item $A,B,C$ sont alignés si et seulement si
  \[\frac{c-a}{b-a}\in\R.\]
\item $(AB)$ et $(AC)$ sont orthogonales si et seulement si
  \[\frac{c-a}{b-a}\in\ii\R.\]
\end{itemize}
\end{proposition}

\begin{remarques}
\remarque Soit $A$, $B$, $C$ trois points d'affixes respectives $a$, $b$, $c\in\C$. Si $A$, $B$, $C$ sont deux à deux distincts, alors
  \begin{eqnarray*}
  \text{$A$, $B$, $C$ sont alignés}
  &\ssi& \frac{c-a}{b-a}\in\R\\
  &\ssi& \frac{c-a}{b-a}=\conj{\p{\frac{c-a}{b-a}}}\\
  &\ssi& (c-a)\conj{(b-a)}=\conj{(c-a)}(b-a).
  \end{eqnarray*}
  On vérifie facilement que même si $A$, $B$, $C$ ne sont pas deux à deux distincts
  \[\text{$A$, $B$, $C$ sont alignés} \quad\ssi\quad (c-a)\conj{(b-a)}=\conj{(c-a)}(b-a).\]
\remarque La proposition précédente étant essentiellement utilisée de cette manière, on pourra tolérer exceptionnellement de l'appliquer, même si $A$, $B$, $C$ ne sont pas deux à deux distincts. Ce genre de \og division par zéro \fg est parfois tolérée en géométrie. Bien entendu, dans tout autre domaine des mathématiques, ces horreurs ne seront pas tolérées.
\end{remarques}

\begin{exoUnique}
\exo Soit $ABCD$ un quadrilatère non croisé. On construit $A_1$ extérieur
au quadrilatère tel que le triangle $BA_1C$ est isocèle et
rectangle en $A_1$. De même pour $B_1,C_1,D_1$. Montrer que les
segments $[A_1C_1]$ et $[D_1B_1]$ ont même longueur et sont
orthogonaux.
\begin{sol}
On suppose que $(ABCD)$ est dans le sens direct. Alors
\[a_1=c+\frac{1}{2}(1+\ii)(b-c), \qquad
  b_1=d+\frac{1}{2}(1+\ii)(c-d),\]
\[c_1=a+\frac{1}{2}(1+\ii)(d-a), \qquad
  d_1=b+\frac{1}{2}(1+\ii)(a-b)\]
donc $a_1-c_1=\ii(d_1-y)$.

Sol. Victor :
$c-a_1=i(b-a_1)$ donc $a_1=\frac{1}{1-i}(c-ib)$. De même $b_1=\frac{1}{1-i}(d-ic)$, $c_1=\frac{1}{1-i}(a-id)$ et $d_1=\frac{1}{1-i}(b-ia)$. D'où $a_1-c_1=i(b_1-d_1)$.
\end{sol}
\end{exoUnique}

\subsection{Les similitudes directes}

\begin{definition}
Soit $\ve{u}$ un vecteur. On appelle \emph{translation} de vecteur $\ve{u}$ l'application qui au point $M$ associe l'unique point $M'$ tel que
\[\ve{MM'}=\ve{u}.\]
\end{definition}

\begin{proposition}
Soit $\ve{u}$ un vecteur d'affixe $u\in\C$. La translation de vecteur $\ve{u}$ transforme le point $M$ d'affixe $z$ en le point $M'$ d'affixe
\[z'=z+u.\]
\end{proposition}

\begin{remarqueUnique}
\remarque Une translation conserve les distances et les angles.
\end{remarqueUnique}

\begin{definition}
Soit $\Omega$ un point du plan et $\rho\in\Rs$. On appelle \emph{homothétie} de centre $\Omega$ et de rapport $\rho$ l'application qui au point $M$ associe l'unique point $M'$ tel que
\[\ve{\Omega M'}=\rho\ve{\Omega M}.\]
\end{definition}

\begin{proposition}
Soit $\Omega$ un point du plan d'affixe $\omega\in\C$ et $\rho\in\Rs$. L'homothétie de centre $\Omega$ et de rapport $\rho$ transforme le point $M$ d'affixe $z$ en le point $M'$ d'affixe
\[z'=\rho(z-\omega)+\omega.\]
\end{proposition}

\begin{remarqueUnique}
\remarque Une homothétie de rapport $\rho\in\Rs$ multiplie les distances par $\abs{\rho}$
	et conserve les angles.
\end{remarqueUnique}

\begin{definition}
Soit $\Omega$ un point du plan et $\theta\in\R$. On appelle \emph{rotation} de centre $\Omega$ et d'angle $\theta$ l'application qui au point $M$ associe
\begin{itemize}
\item $\Omega$ si $M=\Omega$.
\item l'unique point $M'$ tel que
  \[\Omega M'=\Omega M \et \p{\ve{\Omega M}, \ve{\Omega M'}}\equiv\theta\,\cro{2\pi}\]
  sinon.
\end{itemize}
\end{definition}

\begin{proposition}
Soit $\Omega$ un point du plan d'affixe $\omega\in\C$ et $\theta\in\R$. La rotation de centre $\Omega$ et d'angle $\theta$ transforme le point $M$ d'affixe $z$ en le point $M'$ d'affixe
\[z'=\e^{\ii\theta}(z-\omega)+\omega.\]
\end{proposition}

\begin{remarqueUnique}
\remarque Une rotation conserve les distances et les angles.
\end{remarqueUnique}

\begin{exoUnique}
\exo Quelle est l'expression en notation complexe des transformations suivantes~?
\begin{itemize}
\item \textbf{a.} La symétrie centrale de centre $0$.
\item \textbf{b.} L'homothétie de centre $0$ et de rapport $2$.
\item \textbf{c.} L'homothétie de centre $2$ et de rapport $1/2$.
\item \textbf{d.} La composée des deux dernières transformations.
\item \textbf{e.} La rotation de centre $0$ et d'angle $\pi/2$.
\item \textbf{f.} La rotation de centre $1+\ii$ et d'angle $\pi/2$.
\item \textbf{g.} La composée des deux dernières transformations.
\item \textbf{h.} La symétrie orthogonale d'axe $(Ox)$.
\item \textbf{i.} La symétrie orthogonale dont l'axe $\mathcal{D}_\theta$ fait un angle $\theta$ avec l'axe $(Ox)$.
\end{itemize}
\begin{sol}
\begin{itemize}
\item $f(z)=-z$.
\item $h_1(z)=2z$.
\item $h_2(z)=(1/2)(z-2)+2$.
\item $[h_1\circ h_2](z)=2+z$, $[h_2\circ h_1](z)=1+z$.
\item $r_1(z)=\ii z$.
\item $r_2(z)=\ii(z-(1+\ii))+(1+\ii)$.
\item $[r_1\circ r_2](z)=2\ii-z=-(z-\ii)+\ii$ (symétrie de centre $\ii$ ou rotation de centre $\ii$ et d'angle $\pi$), $[r_2\circ r_1](z)=2-z=-(z-1)+1$ (symétrie de centre $1$).
\item $f(z)=\conj{z}$.
\item $f(z)=\e^{2\ii\theta}\conj{z}$.
\end{itemize}
\end{sol}
\end{exoUnique}

\begin{definition}[utile=-3]
Soit $\Omega$ un point du plan, $r\in\RPs$ et $\theta\in\R$. On appelle \emph{similitude} de centre $\Omega$, de rapport $r$ et d'angle $\theta$ la composée (commutative) de l'homothétie de centre $\Omega$ et de rapport $r$ et de la rotation de centre $\Omega$ et d'angle $\theta$.
\end{definition}

\begin{proposition}[utile=-3]
Soit $\Omega$ un point du plan d'affixe $\omega$, $r\in\RPs$ et $\theta\in\R$. La similitude de centre $\Omega$, de rapport $r$ et d'angle $\theta$ transforme le point $M$ d'affixe $z$ en le point $M'$ d'affixe
\[z'=r\e^{\ii\theta}(z-\omega)+\omega.\]
\end{proposition}

\begin{remarques}
\remarque Si $\rho<0$, l'homothétie de centre $\Omega$ et de rapport $\rho$ est
  une similitude de centre $\Omega$, de rapport $\abs{\rho}$ et d'angle $\pi$.
\remarque Une similitude de rapport $r>0$ et d'angle $\theta$ multiplie les distances par
  $r$ et conserve les angles.
\end{remarques}

Dans la suite, on confondra un point et son affixe, un vecteur et son affixe. On identifie ainsi le plan à $\C$.

\begin{definition}
On appelle \emph{similitude directe} toute application $f$ de $\C$ dans $\C$ telle qu'il existe $a\in\Cs$ et $b\in\C$ tels que
\[\forall z\in\C\qsep f(z)=az+b.\]
\end{definition}

\begin{remarqueUnique}
\remarque Les translations et les similitudes de centre $\Omega$ de rapport $r\in\RPs$ et d'angle $\theta\in\R$ sont des similitudes directes. Nous allons voir que ce sont les seules.
\end{remarqueUnique}

\begin{proposition}
La composée de deux similitudes directes est une similitude directe.
\end{proposition}

\begin{proposition}
Soit $f$ une similitude directe et $a\in\Cs$, $b\in\C$ tels que
\[\forall z\in\C\qsep f(z)=az+b.\]
\begin{itemize}
\item Si $a=1$, $f$ est la translation de vecteur $b$.
\item Sinon, il existe $r\in\RPs$ et $\theta\in\R$ tels que $a=r\e^{\ii\theta}$. $f$ admet un unique point fixe $\omega$ et
\[\forall z\in\C\qsep f(z)=r\e^{\ii\theta}(z-\omega)+\omega.\]
Autrement dit $f$ est la similitude de centre $\omega$, de rapport $r$ et d'angle $\theta$.
\end{itemize}
\end{proposition}

\begin{preuve}
Soit $a\in\Cs$, $b\in\C$ et $f$ l'application de $\C$ dans $\C$ définie par
\[\forall z\in\C\qsep f(z)\defeq az+b.\]
Alors $f$ est une similitude. En effet~:
\begin{itemize}
\item Si $a=1$, alors $f(z)=z+b$ donc $f$ est la translation du vecteur d'affixe $b$.
\item Sinon, $f$ admet un unique point fixe. En effet
\begin{eqnarray*}
\forall z\in\C\qsep f(z)=z
&\ssi& az+b = z\\
&\ssi& (1-a)z=b\\
&\ssi& z=\frac{b}{1-a}.
\end{eqnarray*}
On pose $\omega\defeq b/(1-a)$. En mettant $a$ sous forme trigonométrique $a= r\e^{\ii\theta}$, on vérifie que
\[\forall z\in\C\qsep f(z)=r\e^{\ii\theta}\p{z-\omega}+\omega.\]
Donc $f$ est la similitude de centre d'affixe $\omega$, de rapport $r$ et d'angle $\theta$.
\end{itemize}
\end{preuve}

\begin{exoUnique}
\exo {\`A} quelle transformation géométrique correspond la fonction $f : z \mapsto
(3-\ii)+2\ii z$~?
\begin{sol}
C'est la similitude de centre $1+\ii$ d'angle $\pi/2$ et de rapport 2.
\end{sol}
\end{exoUnique}


%END_BOOK

\end{document}
\section{Exercices}
\setcounter{numeroexercice}{1}
\documentclass{magnolia}

\magtex{tex_driver={pdftex}}
\magfiche{document_nom={Exercices sur les nombres complexes},
          auteur_nom={François Fayard},
          auteur_mail={fayard.prof@gmail.com}}
\magexos{exos_matiere={maths},
         exos_niveau={mpsi},
         exos_chapitre_numero={1},
         exos_theme={Nombres Complexes}}
\magmisenpage{misenpage_presentation={tikzvelvia},
         misenpage_format={a4},
         misenpage_nbcolonnes={1},
         misenpage_preuve={non},
         misenpage_sol={non}}
\maglieudiff{}
\magprocess

\begin{document}

%BEGIN_BOOK

\magsection{Le corps des nombres complexes}
\magsubsection{Définition, conjugaison, module}

\exercice{nom={Lieu}}
\begin{questions}
\question On note $f$ la fonction définie sur $\C\setminus\ens{2}$ par
  \[f(z)\defeq\frac{z+1}{z-2}.\]
  Pour quels nombres $z\in\C\setminus\ens{2}$ a-t-on $\abs{f(z)}=1$~? $\Re\p{f(z)}=0$~?
\question On note $g$ la fonction définie sur $\C\setminus\ens{2\ii}$ par
  \[g(z)\defeq\frac{2z-\ii}{z-2\ii}.\]
  Pour quels nombres $z\in\C\setminus\ens{2\ii}$ a-t-on $g(z)\in\R$~? $g(z)\in\U$~?
\end{questions}
\begin{sol}
\begin{questions}
\question $\abs{f(z)}=1$ si et seulement si $z+\conj{z}=1$. On obtient la droite d'équation $x=1/2$.\\
$\Re(f(z))=0$ si et seulement si $\conj{f(z)}=-f(z)$ ce qui conduit à $z\conj{z}-1/2z-1/2\conj{z}+(1/4-(3/2)^2)=0$ d'où :
\[\abs{z-\frac{1}{2}}^2=\p{\frac{3}{2}}^2.\]
On obtient le cercle de centre $1/2$ et de rayon $3/2$, privé de $2$.
\question  $\Im(g(z))=0$ si et seulement si $\conj{g(z)}=g(z)$ ce qui conduit à $\conj{z}=-z$. On obtient donc $\ii\R\setminus\ens{2\ii}$.

$\abs{g(z)}=1$ si et seulement si $|z|=1$.

\end{questions}
\end{sol}

\magsubsection{Inégalité triangulaire}

\exercice{nom={Inégalité}}
Soit $a$ et $b$ deux nombres complexes.
\begin{questions}
\question Montrer que
\[\abs{a}+\abs{b}\leq\abs{a+b}+\abs{a-b}.\]
\question Déterminer les cas d'égalité dans l'inégalité précédente.
\end{questions}
\begin{sol}
\begin{itemize}
\item On peut raisonner par équivalence, tout élever au carré. On est ramené à
\[\p{\abs{a}-\abs{b}}^2+2\abs{a^2-b^2} \geq 0.\]
(je ne suis plus sur des coefficients), ce qui est trivial. Il y a égalité si et seulement si ces termes sont nuls, c'est-à-dire lorsque $a=b$ ou $a=-b$.
\item On peut aussi écrire $2a=\p{a+b}+\p{a-b}$ puis utiliser l'inégalité triangulaire.
Écrire ensuite $2b=\p{a+b}-\p{a-b}$ avant d'utiliser l'inégalité
triangulaire. Sommer ensuite les deux inégalités. Le cas d'égalité est
$a=b$ ou $a=-b$.
\end{itemize}
\end{sol}



\exercice{nom={Somme}}
Soit $z_0,z_1,\ldots,z_n$ des nombres complexes de module 1. Montrer que
  \[\sum_{k=0}^n \frac{z_k}{2^k}\neq 0.\]
On pourra commencer par majorer
\[\abs{\sum_{k=1}^n \frac{z_k}{2^k}}.\]
\begin{sol}
Il suffit de remarquer que
\[\abs{\sum_{k=0}^n \frac{z_k}{2^k}}\geq
  1-\abs{\sum_{k=1}^n \frac{z_k}{2^k}}\geq 1-\sum_{k=1}^n \abs{\frac{z_k}{2^k}} \geq \frac{1}{2^n}\]
\end{sol}

\exercice{nom={Majoration}}
Soit $z\in\C$. On suppose que
\[1+z+z^2+\dots+z^{n-1}=nz^n.\]
Montrer que $\abs{z}\leq 1$.
\begin{sol}
On divise par $z^n$ puis I.T et on raisonne pas l'absurde pour arriver à une contradiction.
\end{sol}

\magsubsection{Puissance entière, binôme de \nom{Newton}}

\exercice{nom={Majorations}}
\begin{questions}
\question Soit $a$ et $b$ deux nombres complexes et $n$ un entier naturel
  non nul. Montrer que si on pose $M\defeq\max(\abs{a},\abs{b})$, alors
  $$\abs{a^n-b^n}\leq nM^{n-1}\abs{a-b}.$$
\question Soit $z\in\C$ et $n$ un entier non nul. Montrer que
  $$\abs{\sum_{k=0}^n \frac{z^k}{k!} - \p{1+\frac{z}{n}}^n}\leq
    \sum_{k=0}^n \frac{\abs{z}^k}{k!}-\p{1+\frac{\abs{z}}{n}}^n.$$
\end{questions}

\begin{sol}
\begin{questions}
\question Formule $a^n-b^n$.
\question Binôme de Newton + on réunit les deux sommes puis I.T et le terme devant $|z|^k$ est positif puis on fais le chemin inverse.
\end{questions}
\end{sol}

\magsection{Forme trigonométrique}
\magsubsection{Exponentielle $\ii\theta$}

\magsubsection{Application à la trigonométrie}

\exercice{nom={Linéarisation}}
Linéariser l'expression $\cos^2 x \sin^3 x$.
\begin{sol}
\[-\frac{1}{16}\sin(5x)+\frac{1}{8}\sin x+\frac{1}{16}\sin(3x)\]
\end{sol}

\exercice{nom={Calcul de sommes trigonométriques}}
On se donne un réel $x$.
\begin{questions}
\question Calculer les sommes suivantes
  $$\sum_{k=0}^n \cos\p{kx}, \qquad \sum_{k=0}^n \binom{n}{k}\cos\p{kx},
    \qquad \sum_{k=0}^n k\cos\p{kx}.$$
\question 
  \begin{questions}
  \question On suppose que $x\not\in 2\pi\Z$. Montrer que
    $$\sum_{k=0}^{n} \sum_{l=-k}^{k} \e^{\ii lx}=
      \p{\frac{\sin\p{\frac{n+1}{2}x}}{\sin\p{\frac{x}{2}}}}^2.$$
  \question Calculer la valeur de cette somme pour $x\in 2\pi\Z$ de deux
    manières distinctes.
  \end{questions}
\end{questions}
\begin{sol}
\begin{questions}
\question Soit $x\in\R$ et $n\in\N$. Alors
 \[\sum_{k=0}^n \cos(kx)=\begin{cases}
   \frac{\cos\p{\frac{nx}{2}}\sin\p{\frac{n+1}{2}x}}{\sin\p{\frac{x}{2}}} & \text{si $x\not\equiv 0\ [2\pi]$}\\
   n+1 & \text{si $x\equiv 0\ [2\pi]$}
   \end{cases}\]
\[\sum_{k=0}^n \binom{n}{k}\cos\p{k\theta}=2^n \cos^n\p{\frac{\theta}{2}}\cos\p{\frac{n\theta}{2}}\]
Pour le dernier, il faut dériver la somme des $\sin(kx)$.
\[\sum_{k=0}^n k\cos(kx)=\begin{cases}
  \frac{n\sin\p{\frac{x}{2}}\cos\p{\frac{nx}{2}}\sin\p{\frac{n+1}{2}x}+(n+1)\sin\p{\frac{x}{2}}\sin\p{\frac{nx}{2}}\cos\p{\frac{n+1}{2}x}-\cos\p{\frac{x}{2}}\sin\p{\frac{nx}{2}}\sin\p{\frac{n+1}{2}x}}{2\sin^2\p{\frac{x}{2}}} & \text{si $x\not\equiv 0\ [2\pi]$}\\
  \frac{n(n+1)}{2} & \text{si $x\equiv 0\ [2\pi]$}
  \end{cases}\]
\question Si $x\not\equiv 0\ [2\pi]$
  \begin{eqnarray*}
  \sum_{k=0}^n \sum_{l=-k}^k \e^{\ii lx}
  &=& \frac{1}{\sin\p{\frac{x}{2}}}\sum_{k=0}^n \sin\cro{\p{k+\frac{1}{2}}x}\\
  &=& \cro{\frac{\sin\p{\frac{n+1}{2}x}}{\sin\p{\frac{x}{2}}}}^2
  \end{eqnarray*}
  Sinon
  \[\sum_{k=0}^n \sum_{l=-k}^k \e^{\ii lx}=\sum_{k=0}^n (2k+1)=(n+1)^2.\]
  On obtient le même résultat en effectuant une limite.
\end{questions}
\end{sol}

\magsubsection{Forme trigonométrique}


\exercice{nom={Mise sous forme trigonométrique}}
Soit $\theta\in\R$. On pose
$$z=-\sin 2\theta+2\ii\cos^2 \theta.$$
\begin{questions}
\question Déterminer le module et l'argument de $z$.
\question Donner une condition nécessaire et suffisante sur $\theta$ pour que
  $z$ et $z-1$ aient même module.
\end{questions}
\begin{sol}
$\quad$
\begin{questions}
\question Factoriser par $2\cos\theta$. On trouve
  $z=2\cos\theta e^{\ii\p{\theta+\frac{\pi}{2}}}$. Discuter selon le signe de
  $\cos\theta$.
  \question On peut soit calculer le module au carré, soit raisonner
    graphiquement. Il faut $\Re z=1/2$ soit
    \[x\equiv-\frac{\pi}{12}\ [\pi] \ou x\equiv -\frac{5\pi}{12}\ [\pi]\]
\end{questions}
\end{sol}

\exercice{nom={Calculs}}
Mettez les nombres complexes suivants sous forme trigonométrique généralisée
$$\textbf{a.}\ \p{\frac{1+\ii\sqrt{3}}{1-\ii}}^{20}, \qquad \textbf{b.}\ \p{1+\e^{\ii\theta}}^n, \qquad
\textbf{c.}\ \e^{\ii\theta}+\e^{\ii\theta'},$$
$$\textbf{d.}\ \frac{a+b}{a-b} \quad \text{et} \quad \textbf{e.}\ \frac{a+b}{1-ab} \qquad
  \text{où $a\defeq \e^{\ii\theta}$ et $b\defeq \e^{\ii\theta'}$}.$$
\begin{sol}
\[2^{10} \e^{-\ii\frac{\pi}{3}}, \qquad 2^n\cos^n\p{\frac{\theta}{2}}\e^{\ii\frac{n\theta}{2}},
  \qquad 2\cos\p{\frac{\theta-\theta'}{2}}\e^{\ii\frac{\theta+\theta'}{2}}\]
\[-\ii\cotan\p{\frac{\theta-\theta'}{2}}, \qquad
  \ii\frac{\cos\p{\frac{\theta-\theta'}{2}}}{\sin\p{\frac{\theta+\theta'}{2}}}\]
\end{sol}

\magsubsection{Exponentielle complexe}








\magsection{Racines d'un nombre complexe}
\magsubsection{L'équation du second degré}

\exercice{nom={L'équation du second degré}}
Résoudre sur $\C$ les équations suivantes
$$\textbf{a.}\ z^2=-7+24\ii, \qquad \textbf{b.}\ z^2=-3-4\ii, \qquad \textbf{c.}\ z^2+z+1=0,$$
$$\textbf{d.}\ z^2-2(2+\ii)z+6+8\ii=0, \qquad \textbf{e.}\ \ii z^2+(4\ii -3)z+\ii-5=0,$$
$$\textbf{f.}\ z^4+2z^3+z^2+2z+1=0 \quad \text{\it On pourra poser $u\defeq z+\frac{1}{z}$},$$
$$\textbf{g.}\ \begin{cases}
  z_1 z_2=\frac{1}{2} &\\
  z_1+2z_2 = \sqrt{3}, &
  \end{cases}
  \qquad\textbf{h.}\ 
  \begin{cases}
  z_1 + z_2=1 &\\
  z_1^2+z_2^2 = 3. &
  \end{cases}
$$
\begin{sol}
\[\ens{\pm(3+4\ii)}, \qquad \ens{\pm(1-2\ii)}, \qquad \ens{\jj,\jj^2}\]
\[\ens{3-\ii,1+3\ii}, \qquad \ens{-3-2\ii,-1-\ii}\]
\[\ens{\frac{-(1+\sqrt{2})\pm\sqrt{2\sqrt{2}-1}}{2},\frac{-1+\sqrt{2}\pm \ii \sqrt{2\sqrt{2}+1}}{2}}\]
\[\ens{\p{\frac{\sqrt{3}+\ii}{2},\frac{\sqrt{3}-\ii}{4}}, \p{\frac{\sqrt{3}-\ii}{2},\frac{\sqrt{3}+\ii}{4}}} \qquad
  \frac{1\pm\sqrt{5}}{2}\]
\end{sol}

\exercice{nom={Modules et arguments des racines d'un trinôme}}
Soit $u$ un réel tel que $\abs{u}<\pi$. Calculer les modules et arguments de
chacune des racines de l'équation
\[z^2-2z\p{\cos u+\ii\sin u}+2\ii\sin u\p{\cos u+\ii\sin u}=0.\]
\begin{sol}
On trouve
\[\e^{\ii u}-1=2\sin\p{\frac{u}{2}}\e^{\ii\frac{(u+\pi)}{2}} \et
  \e^{\ii u}+1=2\cos\p{\frac{u}{2}}\e^{\ii\frac{u}{2}}.\]
\end{sol}
%% \exercice{nom={Un système algèbrique}}
%% Pour quel(s) couple(s) d'entiers $(n,p)$, les équations suivantes
%% ont-elles une racine commune ?
%% $$z^n=1 \qquad \text{et} \qquad (1+z)^p=1$$
%% {\it On raisonnera d'abord par implication}

\exercice{nom={Trinôme dont les racines ont même module}}
Soit $a,b\in\C$. Montrer que les racines de $z^2+az+b=0$ ont même module si et seulement
si il existe $\lambda\in\interf{0}{4}$ tel que $a^2=\lambda b$.

\begin{sol}
Il faut et il suffit qu'il existe $\lambda\in\interf{0}{4}$ tel que
$a^2=\lambda b$. En effet~:
\begin{itemize}
\item Si les racines sont $z_1=\rho e^{i\theta_1}$ et $z_2=\rho e^{i\theta_2}$,
	alors
	\[z_1z_2=\rho^2 \e^{\ii\p{\theta_1+\theta_2}} \et
		z_1+z_2=\rho \e^{\ii\frac{\theta_1+\theta_2}{2}}
		2\cos\p{\frac{\theta_1-\theta_2}{2}}\]
	donc \[\p{z_1+z_2}^2=4\cos^2\p{\frac{\theta_1-\theta_2}{2}} ab\]
	donc $\p{z_1+z_2}^2=\lambda z_1z_2$ donc $a^2=\lambda b$.
\item Réciproquement, si $a^2=\lambda b$ avec $\lambda\in\interf{0}{4}$ alors le discriminant du trinôme vaut $\Delta=a^2-4b=(\lambda-4)b=\p{ib'\sqrt{4-\lambda}}^2$ avec $b'$ une racine carrée de $b$. Les racines du trinôme sont donc $$z_1=\frac{-a+ib'\sqrt{4-\lambda}}{2} \quad ; z_2=\frac{-a-ib'\sqrt{4-\lambda}}{2}.$$
Ainsi, $$|z_1|^2=\frac{1}{4}\p{|a|^2+(4-\lambda)|b'|^2+2\Re\p{-ai\conj{b'}\sqrt{4-\lambda}}}$$ et $$|z_2|^2=\frac{1}{4}\p{|a|^2+(4-\lambda)|b'|^2-2\Re\p{-ai\conj{b'}\sqrt{4-\lambda}}}$$
Or, $$(-ai\conj{b'}\sqrt{4-\lambda})^2=-a^2\conj{b'}^2(4-\lambda)=-\lambda b\conj{b}(4-\lambda)\in\RM.$$
Donc $-ai\conj{b'}\sqrt{4-\lambda}$ est imaginaire pur, donc sa partie réelle et le résultat en découle.
\end{itemize}
\end{sol}

\magsubsection{Racines $n$-ièmes}





\exercice{nom={Équations}}
Résoudre les équations suivantes sur $\C$
$$\textbf{a.}\ z^5=-1, \qquad \textbf{b.}\ z^6=\frac{-4}{1+\ii\sqrt{3}}, \qquad \textbf{c.}\ z^3=\conj{z},$$
$$\textbf{d.}\ (z+\ii)^n=(z-\ii)^n, \qquad\textbf{e.}\ 
  1+\frac{z+\ii}{z-\ii}+\p{\frac{z+\ii}{z-\ii}}^2+\p{\frac{z+\ii}{z-\ii}}^3=0.$$
\begin{sol}
\[-\omega^k, \qquad \sqrt[6]{2}\e^{\ii\frac{\pi}{9}}\omega^k, \qquad \ens{0,1,-1,\ii,-\ii}\]
\[\cotan\p{\frac{k\pi}{n}}, \qquad (?)\] 
\begin{questions}
\question $\p{-1}^5=-1$ donc
  \begin{eqnarray*}
\forall z\in\C\qsep z^5=-1
&\ssi& z^5=(-1)^5\\
&\ssi& \p{\frac{z}{-1}}^5=1\\
&\ssi& \p{-z}^5=1\\
&\ssi& \exists k\in\intere{0}{4}\qsep -z=\omega^k, \quad\text{où $\omega\defeq\e^{\ii 2\pi/5}$}\\
&\ssi& \exists k\in\intere{0}{4}\qsep z=-\omega^k
\end{eqnarray*}
L'ensemble des solutions est donc
\[\ensim{-\omega^k}{k\in\intere{0}{4}}.\]
\question 0 est solution de l'équation. On recherche donc les solutions sur $\Cs$. Soit $z\in\Cs$. Alors, il existe $r>0$ et $\theta\in\R$ tel que $z=r\e^{\ii\theta}$. Donc
  \begin{eqnarray*}
z^3=\conj{z}
&\ssi& \p{r\e^{\ii\theta}}^3=\conj{r\e^{\ii\theta}}\\
&\ssi& r^3 \e^{3\ii\theta}=r\e^{-\ii\theta}\\
&\ssi& r^2 \e^{4\ii\theta}=1\\
&\ssi& r^2=1 \et 4\theta\equiv 0\ [2\pi]\\
&\ssi& r=1 \et \theta\equiv 0\ \cro{\frac{\pi}{2}}\\
&\ssi& z=1 \ou z=\ii \ou z=-1 \ou z=-\ii
  \end{eqnarray*}
L'ensemble des solutions est donc $\ens{0,1,\ii,-1,-\ii}$.
\question On a
  \[\frac{-4}{1+\ii\sqrt{3}}=\frac{-4}{2\e^{\ii \frac{\pi}{3}}}=-2\e^{-\ii\frac{\pi}{3}}=2\e^{\ii\frac{2\pi}{3}}.\]
  On pose
  \[z_0\defeq\sqrt[6]{2}\e^{\ii\frac{\pi}{6}}.\]
  Alors $z_0^6=-4/(1+\ii\sqrt{3})$ donc
  \begin{eqnarray*}
  \forall z\in\C\qsep z^6=\frac{-4}{1+\ii\sqrt{3}}
  &\ssi& z^6 = z_0^6\\
  &\ssi& \p{\frac{z}{z_0}}^6 = 1\\
  &\ssi& \exists k\in\intere{0}{5}\qsep \frac{z}{z_0}=\omega^k, \quad\text{où $\omega\defeq\e^{\ii\frac{2\pi}{6}}$}\\
&\ssi& \exists k\in\intere{0}{5}\qsep z=z_0 \omega^k
  \end{eqnarray*}
  L'ensemble des solutions est donc
\[\ensim{\sqrt[6]{2}\e^{\ii\frac{\pi}{9}}\omega^k}{k\in\intere{0}{5}}.\]
\question $\ii$ n'est pas solution de l'équation. On cherche donc les racines sur $\C\setminus\ens{\ii}$.
\begin{eqnarray*}
\forall z\in\C\setminus\ens{\ii}\qsep \p{z+\ii}^n=(z-\ii)^{n}
&\ssi& \p{\frac{z+\ii}{z-\ii}}^n=1\\
&\ssi& \exists k\in\intere{0}{n-1}\qsep \frac{z+\ii}{z-\ii}=\omega^k, \quad\text{où $\omega\defeq\e^{\ii\frac{2\pi}{n}}$}\\
&\ssi& \exists k\in\intere{0}{n-1}\qsep \p{1-\omega^k}z=-\ii(1+\omega^k)\\
&    & \text{En effet, s'il existe $k\in\intere{0}{n-1}$ tel que $\p{1-\omega^k}z=-\ii(1+\omega^k)$,}\\
&    & \text{alors $k\neq 0$ car sinon on aurait $0\times z=-2\ii$.}\\
&\ssi& \exists k\in\intere{1}{n-1}\qsep \p{1-\omega^k}z=-\ii(1+\omega^k)\\
&\ssi& \exists k\in\intere{1}{n-1}\qsep z=-\ii\cdot\frac{1+\omega^k}{1-\omega^k}\\
&\ssi& \exists k\in\intere{1}{n-1}\qsep z=-\ii\cdot\frac{1+\e^{\ii\frac{k2\pi}{n}}}{1-\e^{\ii\frac{k2\pi}{n}}}\\
&\ssi& \exists k\in\intere{1}{n-1}\qsep z=-\ii\cdot\frac{\e^{\ii\frac{k\pi}{n}}\p{\e^{-\ii\frac{k\pi}{n}}+\e^{\ii\frac{k\pi}{n}}}}{\e^{\ii\frac{k\pi}{n}}\p{\e^{-\ii\frac{k\pi}{n}}-\e^{\ii\frac{k\pi}{n}}}}\\
&\ssi& \exists k\in\intere{1}{n-1}\qsep z=-\ii\cdot\frac{2\cos\p{\frac{k\pi}{n}}}{-2\ii\sin\p{\frac{k\pi}{n}}}\\
&\ssi& \exists k\in\intere{1}{n-1}\qsep z=\cotan\p{\frac{k\pi}{n}}
\end{eqnarray*}
L'ensemble des solutions est donc
\[\ensim{\cotan\p{\frac{k\pi}{n}}}{\intere{1}{n-1}}.\]
\question Soit $z\in\C\setminus\ens{\ii}$. Alors
  \begin{eqnarray*}
  \frac{z+\ii}{z-\ii}=1
  &\ssi& z+\ii=z-\ii\\
  &\ssi& 2\ii=0, \quad\text{ce qui est faux.}
  \end{eqnarray*}
  Donc
  \begin{eqnarray*}
  1+\frac{z+\ii}{z-\ii}+\p{\frac{z+\ii}{z-\ii}}^2+\p{\frac{z+\ii}{z-\ii}}^3=0
  &\ssi& \frac{1-\p{\frac{z+\ii}{z-\ii}}^4}{1-\frac{z+\ii}{z-\ii}}=0, \quad\text{car $\dfrac{z+\ii}{z-\ii}\neq 1$}\\
  &\ssi& \p{\frac{z+\ii}{z-\ii}}^4=1\\
  &\ssi& \exists k\in\intere{0}{3}\qsep \frac{z+\ii}{z-\ii}=\omega^k, \quad\text{avec $\omega=\e^{\ii\frac{2\pi}{4}}=\ii$}\\
  &\ssi& \exists k\in\intere{0}{3}\qsep \frac{z+\ii}{z-\ii}=\ii^k\\
  &\ssi& \exists k\in\intere{0}{3}\qsep (1-\ii^k)z=-\ii(1+\ii^k)\\
  &\ssi& \exists k\in\intere{1}{3}\qsep (1-\ii^k)z=-\ii(1+\ii^k)\\
  &\ssi& \exists k\in\intere{1}{3}\qsep z=-\ii\cdot\frac{1+\ii^k}{1-\ii^k}\\
  &\ssi& \text{$z=0$ ou $z=1$ ou $z=-1$}.
  \end{eqnarray*}
  L'ensemble des solutions est donc $\ens{0,1,-1}$.
\end{questions} 
\end{sol}

\exercice{nom={Équation}}
Soit $n\in\Ns$. Résoudre l’équation \[\p{z^2+1}^n=(z-\ii)^{2n}.\]
\begin{sol}
Soit $n\in\Ns$. Alors
\begin{eqnarray*}
\forall z\in\C\qsep \p{z^2+1}^n=(z-\ii)^{2n}
&\ssi& \p{z^2-\ii^2}^n=(z-\ii)^{2n}\\
&\ssi& \p{z-\ii}^n\p{z+\ii}^n=(z-\ii)^{2n}\\
&\ssi& \p{z-\ii}^n\cro{\p{z+\ii}^n-(z-\ii)^{n}}=0\\
&\ssi& z=\ii \ou \p{z+\ii}^n=(z-\ii)^{n}
\end{eqnarray*}
Puisque $\ii$ est solution de l'équation d'origine, on cherche les solutions de $\p{z+\ii}^n=(z-\ii)^{n}$ sur $\C\setminus\ens{\ii}$. On a donc
\begin{eqnarray*}
\forall z\in\C\setminus\ens{\ii}\qsep \p{z+\ii}^n=(z-\ii)^{n}
&\ssi& \p{\frac{z+\ii}{z-\ii}}^n=1\\
&\ssi& \exists k\in\intere{0}{n-1}\qsep \frac{z+\ii}{z-\ii}=\omega^k, \quad\text{où $\omega=\e^{\ii\frac{2\pi}{n}}$}\\
&\ssi& \exists k\in\intere{0}{n-1}\qsep \p{1-\omega^k}z=-\ii(1+\omega^k)\\
&\ssi& \exists k\in\intere{1}{n-1}\qsep \p{1-\omega^k}z=-\ii(1+\omega^k)\\
&    & \text{En effet, si il existe $k\in\intere{0}{n-1}$ tel que $\p{1-\omega^k}z=-\ii(1+\omega^k)$,}\\
&    & \text{alors $k\neq 0$ car sinon on aurait $0\times z=-2\ii$.}\\
&\ssi& \exists k\in\intere{1}{n-1}\qsep z=-\ii\cdot\frac{1+\omega^k}{1-\omega^k}\\
&\ssi& \exists k\in\intere{1}{n-1}\qsep z=-\ii\cdot\frac{1+\e^{\ii\frac{k2\pi}{n}}}{1-\e^{\ii\frac{k2\pi}{n}}}\\
&\ssi& \exists k\in\intere{1}{n-1}\qsep z=-\ii\cdot\frac{\e^{\ii\frac{k\pi}{n}}\p{\e^{-\ii\frac{k\pi}{n}}+\e^{\ii\frac{k\pi}{n}}}}{\e^{\ii\frac{k\pi}{n}}\p{\e^{-\ii\frac{k\pi}{n}}-\e^{\ii\frac{k\pi}{n}}}}\\
&\ssi& \exists k\in\intere{1}{n-1}\qsep z=-\ii\cdot\frac{2\cos\p{\frac{k\pi}{n}}}{-2\ii\sin\p{\frac{k\pi}{n}}}\\
&\ssi& \exists k\in\intere{1}{n-1}\qsep z=\cotan\p{\frac{k\pi}{n}}
\end{eqnarray*}
Les solutions de l'équation $\p{z^2+1}^n=(z-\ii)^{2n}$ sont donc $\ii$ et les
\[\cotan\p{\frac{k\pi}{n}}\]
pour $k\in\intere{1}{n-1}$.
\end{sol}

\exercice{nom={Relation trigonométrique}}
En considérant les racines 11-ièmes de $1$, montrer que
  $$\cos\p{\frac{\pi}{11}}+\cos\p{\frac{3\pi}{11}}+\cos\p{\frac{5\pi}{11}}+
    \cos\p{\frac{7\pi}{11}}+\cos\p{\frac{9\pi}{11}}=\frac{1}{2}.$$

\exercice{nom={Calcul avec $\jj$}}
  \begin{questions}
  \question Calculer
    $$\p{a+b\jj+c\jj^2}\p{a+b\jj^2+c\jj}.$$
  \question Sans effectuer de développement, retrouver le fait que cette
    expression ne change pas lorsqu'on échange deux variables.
  \end{questions}

\exercice{nom={Autour des racines de l'unité}}
Soit $\p{\omega_k}_{0\leq k\leq n-1}$ les racines n-ièmes de l'unité.
  Calculer pour tout entier $p\in\Z$
  $$\sum_{k=0}^{n-1} \omega_k^p \qquad \text{et} \qquad
    \prod_{k=0}^{n-1} \omega_k.$$

\begin{sol}
\begin{questions}
\question On pose $\omega=\e^{\ii 2\pi/11}$. Alors
  \[\sum_{k=0}^{10} \omega^k=0.\]
  Donc
  \[1+\p{\omega+\omega^{10}}+\p{\omega^2+\omega^9}+\p{\omega^3+\omega^8}+\p{\omega^4+\omega^7}+\p{\omega^5+\omega^6}=0.\]
  En utilisant le fait que $\omega^{11-k}=\conj{\omega^k}$ et le fait que $\omega^k+\conj{\omega^k}=2\Re\p{\omega^k}$, on en déduit que
  \[1+2\cos\p{\frac{2\pi}{11}}+2\cos\p{\frac{4\pi}{11}}+2\cos\p{\frac{6\pi}{11}}+2\cos\p{\frac{8\pi}{11}}+2\cos\p{\frac{10\pi}{11}}=0.\]
  Enfin, en utilisant le fait que $\cos(\pi-x)=-\cos(x)$, on obtient
  \[\cos\p{\frac{\pi}{11}}+\cos\p{\frac{3\pi}{11}}+\cos\p{\frac{5\pi}{11}}+
    \cos\p{\frac{7\pi}{11}}+\cos\p{\frac{9\pi}{11}}=\frac{1}{2}.\]
\question
  \begin{questions}
  \question En développant, on obtient, en utilisant le fait que $\jj^3=1$ puis que $1+\jj+\jj^2=0$
    \begin{eqnarray*}
    \p{a+b\jj+c\jj^2}\p{a+b\jj^2+c\jj}
    &=& a^2+b^2+c^2+\p{\jj+\jj^2}\p{ab+ac+bc}\\
    &=& a^2+b^2+c^2-(ab+ac+bc).
    \end{eqnarray*}
    Cette expression est bien entendu symétrique en $a$, $b$, $c$.
  \question Si on échange $a$ et $b$, on obtient, en utilisant le fait que $\jj^3=1$
    \begin{eqnarray*}
    \p{b+a\jj+c\jj^2}\p{b+a\jj^2+c\jj}
    &=& \jj\p{a+b\jj^2+c\jj}\jj^2\p{a+b\jj+c\jj^2}\\
    &=& \p{a+b\jj+c\jj^2}\p{a+b\jj^2+c\jj}.
    \end{eqnarray*}
    L'expression est donc inchangée si on permute $a$ et $b$. On vérifie que si l'on permute $a$ et $c$ ou $b$ et $c$ l'expression est aussi inchangée.
  \end{questions}
\question Soit $n\in\Ns$. On pose $\omega\defeq\e^{\ii 2\pi/n}$. Soit $p\in\Z$. Alors
\[\sum_{k=0}^{n-1} \omega_k^p= \sum_{k=0}^{n-1} \p{\omega^k}^p = \sum_{k=0}^{n-1} \omega^{kp} = \sum_{k=0}^{n-1} \p{\omega^p}^k\]
Or
\[\omega^p=1 \quad\ssi\quad \p{\e^{\ii\frac{2\pi}{n}}}^p=1 \quad\ssi\quad \e^{\ii\frac{2\pi p}{n}}=1 \quad\ssi\quad \frac{2\pi p}{n}\equiv 0\ [2\pi]\quad\ssi\quad p\equiv 0\ [n].\]
Si $p$ est un multiple de $n$, alors
\[\sum_{k=0}^{n-1} \omega_k^p= \sum_{k=0}^{n-1} 1 = n.\]
Sinon
\[\sum_{k=0}^{n-1} \omega_k^p=\sum_{k=0}^{n-1} \p{\omega^p}^k=\frac{1-\p{\omega^p}^n}{1-\omega^p}=\frac{1-\p{\omega^n}^p}{1-\omega^p}=0.\]
Enfin
\[
\prod_{k=0}^{n-1} \omega_k
= \prod_{k=0}^{n-1} \omega^k = \omega^{\sum_{k=0}^{n-1} k} = \omega^{\frac{n(n-1)}{2}}=
\e^{\ii\frac{2\pi}{n}\cdot\frac{n(n-1)}{2}}=\e^{\ii\pi(n-1)}=(-1)^{n-1}=\begin{cases}
1 & \text{si $n\equiv 1\ [2]$}\\
-1 & \text{si $n\equiv 0\ [2]$.}
\end{cases}
\]
\end{questions}
\end{sol}


\magsection{Nombres complexes et géométrie plane}

\magsubsection{Le plan complexe}

\exercice{nom={Caractérisation d'un triangle équilatéral}}
Soit $A$, $B$ et $C$ trois points du plan d'affixes respectives $a$, $b$ et $c$.
\begin{questions}
\question On suppose que $B\neq A$. Montrer que $ABC$ est équilatéral si et seulement si
  \[\frac{c-a}{b-a}=\e^{\frac{\ii\pi}{3}} \quad\text{ou}\quad \frac{c-a}{b-a}=\e^{-\frac{\ii\pi}{3}}.\]
\question En déduire que $ABC$ est équilatéral si et seulement si
  \[a^2+b^2+c^2=ab+ac+cb.\]
\end{questions}


\exercice{nom={Triangles équilatéraux}}
Soit $(ABC)$ et $(ADE)$ deux triangles équilatéraux directs et $(ACFD)$ un parallélogramme. Montrer que $(BFE)$ est équilatéral direct.
\begin{sol}
On peut écrire toutes les hypothèses (3 égalités par triangle du type $\frac{c-a}{b-a}=\e^{\frac{\ii\pi}{3}}$ et $d-a=f-c$ pour le parallélogramme).
Or, $f-e=d-e+c-a=(a-e)e^{i\pi/3}+(b-a)e^{i\pi/3}=(b-e)e^{i\pi/3}$ ce qui signifie que $(EBF)$ est équilatéral direct.
\end{sol}

\exercice{nom={Algèbre et géométrie}}
A tout nombre complexe $z\neq 4$, on associe le nombre \[z'=\frac{\ii z-4}{z-4}\] et on note
$\mathcal{C}$ l'ensemble des points $M$ d'affixe $z$ du plan tels que $z'$ est réel. Déterminer $\mathcal{C}$ par une méthode algébrique puis par une méthode géométrique.





\exercice{nom={Lieu}}
Soit le point $A$ d'affixe $2$ et le point $B$ d'affixe $-2$. À tout point $M$ d'affixe $z$, autre que $A$, on associe le point $M'$ d'affixe \[z'=\dfrac{2z-4}{\bar{z}-2}.\]
\begin{questions}
\question Déterminer $\abs{z}$. Que peut-on en déduire pour $M'$~?
\question Déterminer l'ensemble $\mathcal{E}$ des points d'affixe $z$ tels que $M'=B$.
\question Pour tout point $M$ de $\mathcal{E}$ et distincts de $A$ et $B$, que peut-on dire de \[\frac{z-2}{z'-2}~?\] Interpréter géométriquement ce résultat et en déduire une construction de $M'$.
\end{questions}

\exercice{nom={Triangles}}
\begin{questions}
\question On considère les points $A$, $B$ et $C$ d'affixes respectives $1$, $z$ et $\ii z$. Déterminer l'ensemble des points $B$ pour lesquels $A$, $B$ et $C$ sont alignés.
\question On considère les points $E$, $F$ et $G$ non alignés, d'affixes respectives, $z_{1}$, $z_{2}$ et $z_{3}$. Déterminer une condition nécessaire et suffisante pour que le triangle $(EFG)$ soit rectangle isocèle en $E$.
\end{questions}
\begin{sol}
\begin{questions}
\question $\dfrac {c-a}{b-a}\in \R$ ssi $z\conj{z}-\dfrac{1+i}{2}z-\dfrac{1-i}{2}\conj{z}=0$. Il s'agit donc du cercle de centre $(1-i)/2$ et de rayon $1/\sqrt 2$.
\question $\dfrac{g-e}{f-e}=i$.
\end{questions}
\end{sol}

\magsubsection{Les similitudes directes}

\exercice{nom={Similitudes}}
\begin{questions}
\question Caractériser géométriquement la similitude
  \[z\mapsto 2(1+\ii)z-7-4\ii.\]
\question Déterminer l'expression complexe de la rotation de centre $1+\ii$ et d'angle $\pi/4$.
\question On note $r$ la rotation de centre $2+\ii$ et d'angle $\pi/2$ et $s$ la symétrie centrale de centre $1-\ii$. Caractériser géométriquement $s\circ r$.
\question On note $r$ la rotation de centre $\ii$ et d'angle $\pi/3$ et $r'$ la rotation de centre $2\ii$ et d'angle $-\pi/3$. Caractériser géométriquement $r'\circ r$.
\end{questions}

\begin{sol}
\begin{questions}
\question Soit $f$ la fonction définie sur $\C$ par
\[\forall z\in\C\qsep f(z)\defeq 2(1+\ii)z-7-4\ii.\]
Alors
\begin{eqnarray*}
\forall z\in\C\qsep f(z)=z
&\ssi& 2(1+\ii)z-7-4\ii=z\\
&\ssi& (1+2\ii)z=7+4\ii\\
&\ssi& z=\frac{7+4\ii}{1+2\ii}\\
&\ssi& z=\frac{(7+4\ii)(1-2\ii)}{5}\\
&\ssi& z=3-2\ii
\end{eqnarray*}
donc $f$ est une similitude directe de centre $3-2\ii$. Or
\[2(1+\ii)=2\sqrt{2}\e^{\ii\frac{\pi}{4}}\]
donc le rapport de cette similitude est $2\sqrt{2}$ et son angle est $\pi/4$.
\question Si $g$ est l'expression complexe de la rotation de centre $1+\ii$ et d'angle $\pi/4$, alors
\[\forall z\in\C\qsep g(z)=\e^{\ii\frac{\pi}{4}}\cro{z-(1+\ii)}+(1+\ii).\]
\question On a 
\[\forall z\in\C\qsep r(z)=\ii\cro{z-(2+\ii)}+(2+\ii)\]
et
\[\forall z\in\C\qsep s(z)=-\cro{z-(1-\ii)}+(1-\ii).\]
On en déduit que
\[\forall z\in\C\qsep (s\circ r)(z)=s(r(z))=-\ii z-1-\ii.\]
Cette fonction admet un point fixe qui est $-1$. c'est donc une rotation de centre $-1$ et d'angle $-\pi/2$.
\question On a
\[\forall z\in\C\qsep r(z)=\e^{\ii\frac{\pi}{3}}\p{z-\ii}+\ii\]
et
\[\forall z\in\C\qsep r'(z)=\e^{-\ii\frac{\pi}{3}}\p{z-2\ii}+2\ii.\]
Donc
\[\forall z\in\C\qsep (r'\circ r)(z)=z-\frac{\sqrt {3}+\ii}{2}.\]
C'est donc une translation de vecteur d'affixe $-(\sqrt{3}+\ii)/2$.
\end{questions}
\end{sol}

\begin{sol}
\begin{questions}
\question $ABC$ est équilatéral si et seulement si il est isocèle de sommet $A$ et que l'angle en $A$ soit égal à $\pm\pi/3$ (direct ou indirect), d'où le résultat.
\question On pose
  \[P\defeq\p{X-\e^{\ii\frac{\pi}{3}}}\p{X-\e^{-\ii\frac{\pi}{3}}}=X^2-\p{\e^{\ii\frac{\pi}{3}}+\e^{-\ii\frac{\pi}{3}}}X+1=X^2-X+1.\]
Donc
\begin{eqnarray*}
\text{$ABC$ est équilatéral} &\ssi & \frac{c-a}{b-a}=\e^{\frac{\ii\pi}{3}} \quad\text{ou}\quad \frac{c-a}{b-a}=\e^{-\frac{\ii\pi}{3}}\\
&\ssi & P\p{\frac{c-a}{b-a}}=0\\
&\ssi &a^2+b^2+c^2=ab+ac+cb.
\end{eqnarray*}
\end{questions}
\end{sol}



%END_BOOK

\end{document}


\chapter{Logique, ensembles}
\setcounter{numeroexercicecours}{1}
\documentclass{magnoliaold}

\magtex{tex_driver={pdftex},
        tex_packages={slashbox,epigraph,xypic}}
\magfiche{document_nom={Cours sur les ensembles},
          auteur_nom={François Fayard},
          auteur_mail={fayard.prof@gmail.com}}
\magcours{cours_matiere={maths},
          cours_niveau={mpsi},
          cours_chapitre_numero={2},
          cours_chapitre={Logique, Ensembles}}
\magmisenpage{}
\maglieudiff{}
\magprocess

\begin{document}

%BEGIN_BOOK
\setlength\epigraphwidth{.7\textwidth}
\epigraph{\og Si la logique est l'hygiène du mathématicien, ce n'est pas elle qui lui fournit sa
  nourriture; le pain quotidien dont il vit, ce sont les grands problèmes.\fg}{--- \textsc{André Weil (1906-1998)}}
\setlength\epigraphwidth{.6\textwidth}
\epigraph{\og  Sur l’enseigne du barbier du village, on peut lire~:~Je rase tous les hommes du
  village qui ne se rasent pas eux-mêmes, et seulement ceux-là. Savez-vous qui rase
  le barbier ? \fg}{--- \textsc{Bertrand Russel (1872--1970)}}
\setlength\epigraphwidth{.5\textwidth}
\epigraph{\og J'aimais et j'aime encore les mathématiques pour elles-mêmes comme n'admettant
  pas l'hypocrisie et le vague, mes deux bêtes d'aversion.\fg}{--- \textsc{Stendhal (1783-1842)}}
\hidemesometimes{
    \bigskip
    \hfill\includegraphics[width=0.3\textwidth]{../../Commun/Images/maths-cours-shadok.png}}

\magtoc

\section{Éléments de logique}

\subsection{Assertion, prédicat}

\begin{definition}[utile=-3]
\begin{itemize}
\item On appelle \emph{assertion} toute phrase mathématique à laquelle
  on peut attribuer une et une seule valeur de vérité~: vrai ou faux.
\item Soit $E$ un ensemble. On appelle \emph{prédicat} sur $E$ toute phrase mathématique
  dont la valeur de vérité dépend d'un élément $x\in E$.
\end{itemize}
\end{definition}

\begin{exemples}
\exemple \og 7 est un nombre premier \fg est une assertion vraie.
  L'assertion \og 7 est divisible par 3 \fg est fausse.
\exemple $P(x)\defeq$ \og $x$ est rationnel \fg est un prédicat sur $\R$. $P\p{3/4}$
  est vrai alors que $P(\sqrt{2})$ est faux.
\exemple $P\p{a,b,c}\defeq$ \og $a^2+b^2=c^2$ \fg est un prédicat sur $\N^3$.
\exemple \og L'ensemble des nombres premiers est infini \fg est une assertion
  vraie. L'assertion \og Il existe une infinité de nombres premiers $p$ tels
  que $p+2$ est premier \fg est une assertion dont on pense qu'elle est vraie. Mais aujourd'hui, personne n'en a fait la preuve.
\end{exemples}

\begin{remarques}
\remarque Deux principes fondamentaux gouvernent les valeurs de vérité des assertions.
  \begin{itemize}
  \item Le \emph{principe de non-contradiction}~: Une assertion ne peut être à la fois vraie et fausse.
  \item Le \emph{principe du tiers exclu}~: Une assertion qui n'est pas vraie est fausse.
  \end{itemize}
\remarque Si $P$ est un prédicat, on dit que $P$ est vrai lorsque, quel que
  soit $x\in E$, $P(x)$ est vraie. Dire que $P$ n'est pas vrai
  signifie qu'il existe $x\in E$ tel que $P(x)$ est faux.
\end{remarques}

\begin{definition}[utile=-3]
$\quad$
\begin{itemize}
\item Le \emph{quantificateur universel} $\forall$ signifie \og pour tout \fg
\item Le \emph{quantificateur existentiel} $\exists$ signifie \og il existe (au moins) un \fg
\end{itemize}
\end{definition}

\begin{remarqueUnique}
\remarque On trouve parfois le quantificateur $\exists !$ qui signifie
  \og il existe un unique \fg. 
\end{remarqueUnique}

\begin{exos}
\exo Les assertions suivantes sont-elles vraies~?
  \begin{questions}
  \question $\forall y\in\R \qsep \exists x\in\R \qsep x+y\geq 0$.
  \question $\exists x\in\R \qsep \forall y\in\R \qsep x+y\geq 0$.
  \question $\exists x\in\R \qsep \forall y\in\R \qsep y^2\geq x$.
  \end{questions}
\exo Déterminer les $x\in\R$ tels que
  \[\forall n\in\N \qsep x^{n+2}\leq x^{n+1}+x^n.\]
\end{exos}

\begin{sol}
$x^n(x^2-x-1)\leq 0$. Si $x\geq 0$, ok.
Si $x<0$, on doit avoir $x^2-x-1=0$, d'où $S= \interf{0}{\frac{1+\sqrt{5}}{2}}\cup\ens{\frac{1-\sqrt{5}}{2}}$.
\end{sol}

\begin{definition}[utile=-3]
Soit $P$ et $Q$ deux assertions.
\begin{itemize}
\item On définit l'assertion $\p{\non\  P}$ comme étant vraie
  lorsque $P$ est fausse et fausse lorsque $P$ est vraie.
\item On définit l'assertion $\cro{P \et Q}$ comme étant vraie
  lorsque $P$ et $Q$ sont vraies et fausse sinon.
\item On définit l'assertion $\cro{P \ou Q}$ comme étant vraie
  lorsqu'au moins l'une des deux assertions est vraie, et fausse sinon.
\end{itemize}
\end{definition}

\begin{remarques}
\remarque Les valeurs de vérité de ces assertions sont données
  par les tables suivantes.
  \begin{center}
  \begin{minipage}{0.3\linewidth}
  \begin{center}
  \begin{tabular}{|c|c|c|}
  \hline
  P & V & F \\
  \hline
  \non\ P &  F & V \\
  \hline
  \end{tabular}
  \[\non\ P\]
  \end{center}
  \end{minipage}
  \begin{minipage}{0.3\linewidth}
  \begin{center}
  \begin{tabular}{|c|c|c|}
  \hline
  \backslashbox{P}{Q} & V & F \\
  \hline
  V & V & F\\
  \hline
  F & F & F\\
  \hline
  \end{tabular}
  \[P \et Q\]
  \end{center}
  \end{minipage}
  \begin{minipage}{0.3\linewidth}
  \begin{center}
  \begin{tabular}{|c|c|c|}
  \hline
  \backslashbox{P}{Q} & V & F \\
  \hline
  V & V & V\\
  \hline
  F & V & F\\
  \hline
  \end{tabular}
  \[P \ou Q\]
  \end{center}
  \end{minipage}

  \end{center}
\remarque Lorsque le menu d'un restaurant vous propose \og fromage ou dessert
  \fg, le \og ou \fg est employé au sens strict (on dit aussi exclusif); il
  n'est pas possible d'avoir les deux. En mathématiques, le \og ou \fg est
  employé au sens large (on dit aussi inclusif). Lorsqu'on dit qu'un
  entier naturel $n$ est divisible par 2 ou par 3, il peut très bien être divisible par 2 et par 3.
\end{remarques}

\subsection{Implication, équivalence}

\begin{definition}[utile=-3]
Soit $P$ et $Q$ deux assertions. On définit l'assertion $P\implique Q$ comme
étant fausse lorsque $P$ est vraie et $Q$ est fausse, et vraie sinon.
\end{definition}

\begin{remarques}
\remarque Montrer $P\implique Q$ revient à prouver que si $P$ est
  vraie, alors $Q$ est vraie.
\remarque Si $P$ et $Q$ sont deux prédicats sur $E$, $P\implique Q$
  signifie que $Q(x)$ est vraie dès que $P(x)$ est vraie.
  Si c'est le cas, on écrit
  \[\forall x\in E \qsep P(x) \implique Q(x)\]
  et on dit que $P$ est une condition suffisante pour $Q$ ou que $Q$ est une
  condition nécessaire pour $P$.
\end{remarques}

\begin{exos}
\exo Dans les exemples suivants, dites si le prédicat $P$ est une
  condition nécessaire ou une condition suffisante pour $Q$.
  \begin{itemize}
  \item $E=\R$, $P(x)\defeq$ \og $x\in\Q$ \fg et $Q(x)\defeq$ \og $x^2\in\Q$ \fg. 
  \item $E$ est l'ensemble des triangles du plan, $P(T)\defeq$ \og
    $T$ est isocèle \fg et $Q(T)\defeq$ \og $T$ est équilatéral \fg.
  \item $E=\R^2$, $P\p{x,y}\defeq$ \og $x\equiv y\ [2\pi]$ \fg et $Q\p{x,y}\defeq$
    \og $x\equiv y\ [\pi]$ \fg.
  \end{itemize}
\exo Montrer que
  \[\forall x,y\in\R \qsep\cro{xy>0 \et x+y>0}\quad\implique\quad
    \cro{x>0 \et y>0}.\]
\exo Montrer que
  \[\forall x\in\R \qsep
    \cro{\forall \epsilon\in\RPs \qsep\abs{x}\leq\epsilon} \quad\implique\quad
    x=0.\]
\end{exos}

\begin{proposition}[utile=-3, nom={Modus Ponens}]
Soit $P$ et $Q$ deux assertions. Si $P$ et $P\implique Q$ sont vraies, alors
$Q$ est vraie.   
\end{proposition}

\begin{remarqueUnique}
\remarque En pratique, on utilise cette proposition lorsque $P$ et $Q$ sont des
  prédicats. Si $P\implique Q$ est vrai et $x$ est un élément de $E$ tel que
  $P(x)$ est vrai, alors $Q(x)$ est vrai. Dans ce cadre, on dit que
  $P\implique Q$ est un théorème. Vérifier les hypothèses
  du théorème revient à vérifier que $P(x)$ est vrai et appliquer le
  théorème nous permet de conclure que $Q(x)$ est vrai.
  Traduisons mathématiquement le raisonnement suivant~: \og Socrate
  est un homme. Puisque tous les hommes sont mortels, alors Socrate est
  mortel \fg. Si $P(x)\defeq$ \og $x$ est un homme \fg et $Q(x)\defeq$ \og $x$ est
  mortel \fg, alors l'énoncé \og Tous les hommes sont mortels \fg s'écrit
  \[\forall x\in U \qsep P(x) \implique Q(x).\]
  Puisque Socrate est un homme ($P\p{\text{Socrate}}$ est vrai), on en déduit
  que Socrate est mortel ($Q\p{\text{Socrate}}$ est vrai).
\end{remarqueUnique}

\begin{exoUnique}
\exo Soit $a,b\in\R$ tels que
  \[\forall x\in\R \qsep x<a \quad\implique\quad x\leq b.\]
  Montrer que $a\leq b$.
\end{exoUnique}

\begin{sol}
Absurde + on cale un $c$ strictement entre $a$ et $b$.
\end{sol}

\begin{definition}[utile=-3]
Soit $P$ et $Q$ deux assertions. On définit l'assertion $P\ssi Q$ comme
étant vraie lorsque $P$ et $Q$ ont même valeur de vérité, et fausse sinon.
\end{definition}

\begin{remarques}
\remarque Les valeurs de vérité des assertions $P \implique Q$ et
  $P \ssi Q$ sont regroupées dans les tableaux suivants.
  \begin{center}
  \begin{minipage}{0.3\linewidth}
  \begin{center}
  \begin{tabular}{|c|c|c|}
  \hline
  \backslashbox{P}{Q} & V & F \\
  \hline
  V &  V & F \\
  \hline
  F &  V & V \\
  \hline
  \end{tabular}
  \[P \implique Q\]
  \end{center}
  \end{minipage}
  \begin{minipage}{0.3\linewidth}
  \begin{center}
  \begin{tabular}{|c|c|c|}
  \hline
  \backslashbox{P}{Q} & V & F \\
  \hline
  V & V & F \\
  \hline
  F &  F & V\\
  \hline
  \end{tabular}
  \[P \ssi Q\]
  \end{center}
  \end{minipage}
  \end{center}
\remarque Les assertions $P\ssi Q$ et $Q\ssi P$ ont même valeur de vérité; on
  dit que la relation d'équivalence est symétrique.
\remarque Si $P$ et $Q$ sont deux prédicats sur $E$, dire que $P\ssi Q$
  est vrai signifie que $Q(x)$ et $P(x)$ ont même valeur de vérité
  quel que soit $x\in E$. Si c'est le cas, on écrit
  \[\forall x\in E \qsep P(x) \ssi Q(x)\]
  et on dit que $P$ est une condition nécessaire et suffisante pour $Q$.
\end{remarques}

\begin{proposition}[utile=-3]
Soit $P$ et $Q$ deux assertions. Alors $P\ssi Q$ et $[\p{P\implique Q}$ et
$\p{Q \implique P}]$ ont même valeur de vérité.
\end{proposition}

\begin{remarqueUnique}
\remarque Pour démontrer que $P\ssi Q$, on pourra démontrer que
  $P\implique Q$, puis que $Q\implique P$; on dit qu'on raisonne par
  double implication.
\end{remarqueUnique}

\begin{exoUnique}
\exo Soit $\lambda\in\R$ et $f$ la fonction définie sur $\R$ par
  \[\forall x\in\R \qsep f(x)\defeq \sin(\lambda x).\]
  Donner une condition nécessaire et suffisante sur $\lambda$ pour que
  $f$ soit $2\pi$-périodique.
% \exo Soit $a,b,c\in\R$ avec $a\neq 0$. On note $P\defeq aX^2+bX+c$. Donner une
%   condition nécessaire et suffisante sur $a,b,c$ pour qu'il existe $x,y\in\R$
%   tels que $P(x)=y$ et $P(y)=x$.
\end{exoUnique}

\begin{sol}
Ok pour $\lambda=0$
Si $f$ est $2\pi$-per, prendre $x=\pi/2\lambda$ et on obtient $\lambda\in\Z$. Réciproque ok.
\end{sol}

\begin{proposition}[utile=-3]
Soit $P$, $Q$, $R$ trois assertions. Alors
\begin{eqnarray*}
\cro{P \et \p{Q \ou R}} &\ssi& \cro{\p{P \et Q} \ou \p{P \et R}},\\
\cro{P \ou \p{Q \et R}} &\ssi& \cro{\p{P \ou Q} \et \p{P \ou R}}.
\end{eqnarray*}
\end{proposition}

\begin{preuve}

~~
\begin{center}
  \begin{tabular}{|c|c|c|c|c|c|c|c|}
  \hline
  P & Q & R & Q$\lor$ R & P$\wedge$ (Q$\lor$ R) & P$\wedge$ Q & P $\wedge$ R & (P$\wedge$ Q)$\lor$(P$\wedge$ R) \\
  \hline
  V &V & V &V &V &V &V &V \\
  \hline
  V & V & F & V &V &V &F & V \\
  \hline
  V& F & F & F &F &F &F & F \\
  \hline
  V & F & V & V &V &F &V & V \\
  \hline
  F & V & V & V &F &F &F & F \\
  \hline
  F & V & F & V &F &F &F & F \\
  \hline
  F &F &F & F &F &F &F & F \\
  \hline
  F &F &V & V &F &F &F & F \\
  \hline
  \end{tabular}
  Preuve de la première équivalence
  
\end{center}
\end{preuve}

\begin{proposition}[utile=-3, nom={Lois de \nom{Morgan}}]
Soit $P$ et $Q$ deux assertions. Alors
\begin{eqnarray*}
\non\p{P \et Q} &\ssi& \cro{\p{\non\  P} \ou \p{\non\  Q}},\\
\non\p{P \ou Q} &\ssi& \cro{\p{\non\  P} \et \p{\non\  Q}},\\
\non\p{\non\  P} &\ssi& P.
\end{eqnarray*}
\end{proposition}

\begin{proposition}[utile=-3, nom={Raisonnement par contraposée}]
Soit $P$ et $Q$ deux assertions. Alors
\[\cro{P \implique Q} \quad\ssi\quad \cro{\non\  Q \implique \non\  P}.\]
\end{proposition}

\begin{remarqueUnique}
\remarque Lorsque l'on démontre $\cro{\p{\non\  Q} \implique \p{\non\  P}}$ pour montrer
  que $\cro{P\implique Q}$, on dit que l'on raisonne par contraposée.
\end{remarqueUnique}

\begin{exoUnique}
\exo Supposons que l'on ait montré que $\pi^2$ est irrationnel. Peut-on en
  déduire que $\pi$ est irrationnel~?
\end{exoUnique}

\begin{proposition}[utile=-3]
Soit $P$ et $Q$ deux assertions. Alors
\[\cro{\non\p{P \implique Q}} \quad\ssi\quad \cro{P \et \p{\non\ Q}}.\]  
\end{proposition}


\begin{proposition}[utile=-3]
Soit $P$ un prédicat sur l'ensemble $E$. Alors
\begin{eqnarray*}
\non\cro{\forall x\in E \qsep P(x)} &\ssi&
\cro{\exists x\in E \qsep \non\p{P(x)}},\\
\non\cro{\exists x\in E \qsep P(x)} &\ssi&
\cro{\forall x\in E \qsep \non\p{P(x)}}.
\end{eqnarray*}
\end{proposition}

\begin{exoUnique}
\exo Soit $f:\R\to\R$. Écrire les phrases suivantes avec des
  quantificateurs. En déduire leur négation. \begin{center}
  \og $f$ est majorée \fg,
  \og $f$ est croissante \fg, \og $f$ est décroissante \fg. \end{center}
\end{exoUnique}
% \subsection{Principes de démonstration}

% Utiliser dans une preuve une hypothèse de la forme $\cro{P \et Q}$.



% Prouver une proposition de la forme $\cro{P \et Q}$.



% Utiliser dans un raisonnement une proposition de la forme $\cro{P \ou Q}$.

% \begin{exos}
% \exo Soit $a,b\in\R$. Montrer que
%    \[\max\p{a,b}=\frac{a+b+\abs{a-b}}{2} \et
%      \min\p{a,b}=\frac{a+b-\abs{a-b}}{2}\]
% \exo Si $n$ est la somme de deux carrés d'entiers, montrer que le reste
%   de la division euclidienne de $n$ par 4 n'est jamais égal à 3.
%   \begin{sol}
%   Traiter les cas où~:
%   \begin{itemize}
%   \item $a$ et $b$ sont pairs~: on trouve 0.
%   \item $a$ est pair, $b$ est impair~: on trouve 1. Le cas où $a$ est impair
%     et $b$ est pair se montre de même.
%   \item $a$ et $b$ sont impairs.
%   \end{itemize}
%   \end{sol}
% \end{exos}

% Prouver une proposition de la forme $P \ou Q$.

% Le principe à retenir est le suivant~:
% \begin{itemize}
% \item Si $P$ est vraie, alors $\cro{P \ou Q}$ est vraie.
% \item Sinon, on peut rajouter $\p{\non\ P}$ aux hypothèses et on cherche à
%   montrer que $Q$ est vraie.
% \end{itemize}

% \begin{exos}
% \exo Montrer qu'il existe deux irrationnels strictement positifs $x$ et $y$
%   tels que $x^y$ soit rationnel.
% \end{exos}


\section{Ensemble}
\subsection{Ensemble, élément}

\begin{definition}[utile=-3]
Les notions d'\emph{ensemble}, d'\emph{élément} et d'\emph{appartenance} sont des notions
premières en mathématiques que l'on ne définit pas. Intuitivement, un ensemble
est une collection d'objets mathématiques appelés éléments. La notation $x\in E$
signifie que l'élément $x$ appartient à l'ensemble $E$.
\end{definition}

\begin{remarqueUnique}
\remarque Si $x_1,\ldots,x_n$ sont des objets mathématiques, l'ensemble
  constitué de ces éléments est noté $\ens{x_1,\ldots,x_n}$.
\end{remarqueUnique}

\begin{definition}[utile=-3]
Soit $A$ et $B$ deux ensembles. On dit que $A$ est \emph{inclus} dans $B$ et on note
$A\subset B$ lorsque
\[\forall x\in A \qsep x\in B.\]  
\end{definition}

\begin{proposition}[utile=-3]
Deux ensembles $A$ et $B$ sont égaux lorsqu'ils possèdent les mêmes
éléments, c'est-à-dire lorsque
\[A\subset B \et B\subset A.\]
\end{proposition}

\begin{remarqueUnique}
\remarque En particulier $\ens{0,1}=\ens{1,0}$ et $\ens{0,0,1}=\ens{0,1}$.
\end{remarqueUnique}

\begin{definition}[utile=-3]
Soit $E$ un ensemble. On appelle partie de $E$ tout ensemble $A$ inclus dans $E$.
L'ensemble des parties de $E$ est noté $\mathcal{P}(E)$.
\end{definition}

\begin{remarqueUnique}
\remarque Un même objet mathématique peut très bien, selon le contexte, être un élément
  ou un ensemble. Par exemple, l'ensemble $\N$ est un élément de $\mathcal{P}(\R)$.
\end{remarqueUnique}

\begin{exoUnique}
\exo Déterminer $\parties{\ens{1,2}}$,
  $\parties{\parties{\emptyset}}$ et $\parties{\parties{\parties{\emptyset}}}$.
\end{exoUnique}

\begin{sol}
$\parties{\ens{1,2}}=\ens{\emptyset,\ens{1},\ens{2},\ens{1,2}}$. $\parties{\emptyset}=\ens{\emptyset}$ donc $\parties{\parties{\emptyset}}=\ens{\emptyset, \ens{\emptyset}}$. Pour finir, $\parties{\parties{\emptyset}}=\ens{\emptyset,\ens{\emptyset},\ens{\emptyset,\ens{\emptyset}},\ens{\ens{\emptyset}}}$.

\end{sol}

\subsection{Opérations élémentaires}

\begin{definition}[utile=-3]
Soit $E$ un ensemble et $P$ un prédicat sur $E$. On définit
\[\enstq{x\in E}{P(x)}\]
comme l'ensemble des éléments $x$ de $E$ tels que $P(x)$ est vrai. C'est une partie
de $E$.
\end{definition}

\begin{definition}[utile=-3]
Soit $A$ et $B$ deux parties de $E$. On définit
\[A\cap B\defeq\enstq{x\in E}{x\in A \et x\in B}, \qquad
  A\cup B\defeq\enstq{x\in E}{x\in A \ou x\in B},\]
\[\overline{A}\defeq\enstq{x\in E}{x\not\in A}.\]
\end{definition}

\begin{remarques}
\remarque Le complémentaire de $A$ dans $E$ est aussi noté $A^c$.
\remarque On dit que deux ensembles $A$ et $B$ sont \emph{disjoints} lorsque
  $A\cap B=\emptyset$.
\end{remarques}

\begin{proposition}[utile=-3]
Soit $A, B, C$ trois parties de $E$. Alors
\begin{eqnarray*}
A\cap\p{B\cup C}&=&\p{A\cap B}\cup\p{A\cap C},\\
A\cup\p{B\cap C}&=&\p{A\cup B}\cap\p{A\cup C}.
\end{eqnarray*}
\end{proposition}

\begin{preuve}
Avec la proposition 1.3, on raisonne par équivalence avec $x\in A\cap\p{B\cup C}$.
\end{preuve}

\begin{proposition}[utile=-3, nom={Lois de \nom{Morgan}}]
Soit $A$ et $B$ deux parties de $E$. Alors
\begin{eqnarray*}
\overline{A\cap B}&=&\overline{A}\cup \overline{B},\\
\overline{A\cup B}&=&\overline{A}\cap \overline{B},\\
\overline{\overline{A}}&=&A.
\end{eqnarray*}
\end{proposition}

\begin{preuve}
Avec la proposition 1.4, on raisonne par équivalence avec $x\in \overline{A\cap B}$.
\end{preuve}

\begin{exos}
\exo Soit $A$ et $B$ deux parties d'un même ensemble. Montrer que
  $\parties{A\cap B}=\parties{A}\cap\parties{B}$.
\exo Soit $A$, $B$, $C$ trois parties d'un ensemble $E$ non vide.
  \begin{questions}
  \question Si $A\cup B=A\cup C$, a-t-on $B=C$~?
  \question Si $A\cup B=A\cap B$, a-t-on $A=B$~?
  \question Montrer que si $A\cup B=A\cup C$ et $A\cap B=A\cap C$, alors $B=C$.
  \question Montrer que si $A\cup B=E$ et $A\cap B=\emptyset$, alors $A=\overline{B}$ et
  $B=\overline{A}$.
  \end{questions}
\end{exos}

\begin{definition}[utile=-3]
Soit $A$ et $B$ deux parties de $E$. On définit
\[A\setminus B\defeq\enstq{x\in E}{x\in A \et x\not\in B}.\]
\end{definition}

\begin{remarques}
\remarque L'ensemble $A\setminus B$ se lit \og $A$ privé de $B$ \fg.
\remarque Si $A$ est une partie de $E$, $\bar{A}=E\setminus A$.
\end{remarques}

\begin{definition}
Soit $A$ et $B$ deux ensembles. On définit $A\times B$ comme l'ensemble
des \emph{couples} $\p{a,b}$ avec $a\in A$ et $b\in B$. Par définition, deux couples
$\p{a_1,b_1}, \p{a_2,b_2}\in A\times B$ sont égaux lorsque $a_1=a_2$ et $b_1=b_2$.
\end{definition}

\begin{definition}[utile=-3]
\begin{itemize}
\item Si $A_1,\ldots,A_n$ sont $n$ ensembles, on définit
  $A_1\times\cdots\times A_n$ comme l'ensemble des $n$-uplets
  $\p{a_1,\ldots,a_n}$ avec $a_1\in A_1,\ldots,a_n\in A_n$. Par définition, deux $n$-uplets
  $\p{a_1,\ldots,a_n}, \p{b_1,\ldots,b_n}\in A_1\times\cdots\times A_n$ sont égaux lorsque
  \[\forall k\in\intere{1}{n} \qsep a_k=b_k.\]
\item Si A est un ensemble et $n\in\N$, on définit $A^n$ comme
  \[A^n\defeq\underbrace{A\times A\times\cdots\times A}_{\text{$n$ fois $A$}}.\]
\end{itemize}
\end{definition}

\begin{remarques}
\remarque $A^1$ est l'ensemble des 1-uplets $(a)$, pour $a\in A$; on confondra cet
  ensemble avec $A$. Quant à $A^0$, c'est l'ensemble qui contient un unique élément,
  le $0$-uplet $()$.
\remarque Pour énoncer qu'un prédicat portant sur deux variables est vrai, on peut écrire
  \og $\forall x\in A\qsep \forall y\in A\qsep P(x,y)$ \fg.
  On condense cependant souvent cette phrase en \og $\forall (x,y)\in A^2\qsep P(x,y)$ \fg ou
  en \og $\forall x,y\in A\qsep P(x,y)$\fg.
  % Daniel Perrin le fait
\end{remarques}

\section{Application}

\subsection{Définition, exemples}
\begin{definition}[utile=-3]
Soit $E$ et $F$ deux ensembles. Une \emph{application} $f$ de $E$ dans $F$ associe
à tout élément $x\in E$ un unique élément $f(x)\in F$, appelé image de $x$ par
$f$. On note
\[\dspappli{f}{E}{F}{x}{f(x).}\]
On dit que $E$ est le \emph{domaine} de $f$ et que $F$ est son \emph{codomaine}. L'ensemble des
applications de $E$ dans $F$ est noté $\mathcal{F}\p{E,F}$.
\end{definition}


\begin{remarques}
\remarque Deux applications sont égales lorsqu'elles ont même domaine et codomaine et qu'elles prennent la
  même valeur en chaque point de ce domaine.
\remarque On utilise aussi les expressions \og ensemble de départ \fg et \og ensemble d'arrivée \fg
  d'une application pour désigner respectivement son domaine et son codomaine.
\remarque \og application \fg et \og fonction \fg sont synonymes. L'usage veut cependant que l'on
  réserve le mot \og fonction \fg aux applications dont le domaine et le codomaine sont des parties de $\C$.
\remarque L'ensemble $\mathcal{F}\p{E,F}$ est aussi noté $F^E$.
\remarque Pour les fonctions usuelles, il arrive qu'on omette les parenthèses et
  qu'on écrive $\sin x$ au lieu de $\sin(x)$. Cependant, on ne se permettra pas de faire cela avec les
  autres fonctions.
\end{remarques}

\begin{definition}[utile=-3]
  Si $f$ est une application de $E$ dans $F$, on appelle \emph{graphe} de $f$ l'ensemble
  \[\enstq{\p{x,y}\in E\times F}{f(x)=y}.\]
  \end{definition}
  

\begin{definition}
Soit $A$ une partie de $E$. On appelle \emph{fonction caractéristique} de $A$ et on note $\mathds{1}_A$ l'application de $E$ dans $\ens{0,1}$ définie
par
  \[\forall x\in E \qsep \mathds{1}_A(x)\defeq
    \begin{cases}
    1 & \text{si $x\in A$}\\
    0 & \text{sinon.}
    \end{cases}\]
\end{definition}

\begin{remarques}
\remarque Deux parties $A$ et $B$ de $E$ sont égales si et seulement si $\mathds{1}_{A}=\mathds{1}_{B}$.
\remarque Si $A$ et $B$ sont deux parties de $E$, alors
  \[\forall x\in E\qsep \mathds{1}_{A\cap B}(x)=\mathds{1}_{A}(x)\mathds{1}_{B}(x) \et \mathds{1}_{A\cup B}(x)=\max\p{\mathds{1}_{A}(x), \mathds{1}_{B}(x)}.\]
\end{remarques}

%% Remarques :
%% 1) 
%% 2) 
%% 3) Lorsqu'on définit une fonction f:E->F, il faudra toujours bien veiller à
%%    ce que :
%%    - f(x) a un sens pour tout x dans E
%%    - f(x) est dans F
%%     -f(x) ne dépend que de x : (problème avec exp(i theta)->theta
%%
%% Exemples :
%% 1) f : R -> R   x -> 1/(1+x^2)
%% 2) Si A est une partie de E, on définit
%%    f: P(E) -> P(E)   X -> X \cap A
%% 3) f : F(R,R) -> F(RPs,RPs)  g -> exp o g o ln


\begin{definition}[utile=-3]
Soit $f:E\to F$ et $y\in F$. On appelle \emph{antécédent} de $y$ tout élément $x\in E$
tel que $f(x)=y$.  
\end{definition}

%% Exemples :
%% 1) Soit f : R -> R  x -> sin(x)
%%    Calcul des antécédents du réel y

\begin{exoUnique}
\exo Soit $f$ l'application de $\R^2$ dans $\R^2$ qui au couple $(x,y)$
  associe le couple $(x+2y,xy)$. Déterminer les antécédents de $(3,1)$.
  \begin{sol}
  Ce sont $(1,1)$ et $(2,1/2)$.
  \end{sol}
\end{exoUnique}

\begin{definition}[utile=-3]
Soit $f:E\to F$.
\begin{itemize}
\item Soit $B$ une partie de $F$. On appelle image réciproque de $B$
  et on note $f^\applirec (B)$ l'ensemble des éléments de $E$
  dont l'image par $f$ est dans $B$.
  \[f^\applirec (B)\defeq\enstq{x\in E}{f(x)\in B}\]
\item Soit $A$ une partie de $E$. On appelle image directe de $A$ 
  et on note $f(A)$ l'ensemble des éléments de $F$ qui sont
  image d'un élément de $A$ par $f$.
  \[f(A)\defeq\enstq{y\in F}{\exists x\in A \qsep f(x)=y}\]  
  L'ensemble $f(E)$ est appelé image de $f$ et noté $\im f$.
\end{itemize}
\end{definition}

\begin{remarqueUnique}
\remarque L'ensemble image $f(A)$ est aussi noté
  \[\ensim{f(x)}{x\in A}.\]
\end{remarqueUnique}

\begin{exos}
\exo Soit $f$ la fonction
  \[\dspappli{f}{\R}{\R}{x}{\sin x}\]
  Calculer $f^{-1}(f(\ens{\pi/2}))$ et $f(f^{-1}(\ens{0, 2}))$.
\exo Soit $f$ la fonction de $\C\setminus\ens{\ii}$ dans $\C$ qui à $z$
  associe $\frac{z+\ii}{z-\ii}$. Calculer %$f\p{\U\setminus\ens{i}}$ et  
  $f^\applirec \p{\U}$.
  \begin{sol}
  On a
  \[f\p{e^{i\theta}}=-i\cotan\p{\frac{\theta}{2}-\frac{\pi}{4}}\]
  donc $f\p{\U\setminus\ens{i}}=i\R$.) On trouve $f^\applirec \p{\U}=\R$
  (utiliser que la norme au carré fait $1$ et passer par le conjugué).
  \end{sol}
\exo Soit $f$ la fonction de $\R$ dans $\R$ définie par
  \[\forall x\in\R \qsep f(x)\defeq\frac{x}{1+x^2}.\]
  En lisant le tableau de variations de $f$, intuiter $f(\R)$, puis
  prouver rigoureusement ce résultat.
  \begin{sol}
  On trouve $\interf{-1/2}{1/2}$ (étude de fonctions avec dérivée et tableau de variations).
  \end{sol}
% \exo Soit $f$ l'application de $\R$ dans $\R$ définie par
%   \[\forall x\in\R \qsep f(x)=\sqrt{3}\sin x+\cos x\]
%   Déterminer $f^\applirec\p{\interf{0}{1}}$.
\exo Soit $f$ une application de $E$ dans $F$. Si $A$ est une partie de $E$,
  comparer $f^{-1}(f(A))$ et $A$. De même, si $B$ est une partie de $F$,
  comparer $f(f^{-1}(B))$ et $B$.  
 \begin{sol}
 $f^{-1}(f(A))\supset A$ et $f(f^{-1}(B))\subset B$.
 \end{sol}
\end{exos}

%% Exemples :
%% 1) Étudier l'image réciproque par f(x)=x^2 de [-1,1] puis l'image directe par
%%    f de [-1,1]

\begin{definition}[utile=-3]
Soit $f$ une application de $E$ dans $F$.
\begin{itemize}
\item Si $A$ est une partie de $E$, l'application
  \[\dspappli{\restri{f}{A}}{A}{F}{x}{f(x)}\]
  est appelée \emph{restriction} de $f$ à $A$.
\item On dit qu'une application $g$ est un
  \emph{prolongement} de $f$ lorsque $f$ est une restriction de $g$.
\item Si $B$ est une partie de $F$ telle que
  \[\forall x\in E \qsep f(x)\in B\]
  l'application
  \[\dspappli{\corestri{f}{B}}{E}{B}{x}{f(x)}\]
  est appelée \emph{corestriction} de $f$ à $B$.
\end{itemize}
\end{definition}

\begin{remarqueUnique}
\remarque Soit $f:E\to F$, $A$ une partie de $E$ et $B$ une partie de $F$
  telles que
  \[\forall x\in A\qsep f(x)\in B.\]
  Alors, on peut définir l'application
  \[\dspappli{\restricorestri{f}{A}{B}}{A}{B}{x}{f(x)}\]
  appelée restriction de $f$ à $A$, corestreinte à $B$.
\end{remarqueUnique}


%% Exemples :
%% 1) l'application f : [-Pi/2,Pi/2] -> [-1,1]  x -> sin(x) est la restriction
%%    de l'application de g à [-Pi/2,Pi/2] corestreinte à [-1,1] où g est
%%    l'application R -> R  x -> sin(x)
%% 2) l'application f : R -> R  x -> (sin x)/x si x \neq 0 et x-> 1 si x=0
%%    est un prolongement de x->sin x/x. C'est d'ailleurs le seul prolongement
%%    continu en 0.

\begin{definition}[utile=-3]
Soit $f:E\to F$ et $g:F\to G$. On définit l'application $g\circ f$ de $E$ dans $G$ par
\[\forall x\in E\qsep (g\circ f)(x)\defeq g\p{f(x)}.\]
\end{definition}

\begin{remarqueUnique}
\remarque Si $A$ est une partie de $E$, alors $(g\circ f)(A)=g(f(A))$. De même, si
  $B$ est une partie de $G$, alors \[(g\circ f)^{-1}(B)=f^{-1}(g^{-1}(B)).\]
\end{remarqueUnique}

%% Exemple :
%% 1) Si f : R -> RP x -> x^2 et g : RP -> R x -> sqrt(x)
%%    alors (f o g) : RP -> RP  x -> x
%%    et    (g o f) : R  -> R   x -> abs(x) 
%%
%% Remarque :
%% 1) si f : E -> B et g : F -> G, et que B\subset F, on peut composer gof
%%    il suffit de restreindre g

\begin{proposition}[utile=-3]
Soit $f:E\to F$, $g:F\to G$ et $h:G\to H$. Alors
\[\p{h\circ g}\circ f=h\circ\p{g\circ f}.\]
On note cette application $h\circ g\circ f$.
\end{proposition}

\subsection{Application injective, surjective, bijective}

\begin{definition}[utile=-3]
Soit $f:E\to F$. On dit que $f$ est \emph{injective} lorsque
\[\forall x_1,x_2\in E \qsep f\p{x_1}=f\p{x_2} \implique x_1=x_2\]
c'est-à-dire lorsque tout élément de $F$ a au plus un antécédent.
\end{definition}

\begin{exos}
\exo Soit $f$ une fonction de $\R$ dans $\R$. Montrer que si $f$ est
  strictement monotone alors elle est injective. La réciproque est-elle vraie~?
\exo Soit $f$ une fonction de $\R$ dans $\R$ telle que :
  $\forall x,y\in\R \qsep \abs{f(x)-f(y)}\geq\abs{x-y}$. Montrer qu'elle est
  injective.
\exo Soit $\phi$ l'application qui à la fonction $f$ de
  $[-1,1]$ dans $\R$ associe la fonction $\phi(f)$ de $\R$ dans $\R$
  définie par
  \[\forall x\in\R \qsep \cro{\phi(f)}(x)\defeq f\p{\sin x}\]
  Montrer que $\phi$ est injective.
  \begin{sol}
  $\phi(f)=\phi(g)$ implique $f\p{\sin x}=g\p{\sin x}, \forall x\in\R$ donc $f$ et $g$ sont égales étant donné les valeurs prises par $\sin$.
  \end{sol}
\exo Soit $E$ un ensemble et $A$ une partie de $E$. Donner une condition
  nécessaire et suffisante sur $A$ pour que
  \[\dspappli{\phi}{\parties{E}}{\parties{E}}{X}{X\cap A}\]
  soit injective.
  \begin{sol}
  $A=E$, sinon on raisonne par l'absurde en prenant un élément $b$ de $E\setminus A$ et $X=A$, $Y=A\cup \ens{b}$.
  \end{sol}
\end{exos}

%% Exemple :
%% 1) f : RP -> R x->x^2, alors f est injective
%% 2) f : R -> R^2  x -> (x,0) est injective
%% 3) l'application phi qui a une fonction f de [-1,1] dans R associe le fonction
%%    g de R dans R définie par g(x)=f(sin(x)) est une application injective.  
%%
%% En pratique :
%% Pour montrer qu'une application est injective :
%% 1) On montre que l'équation f(x)=y a au plus une solution : pour cela, on
%%    peut raisonner par implication (ou par analyse)
%% 2) On montre que f(x_1)=f(x_2) implique x_1=x_2

\begin{definition}[utile=-3]
Soit $f:E\to F$. On dit que $f$ est \emph{surjective} lorsque
\[\forall y\in F \qsep \exists x\in E \qsep f(x)=y\]
c'est-à-dire lorsque tout élément de $F$ a au moins un antécédent.
\end{definition}

%% Exemple :
%% 1) f : R -> RP   x -> x^2 est surjective
%% 2) f : R^2 -> R   (x,y) -> x est surjective
%% 3) Si D est une droite du plan P, l'application f de P dans D qui au point M
%%    associe son projeté orthogonal sur D est surjective 
%%
%% En pratique :
%% 1) On montre que l'équation f(x)=y a au moins une solution
%%
%% Remarque :
%% 1) f est surjective si et seulement si f(E)=F

\begin{proposition}[utile=-3]
Une application $f:E\to F$ est surjective si et seulement si $\im f=F$.
\end{proposition}

\begin{exos}
% \exo Soit $f$ le fonction de $\R$ dans $\R$ définie par~:
%   $\forall x\in\R \qsep f(x)=x^3-x$. $f$ est-elle injective~? surjective~?
%   Calculer $f\p{\RP}$ et $f^\applirec \p{\interf{-6}{6}}$.
\exo L'application
  \[\dspappli{\phi}{\mathcal{F}\p{\R,\R}}{\R}{f}{f\p{0}}\]
  est-elle injective~? surjective~?
% \exo Soit $f:E\to F$ et $g:E\to G$. On définit l'application $\phi$ de
%   $E$ dans $F\times G$ par \[\forall x\in E \qsep \phi(x)\defeq \p{f(x),g(x)}.\]
%   Que dire des assertions \og $\phi$ est injective si et seulement si $f$
%   et $g$ le sont \fg et \og $\phi$ est surjective si et seulement si $f$ et
%   $g$ le sont \fg~?
%   \begin{sol}
%   Faux. Si $f=\id$ et $g=1$, $\phi$ est injective sans que $g$ le soit. Il suffit en fait que l'une des deux soit injective pour que $\phi$ le soit.
%   Pour la deuxième, c'est l'inverse, si $\phi$ est surjective, alors $f$ et $g$ le sont. Mais si $f=\id=g$, $\phi$ n'est pas surjective.
%   \end{sol}
\end{exos}

\begin{definition}[utile=-3]
On dit qu'une application $f:E\to F$ est \emph{bijective} lorsqu'elle est injective et
surjective, c'est-à-dire lorsque tout élément de $F$ possède un unique antécédent.
\end{definition}

\begin{exos}
\exo Montrer que la fonction $f$ qui à $x$ associe $\frac{1+\ii x}{1-\ii x}$
  réalise une bijection de $\R$ dans $\U\setminus\ens{-1}$.
  \begin{sol}
  $\displaystyle z=\frac{1+ix}{1-ix}\Longleftrightarrow x=-i \frac{z-1}{z+1}$, ce qui est bien dans $\R$, on le constate en prenant le conjugué et en utilisant que $\conj{z}=1/z$.
  \end{sol}
\exo Montrer que l'application
  \[\dspappli{f}{\N^2}{\N}{\p{a,b}}{2^a\p{2b+1}-1}\]
  est bijective.
\exo Soit $X$ un ensemble et $f:X^2\to X$ une bijection. Montrer que
  \[\dspappli{g}{X^3}{X}{(x,y,z)}{f\p{x,f(y,z)}}\]
  est bijective.
\end{exos}


%% Exemples :
%% 1) Si n>=1, f : RP -> RP  x -> x^n  est bijective
%% 2) f : RPs -> R  x -> ln(x) est bijective
%% 3) f : R -> RPs  x -> exp(x) est bijective
%% 4) f : [-Pi/2,Pi/2] -> [-1,1] -> sin(x) est bijective
%% 5) L'application de P(E) dans P(E) qui a X associe son complémentaire est
%%    bijective.
%% 6) L'apllication phi de F(RP,RP) dans F(R,R) qui à la fonction f associe la
%%    fonction g définie par g(x)=ln(f(exp(x))) est bijective.
%%
%% En pratique :
%% 1) On montre que l'équation f(x)=y a une et une seule solution
%% 2) On montre que f est injective, puis que f est surjective

\begin{proposition}[utile=-3]
\begin{itemize}
\item La composée de deux applications injectives est injective.
\item La composée de deux applications surjectives est surjective.
\item La composée de deux applications bijectives est bijective.
\end{itemize}
\end{proposition}

\begin{exos}
\exo Soit $f:E\to F$ et $g:F\to G$. Montrer que si $g\circ f$ est injective,
  alors $f$ est injective. De même, montrer que si $g\circ f$ est surjective,
  alors $g$ est surjective.
\exo Est-il vrai que si $g\circ f$ est bijective, $f$ et $g$ le sont~?
\begin{sol}
Contre-exemple : $f:\ens{0,1}\to\R$ telle que $f(x)=x$, $g:\R\to\ens{0,1}$ telle que $g(x)=x$ si $x\in \ens{0,1}$ et $g(x)=0$ sinon.
\end{sol}
\end{exos}

\begin{definition}[utile=-3]
Soit $E$ un ensemble. On appelle \emph{identité} et on note $\id_E$
l'application de $E$ dans $E$ définie par
\[\forall x\in E \qsep \id_E(x)\defeq x.\]
Si $f$ est une application de $E$ dans $F$
\[f\circ \id_E=f \et \id_F\circ f=f.\]
\end{definition}

\begin{proposition}[utile=-3]
Soit $f$ une application de $E$ dans $F$.
\begin{itemize}
\item L'application $f$ est bijective si et seulement si il existe une
  application $g:F\to E$ telle que
  \[g\circ f=\id_E \et f\circ g=\id_F.\]
  Si tel est le cas, $g$ est unique; on l'appelle \emph{bijection réciproque} de $f$
  et on la note $f^{-1}$.
\item Si $f:E\to F$ est bijective, $f^{-1}$ est bijective et $\p{f^{-1}}^{-1}=f$.
\end{itemize}
\end{proposition}

\begin{preuve}
On raisonne par double implication. 
\begin{itemize}
\item [$\bullet$]
Tout d'abord, si il existe une application $g:F\to E$ telle que \[g\circ f=\id_E \et f\circ g=\id_F.\]
Comme $\id_E$ est injective, $g\circ f$ est injective donc $f$ est injective (cf. ex. précédent). De même, comme $\id_F$ est surjective, $f\circ g $ est surjective donc $f$ est surjective.
Finalement, $f$ est bijective.
\item [$\bullet$] Si $f$ est bijective. Fixons $y\in F$, il existe un unique $x$ dans $E$ tel que $f(x)=y$. Notons-le $g(y)$. On définit ainsi une fonction $g:F\to E$ vérifiant $f(g(y))=y, \forall y \in F$. $g$ associe donc à tout élément de $F$ son antécédent par $f$. Ainsi, pour tout $x\in E$, $g$ associe à $f(x)$ son antécédent par $f$, c'est-à-dire $x$, donc $g(f(x))=x$.
\end{itemize}
Pour l'unicité, supposons qu'il existe $h$ qui réalise la même chose, on a donc $h\circ f=\id_E$, on compose par $g$ à droite et on obtient $h=g$.

Enfin, si $f$ est bijective, il existe donc $f^{-1} : F\to E$ telle que \[f^{-1}\circ f=\id_E \et f\circ f^{-1}=\id_F.\] donc en appliquant la première partie du théorème à $f^{-1}$ avec $f$ qui joue le rôle de $g$, on obtient bien que $f^{-1}$ est bijective et $\p{f^{-1}}^{-1}=f$.

\end{preuve}

\begin{remarques}
\remarque Soit $f$ une bijection de $E$ dans $F$. Quel que soit $y\in F$, si $x\in E$
  est tel que $f(x)=y$, alors $f^{-1}(y)=x$.
\remarque La fonction $\ln$ de $\RPs$ dans $\R$ est une bijection et sa
  bijection réciproque est la fonction $\exp$ de $\R$ dans $\RPs$.
\end{remarques}

\begin{exos}
\exo Montrer que l'application
  \[\dspappli{f}{\Z^2}{\Z^2}{(x,y)}{(2x+y,5x+3y)}\]
  est bijective et calculer $f^{-1}$.
	\begin{sol}
$f^{-1}(a, b)=(3a-b,-5a+2b)$
	\end{sol}
\exo Soit $f$ une bijection de $\R$ dans $\R$. Montrer que si $f$ est
  strictement croissante, il en est de même pour $f^{-1}$. Que dire si $f$
  est impaire~?
  \begin{sol}
  Si $f$ est impaire, on regarde $f(f^{-1}(-x))=-x$ et $f(-f^{-1}(x))=-f(f^{-1}(x))=-x$ donc $f^{-1}(-x)=-f^{-1}(x)$ par injectivité.
  
  Une fonction paire ne peut pas être bijective.
  \end{sol}
% \exo Soit $f:F\to E$ et $g:G\to E$ deux applications. Montrer qu'il existe
%   une application $h:G\to F$ telle que $g=f\circ h$ si et seulement si
%   $g(G)\subset f(F)$.
\end{exos}

%% Exemples : calcul de la bijection réciproque de
%% 1) Si n>=1, f : RP -> RP  x -> x^n  est bijective
%% 2) f : RPs -> R  x -> ln(x) est bijective
%% 3) f : R -> RPs  x -> exp(x) est bijective
%% 4) f : [-Pi/2,Pi/2] -> [-1,1] -> sin(x) est bijective
%% 5) L'application de P(E) dans P(E) qui a X associe son complémentaire est
%%    bijective.
%% 6) L'apllication phi de F(RP,RP) dans F(R,R) qui à la fonction f associe la
%%    fonction g définie par g(x)=ln(f(exp(x))) est bijective.
%%
%% En pratique :
%% 1) On résout l'équation f(x)=y. L'unique solution en x est f^{-1}(y)
%%
%% Applications :
%% 1) Changement de variable
%%    \forall x\in E_1 f_1(x)=f_2(x) \ssi
%%      \forall t\in E_2 f_1\p{x(t)}=f_2\p{x(t)}
%%    pour cela, il suffit que x soit surjective (exemple avec sinus).
%%    Le plus souvent le changement de variable est bijectif (exemple avec
%%    l'exponentielle)

\begin{proposition}[utile=-3]
Soit $f:E\to F$ et $g:F\to G$ deux applications bijectives. Alors $g\circ f$
est bijective et
\[\p{g\circ f}^{-1}=f^{-1} \circ g^{-1}.\]
\end{proposition}

\subsection{Famille}

Si $E$ est un ensemble, il est courant de se donner $n$ éléments
$f_1,\ldots,f_n$ de $E$. Cela revient à définir une application
\[\dspappli{f}{\intere{1}{n}}{E}{i}{f(i)}\]
où l'on pose $f(i)\defeq f_i$ pour tout $i\in\intere{1}{n}$. Nous dirons que $f$ est une
famille d'éléments de $E$ indexée par $\intere{1}{n}$. On peut généraliser ce principe
et construire des familles indexées par un ensemble quelconque. Par exemple, on peut considérer
l'application $f$ de $\R$ dans $\mathcal{F}(\R, \R)$, qui à $\lambda\in\R$ associe la
fonction $f_\lambda:\R\to\R$ définie par
\[\forall x\in\R\qsep f_\lambda(x)\defeq\e^{\lambda x}.\]
On a ainsi défini une famille d'éléments de $\mathcal{F}(\R, \R)$ indexée par $\R$.

\begin{definition}[utile=-3]
Soit $E$ un ensemble et $I$ un ensemble, appelé ensemble d'indices. On appelle
\emph{famille d'éléments de $E$ indexée par $I$} toute application
\[\dspappli{f}{I}{E}{i}{f_i.}\]
Cette application est notée $\p{f_i}_{i\in I}$. L'ensemble des familles
d'éléments de $E$ indexées par $I$ est noté $E^I$.
\end{definition}

\begin{remarques}
\remarque Une famille d'éléments de $E$ indexée par $\N$ est une suite d'éléments de $E$.
\remarque On appelle sous-famille d'une famille $\p{f_i}_{i\in I}$ toute famille
  de la forme $\p{f_i}_{i\in J}$ où $J$ est une partie de $I$.
\remarque Si $A$ est un ensemble, on dit qu'une famille $(f_i)_{i\in I}$ est la famille des éléments de $A$ lorsque $f$ est une bijection de $I$ dans $A$. Le fait de parler de \og la\fg famille des éléments de $A$ est un abus de langage, car cette famille n'est pas unique.
% Par abus de langage, bien qu'elle ne soit pas unique, on parlera souvent de la famille des éléments d'un ensemble pour désigner une telle famille.
% Si $A$ est un ensemble, on appelle famille des éléments de $A$
%   l'application
%   \[\dspappli{f}{A}{A}{a}{f_a=a}\]
%   que l'on note $(f_a)_{a\in A}$ ou plus simplement $(f_i)_{i\in I}$ (où $I=A$, ce
%   que l'on s'empresse d'oublier).
\end{remarques}

%% Exemples :
%% 1) une suite  d'éléments de E est une famille indexée par N
%% 2) (f_a)_(a\in\R) la famille indexée par \R des fonctions de \R dans \C définies
%%    par f_a(x)=exp(iax)
%% 3) La famille ([-1/n,1/n])_(n\in\Ns)
%% Exemples :
%% 1) la famille (f_n)_(n\in\N) définie par f_n(x)=exp(inx) est une sous-famille
%%    de la famille définie plus haut.

\begin{definition}[utile=-3]
Soit $E$ un ensemble et $\p{A_i}_{i\in I}$ une famille de parties de $E$. On
définit alors
\[\bigcap\limits_{i\in I} A_i\defeq \enstq{x\in E}{\forall i\in I \qsep x\in A_i},\]
\[\bigcup\limits_{i\in I} A_i\defeq \enstq{x\in E}{\exists i\in I \qsep x\in A_i}.\]
\end{definition}

%% Exemples :
%% 1) intersection des [-1/n,1/n]

\begin{exoUnique}
% \exo Soit $f:E\to F$ et $\p{A_i}_{i\in I}$ une famille de parties de
%   $E$. Comparer
%   \[f\p{\bigcap\limits_{i\in I} A_i} \et \bigcap\limits_{i\in I} f\p{A_i}\]
\exo Soit $f:E\to E$. On définit $f^n$ pour tout $n\in\N$ par
  \[f^0\defeq \id_E \et \cro{\forall n\in\N \qsep f^{n+1}\defeq f\circ f^n}\]
  Soit $A$ une partie de $E$. Pour tout $n\in\N$, on pose $A_n\defeq f^n(A)$.
  Enfin, on pose $B\defeq \cup_{n\in\N} A_n$. Montrer que $A\subset B$ et que
  $f(B)\subset B$.
%   De plus, montrer que si $C$ est une partie de $E$ telle que
%   $A\subset C$ et $f(C)\subset C$, alors $B\subset C$.
\end{exoUnique}

\begin{proposition}[utile=-3]
Soit $E$ un ensemble et $\p{A_i}_{i\in I}$ une famille de parties de $E$. Alors
\[\overline{\bigcap\limits_{i\in I} A_i}=\bigcup\limits_{i\in I} \overline{A_i} \et
  \overline{\bigcup\limits_{i\in I} A_i}=\bigcap\limits_{i\in I} \overline{A_i}.\]
\end{proposition}

\begin{definition}[utile=-3, nom=Partition]
Soit $E$ un ensemble et $\p{A_i}_{i\in I}$ une famille de parties de $E$.
On dit que $\p{A_i}_{i\in I}$ est une \emph{partition} de $E$ lorsque
\[E = \bigcup_{i\in I} A_i \qquad\et\qquad \cro{\forall i,j\in I\qsep i\neq j\implique
  A_i\cap A_j=\emptyset}.\]
\end{definition}

\begin{remarques}
\remarque La définition de partition peut varier d'un cours à l'autre.
  Dans certains cours, on demande en plus que les $A_i$ soient non vides; on appelle
  alors \emph{recouvrement disjoint} ce que nous appelons ici partition.
\remarque La notion de partition a été définie à l'aide de familles. Mais on peut aussi
  la définir de manière ensembliste; on dit qu'une partie $\mathcal{P}$ de
  $\mathcal{P}(E)$ est une \emph{partition (au sens ensembliste)} de $E$ lorsque  
  \begin{itemize}
  \item $\forall x\in E\qsep \exists A\in\mathcal{P}\qsep x\in A$.
  \item $\forall A_1, A_2\in\mathcal{P}\qsep A_1\neq A_2\implique A_1\cap A_2=\emptyset$.
  \item $\forall A\in\mathcal{P}\qsep A\neq\emptyset$. 
  \end{itemize}
  Remarquons que dans la définition ensembliste, on demande à ce que les ensembles appartenant à $\mathcal{P}$
  soient non vides.
\end{remarques}

\begin{exoUnique}
\exo Déterminer les partitions au sens ensembliste de $E\defeq\ens{1,2,3}$.
\end{exoUnique}

\section{Relation binaire}

\begin{definition}[utile=-3]
Soit $E$ un ensemble. On appelle \emph{relation binaire} sur $E$ tout prédicat
$\mathcal{R}$ défini sur $E\times E$. Si $x$ et $y$ sont deux
éléments de $E$ et $\mathcal{R}\p{x,y}$ est vrai, on écrit $x\mathcal{R} y$.
\end{definition}

\begin{definition}[utile=-3]
On dit qu'une relation binaire $\mathcal{R}$ sur $E$ est
\begin{itemize}
\item \emph{réflexive} lorsque
  \[\forall x\in E \qsep x\mathcal{R} x.\]
\item \emph{transitive} lorsque
  \[\forall x,y,z\in E \qsep \cro{x\mathcal{R} y \et y\mathcal{R} z} \implique
    x\mathcal{R} z.\]
\item \emph{symétrique} lorsque
  \[\forall x,y\in E \qsep x\mathcal{R} y \implique y\mathcal{R} x.\]
\item \emph{antisymétrique} lorsque
  \[\forall x,y\in E \qsep \cro{x \mathcal{R}y \et y\mathcal{R} x} \implique
    x=y.\]
\end{itemize}
\end{definition}

\subsection{Relation d'ordre}

\begin{definition}[utile=-3]
On dit qu'une relation binaire $\preceq$ est une
\emph{relation d'ordre} lorsqu'elle est
\begin{itemize}
\item réflexive~: $\forall x\in E \qsep x\preceq x.$
\item transitive~: $\forall x,y,z\in E \qsep \cro{x\preceq y \et y\preceq z}
  \implique x\preceq z.$
\item antisymétrique~: $\forall x,y\in E \qsep \cro{x\preceq y \et y\preceq x}
  \implique x=y.$
\end{itemize}
On appelle \emph{ensemble ordonné} tout ensemble muni d'une relation d'ordre.
\end{definition}

\begin{remarques}
\remarque La relation $\leq$ est une relation d'ordre sur $\R$.  La relation
  $\leq$ définie sur $\mathcal{F}\p{\R,\R}$ par
  \[\forall f,g\in\mathcal{F}\p{\R,\R} \qsep f\leq g \quad\ssi\quad
    \cro{\forall x\in\R \qsep f(x)\leq g(x)}\]
  est une relation d'ordre sur $\mathcal{F}\p{\R,\R}$.
 \remarque Si $E$ est un ensemble,
  la relation d'inclusion est une relation d'ordre sur $\parties{E}$. 
\remarque Si $\preceq$ est une relation d'ordre sur $E$, la relation
  $\succeq$ définie par
  \[\forall x,y\in E \qsep x\succeq y \quad\ssi\quad y\preceq x\]
  est une relation d'ordre appelée relation d'ordre opposée à la première.
\remarque La relation $<$ n'est pas une relation d'ordre sur $\R$ car elle
  n'est pas réflexive.
\end{remarques}

\begin{exoUnique}
\exo Montrer que la relation | définie sur $\N$ par
  \[\forall a,b\in\N \qsep a|b \quad\ssi\quad \cro{\exists k\in\N \qsep b=ka}\]
  est une relation d'ordre sur $\N$.
\end{exoUnique}

%% Exemples :
%% 1) relation d'ordre <= sur R
%% 2) relation d'ordre <= sur F(R,R)
%% 3) inclusion sur P(E)
%% 4) divisibilité sur N
%%
%% Remarques :
%% 1) si <= est une relation d'ordre, >= est aussi une relation d'ordre appelée
%%    relation d'ordre opposée
%% 2) La relation < n'est pas une relation d'ordre sur R car elle n'est pas
%%    réflexive

\begin{definition}[utile=-3]
On dit qu'une relation d'ordre $\preceq$ est totale lorsque
\[\forall x,y\in E \qsep x\preceq y \ou y\preceq x.\]
\end{definition}

\begin{remarqueUnique}
\remarque La relation d'ordre $\leq$ est totale sur $\R$. Par contre, les
  relations $\leq$ sur $\mathcal{F}\p{\R,\R}$, $\subset$ sur $\parties{E}$ et
  | sur $\N$ ne sont pas totales.
\end{remarqueUnique}

%% Exemples :
%% Voir exemples précédents
%% 1) est totale
%% 2) n'est pas totale : f(x)=0 et g(x)=x
%% 3) n'est pas totale dès que E possède plus de deux éléments
%% 4) n'est pas totale : 2 et 3

\begin{definition}[utile=-3]
Soit $\p{E,\preceq}$ un ensemble ordonné et $A$ une partie de $E$.
\begin{itemize}
\item On dit que $M\in E$ est un \emph{majorant} de $A$ lorsque
  \[\forall a\in A \qsep a\preceq M.\]
\item On dit que $m\in E$ est un \emph{minorant} de $A$ lorsque
  \[\forall a\in A \qsep m\preceq a.\]
\end{itemize}
\end{definition}

\begin{exoUnique}
\exo Soit $c>0$. On définit la relation $\preceq$ sur $\R^2$ par
  \[\forall (x,t),(x',t')\in\R^2 \qsep (x,t)\preceq (x',t')\quad\ssi\quad \abs{x'-x}\leq c\cdot (t'-t).\]
  Vérifier que c'est une relation d'ordre. Dessiner l'ensemble des majorants et
  des minorants d'un couple $(x_0,t_0)$. L'ordre est-il total~?
  \begin{sol}
  L'ensemble des majorants et des minorants forme un cône (on peut le tracer dans un repère $(x,t)$ pour une valeur donnée de $t_0$). L'ordre n'est pas
  total car $(0,0)$ et $(1,0)$ ne sont pas comparables.
  \end{sol}
\end{exoUnique}

%% Exemples :
%% 1) Si A=[0,1[, l'ensemble des majorants de A est [1,+infty[
%% 2) Z n'admet ni majorant ni minorant dans R
%% 3) Si (A_i)_(i\in I), un majorant de {A_i : i\in I} est l'union, et un minorant
%%    de cet ensemble est l'intersection.
%% 4) Si a,b sont des entiers naturels, A={a,b} est majoré par ab. ppcm(a,b) est
%%    un autre majorant (pour la relation d'ordre de divisibilité)

\begin{definition}[utile=-3]
Soit $\p{E,\preceq}$ un ensemble ordonné et $A$ une partie de $E$.
\begin{itemize}
\item On dit que $A$ admet un \emph{plus grand élément} lorsqu'il existe un majorant
  de $A$ appartenant à $A$. Si un tel élément existe, il est unique et on
  l'appelle plus grand élément de $A$.
\item On dit que $A$ admet un \emph{plus petit élément} lorsqu'il existe un minorant
  de $A$ appartenant à $A$. Si un tel élément existe, il est unique et on
  l'appelle plus petit élément de $A$.
\end{itemize}
\end{definition}

\begin{remarques}
\remarque Muni de l'ordre usuel, $\interfo{0}{1}$ admet un plus petit élément 0
  mais n'admet pas de plus grand élément. Muni de la relation de divisibilité,
  $\ens{2,3}$ n'admet ni de plus grand ni de plus petit élément.
\remarque Un ensemble admettant un plus petit ou un plus grand élément est non
   vide.
\remarque Si $E$ est totalement ordonné et $A$ est une partie finie non vide de
  $E$, alors il admet un plus petit et un plus grand élément.
\end{remarques}

\subsection{Relation d'équivalence}

\begin{definition}[utile=-3]
On dit qu'une relation binaire $\mathcal{R}$ sur $E$ est une
\emph{relation d'équivalence} lorsqu'elle est
\begin{itemize}
\item réflexive~: $\forall x\in E \qsep x\mathcal{R} x.$
\item transitive~: $\forall x,y,z\in E \qsep \cro{x\mathcal{R} y \et
  y\mathcal{R} z} \implique x\mathcal{R} z.$
\item symétrique~: $\forall x,y\in E \qsep x\mathcal{R} y \implique
  y\mathcal{R} x.$
\end{itemize}
\end{definition}

\begin{remarqueUnique}
\remarque Si $E$ est un ensemble quelconque, la relation d'égalité est une
  relation d'équivalence. Si $n\in\Ns$, la relation $\mathcal{R}$ définie sur $\Z$ par
  \og $\forall a,b\in\Z \qsep a\mathcal{R} b \quad\ssi\quad a\equiv b\ [n]$ \fg
  est une relation d'équivalence. De même, si $f$ est une application
  de $E$ dans $F$, la relation $\mathcal{R}$ définie sur $E$ par
  \og $\forall x,y\in E \qsep x\mathcal{R}y \quad\ssi\quad f(x)=f(y)$ \fg
  est une relation d'équivalence.
\end{remarqueUnique}

\begin{exoUnique}
\exo Soit $E$ un ensemble. Montrer que la relation $\mathcal{R}$ définie
  sur $\parties{E}$ par
  \[\forall A,B\in\parties{E} \qsep A\mathcal{R} B \quad\ssi\quad
    \text{\og Il existe une bijection de $A$ dans $B$. \fg}\]
  est une relation d'équivalence.
\end{exoUnique}

%% Exemples :
%% 1) l'égalité sur E
%% 2) Si n\in\N la relation aRb \ssi a=b+kn est une relation d'équivalence sur
%%    \Z
%% 3) Si f : E -> F, la relation f(x)=f(y) est une relation d'équivalence sur E.
%% 4) La relation << est dans la même classe que >> est une relation
%%    d'équivalence sur l'ensemble des élèves de Janson de Sailly.

\begin{definition}[utile=-3]
Soit $\mathcal{R}$ une relation d'équivalence sur $E$ et $x\in E$. On appelle
\emph{classe d'équivalence de $x$} et on note ${\rm Cl}(x)$ l'ensemble des
éléments de $E$ en relation avec $x$
\[{\rm Cl}(x)\defeq\enstq{y\in E}{x\mathcal{R}y}.\]
On dit qu'une partie $A$ de $E$ est \emph{une classe d'équivalence} lorsqu'il existe
$x\in E$ tel que $A={\rm Cl}(x)$.
\end{definition}

\begin{remarqueUnique}
\remarque Si $x,y\in E$, alors
  \[{\rm Cl}(x)={\rm Cl}(y) \quad\ssi\quad x \mathcal{R} y.\]
\end{remarqueUnique}

% Soit $E$ un ensemble et $\p{A_i}_{i\in I}$ une famille de parties de $E$.
% \begin{itemize}
% \item On dit que $\p{A_i}_{i\in I}$ est un \emph{recouvrement disjoint} de $E$ lorsque
% \[\cro{\forall i,j\in I \qsep A_i\cap A_j\neq\emptyset
%   \implique i=j}  \et \bigcup\limits_{i\in I} A_i=E.\]
% \item On dit que $\p{A_i}_{i\in I}$ est une \emph{partition} de $E$ lorsque c'est un recouvrement disjoint et
%   que
%   \[\forall i\in I\qsep A_i\neq\emptyset.\]
% \end{itemize}

\begin{proposition}[utile=-3]
Soit $\mathcal{R}$ une relation d'équivalence sur $E$. Alors, l'ensemble des
classes d'équivalence réalise une partition de $E$.
\end{proposition}

\begin{exoUnique}
\exo Soit $n\in\Ns$. Déterminer le nombre de classes d'équivalence sur $\Z$
  pour la relation de congruence modulo $n$.
\end{exoUnique}
  
  

\section{L'ensemble des entiers naturels}

Dans ce cours, nous ne chercherons pas à construire l'ensemble des entiers naturels.
Nous nous limiterons à la définition intuitive suivante.
\[\N\defeq\ens{0,1,2,3,\ldots}\]
Nous supposerons aussi définies les opérations usuelles $+$ et $\times$ ainsi que la
relation d'ordre totale $\leq$. Nous admettrons enfin la proposition suivante.

\begin{proposition}[utile=-3]
Toute partie non vide de $\N$ admet un plus petit élément.
\end{proposition}
  
\begin{proposition}[utile=-3]
Toute partie non vide majorée de $\N$ admet un plus grand élément.
\end{proposition}

\subsection{Récurrence}


\begin{preuve}
Admettons le premier point pour démontrer le second. Soit $A$ une partie non vide majorée de $\N$ et $B$ l'ensemble de ses majorants ($B\neq \emptyset$). $B$ admet donc un ppe. Notons-le $m$. Montrons que $m\in A$.

Supposons un instant que $m\notin A$. Comme $m$ est un majorant de $A$, on a $\forall x \in A, x<m$. Ainsi, $m>0$, sinon $A$ serait vide. Mais alors, $\forall x \in A, x\leq m-1$ donc $m-1$ est un majorant de $A$. CONTRADICTION. Donc $m$ est le pge de $A$.
\end{preuve}

\begin{proposition}[utile=-3, nom={Principe de récurrence}]
Soit $A$ une partie de $\N$ telle que
\begin{itemize}
\item $0\in A$,
\item $\forall n\in\N\qsep n\in A \implique n+1 \in A.$
\end{itemize}
Alors $A=\N$.
\end{proposition}

\begin{preuve}
Procédons par l'absurde. Supposons le contraire : $\exists n\in \N, n\notin A$.
Posons $B=\ens{n\in \N | n\notin A}$. $B$ est non vide donc admet un ppe. Notons-le $p$. $p\neq 0$ car $0\notin B$. Posons $q=p-1 \in \N$. $q\in A$ car il ne peut pas être dans $B$. Mais alors, on conclut d'après la deuxième hypothèse.

\end{preuve}

\begin{remarques}
\remarque Cette proposition est au cœur du principe de récurrence. Si $\mathcal{H}$ est un prédicat sur $\N$ tel que
\begin{itemize}
\item $\mathcal{H}_0$ est vraie,
\item $\forall n\in\N\qsep \mathcal{H}_n \implique \mathcal{H}_{n+1}$,
\end{itemize}
alors $\mathcal{H}_n$ est vraie pour tout $n\in\N$. Il suffit pour démontrer cela d'appliquer la proposition précédente à \[A\defeq\enstq{n\in\N}{\text{$\mathcal{H}_n$ est vraie}}.\]
\remarque Le principe de récurrence double est une conséquence du principe de récurrence. En effet, si $\mathcal{H}$ est un prédicat sur $\N$ tel que
\begin{itemize}
\item $\mathcal{H}_0$ et $\mathcal{H}_1$ sont vraies,
\item $\forall n\in\N\qsep \cro{\mathcal{H}_n \et \mathcal{H}_{n+1}}\implique \mathcal{H}_{n+2}$,
\end{itemize}
alors $\mathcal{H}_n$ est vraie pour tout $n\in\N$. Il suffit pour cela de remarquer que le prédicat $\mathcal{P}$ défini sur $\N$ par \[\forall n\in\N\qsep \mathcal{P}_n\defeq\text{\og} \mathcal{H}_n\text{ et }\mathcal{H}_{n+1}\text{ sont vraies} \text{\fg}\]
vérifie le principe de récurrence.
\remarque De même, le principe de récurrence forte est une conséquence du principe de récurrence.
  En effet, si $\mathcal{H}$ est un prédicat sur $\N$ tel que
\begin{itemize}
\item $\mathcal{H}_0$ est vraie,
\item $\forall n\in\N\qsep \cro{\mathcal{H}_0 \et\ldots\et \mathcal{H}_n}\implique \mathcal{H}_{n+1}$,
\end{itemize}
  alors $\mathcal{H}_n$ est vraie pour tout $n\in\N$. Il suffit pour cela de remarquer que le prédicat $\mathcal{P}$ défini sur $\N$ par 
\[\forall n\in\N\qsep \mathcal{P}_n\defeq\text{\og} \mathcal{H}_0, \mathcal{H}_1, \ldots, \mathcal{H}_{n}\text{ sont vraies} \text{\fg}.\]
vérifie le principe de récurrence.
\end{remarques}

\begin{exos}
\exo Montrer que pour tout $n\in\N$, $4^n+2$ est un multiple de 3.
\exo Soit $(u_n)$ la suite définie par
  \[u_0\defeq 1\qsep u_1\defeq 1, \et \forall n\in\N\qsep u_{n+2}\defeq u_{n+1}+\frac{2}{n+2}u_n.\]
  Montrer que pour tout $n\in\Ns$, $u_n\leq n^2$.
\exo
  % On admet que tout entier supérieur ou égal à 2 admet un diviseur premier.
  Montrer que tout entier $n\in\Ns$ s'écrit comme produit
  de nombres premiers.
% \exo Soit $f:\N\to\N$ une fonction surjective telle que
%   \[\forall n\in\N\qsep f(n)\geq n.\]
%   Montrer que $f=\id$.
\end{exos}

% \begin{exoUnique}
% \exo Montrer que
%   \[\forall n\geq 2\qsep \frac{1}{\sqrt{4n+1}}\leq\frac{1\times 3\times 5\times\cdots\times(2n-1)}{2\times 4\times 6\times\cdots\times (2n)}\leq \frac{1}{\sqrt{3n+1}}.\]
% \end{exoUnique}

% \begin{sol}
% Récurrence simple et on veut vérifier si $\dfrac{1}{\sqrt{4n+5}}\leq \dfrac{1}{\sqrt{4n+1}}\dfrac{2n+1}{2n+2}$. On passe tout au numérateur, on élève au carré et on développe et ça marche. Même chose de l'autre côté.
% \end{sol}

\subsection{Définition par récurrence}

\begin{proposition}
Soit $E$ un ensemble, $f\in\mathcal{F}(E,E)$  et $x\in E$. Alors, il existe une unique
suite $(u_n)$ d'éléments de $E$ telle que
\[u_0=x \quad\et\quad \forall n\in\N\qsep u_{n+1}=f(u_n).\]
\end{proposition}

\begin{remarqueUnique}
\remarque Il arrive qu'au lieu d'avoir une fonction $f\in\mathcal{F}(E,E)$,
  on ait une fonction $f\in\mathcal{F}(A, E)$ où $A$ est une partie de $E$. Si $x\in A$,
  et que l'on souhaite prouver l'existence d'une unique suite $(u_n)$ telle que
  \[u_0=x \quad\et\quad \forall n\in\N\qsep u_{n+1}=f(u_n),\]
  nous sommes face à un problème bien plus délicat. En effet, l'existence d'une
  telle suite n'est pas garantie. 
  % Si on trouve une partie $B$ de $A$ contenant $x$ et stable par $f$, il suffit
  % d'appliquer la proposition précédente à l'application $f$ induite à $B$.
  % Mais sans une telle partie $B$, il est impossible de prouver l'existence d'une telle suite.
  Par exemple, il n'existe pas de suite $(u_n)$ telle que
  \[u_0=\frac{1}{2} \quad\et\quad \forall n\in\N\qsep u_{n+1}=\frac{3 u_n-2}{u_n-1}.\]
  En effet, si tel était le cas,
  on aurait $u_1=1$ donc $u_2=(3u_1-2)/(u_1-1)$ ne serait pas défini. On ne peut tout
  simplement pas appliquer la proposition précédente, car la fonction
  \[\dspappli{f}{\R\setminus\ens{1}}{\R}{x}{\frac{3x-2}{x-1}}\]
  n'est pas définie sur $\R$ mais sur $\R\setminus\ens{1}$.
\end{remarqueUnique} 

\begin{definition}
Soit $E$ un ensemble, $A$ une partie de $E$ et $f\in\mathcal{F}(A, E)$. On dit qu'une
partie $B$ de $A$ est \emph{stable} par $f$ lorsque
\[\forall x\in B\qsep f(x)\in B.\]
\end{definition}

\begin{remarques}
\remarque Si $B$ est stable par $f$, il est possible de considérer la restriction de
  $f$ à $B$, corestreinte à $B$. On parle alors d'application \emph{induite} à $B$.
\remarque Si $A$ est une partie de $E$, $f\in\mathcal{F}(A, E)$ et $x\in A$,
  pour prouver l'existence d'une unique suite $(u_n)$ telle que
\[u_0=x \quad\et\quad \forall n\in\N\qsep u_{n+1}=f(u_n),\]
  il suffit de trouver une partie $B$ de $A$, stable par $f$ et telle que
  $x\in B$.
  % Par contre, si $x\in\interfo{2}{+\infty}$, il existe
  % bien une et une seule suite $(u_n)$ telle que
  % \[u_0=x \quad\et\quad \forall n\in\N\qsep u_{n+1}=\frac{3x-2}{x-1}\]
  % car~: $\forall x\geq 2\qsep f(x)\geq 2$. Autrement dit, $B\defeq \interfo{2}{+\infty}$ est stable par $f$.
\end{remarques}

\begin{exoUnique}
\exo Soit $x\in\interfo{2}{+\infty}$. Montrer qu'il existe une unique suite $(u_n)$ telle que
  \[u_0= x \quad\et\quad \forall n\in\N\qsep u_{n+1}= \frac{3 u_n-2}{u_n-1}.\]
\end{exoUnique}



% \subsection{Ensembles finis}

% \begin{definition}[utile=-3]
% On dit qu'un ensemble $E$ est \emph{fini} lorsqu'il existe $n\in\N$ tel que $\intere{1}{n}$ et $E$ sont en bijection. Sinon, on dit qu'il est \emph{infini}.
% \end{definition}

% \begin{proposition}[utile=-3]
% Si $E$ est un ensemble fini alors il existe un unique $n\in\N$ pour lequel il existe une bijection de $\intere{1}{n}$ dans $E$. Cet entier est appelé \emph{cardinal} de $E$ et noté  $\card(E)$, $|E|$, ou $\# E$.
% \end{proposition}

% \begin{preuve}
% S'il en existe un autre, cela donnera une bijection de $\intere{1}{n}$ dans $\intere{1}{m}$. Et si $n\neq m$, c'est impossible. En effet, démontrons cela par récurrence sur $n\in \N$ avec $\mathcal{H}n$ : "$\forall m\in \Ns$, s'il y a une bijection de $\intere{1}{n}$ dans $\intere{1}{m}$ alors $m=n$."

% On initialise à $1$ et pour l'hérédité, on considère la restriction où on enlève $n+1$ au départ et $f(n+1)$ à l'arrivée.
% \end{preuve}

% \begin{remarqueUnique}
% \remarque Soit $a,b\in\Z$ tels que $a\leq b$. Alors $\intere{a}{b}$ est fini et $\card(\intere{a}{b})=b-a+1$.
% \end{remarqueUnique}

% \begin{proposition}[utile=3] Soient $E$ et $F$ deux ensembles finis.
% \begin{itemize}
% \item Il existe une injection de $E$ dans $F$ si et seulement si $\card (E)\leq\card (F)$.
% \item Il existe une surjection de $E$ dans $F$ si et seulement si $\card (E)\geq\card (F)$.
% \item Il existe une bijection de $E$ dans $F$ si et seulement si $\card (E)=\card (F)$.
% \end{itemize}
% \end{proposition}

% \begin{preuve}
% Il suffit de considérer $f(E)$ dans un sens et de la construire dans l'autre.
% \end{preuve}


% \begin{remarqueUnique}
% \remarque Du premier point découle le principe des tiroirs~: si $f$ est une application d'un ensemble fini $E$ dans un ensemble fini $F$ de cardinal strictement inférieur, alors $f$ n'est pas injective. Par exemple, si on possède 4 T-shirts que l'on a rangés dans une commode possédant 3 tiroirs, alors deux T-shirts se trouvent dans un même tiroir.
% \end{remarqueUnique}

% \begin{exoUnique}
% \exo On estime que le nombre de cheveux qu'une personne a sur la tête est inférieur à $1\ 000\ 000$. Montrer qu'il existe deux français qui ont autant de cheveux sur la tête.
% \end{exoUnique}

% \begin{proposition}[utile=3]
% Soit $E$ un ensemble fini et $F \subset E$. Alors $F$ est fini et $\card (F)
% \leq \card (E)$. De plus, il y a égalité si et seulement si $E = F$.
% \end{proposition}

% \begin{preuve}
% On peut le faire par récurrence sur le cardinal de $E$. Dans l'hérédité, ou bien $F=E$ ou bien on prend $a\in E\setminus F$ et on raisonne sur $E\setminus\ens{a}$.
% \end{preuve}

% \begin{remarqueUnique}
% \remarque En particulier, un ensemble fini ne peut pas être en bijection avec une partie stricte de lui-même. Remarquons que c'est faux pour les ensembles infinis. Par exemple
% \[\dspappli{f}{\N}{\Ns}{n}{n+1}\]
% est une bijection alors que $\Ns$ est une partie stricte de $\N$.
% \end{remarqueUnique}

% \begin{proposition}[utile=3]
% Soit $E$ et $F$ deux ensembles finis et $f$ une application de $E$ dans $F$.
% \begin{itemize}
% \item Si $f$ est injective et $\card(E)=\card(F)$, alors $f$ est bijective.
% \item Si $f$ est surjective et $\card(E)=\card(F)$, alors $f$ est bijective.
% \end{itemize}
% Autrement dit, si $\card(E)=\card(F)$
% \[\text{$f$ est injective} \quad\ssi\quad \text{$f$ est bijective}
%    \quad\ssi\quad \text{$f$ est surjective}.\]
% \end{proposition}

% \begin{preuve}
% On utilise $f$ injective ssi $|f(E)|=|E|$ et $f$ surjective ssi $|f(E)|=|F|$.
% \end{preuve}

%END_BOOK

\end{document}
\section{Exercices}
\setcounter{numeroexercice}{1}
\documentclass{magnolia}

\magtex{tex_driver={pdftex}}
\magfiche{document_nom={Logique, Ensembles},
          auteur_nom={François Fayard},
          auteur_mail={fayard.prof@gmail.com}}
\magexos{exos_matiere={maths},
         exos_niveau={mpsi},
         exos_chapitre_numero={2},
         exos_theme={Logique, Ensembles}}
\magmisenpage{}
\maglieudiff{}
\magprocess

\begin{document}

%BEGIN_BOOK
\magsection{Éléments de logique}
\magsubsection{Assertion, prédicat}

% \exercice{nom={Commentaires sportifs}}
% En utilisant le fait que $\non(\non A)$ et $A$ ont même valeur de vérité, simplifier
% la phrase
% \begin{center}
% \og Il ne se trouvera aucun sportif pour nier que le contraire n'eût été
% immérité. \fg.
% \end{center}
% \begin{sol}
% Il y aura un sportif pour dire que ce fut mérité.
% \end{sol}

\magsubsection{Implication, équivalence}

% \exercice{nom={Liens logiques}}
% Déterminer les liens d'implication et d'équivalence entre les assertions
% suivantes (on ne cherchera pas à savoir si les assertions suivantes sont
% vraies ou fausses)~:
% \begin{questions}
% \question Tous les australopithèques sont mortels.
% \question Tous les australopithèques sont immortels.
% \question Aucun australopithèque n'est mortel.
% \question Aucun australopithèque n'est immortel.
% \question Il existe des australopithèques immortels.
% \question Il existe des australopithèques mortels.
% \end{questions}
\exercice{nom={Phrases mathématiques}}
On considère les propositions
\begin{enumerate}
\item $\forall x\in\R \qsep \cro{\p{\forall y\in\R \qsep xy=0} \implique x=0}$
\item $\forall x\in\R \qsep \forall y\in\R \qsep \cro{xy=0 \implique x=0}.$
\end{enumerate}
Sont-elles vraies ou fausses~? Bien entendu, on justifiera.

\exercice{nom={Quantificateurs}}
Écrire avec des quantificateurs les assertions suivantes, où $(u_n)$ désigne une
suite réelle et $f$ désigne une fonction de $\R$ dans $\R$.
\begin{questions}
\question La suite $u$ est majorée.
\question La suite $u$ n'est pas majorée.
\question La fonction $f$ est nulle.
\question La fonction $f$ n'est pas nulle.
\question La fonction $f$ n'est pas croissante.
\question La fonction $f$ est périodique.
\question La fonction $f$ n'est pas périodique.
\question La fonction $f$ n'est pas paire.
\question La fonction $f$ n'est pas bornée.
\end{questions}
\begin{sol}
$\quad$
\begin{questions}
\question $\exists M\in\R \qsep \forall n\in\N \qsep u_n\leq M$
\question $\forall M\in\R \qsep \exists n\in\N \quad u_n > M$
\question $\forall x\in\R \quad f(x)=0$
\question $\exists x\in\R \quad f(x)\neq 0$
\question $\non\p{\forall x\in\R \quad \forall y\in\R \quad x\leq y \implique
  f(x)\leq f(y)}$\\
  $\exists x\in\R \quad \exists y\in\R \quad x\leq y \et f(x)>f(y)$
\question $\exists T\in\Rs \quad \forall x\in\R \quad f\p{x+T}=f(x)$\\
  $\forall T\in\Rs \quad \exists x\in\R \quad f\p{x+T}\neq f(x)$
\question $\non\p{\forall x\in\R \quad f\p{-x}=f(x)}$\\
  $\exists x\in\R \quad f\p{-x}\neq f(x)$
\question $\non\p{\exists M\in\RP \quad \forall x\in\R \quad
  \abs{f(x)}\leq M}$\\
  $\forall M\in\RP \quad \exists x\in\R \quad \abs{f(x)}>M$
\end{questions}
\end{sol}

\exercice{nom={Autour des suites}}
Soit $\p{u_n}$ une suite réelle. Montrer l'équivalence des deux propriétés
suivantes
\begin{enumerate}
\item $\forall A\in\R \qsep \exists n\in\N \qsep u_n \geq A.$
\item $\forall A\in\R \qsep \exists n\in\N \qsep u_n > A.$
\end{enumerate}
Que signifient ces énoncés~?

\magsection{Ensemble}
\magsubsection{Ensemble, élément}
\magsubsection{Opérations élémentaires}

\exercice{nom={Ensembles}}
Soit $A$ et $B$ deux parties d'un ensemble $E$.
\begin{questions}
\question Montrer que
   \[A\cap B=A\cup B \implique A=B.\]
\question Montrer que
  \[A\setminus\p{A\setminus B}=A\cap B.\]
\end{questions}

\exercice{nom={Ensembles}}
Soit $E$ un ensemble et $A,B,C\in\mathcal{P}(E)$. Montrer que les 3 assertions suivantes
sont deux à deux équivalentes.
\begin{itemize}
\item $A\setminus B\subset C$.
\item $A\setminus C\subset B$. 
\item $A\subset B\cup C$.
\end{itemize}

% \exercice{nom={Produit cartésien}}
% Soit $E,F,G,H$ quatre parties d'un même ensemble. Comparer les ensembles
% \[(E\times F)\cap(G\times H) \et (E\cap G)\times(F\cap H).\] On fera un dessin,
% puis une preuve.

\exercice{nom={Équation ensembliste}}
Soit $A,B$ deux parties de $E$.
\begin{questions}
\question On souhaite résoudre l'équation $A\cup X=B$ pour $X\in\parties{E}$.
  \begin{questions}
  \question Montrer que si l'équation admet au moins une solution, alors $A\subset B$.
  \question Montrer que si $A\subset B$, l'ensemble des solutions de l'équation est
    \[\ensim{(B\setminus A)\cup T}{T\in\mathcal{P}(A)}.\]
  \question Conclure.
  \end{questions}
\question On souhaite résoudre l'équation $A\cap X=B$ pour $X\in\parties{E}$.
  \begin{questions}
  \question Montrer que si l'équation admet au moins une solution, alors $B\subset A$.
  \question Montrer que si $B\subset A$, l'ensemble des solutions de l'équation est
    \[\ensim{B\cup T}{T\in\mathcal{P}(\bar{A})}.\]
  \question Conclure.
  \end{questions}
\end{questions}


\exercice{nom={Équation ensembliste}}
Soit $E$ un ensemble et $A$ et $B$ deux parties de $E$. Discuter et résoudre
l'équation
\[(A\cap X)\cup(B\cap X^c)=\varnothing.\]

\magsection{Application}
\magsubsection{Définition, exemples}



\exercice{nom={Différence symétrique}}
Soit $E$ un ensemble. Quels que soient $A,B\in\mathcal{P}(E)$, on définit
la différence symétrique $A\Delta B$ entre $A$ et $B$ par
\[A\Delta B\defeq(A\cup B)\setminus(A\cap B).\]
\begin{questions}
\question Déterminer $A\Delta B$ dans les deux exemples suivants.
\begin{itemize}
\item $E\defeq\ens{1,2,3,4}$, $A\defeq\ens{1, 2}$, $B\defeq\ens{1, 3}$.
\item $E\defeq\R$, $A\defeq\interof{-\infty}{2}$, $B\defeq\interfo{1}{+\infty}$.
\end{itemize}
\question Montrer que
\[\forall A,B\in\mathcal{P}(E)\qsep A\Delta B=\p{A\cap\overline{B}}\cup\p{B\cap\overline{A}}.\]
\question Montrer que pour tout $(A,B)\in\mathcal{P}(E)^2$
\[\mathds{1}_{A\Delta B}=\mathds{1}_A + \mathds{1}_B - 2 \mathds{1}_A \mathds{1}_B.\]
\question En déduire que la loi $\Delta$ est associative sur $\mathcal{P}(E)$,
  c'est-à-dire que
	\[\forall A,B,C\in\mathcal{P}(E)\qsep A\Delta(B\Delta C)=(A\Delta B)\Delta C.\]
\end{questions}
\magsubsection{Application injective, surjective, bijective}

\exercice{nom={Fonction de $\Cs$ dans $\C$}}
Soit $f$ la fonction
\[\dspappli{f}{\Cs}{\C}{z}{z-\frac{1}{z}}\]
\begin{questions}
\question Montrer que $f$ est surjective mais non injective.
\question Déterminer $f^\applirec(\ii\R)$.
\question Déterminer $f\p{\U}$.
\end{questions}
\begin{sol}
$\quad$
\begin{questions}
\question La fonction n'est pas injective car $f(1)=f(-1)=0$.
\question $f(\U)=\enstq{i\alpha}{\alpha\in\interf{-2}{2}}$.
\question On a $f(z)\in i\R \ssi (z+\conj{z})(z\conj{z}-1)=0$. Donc
  $f^{-1}(J)=\U\cup i\Rs$.
\end{questions}
\end{sol}

\exercice{nom={Fonction de $\C^2$ dans $\C^2$}}
Soit $f$ la fonction
\[\dspappli{f}{\C^2}{\C^2}{\p{u,v}}{\p{u^2+v^2,uv}}\]
\begin{questions}
\question $f$ est-elle injective~?
\question $f$ est-elle surjective~?
\question Déterminer les antécédents de $\p{3-2\ii,3+\ii}$ par $f$.
\end{questions}
\begin{sol}
$\quad$
\begin{questions}
\question Non, car $f\p{0,1}=f\p{1,0}$.
\question Oui, car les équations de degré 2 admettent toujours une solution
  sur $\C$.
\question $\p{u,v}=\p{-1-i,-2+i}$, $\p{u,v}=\p{1+i,2-i}$ et deux autres
  solutions obtenues par échange de $u$ et $v$.
\end{questions}
\end{sol}


\exercice{nom={Injection, surjection}}
Soit $E$ et $F$ deux ensembles, $f:E\to F$ et $g:F\to E$.
\begin{questions}
\question Montrer que si $f\circ g\circ f=f$ et que $f$ est injective, alors $g$ est surjective.
\question Montrer que si $g\circ f\circ g=g$ et que $g$ est surjective, alors $f$ est
injective.
\end{questions}

\exercice{nom={Image directe, image réciproque}}
Soit $f:E\to F$, $A$ une partie de $E$ et $B$ une partie de $F$. Montrer que
\[f\p{A\cap f^\applirec (B)}=f(A)\cap B.\]
\begin{sol}
Soit $f$ une application de $E$ dans $F$, $A$ une partie de $E$ et $B$ une partie de $F$. Montrons que 
\[f\p{A\cap f^\applirec (B)}=f(A)\cap B.\]
On procède par double inclusion.
\begin{itemize}
\item $f(A\cap f^\applirec (B))\subset f(A)\cap B$. Soit $y\in f(A\cap f^\applirec (B))$. Montrons que $y\in f(A)\cap B$. Il existe $x\in A\cap f^\applirec (B)$ tel que $y=f(x)$. Puisque $x\in A$, on en déduit que $f(x)\in f(A)$. Puisque $x\in f^\applirec (B)$, on en déduit que $f(x)\in B$. Donc $y\in f(A)\cap B$.
\item $f(A)\cap B\subset f(A\cap f^\applirec (B))$. Soit $y\in f(A)\cap B$. Montrons que $y\in f(A\cap f^\applirec (B))$. Puisque $y\in f(A)$, il existe $x\in A$ tel que $y=f(x)$. Puisque $f(x)=y\in B$, on en déduit que $x\in f^\applirec(B)$. Donc $x\in A\cap f^\applirec(B)$. Donc $y\in f(A\cap f^\applirec(B))$.
\end{itemize}
Par double inclusion, $f(A\cap f^\applirec (B))= f(A)\cap B$.
\end{sol}


\exercice{nom={Une bijection de $\interf{0}{1}$ dans $\intero{0}{1}$}}
Montrer que la fonction $f:\interf{0}{1} \rightarrow \intero{0}{1}$ définie
par
\[\forall x\in\interf{0}{1} \qsep f(x)\defeq
  \begin{cases}
  \frac{1}{2} & \text{si $x=0$}\\
  \frac{1}{n+2} & \text{si $x=\frac{1}{n}$ avec $n\in\Ns$}\\
  x  &\text{sinon}
  \end{cases}\]
est une bijection.
\begin{sol}
Le plus simple est de donner la bijection réciproque
\[\forall x\in\intero{0}{1} \quad g(x)=
  \begin{cases}
  0 & \text{si $x=\frac{1}{2}$}\\
  \frac{1}{n-2} & \text{si $x=\frac{1}{n}$ avec $n\geq 3$}\\
  x  &\text{sinon}
  \end{cases}\]
\end{sol}

\exercice{nom={Application ensembliste}}
Soit $A$ et $B$ deux parties de $E$ et
\[\dspappli{f}{\parties{E}}{\parties{A}\times\parties{B}}{X}{\p{X\cap A,X\cap B}.}\]
\begin{questions}
\question Montrer que $f$ est injective si et seulement si $A\cup B=E$.
\question Montrer que $f$ est surjective si et seulement si $A\cap B=\emptyset$.
\question On suppose que $f$ est bijective. Calculer $f^{-1}$.
\end{questions}
\begin{sol}
$\quad$
\begin{questions}
\question Montrons que $f$ est injective si et seulement si $A\cup B=E$.
  \begin{itemize}
  \item Supposons que $A\cup B=E$. Montrons que $f$ est injective. Soit $X_1,X_2\in\parties{E}$ tels que $f(X_1)=f(X_2)$. Montrons que $X_1=X_2$. On a $(X_1\cap A, X_1\cap B)=(X_2\cap A, X_2\cap B)$, donc $X_1\cap A=X_2\cap A$ et $X_1\cap B=X_2\cap B$. Donc
  \begin{eqnarray*}
  X_1&=&X_1\cap E=X_1\cap(A\cup B)=(X_1\cap A)\cup(X_1\cap B)\\
     &=&(X_2\cap A)\cup(X_2\cap B)=X_2\cap(A\cup B)\\
     &=&X_2\cap E=X_2.
  \end{eqnarray*}
  \item Supposons que $f$ est injective. Montrons que $A\cup B=E$. On a
  \[f(A\cup B)=((A\cup B)\cap A, (A\cup B)\cap B)=(A,B)\]
  et
  \[f(E)=(E\cap A, E\cap B)=(A, B)\]
  donc $f(E)=f(A\cup B)$. Or $f$ est injective, donc $E=A\cup B$.
  \end{itemize}
\question Montrons que $f$ est surjective si et seulement si $A\cap B=\emptyset$.
\begin{itemize}
\item Supposons de $A\cap B=\emptyset$. Montrons que $f$ est surjective. Soit $(U,V)\in\parties{A}\times\parties{B}$. On pose $X\defeq U\cup V$. Alors
\[f(X)=(X\cap A, X\cap B).\]
On a $X\cap A=(U\cup V)\cap A=(U\cap A)\cup(V\cap A)$. Or $U\cap A=U$ car $U\subset A$ et $V\cap A=\emptyset$ car $V\subset B$ donc $V\cap A\subset B\cap A=\emptyset$. Donc $X\cap A=U\cup\emptyset=U$. De même $X\cap B=V$. Donc $f(X)=(U, V)$.
\item Supposons que $f$ est surjective. Montrons que $A\cap B=\emptyset$. On pose $Y\defeq (A\cap B,\emptyset)\in\parties{A}\times\parties{B}$. Puisque $f$ est surjective, il existe une partie $X$ de $E$ telle que $f(X)=(A\cap B,\emptyset)$. Alors $X\cap A=A\cap B$ et $X\cap B=\emptyset$. Alors
\[A\cap B=(A\cap B)\cap B=(X\cap A)\cap B=(X\cap B)\cap A=\emptyset\cap A=\emptyset.\]
\end{itemize}
\question $f$ est bijective si et seulement si elle est injective et surjective, c'est-à-dire si et seulement si $A\cap B=\emptyset$ et $A\cup B=E$, c'est-à-dire si et seulement si $B=\overline{A}$. Supposons que $f$ est bijective. Alors, si $(U,V)\in\parties{A}\times\parties{B}$, on a vu dans la seconde question que $X\defeq U\cup V$ est un antécédent de $(U,V)$. Or $f^{-1}$ est l'application qui à $(U,V)\in\parties{A}\times\parties{B}$ associe son unique antécédent par $f$. Donc
\[\forall (U,V)\in\parties{A}\times\parties{B}\qsep f^{-1}(U,V)=U\cup V.\]
\end{questions}
\end{sol}

\exercice{nom={Application fonctionnelle}}
Soit $f$ une application de $E$ dans $E$ et $\phi$ l'application
\[\dspappli{\phi}{\mathcal{F}\p{E,E}}{\mathcal{F}\p{E,E}}{g}{g\circ f.}\]
Montrer que $\phi$ est bijective si et seulement si $f$ l'est.
\begin{sol}
On suppose que $f$ est bijective. Montrons que $\phi$ est bijective. On définit
\[\dspappli{\psi}{\mathcal{F}\p{E,E}}{\mathcal{F}\p{E,E}}{g}{g\circ f^{-1}.}\]
Alors, quel que soit $g\in\mathcal{F}\p{E,E}$
\[(\phi\circ\psi)(g)=\phi(\psi(g))=\phi(g\circ f^{-1})=g\circ f^{-1}\circ f=g\circ\id=g\]
\[(\psi\circ\phi)(g)=\psi(\phi(g))=\psi(g\circ f)=g\circ f\circ f^{-1}=g\circ\id=g.\]
Donc $\phi\circ\psi=\id$ et $\psi\circ\phi=\id$. Donc $\phi$ est bijective.\\

Réciproquement, supposons que $\phi$ est bijective. Montrons que $f$ est bijective. On commence par prouver que $f$ est injective. Puisque $\phi$ est surjective, il existe $g\in\mathcal{F}(E,E)$ tel que $\id=g\circ f$. Soit $x_1,x_2\in E$ tels que $f(x_1)=f(x_2)$. Alors $g(f(x_1))=g(f(x_2))$, donc $x_1=x_2$.\\
Montrons désormais que $f$ est surjective, c'est-à-dire que $\im f=E$. Si $E$ est vide ou ne possède qu'un élément, c'est évident. On suppose donc que $E$ possède au moins deux éléments. Pour montrer que $\im f=E$, on raisonne par l'absurde et on suppose que $\im f$ est strictement inclus dans $E$. Puisque $E$ possède au moins 2 éléments, il existe $y,y'\in E$ tels que $y\neq y'$. On définit alors les fonctions $g_1$ et $g_2$ de $E$ dans $E$ par
\[\forall x\in E\qsep g_1(x)=\begin{cases}
x &\text{si $x\in\im f$}\\
y &\text{sinon}
\end{cases} \et
g_2(x)=\begin{cases}
x &\text{si $x\in\im f$}\\
y' &\text{sinon.}
\end{cases}
\]
Alors $g_1\neq g_2$ car $\im f$ est strictement inclus dans $E$ et $y\neq y'$. Pourtant
\[\forall x\in E\qsep (g_1\circ f)(x)=(g_2\circ f)(x)\]
donc $\phi(g_1)=\phi(g_2)$. C'est absurde, car $\phi$ est injective.\\
Voici une autre preuve de la surjectivité de $f$, plus courte mais surement moins naturelle.
Montrons désormais que $f$ est surjective. On a $f\circ g\circ f=f\circ(g\circ f)=f\circ\id=f=\id\circ f$. Donc $\phi(f\circ g)=\phi(\id)$. Puisque $\phi$ est injective, $f\circ g=\id$. On en déduit que $f$ est surjective. En effet, soit $y\in E$. Alors, $(f\circ g)(y)=y$ donc $f(g(y))=y$. En posant $x\defeq g(y)$, on en déduit que $f(x)=y$.\\
\end{sol}

\exercice{nom={\nom{Haskell Curry}}}
Soit $A,B,C$ trois ensembles. Montrer que l'application
\[\dspappli{\phi}{\mathcal{F}(A\times B, C)}{\mathcal{F}(A,\mathcal{F}(B, C))}{f}{
	\dspappli{\phi(f)}{A}{\mathcal{F}(B, C)}{a}{
		\dspappli{[\phi(f)](a)}{B}{C}{b}{f(a,b)}
	}}\]
est bijective.

\exercice{nom={Il n'y a pas de surjection de $E$ dans $\mathcal{P}(E)$}}
Montrons que si $E$ est un ensemble, il n'existe pas de surjection de $E$ dans $\mathcal{P}(E)$. On raisonne par l'absurde et on suppose qu'il existe une surjection $\phi$ de $E$ dans $\mathcal{P}(E)$.
Conclure a une absurdité en considérant
\[A\defeq\enstq{ x \in E}{x \not\in \phi(x)}.\]
\begin{sol}
Montrons que si $E$ est un ensemble, il n'existe pas de surjection de $E$ dans $\mathcal{P}(E)$. On raisonne par l'absurde et on suppose qu'il existe une surjection $\phi$ de $E$ dans $\mathcal{P}(E)$. On pose
\[A\defeq\enstq{ x \in E}{x \not\in \phi(x)}.\]
Alors $A\in\parties{E}$. Puisque $\phi$ est surjective, il existe $x\in E$ tel que $A=\phi(x)$.
\begin{itemize}
\item Supposons que $x\in A$. Alors $x\not\in \phi(x)=A$. C'est impossible.
\item Supposons que $x\not\in A$. Alors $x\not\in\phi(x)$, donc $x\in A$. C'est impossible.
\end{itemize}
On aboutit donc à une absurdité. Donc $\phi$ n'est pas surjective.
\end{sol}

\begin{sol}
$A\in \mathcal{P}(E)$. Donc il existe $y\in E$ tel que $A=\phi(y)$. Si $y\in A, y\notin \phi(y)=A$, absurde, donc $y\notin A=\phi(y)$ donc $y\in A$. CONTRADICTION.
\end{sol}


% \exercice{nom={Injection, surjection et composition}}
% Dans cet exercice, les applications sont toutes définies sur des ensembles tels que
% les compositions indiquées existent. Chaque question est indépendante.
% \begin{questions}
% \question Montrer que si $g\circ f$ est injective, alors $f$ est injective.
% \question Montrer que si $g\circ f$ est surjective, alors $g$ est surjective.
% \end{questions}

\exercice{nom={Composition, injection et surjection}}
\begin{questions}
\question Soit $A$, $B$, $C$, $D$ quatre ensembles, $f:A\rightarrow B$,
  $g:B\rightarrow C$, $h:C\rightarrow D$ trois applications telles que $g\circ f$
  et $h\circ g$ sont bijectives. Montrer que $f$, $g$ et $h$ sont bijectives.
\question Soit $X$, $Y$, $Z$ trois ensembles et  $f:X\rightarrow Y$,
  $g:Y\rightarrow Z$, $h:Z\rightarrow X$ trois applications. On forme les
  applications composées $h\circ g\circ f$, $g\circ f\circ h$,
  $f\circ h\circ g$. On suppose que deux d'entre elles sont surjectives et la
  troisième injective. Montrer qu'alors $f$, $g$ et $h$ sont bijectives.
\end{questions}
\begin{sol}
\begin{questions}
\question  Soit $A,B,C,D$ quatre ensembles, $f:A\rightarrow B$,
  $g:B\rightarrow C$, $h:C\rightarrow D$ trois applications telles que $g\circ f$
  et $h\circ g$ sont bijectives. Montrons que $f$, $g$ et $h$ sont bijectives. Commençons par montrer que $g$ est bijective.
  \begin{itemize}
  \item $h\circ g$ est bijective, donc injective. Montrons que $g$ est injective. Soit $b_1,b_1\in B$ tels que $g(b_1)=g(b_2)$. Alors $h(g(b_1))=h(g(b_2))$, donc $(h\circ g)(b_1)=(h\circ g)(b_2)$. Or $h\circ g$ est injective, donc $b_1=b_2$.
  \item $g\circ f$ est bijective, donc surjective. Montrons que $g$ est surjective. Soit $c\in C$. Puisque $g\circ f$ est surjective, il existe $a\in A$ tel que $c=(g\circ f)(a)$. Donc $c=g(f(a))$. On pose $b\defeq f(a)$. Alors $c=g(b)$.
  \end{itemize}
  Donc $g$ est bijective. Donc $f=g^{-1} \circ (g\circ f)$ est bijective comme composée de deux bijections. De même $h=(h\circ g)\circ g^{-1}$ est bijective comme composée de deux bijections.
\question Soit $X$, $Y$, $Z$ trois ensembles et  $f:X\rightarrow Y$,
  $g:Y\rightarrow Z$, $h:Z\rightarrow X$ trois applications. On forme les
  applications composées $h\circ g\circ f$, $g\circ f\circ h$,
  $f\circ h\circ g$. On suppose que $h\circ g\circ f$ et $g\circ f\circ h$ sont surjectives et que $f\circ h\circ g$ est injective. Montrons que $f$, $g$ et $h$ sont bijectives. Puisque $g\circ f\circ h$ est surjective, on en déduit que $g$ est surjective. De même, $f\circ h\circ g$ est injective donc $g$ est injective. Donc $g$ est bijective. Donc $f\circ h=(f\circ h\circ g)\circ g^{-1}$ est injective comme composée de deux injections. Donc $h$ est injective. Or $h\circ g\circ f$ est surjective donc $h$ est surjective. Donc $h$ est bijective. Donc $f=g^{-1}\circ h^{-1}\circ(h\circ g\circ f)$ est surjective comme composée de surjections. De même $f=(f\circ h\circ g)\circ g^{-1} \circ h^{-1}$ est injective comme composée d'injections. Donc $f$ est bijective. Les autres cas se démontrent de la même manière.
\end{questions}
\end{sol}

\exercice{nom={Inversion à droite, à gauche d'une application}}
Soit $f$ une application de $A$ dans $B$.
\begin{questions}
\question
  \begin{questions}
  \question Montrer que $f$ est surjective si et seulement si il existe une
    application $g$ de $B$ dans $A$ telle que $f\circ g=\id_{B}$.
  \question Dans le cas où $f$ est surjective, montrer que $g$ est unique si
    et seulement si $f$ est bijective.
  \end{questions}
\question
  \begin{questions}
  \question Montrer que $f$ est injective si et seulement si il existe une
    application $g$ de $B$ dans $A$ telle que $g\circ f=\id_{A}$.
  \question Dans le cas où $f$ est injective, montrer que $g$ est unique si
    et seulement si $f$ est bijective.
\end{questions}
\end{questions}

\begin{sol}
Soit $f$ une application de $A$ dans $B$.
\begin{questions}
\question
  \begin{questions}
  \question 
  \begin{itemize}
  \item[$\bullet$]S'il existe une application $g$ de $B$ dans $A$ telle que $f\circ g=\id_{B}$. 
%  \textit{(Alors $f\circ g$ est surjective. On a déjà vu dans un exercice du cours que cela implique que $f$ est surjective. Redémontrons-le ici dans ce cas particulier.)}
   Soit donc $y\in B$, on a $y=(f\circ g) (y)=f(g(y))$ donc $y\in \im f$, ce qui signifie bien que $f$ est surjective.
  
  \item[$\bullet$]Réciproquement, si $f$ est surjective. Fixons $y\in B$, il existe $x$ dans $A$ tel que $f(x)=y$. Notons-le $g(y)$. On définit ainsi une fonction $g:B\to A$ vérifiant $f(g(y))=y, \forall y \in B$. On a donc bien $f\circ g=\id_{B}$.
  \end{itemize}
  \question On suppose dans cette question $f$ surjective.
  
  \begin{itemize}
  \item[$\bullet$] Supposons que $f$ est bijective et montrons que $g$ est unique.
  S'il existe $g$ et $h$ deux fonctions de $B$ dans $A$ telles que $f\circ g=f\circ h=\id_{B}$. En composant à gauche par $f^{-1}$, on obtient $g=h$. 
  
  \item[$\bullet$] Supposons que $g$ est unique et montrons que $f$ est bijective.
  Il reste à montrer que $f$ est injective.\\
   Procédons par l'absurde et supposons un instant que $f$ n'est pas injective. Soit alors $y\in B$ qui admet plusieurs antécédents par $f$. On a vu dans la construction de $g$ à la question précédente que $g$ envoie une image de $f$ sur un de ses antécédents. $g(y)$ est donc un antécédent de $y$. Notons $z\neq g(y)$ tel que $f(z)= y$. On peut alors définir $h$ qui est égale à $g$ sur $B\setminus\ens{y}$ et telle que $h(y)=z$.
  Ainsi, $\forall u \in B\setminus\ens{y}, (f\circ h)(u)=f(h(u))=f(g(u))=u$ et $(f\circ h)(y)=f(h(y))=f(z)=y$. Finalement, $f\circ h=\id_{B}$, ce qui contredit l'unicité de $g$.\\
  $f$ est donc bien injective et donc bijective.
  
  
  \end{itemize}
  \end{questions}
\question
  \begin{questions}
  \question \begin{itemize}
  \item[$\bullet$]S'il existe une application $g$ de $B$ dans $A$ telle que $g\circ f=\id_{A}$. 
  %\textit{(Alors $g\circ f$ est injective. On a déjà vu dans un exercice du cours que cela implique que $f$ est injective. Redémontrons-le ici dans ce cas particulier.)}
   Soit donc $x_1,x_2$ dans $A$ tels que $f(x_1)=f(x_2)$, alors on a $x_1=\id(x_1)=(g\circ f) (x_1)=g(f(x_1))=g(f(x_2))=(g\circ f) (x_2)=\id(x_2)=x_2$ ce qui signifie bien que $f$ est injective.\\
  
  \item[$\bullet$]Réciproquement, si $f$ est injective, considérons $y$ un élément de $f(A)$, $y$ admet dans $A$ un unique antécédent $x$ par $f$. Notons-le $g(y)$. Alors, on a $\forall x \in A$, $(g\circ f)(x)=g(f(x))=x$. $g$ est alors bien défini sur $f(A)$ et on peut la définir comme on le souhaite sur $B\setminus f(A)$.
  \end{itemize}
  \question On suppose dans cette question $f$ surjective.
  \begin{itemize}
  \item[$\bullet$]Supposons que $f$ est bijective et montrons que $g$ est unique.
  S'il existe $g$ et $h$ deux fonctions de $B$ dans $A$ telles que $g\circ f=h\circ f=\id_{A}$. En composant à droite par $f^{-1}$, on obtient $g=h$. 
  \item[$\bullet$]Supposons que $g$ est unique et montrons que $f$ est bijective.
  Il reste à montrer que $f$ est surjective.\\
  Procédons par l'absurde et supposons un instant que $f$ n'est pas surjective. Il existe donc $y\in B\setminus f(A)$. On peut alors définir $h$ qui est égale à $g$ sur $B\setminus\ens{y}$ et telle que $h(y)\neq g(y)$. On a bien $\forall x\in A$, $(h\circ f)(x)=h(f(x))=g(f(x))=x$ car $f(x)\neq y$. Finalement, $h\circ f=\id_{A}$, ce qui contredit l'unicité de $g$.\\
  $f$ est donc bien surjective et donc bijective.
  
  
  \end{itemize}
\end{questions}
\end{questions}
\end{sol}

\exercice{nom={Image directe et réciproque}}
Soit $f$ une application de $A$ dans $B$.
\begin{questions}
\question Montrer que $f$ est surjective si et seulement si quelle que soit la
  partie $Y$ de $B$, on a $f(f^\applirec(Y))=Y$.
\question Montrer que $f$ est injective si et seulement si quelle que soit la
  partie $X$ de $A$, on a $f^\applirec\p{f(X)}=X$.
\end{questions}
\begin{sol}
Soit $f$ une application de $A$ dans $B$.
\begin{questions}

  \question Montrons que $f$ est surjective si et seulement si, quelle que soit la partie $Y$ de $B$, on a $f(f^\applirec(Y))=Y$.
  \begin{itemize}
  \item On suppose que $f$ est surjective. Soit $Y$ une partie de $B$. Montrons que $f(f^\applirec(Y))=Y$. On procède par double inclusion.
    \begin{itemize}
    \item $f(f^\applirec(Y))\subset Y$. Soit $y\in f(f^\applirec(Y))$. Alors il existe $x\in f^\applirec(Y)$ tel que $y=f(x)$. Puisque $x\in f^\applirec(Y)$, alors $f(x)\in Y$. Donc $y \in Y$.
    \item $Y\subset f(f^\applirec(Y))$. Soit $y\in Y$. Puisque $f$ est surjective, il existe $x\in A$ tel que $y=f(x)$. Comme $f(x)\in Y$, on en déduit que $x\in f^\applirec(Y)$. Donc $y\in f(f^\applirec(Y))$.
    \end{itemize}
    Donc $Y=f(f^\applirec(Y))$.
  \item On suppose que pour toute partie $Y$ de $B$, $f(f^\applirec(Y))=Y$. Montrons que $f$ est surjective. Pour $Y\defeq B$, $f(f^\applirec(B))=B$. Soit $y\in B$. Il existe donc $x\in f^\applirec(B)$ tel que $y=f(x)$.
  \end{itemize}
\question Montrons que $f$ est injective si et seulement si, quelle que soit la partie $X$ de $A$, on a $f^\applirec\p{f(X)}=X$.
  \begin{itemize}
  \item On suppose que $f$ est injective. Soit $X$ une partie de $A$. Montrons que $f^\applirec\p{f(X)}=X$. On procède par double inclusion.
  \begin{itemize}
  \item $X\subset f^\applirec\p{f(X)}$. Soit $x\in X$. Montrons que $x\in f^\applirec\p{f(X)}$. Puisque $x\in X$, $f(x)\in f(X)$. Donc $x\in f^\applirec\p{f(X)}$.
  \item $f^\applirec\p{f(X)}\subset X$. Soit $x\in f^\applirec\p{f(X)}$. Alors $f(x)\in f(X)$. Il existe donc $x'\in X$ tel que $f(x)=f(x')$. Or $f$ est injective donc $x=x'$. Donc $x\in X$.
  \end{itemize}
  Donc $X= f^\applirec\p{f(X)}$.
\item On suppose que quelle que soit la partie $X$ de $A$, on a $f^\applirec\p{f(X)}=X$. Montrons que $f$ est injective. Soit $x_1,x_2\in A$ tels que $f(x_1)=f(x_2)$. Montrons que $x_1=x_2$. On pose $X\defeq\ens{x_1}$. Alors $f(X)=\ens{f(x_1)}$, donc $f(x_2)\in f(X)$. Donc $x_2\in f^\applirec\p{f(X)}$. Or $f^\applirec\p{f(X)}=X$, donc $x_2\in X$. Donc $x_2=x_1$.
  \end{itemize}
\end{questions}
\end{sol}

% \exercice{nom={Identification d'ensembles par bijection}}
% \begin{questions}
% \question Soit $A$ et $A'$ deux ensembles tels qu'il existe une bijection
%   $f_1$ de $A$ dans $A'$, $B$ et $B'$ deux ensembles tels qu'il existe une
%   bijection $f_2$ de $B$ dans $B'$. On définit l'application $\phi$ de
%   $\mathcal{F}(A,B)$ dans $\mathcal{F}(A',B')$ qui à l'application $f$ associe
%   l'application $f_2\circ f\circ f_1^{-1}$.
%   \begin{questions}
%   \question Montrer que $\phi$ est injective.
%   \question Montrer que $\phi$ est surjective.
%   \question Calculer $\phi^{-1}$.
%   \question Montrer que l'image par $\phi$ d'une application injective est
%     injective.
%   \question Montrer que l'image par $\phi$ d'une application surjective est
%     surjective.
%   \question Que dire d'une application dont l'image par $\phi$ est injective~?
%     surjective~?
%   \end{questions}
% \question Soit $A$ et $A'$ deux ensembles, $f_1$ une application de $A'$ dans
%   $A$, $B$ et $B'$ deux ensembles, $f_2$ une application de $B$ dans $B'$. On
%   définit l'application $\phi$ de $\mathcal{F}(A,B)$ dans $\mathcal{F}(A',B')$
%   qui à l'application $f$ associe l'application $f_2\circ f\circ f_1$.
%   \begin{questions}
%   \question Montrer que $\phi$ est injective si et seulement si $f_1$ est
%     surjective et $f_2$ est injective.
%   \question Montrer que $\phi$ est surjective si et seulement si $f_1$ est
%     injective et $f_2$ est surjective.
%   \end{questions}
% \end{questions}

% \exercice{nom={Équipotence}}
% On dit que deux ensembles $A$ et $B$ sont équipotents lorsqu'il existe une
% bijection entre $A$ et $B$. Montrer que les ensembles suivants sont
% équipotents.
% \begin{questions}
% \question $E\times F$ et $F\times E$.
% \question $(E\times F)^G$ et $E^G \times F^G$.
% \question $(E^F)^G$ et $E^{F\times G}$.
% \question $(\mathcal{P}(E))^F$ et $\mathcal{P}(E\times F)$.
% \end{questions}

% \exercice{nom={Bijection entre $\N$ et $\N^2$}}
% \begin{questions}
% \question Montrer graphiquement qu'il existe une bijection entre $\N$
%   et $\N^2=\N\times\N$.
% \question En déduire, par récurrence sur $n$, qu'il existe une bijection entre
%   $\N$ et $\N^n$.
% \end{questions}

\exercice{nom={Application}}
Soit $E$ un ensemble et $f$ une application de $E$ dans $E$ telle que $f\circ f\circ f=f$. Montrer que $f$ est injective si et seulement si $f$ est surjective.

\exercice{nom={Application}}
Soit $E$ et $F$ deux ensembles, $f$ une application de $E$ dans $F$, et $g$ une application de $F$ dans $E$. On suppose que $f\circ g\circ f$ est bijective. Montrer que $f$ et $g$ sont bijectives.

\magsubsection{Familles}

\exercice{nom={Partition}}
Soit $E$ l'ensemble des fonctions de $\N$ dans $\ens{1,2,3}$. Pour tout
$i\in\ens{1,2,3}$, on pose $A_i\defeq\enstq{f\in E}{f(0)=i}$. Montrer que les $A_i$
forment une partition de $E$.

\magsection{Relation binaire}
\magsubsection{Relation d'ordre}


\exercice{nom={Ordre sur $\Ns$}}
On considère la relation $\mathcal{R}$ définie sur $\Ns$ par
\[\forall n,m\in\Ns\qsep n\mathcal{R}m \quad\ssi\quad \cro{\exists q\in\Ns \qsep m=n^q}.\]
\begin{questions}
\question Montrer que $\mathcal{R}$ est une relation d'ordre sur $\Ns$.
\question Est-ce que $\mathcal{R}$ est totale~?
\end{questions}

\exercice{nom={Plus grand, plus petit élément}}
Montrer que si $E$ est un ensemble ordonné dont l'ordre est total, toute partie
finie de $E$ admet un plus petit et un plus grand élément. Que dire si l'ordre
n'est pas total~?

\exercice{nom={Applications croissantes}}
Soit $\p{E,\preceq}$ un ensemble ordonné. Une application $f$ de $E$ dans $E$
est dite croissante lorsque
\[\forall x,y\in E \qsep x\preceq y \implique f(x)\preceq f(y).\]
\begin{questions}
\question Montrer que la composée de deux applications croissantes est
  croissante.
\question Montrer que si $E$ est totalement ordonné, l'application réciproque
  d'une bijection croissante est croissante.
  %% Trop dur :
  % Ce résultat est-il toujours vrai
  % lorsque $E$ n'est pas totalement ordonné~?
\end{questions}
\begin{sol}
$\quad$
\begin{questions}
\question Soit $f$ et $g$ deux applications croissantes de $E$ dans $E$. Montrons
  que $g\circ f$ est croissante. Soit $x,y\in E$ tels que $x\preceq y$. Alors
  $f(x)\preceq f(y)$ car $f$ est croissante, donc
  $g\p{f(x)}\preceq g\p{f(y)}$ car $g$ est croissante. Donc $g\circ f$
  est croissante.
\question Soit $f$ une bijection croissante de $E$ dans $E$. Montrons que $f^{-1}$
  est croissante. Soit $y_1,y_2\in E$ tels que $y_1\preceq y_2$. Montrons que
  $f^{-1}\p{y_1}\preceq f^{-1}\p{y_2}$. On raisonne par l'absurde et on suppose
  que ce n'est pas le cas. Comme $E$ est totalement ordonné, on en déduit que
  $f^{-1}\p{y_2}\preceq f^{-1}\p{y_1}$. Par croissance de $f$, on obtient
  $f\p{f^{-1}\p{y_2}}\preceq f\p{f^{-1}\p{y_1}}$ donc $y_2\preceq y_1$. Mais
  comme $y_1\preceq y_2$, on en déduit que $y_1=y_2$. Donc
  $f^{-1}\p{y_1}=f^{-1}\p{y_2}$ donc $f^{-1}\p{y_1}\preceq f^{-1}\p{y_2}$ ce qui est
  absurde. Donc $f^{-1}\p{y_1}\preceq f^{-1}\p{y_2}$. En conclusion, $f^{-1}$ est
  croissante.
\end{questions}
\end{sol}

\magsubsection{Relation d'équivalence}


% \exercice{nom={Borne supérieure}}
% Montrer que les relations suivantes sont des relations d'ordre. Sont-elles
% totales~? Les parties non vides finies admettent-elles une borne supérieure~? Si
% oui, calculer cette borne supérieure. Les parties non vides majorées admettent
% elles une borne supérieure~? Si oui, calculer cette borne supérieure.
% \begin{enumerate}
% \question $\mathcal{P}(E)$ muni de la relation d'inclusion.
% \question $\mathcal{C}^0(\R,\R)$ muni de la relation $\leq$
%   (définir cette relation).
% \question $\mathcal{C}^1(\R,\R)$ muni de la relation $\leq$
%   (définir cette relation).
% \question $\mathbb{N}$ muni de la relation de divisibilité.
% \question $\mathbb{Z}$ muni de la relation de divisibilité.
% \end{enumerate}

% \magsection{Relations d'équivalence}

\exercice{nom={Relation sur $\R$}}
On note $\mathcal{R}$ la relation définie sur $\R$ par
\[\forall x,y\in\R\qsep x\mathcal{R}y \quad\ssi\quad x^2 - 2x = y^2 - 2y.\]
\begin{questions}
\question Montrer que $\mathcal{R}$ est une relation d'équivalence sur $\R$.
\question Pour tout $x\in\R$, déterminer la classe d'équivalence modulo $\mathcal{R}$.
\end{questions}

\exercice{nom={Factorisation canonique}}
Soit $f$ une application de $A$ dans $B$ et $\mathcal{R}$ la relation binaire définie sur $A$ par
\[\forall x_1,x_2 \in A \qsep x_1\mathcal{R}x_2 \quad\Longleftrightarrow\quad
  f(x_1)=f(x_2).\]
\begin{questions}
\question Montrer que $\mathcal{R}$ est une relation d'équivalence.
\question Montrer que les classes d'équivalence sont les images réciproques des
  $\{y\}$ pour $y\in f(A)$.
\enonce On appelle ensemble quotient $A/\mathcal{R}$ l'ensemble des classes d'équivalence pour la relation $\mathcal{R}$.
\question Soit $s$ l'application de $A$ dans $A/\mathcal{R}$ qui à $x$ associe la classe de
  $x$. Montrer que $s$ est une surjection appelée surjection canonique.
\question Soit $i$ l'application de $f(A)$ dans $B$ qui à $y$ associe $y$. Montrer
  que $i$ est une injection appelée injection canonique.
\question Montrer qu'il existe une et une seule application $\bar{f}$ de $A/\mathcal{R}$
  dans $f(A)$ telle que $f=i\circ \bar{f}\circ s$. Montrer que $\bar{f}$ est une
  bijection.
% \question Retrouver le lemme des Bergers.
\question Soit $C$ un ensemble et $g$ une application de $A$ dans $C$. Montrer
  qu'il existe une application $\bar{g}$ de $A/\mathcal{R}$ dans $C$ telle que
  $g=\bar{g}\circ s$ si seulement si
  $$\forall x_1,x_2 \in A \qsep x_1\mathcal{R}x_2 \Longrightarrow g(x_1)=g(x_2).$$
\end{questions}

\magsection{L'ensemble des entiers naturels}

\magsubsection{Récurrence}

\exercice{nom={Inégalité}}
Montrer par récurrence que
\[\forall n\in\Ns\qsep \sqrt{1}+\sqrt{2}+\cdots+\sqrt{n}\leq \frac{4n+3}{6}\sqrt{n}.\]

\exercice{nom={Fibonacci}}
On considère la suite de Fibonacci définie par
\[F_0\defeq 0,\quad F_1\defeq 1 \et \forall n\in\N\qsep F_{n+2}\defeq F_{n+1}+F_n.\]
\begin{questions} 
\question Montrer que~: $\forall n\in\N\qsep F_n\geq n-1$. 
\question Montrer que~: $\forall n\in\N\qsep F_nF_{n+2}-F_{n+1}^2=(-1)^{n+1}$. 
\question Montrer que
  \[\forall n\in\N\qsep F_{2n+1}=F_{n+1}^2+F_n^2 \et F_{2n+2}=F_{n+2}^2-F_{n}^2.\]
\end{questions}
% \begin{sol}
\begin{questions}
\question
\question
\begin{eqnarray*}
F_{n+1}F_{n+3}-F_{n+2}^2
&=& F_{n+1}(F_{n+2}+F_{n+1})-(F_n+F_{n+1})^2\\
&=& F_{n+1}F_{n+2}+F_{n+1}^2-F_n^2-2F_n F_{n+1}-F_{n+1}^2\\
&=& F_{n+1}(F_n+F_{n+1}) -F_n^2-2F_n F_{n+1}\\
&=& F_{n+1}^2-F_n(F_n+F_{n+1})\\
&=& F_{n+1}^2-F_n F_{n+2}\\
&=& -(-1)^{n+1} = (-1)^{n+2}
\end{eqnarray*}
\question
\begin{eqnarray*}
F_{2n+4}
&=& F_{2n+3}+F_{2n+2}=F_{n+2}^2+F_{n+1}^2+F_{n+2}^2-F_{n}^2\\
&=& (F_{n+3}-F_{n+1})^2+F_{n+1}^2+F_{n+2}^2-F_{n}^2\\
&=& F_{n+3}^2-2F_{n+3}F_{n+1}+2F_{n+1}^2+F_{n+2}^2-F_{n}^2\\
&=& F_{n+3}^2-2F_{n+2}^2-2(-1)^{n+1}+2F_{n+1}^2+F_{n+2}^2-F_{n}^2\\
&=& F_{n+3}^2+2F_{n+1}^2-F_{n+2}^2-F_{n}^2-2(-1)^{n+1}\\
&=& F_{n+3}^2+2F_{n+1}^2-((F_{n+2}-F_{n})^2+2F_{n+2}F_n)-2(-1)^{n+1}\\
&=& F_{n+3}^2+2F_{n+1}^2-F_{n+1}^2-2F_{n+2}F_n-2(-1)^{n+1}\\
&=& F_{n+3}^2+F_{n+1}^2-2((-1)^n+F_{n+1}^2)-2(-1)^{n+1}\\
&=& F_{n+3}^2-F_{n+1}^2-2((-1)^n+(-1)^{n+1})\\
&=&F_{n+3}^2-F_{n+1}^2
\end{eqnarray*}
\end{questions}
% \end{sol}

\exercice{nom={Inégalité sur Fibonacci}}
On considère la suite de Fibonacci définie par
\[F_0\defeq 1, \quad F_1\defeq 1, \quad\text{et}\quad \forall n\in\N\qsep F_{n+2}\defeq F_{n+1}+F_n.\]
Déterminer les $r\in\RPs$ tels qu'il existe $\alpha\in\RPs$ tel que
\[\forall n\in\N\qsep F_n\leq \alpha r^n.\]

\exercice{nom={Inégalité}}
Montrer que
\[\forall n\in\Ns\qsep \frac{4^n}{2\sqrt{n}}\leq\binom{2n}{n}\leq\frac{4^n}{\sqrt[3]{n}}.\]

% \exercice{nom={Inégalité triangulaire}}
% Montrer que quels que soient $z_1,\ldots,z_n\in\C$
% \[\abs{z_1+\cdots+z_n}\leq\abs{z_1}+\cdots+\abs{z_n}\]
% et que cette inégalité est une égalité si et seulement si il existe $\theta\in\R$ tel que
% \[\forall k\in\intere{1}{n}\qsep z_k=\abs{z_k}\e^{\ii \theta}.\]

\exercice{nom={Injections de $\N$ dans $\N$}}
Déterminer les injections $f$ de $\N$ dans $\N$ telles que
\[\forall n\in\N\qsep f(n)\leq n.\]

\exercice{nom={Équation fonctionnelle}}
Le but de cet exercice est de déterminer les fonctions $f$ de $\N$ dans $\N$ telles que
\[(E) \qquad \forall n\in\N\qsep f(n)+(f\circ f)(n)=2n.\]
\begin{questions}
\question Déterminer une fonction vérifiant $(E)$.
\question Réciproquement, soit $f$ une fonction de $\N$ dans $\N$ vérifiant $(E)$.
\begin{questions}
\question Montrer que $f$ est injective.
\question Montrer que~: $\forall n\in\N\qsep f(n)=n$. Conclure.
\end{questions}
\end{questions}


\exercice{nom={Équation fonctionnelle}}
Montrer qu'il existe une unique bijection $\app{f}{\N}{\N}$ telle que
\[\forall n\in\N\qsep \abs{f(n)-n}=1.\]
Que se passe-t-il si on remplace $\N$ par $\Z$~?

\begin{sol}
Posons $\mathcal{H}_n$ : "$ \dspappli{f_{|\intere{0}{2n-1}}}{\intere{0}{2n-1}}{\intere{0}{2n-1}}{k}{\begin{cases} k+1 \text{ si } k \text{ est pair}\\ k-1 \text{ si } k \text{ est impair}.\end{cases}}$"

Et ça fonctionne bien.

Sur $\Z$, on peut tous les décaler de $1$ dans un sens ou dans l'autre. Ou bien les échanger deux à deux : deux nouvelles façons si on commence par échanger $0$ avec $1$ ou $-1$.
\end{sol}

\exercice{nom={Les crayons de couleur}}
Nous allons démontrer que toute boite de crayons de couleur ne possède que des crayons de la même couleur. Pour cela, on procède par récurrence sur le nombre $n$ de crayons. L'initialisation est évidente car une boite ne contenant qu'un crayon ne possède que des crayons de la même couleur. Pour l'hérédité, supposons que le résultat est vrai pour $n$ crayons et considérons une boite de $n+1$ crayons de couleur. On enlève le premier crayon. Par hypothèse de récurrence, tous les autres crayons ont la même couleur. On replace le premier crayon et on enlève le dernier crayon. De même, tous les autres crayons ont la même couleur. On en déduit que les $n+1$ crayons ont la même couleur.\\

Quelle est l'erreur de ce raisonnement~?

\begin{sol}
L'hérédité ne passe pas de $1$ à $2$, on a aucune intersection entre les deux paquets de crayons...
\end{sol}

\exercice{nom={Les Moines}}
Dans un camp de bouddhistes, on apprend qu'il y a au moins un malade. Cette maladie n'est pas contagieuse ni évolutive (le nombre de malades n'évoluera plus). Afin de préserver une entière pureté et ne pas perturber les méditations, un bouddhiste qui se sait malade doit partir.
Cette maladie se caractérise uniquement par une tâche rouge sur le front. Un symptôme qui leur permet de reconnaître sans hésitation si une personne est malade.
Le problème est qu'il n'y a aucun moyen pour un bouddhiste de se voir. Il n'y a aucun miroir ou autre moyen permettant de voir son propre front. De plus, les moines bouddhistes ont fait le vœu de silence et ne communiquent d'aucune façon. Ils ne font que méditer, lire et ont un esprit très logique.
Ils se réunissent tous une seule fois par jour au lever du soleil pour une méditation commune de 3 heures. Pendant ces trois heures, ils n'ont toujours pas le droit de communiquer entre eux ni de partir avant la fin de la séance commune. Au bout de 5 jours, tous les malades sont partis.
Combien y avait-il de moines malades ?

% \exercice{nom={Les cocus de Babylone}}
% Un jour paraît dans la bonne ville de Babylone un décret royal dont voici
% le texte~:
% \og Par ordre du roi, tout babylonien s'apercevant que sa femme le trompe est
% tenu de la répudier le jour même.  Il y a actuellement au moins un cocu à
% Babylone. \fg.
% Or à cette époque, et dans cette bonne ville de Babylone, les habitants, fort
% rompus aux techniques du raisonnement logique, avaient la caractéristique
% surprenante de connaitre parfaitement toutes les histoires extra-conjugales
% des maisons voisines sans jamais s'apercevoir de quoi que ce soit sous leur
% propre toit. Et, fiers de nature, ils n'auraient pour rien au monde demandé
% à un autre ce qui se passait chez eux~!
% \begin{itemize}
% \item Le premier jour, tout va bien, aucune femme n'est répudiée.
% \item Le deuxième jour, toujours rien. Et ainsi de suite...
% \item Le cent soixante huitième jour, cent soixante huit femmes sont répudiées.
% \end{itemize}
% Expliquez.

% \begin{sol}
%     S'il y a un seul cocu, le problème est simple : connaissant la vie conjugale de tous les autres babyloniens, il sait qu'aucun autre n'est trompé. Or l'édit royal affirme l'existence de cocus à Babylone : c'est donc que lui-même est cocu !
%     L'homme répudiera donc sa femme le premier jour.
%     S'il y a deux cocus à Babylone, chacun des deux connait l'autre. Il n'a donc aucune raison de répudier sa femme le premier jour. En revanche, voyant le soir venu qu'aucune femme n'a été répudiée, il en déduit immédiatement qu'il est cocu lui-même.
%     En effet, il n'y aurait dans le cas contraire qu'un seul cocu à Babylone, et celui-ci aurait répudié sa femme le premier jour.
%     Les deux hommes répudieront donc leur femme le deuxième jour.
%     S'il y a trois cocus à Babylone, chacun connait les deux autres. De son point de vue, il y a donc soit deux, soit trois cocus à Babylone. Il va donc observer ce qui se passe le deuxième jour. Mais chacun des trois fait la même chose, et donc aucun ne répudie sa femme le deuxième jour.
%     Comme il ne se passe rien, chacun des trois en déduit qu'il y a trois cocus à Babylone, dont lui-même, et répudie donc sa femme le troisième jour.
%     Et ainsi de suite. Par récurrence sur n, on montre ainsi que s'il y a n cocus à Babylone, tous répudieront leur femme le n-ème jour !
%     Il y a donc 168 cocus à Babylone.
% \end{sol}

\magsubsection{Définition par récurrence}

\exercice{nom={Suite définie par récurrence}}
Montrer qu'il existe une unique suite $(u_n)$ telle que
\[u_0=1\qsep u_1=1, \et \cro{\forall n\in\N\qsep u_{n+2}=\frac{1}{u_{n+1}}+u_n}.\]

%END_BOOK

\end{document}


\chapter{Compléments d'analyse}
\setcounter{numeroexercicecours}{1}
\documentclass{magnoliaold}

\magtex{tex_driver={pdftex},
        tex_packages={epigraph,slashbox,xypic}}
\magfiche{document_nom={Réels},
          auteur_nom={François Fayard},
          auteur_mail={fayard.prof@gmail.com}}
\magcours{cours_matiere={maths},
          cours_niveau={mpsi},
          cours_chapitre_numero={3},
          cours_chapitre={Compléments d'analyse}}
\magmisenpage{misenpage_presentation={tikzvelvia},
          misenpage_format={a4},
          misenpage_nbcolonnes={1},
          misenpage_preuve={non},
          misenpage_sol={non}}
\maglieudiff{}
\magprocess

\begin{document}

%BEGIN_BOOK
\setlength\epigraphwidth{.8\textwidth}
\epigraph{\og Le calcul infinitésimal est l'apprentissage du maniement des inégalités bien plus que des égalités, et on pourrait le résumer en trois mots~: majorer, minorer, approcher.\fg}{--- \textsc{Jean Dieudonné (1906--1992)}}
\setlength\epigraphwidth{.6\textwidth}
\epigraph{\og Les hommes sont comme les chiffres, ils n'acquièrent de valeur que par leur position.\fg}{--- \textsc{Napoléon Bonaparte (1769--1821)}}
\hidemesometimes{
\hfill\includegraphics[width=0.55\textwidth]{../../Commun/Images/maths-cours-diff_int.png}}
\magtoc

\section{Le corps ordonné $\R$}

\subsection{La relation d'ordre sur $\R$}

%% On utilise a,b,c,d et non x,y,.. pour avoir 4 lettres

\begin{proposition}[utile=-3]
La relation d'ordre $\leq$ définie sur $\R$ possède les propriétés suivantes.
\begin{itemize}
\item Elle est totale.
  \[\forall a,b\in\R \qsep a\leq b \quad\ou\quad b\leq a.\]
\item Elle est compatible avec l'addition.
  \[\forall a,b,c\in\R \qsep a\leq b \quad\implique\quad a+c\leq b+c.\]
\item Elle est compatible avec la multiplication.
  \[\forall a,b\in\R \qsep \cro{0\leq a \et 0\leq b} \quad\implique\quad 0\leq ab.\]
\end{itemize}
\end{proposition}

\begin{preuve}
On admet cette proposition pour démontrer les trois suivantes.
\end{preuve}

\begin{remarques}
\remarque La relation $\leq$ étant antisymétrique sur $\R$, 0 est le seul réel
  à la fois positif et négatif.
  \begin{sol}
  $a\leq 0$ et $0\leq a$ $\Rightarrow$ $a=0$.
  \end{sol}
\remarque Si $a,b\in\R$, la négation de \og $a\leq b$ \fg  est \og $a>b$ \fg.
\remarque Deux réels $a$ et $b$ sont de même signe si et
  seulement si $ab\geq 0$. On dit qu'ils sont de même signe au sens strict
  lorsque $ab>0$.
\begin{sol}
Déjà, si $b\leq 0$, alors d'après (ii) $b+(-b)\leq 0+(-b)$, i.e. $0\leq -b$.
Ainsi, si $0<a$ et $b<0$ alors $0\leq a$ et $b\leq 0$ (donc $0\leq -b$ d'après ce qu'on vient de dire). D'après (iii), cela implique $0\leq a(-b)=-ab$.
Or, si $0\leq c$, alors d'après (ii), $0+(-c)\leq c+(-c)$, i.e. $-c\leq 0$.
Ainsi, $-(-ab)\leq 0$, i.e. $ab\leq 0$. Mais $ab\neq 0$ donc $ab<0$.
On a montré $$0<a \text{ et } b<0 \Longrightarrow ab<0.$$
\end{sol}  
\remarque Quel que soit $a\in\R$, $a^2\geq 0$.
\end{remarques}

\begin{exos}
\exo Soit $a,b$ deux réels positifs. Montrer que
  \[\sqrt{ab}\leq\frac{a+b}{2}.\]
\exo Soit $a,b,c\in\R$ tels que $a^2+b^2+c^2=ab+bc+ca$. Montrer que
  $a=b=c$.
  \begin{sol}
  $(a-b)^2+(a-c)^2+(b-c)^2=...$
  \end{sol}
\end{exos}

\begin{proposition}
\begin{eqnarray*}
\forall a,b,c,d\in\R, & & \cro{a\leq b \et c\leq d}\quad\implique\quad
       a+c\leq b+d\\
\forall a,b,c\in\R, & & \cro{a\leq b \et 0\leq c}\quad\implique\quad ac\leq bc\\
\forall a,b,c,d\in\R, & & \cro{0\leq a\leq b \et 0\leq c\leq d}\quad\implique\quad
       0\leq ac\leq bd\\
\forall a,b\in\R, \quad \forall n\in\N, & & 0\leq a\leq b \quad\implique\quad
  0\leq a^n\leq b^n.
\end{eqnarray*}
\end{proposition}

\begin{preuve}
$a+c\leq b+c\leq b+d$ avec (ii)
$b-a\geq 0$ donc $c(b-a)\geq 0$ avec (iii)
$ac\leq bc \leq bd $ avec le point précédent.
Récurrence avec le point précédent.
\end{preuve}

\begin{remarqueUnique}
\remarque On peut multiplier une inégalité de signe quelconque par un réel
  négatif. Dans ce cas, l'inégalité change de sens.
\begin{sol}
Si [$a\leq b$ et $c\leq 0$], alors [$a\leq b$ et $0\leq -c$] donc $-ac\leq -bc$ d'où d'après (iii) en ajoutant $ac+bc$ des deux côtés, $bc\leq ac$.
\end{sol}  
  
% \remarque On peut multiplier une inégalité de signe quelconque par
%   un réel positif mais pas une inégalité positive par une inégalité de
%   signe quelconque. Par exemple $-1\leq -1$ et $0\leq 1$ mais $0\not\leq -1$.
% \remarque De cette proposition, on peut en déduire que~:
%   \begin{itemize}
%   \item La somme de deux fonctions croissantes est croissante.
%   \item Le produit d'une fonction croissante par un réel positif est
%     croissante.
%   \item Le produit de deux fonctions croissantes positives est croissante
%     positive.
%   \item Si $f$ est une fonction croissante positive et $n\in\N$ alors $f^n$ est
%     croissante positive.
%   \end{itemize}
%   En utilisant le fait qu'une fonction $f$ est croissante si et seulement si
%   $-f$ est décroissante, on en déduit une série d'assertions qui permettent
%   souvent de déterminer la monotonie d'ne fonction élémentaire sans la dériver.
% \remarque La bijection réciproque d'une fonction croissante est croissante.
\end{remarqueUnique}

\begin{exoUnique}
\exo L'assertion suivante est-elle vraie~?
  \[\forall a,b,c,d\in\R \qsep\cro{a\leq b \et 0\leq c\leq d}\implique
    ac\leq bd\]
\end{exoUnique}
\begin{sol}
FAUX $-2\leq -1$ mais $(-2)\times 2 > (-1)\times 5$.
\end{sol}

\begin{proposition}
Soit $a,b\in\R$. Alors
\[0<a\leq b \quad\implique\quad 0<\frac{1}{b}\leq\frac{1}{a}.\]
\end{proposition}

\begin{preuve}
Déjà, si $0\leq a$ alors $0\leq \frac{1}{a}$. En effet, supposons le contraire, i.e. $\frac{1}{a}<0$ alors $\frac{1}{a}\leq 0$ donc $a\frac{1}{a}\leq a\dot 0$ i.e. $1\leq 0$, d'où la contradiction.\\

$a\leq b$ et $0\leq 1/a$ donc $1\leq b/a$. Or $0\leq 1/b$ donc $1/b\leq 1/a$.

\end{preuve}

\begin{proposition}
\begin{eqnarray*}
\forall a,b,c,d\in\R, & & \cro{a\leq b \et c< d}\quad\implique\quad
       a+c< b+d\\
\forall a,b,c\in\R, & & \cro{a< b \et 0< c}\quad\implique\quad ac< bc\\
\forall a,b,c,d\in\R, & & \cro{0\leq a< b \et 0\leq c< d}\quad\implique\quad
       0\leq ac< bd\\
\forall a,b\in\R, \quad \forall n\in\Ns, & & 0\leq a< b \quad\implique\quad
  0\leq a^n< b^n.
\end{eqnarray*}
\end{proposition}

\begin{preuve}
Il faut prouver (ii) et (iii) de la proposition 1.1 avec des inégalités strictes.
(iii) est évidente.
Pour (ii), supposons $a<b$. On cherche à montrer que $a+c<b+c$. Supposons le contraire, $b+c\leq a+c$ donc d'après (ii) $b+c-c\leq a+c-c$ donc $b\leq a$. CONTRADICTION
On peut maintenant prouver la prop comme on avait prouvé la prop. 1.2 avec les nouveaux (ii) et (iii).
\end{preuve}


\begin{definition}[utile=-3]
Soit $a,b\in\R$ avec $a\leq b$. On définit
\[\interf{a}{b}=\enstq{x\in\R}{a\leq x\leq b}, \qquad
  \intero{a}{b}=\enstq{x\in\R}{a< x< b},\]
\[\interof{a}{b}=\enstq{x\in\R}{a< x\leq b}, \qquad
  \interfo{a}{b}=\enstq{x\in\R}{a\leq x< b},\]
\[\interfo{a}{+\infty}=\enstq{x\in\R}{a\leq x}, \qquad
  \intero{a}{+\infty}=\enstq{x\in\R}{a< x},\]
\[\interof{-\infty}{b}=\enstq{x\in\R}{x\leq b}, \qquad
  \intero{-\infty}{b}=\enstq{x\in\R}{x< b}.\]
\end{definition}

\subsection{Valeur absolue}

\begin{definition}[utile=-3]
Pour tout réel $a$, on définit sa \emph{valeur absolue}, notée $\abs{a}$ par
  \[\abs{a}\defeq
    \begin{cases}
    a  & \text{si $a\geq 0$}\\
    -a & \text{si $a<0$.}
    \end{cases}\]
\end{definition}

\begin{remarques}
\remarque Pour tout $a\in\R$, $a^2=\abs{a}^2$.
\remarque Si $a$ et $b$ sont deux réels, on définit la distance de $a$ à $b$,
  notée $\dis{a}{b}$ par
  \[\dis{a}{b}\defeq\abs{a-b}.\]
\end{remarques}

\begin{exoUnique}
\exo Soit $a,b\in\R$. Exprimer $\min(a,b)$ et $\max(a,b)$
  à l'aide de $a$, $b$ et de la valeur absolue.
\end{exoUnique}
\begin{sol}
$\max(a,b)=\dfrac{a+b+\abs{a-b}}{2}$ et $\min(a,b)=\dfrac{a+b-\abs{a-b}}{2}$.
\end{sol}

\begin{proposition}
\begin{eqnarray*}
\forall a\in\R, & & \abs{a}\geq 0\\
\forall a\in\R, & & \abs{a}=0 \ssi a=0\\
\forall a\in\R, & & \abs{-a}=\abs{a}\\
\forall a,b\in\R, & & \abs{ab}=\abs{a}\abs{b}.
\end{eqnarray*}
\end{proposition}
\begin{preuve}
RAS pour les trois premiers points. Disjonction de 2 cas + inversion des rôles pour le quatrième point.
\end{preuve}

\begin{remarques}
\remarque Soit $a\in\R$. Alors, quel que soit $n\in\N$, $\abs{a^n}=\abs{a}^n$. Si de
  plus $a\neq 0$, quel que soit $n\in\Z$, $\abs{a^n}=\abs{a}^n$. 
\remarque De cette proposition, on déduit les résultats suivants sur la distance
  entre deux réels.
  \begin{eqnarray*}
  \forall a,b\in\R, & & \dis{a}{b}\geq 0\\
  \forall a,b\in\R, & & \dis{a}{b}=0 \quad\ssi\quad a=b\\
  \forall a,b\in\R, & & \dis{b}{a}=\dis{a}{b}.
  \end{eqnarray*}
\end{remarques}

\begin{exoUnique}
\exo Soit $a>0$ et $x,y\geq a$. Montrer que
  \[\abs{\sqrt{x}-\sqrt{y}}\leq\frac{1}{2\sqrt{a}}\abs{x-y}.\]
	\begin{sol}
	Suffit de multiplier en haut et en bas par la quantité conjuguée.
	\end{sol}
\end{exoUnique}



\begin{proposition}
Soit $a$ un réel. Alors
\[\abs{a}=\max\ens{a,-a}.\]
\end{proposition}

\begin{remarques}
\remarque En particulier, si $M$ est un réel positif, pour montrer que
  $\abs{a}\leq M$ il suffit de montrer que
  \[a\leq M \et -a\leq M.\] 
% \remarque On n'a pas $a\leq b \implique \abs{a}\leq \abs{b}$. Par exemple
%   $-1\leq 0$ mais $\abs{-1}=1\not\leq\abs{0}=0$.
\remarque Soit $a$ un réel et $M$ un réel positif. Alors
  \[\abs{a}\leq M \quad\ssi\quad -M\leq a\leq M\]
  \[\abs{a}\geq M \quad\ssi\quad \cro{a\leq -M \ou a\geq M}.\]
\end{remarques}

\begin{exoUnique}
\exo Soit $x,y\in\R$. Montrer que $\cos(x)\sin(y)\geq -1$.  
\end{exoUnique}

\begin{proposition}
Soit $a$ et $b$ deux réels. Alors
\[\abs{a+b}\leq \abs{a}+\abs{b}.\]
De plus, l'égalité a lieu si et seulement si $a$ et $b$ sont de même signe.  
\end{proposition}

\begin{preuve}
Ou bien on dit qu'on l'a vu dans $\C$, ou bien mieux dans $\R$, c'est plus facile :

$a\leq \abs{a}$, $b\leq\abs{b}$ donc $a+b\leq |a|+|b|$.

On applique ce résultat à $-a$ et $-b$ puis on dit que la valeur absolue est un des deux cas qu'on vient de prouver. 

Pour le cas d'égalité, il suffit de passer l'égalité au carré.
\end{preuve}

\begin{proposition}
Soit $a$ et $b$ deux réels. Alors
\[\abs{\abs{a}-\abs{b}}\leq\abs{a-b} \et \abs{a+b}\geq \abs{a}-\abs{b}.\]
\end{proposition}

\begin{preuve}
cf. $\C$
\end{preuve}

\begin{remarqueUnique}
\remarque Soit $a$, $b$, $c\in\R$. Alors
  \[\abs{\dis{a}{b}-\dis{b}{c}} \leq \dis{a}{c} \leq
    \dis{a}{b}+\dis{b}{c}.\]
\end{remarqueUnique}

\begin{exoUnique}
\exo Parmi les assertions suivantes, lesquelles sont vraies~?
  \begin{multicols}{2}
  \begin{itemize}
  \item $\forall a,b\in\R \qsep \abs{a-b}\leq\abs{a}-\abs{b}$.
  \item $\forall a,b\in\R \qsep a^2\leq b^2 \ssi \abs{a}\leq\abs{b}$.
  \item $\forall a,b\in\R \qsep \abs{a-b}\leq\abs{a}+\abs{b}$.
  \item $\forall a,b\in\R \qsep a\leq b \implique \abs{a}\leq\abs{b}$.
  \item $\forall a,b\in\R \qsep \abs{a-b}\geq\abs{a}-\abs{b}$.
  \item $\forall a,b\in\R \qsep \abs{a}\leq\abs{b}+\abs{b-a}$.
  \end{itemize}  
  \end{multicols}
\end{exoUnique}
\begin{sol}
\begin{itemize}
  \item FAUX. $a=2$, $b=-1$.
  \item VRAI
  \item VRAI
  \item FAUX. Prendre $b<0$.
  \item VRAI
  \item VRAI
  \end{itemize}  
  \end{sol}

\begin{proposition}
Soit $a_1,\ldots,a_n\in\R$. Alors
\[\abs{\sum_{k=1}^n a_k}\leq \sum_{k=1}^n \abs{a_k}.\]
\end{proposition}

\begin{exoUnique}
\exo Soit $a_1,\ldots,a_n\in\R$ et $\theta_1,\ldots,\theta_n\in\R$. Montrer
  que
  \[\abs{\sum_{k=1}^n a_k \sin(\theta_k)}\leq\sum_{k=1}^n \abs{a_k}.\]
\end{exoUnique}


\subsection{Racine}

\begin{definition}
Soit $a\in\R$ et $n\in\Ns$.
\begin{itemize}
\item Si $n$ est pair et $a\geq 0$, il existe un unique réel positif $x$ tel que
  $x^n=a$. On le note $\sqrt[n]{a}$.
\item Si $n$ est impair, il existe un unique réel $x$ tel que
  $x^n=a$. On le note $\sqrt[n]{a}$.
\end{itemize}
\end{definition}


\begin{remarques}
\remarque Soit $a\in\R$ et $n\in\Ns$.
  \begin{itemize}
  \item Si $n$ est pair et $a\geq 0$, alors
    \[\forall x\in\R\qsep x^n=a \quad\ssi\quad x=\sqrt[n]{a} \ou x=-\sqrt[n]{a}\]
  \item Si $n$ est pair et $a< 0$, l'équation $x^n=a$ n'admet aucune solution sur $\R$.
  \item Si $n$ est impair, alors
    \[\forall x\in\R\qsep x^n=a \quad\ssi\quad x=\sqrt[n]{a}\]
  \end{itemize}
\remarque Soit $n\in\Ns$.
  \begin{itemize}
  \item Si $n$ est pair
  \begin{eqnarray*}
  \forall a\in\RP, & & \p{\sqrt[n]{a}}^n=a,\\
  \forall a\in\R, & & \sqrt[n]{a^n}=\abs{a}.
  \end{eqnarray*}
  \item Si $n$ est impair
  \begin{eqnarray*}
  \forall a\in\R, & & \p{\sqrt[n]{a}}^n=a,\\
                  & & \sqrt[n]{a^n}=a.
  \end{eqnarray*}
  \end{itemize}
\remarque Soit $n\in\Ns$.
  \begin{itemize}
  \item Si $n$ est pair
  \[\forall x,y\in\R\qsep x^n\leq y^n \quad\ssi\quad \abs{x}\leq\abs{y}.\]
  \item Si $n$ est impair
  \[\forall x,y\in\R\qsep x^n\leq y^n \quad\ssi\quad x\leq y.\]
  \end{itemize}
\end{remarques}

\begin{proposition}
Soit $n\in\Ns$.
\begin{itemize}
\item Si $n$ est pair
  \[\forall a,b\in\RP\qsep \sqrt[n]{ab}=\sqrt[n]{a}\sqrt[n]{b}\]
\item Si $n$ est impair
  \[\forall a,b\in\R\qsep \sqrt[n]{ab}=\sqrt[n]{a}\sqrt[n]{b}\]
\end{itemize}
\end{proposition}

\subsection{Partie entière, approximation}

\begin{proposition}
$\R$ possède la propriété
\[\forall x\in\R \qsep \forall \epsilon>0 \qsep \exists n\in\N \qsep n\epsilon \geq x.\]
On dit que $\R$ est \emph{archimédien}.
% \[\forall x\in\R \qsep \exists n\in\N \qsep x \leq n.\]
\end{proposition}

\begin{preuve}
ADMIS

\end{preuve}

\begin{remarqueUnique}
\remarque En particulier, si on note $x$ le volume d'eau de l'océan et $\epsilon$ le volume
que peut contenir une petite cuillère, l'archimédisme de $\R$ nous permet de montrer qu'une
personne (patiente) arrivera à vider l'océan à l'aide de cette petite cuillère.
\end{remarqueUnique}


% \begin{theoreme}
% On a~:
% \[\forall \epsilon>0 \quad \forall x \in\R \quad \exists n\in\N
%   \quad n\epsilon\geq x\]  
% \end{theoreme}

% \begin{remarques}
% \remarque Cette proposition permet de prouver que l'on peut (théoriquement)
%   vider l'océan avec une petite cuillère.
% \end{remarques}

\begin{definition}[utile=-3]
Soit $x\in\R$. Il existe un unique $n\in\Z$ tel que
\[n\leq x<n+1.\]
Cet entier est appelé \emph{partie entière} de $x$ et est noté $\ent{x}$.
\end{definition}

\begin{preuve}
$\quad$
\begin{itemize}
\item \textbf{Unicité~:}
  Soit $x\in\R$. On se donne $n_1,n_2\in\Z$ tels que~:
  \[n_1\leq x<n_1+1 \et n_2\leq x<n_2+1\]
  Alors $x-1<n_1\leq x$ et $x-1<n_2\leq x$ donc $-x\leq -n_2<-\p{x-1}$. En
  sommant ces deux inégalités, on obtient $-1<n_1-n_2<1$. Comme $n_1-n_2\in\Z$,
  on en déduit que $n_1-n_2=0$ donc que $n_1=n_2$.
\item \textbf{Existence~:}
  Soit $x\in\R$. On suppose que $x\geq 0$ et on pose
  \[A=\enstq{n\in\N}{n\leq x}\]
  Alors $A$ est une partie non vide de $\N$ car $0\in A$. De plus, $A$ est
  majorée.
  En effet, par archimédisme de $\R$, il existe $m\in\N$ tel que $m\geq x$. $A$
  est alors majoré par $m$. Puisque $A$ est une partie non vide majorée de $\N$,
  elle admet un plus grand élément $n$. Comme $n\in A$, on en déduit que
  $n\leq x$ et comme $n+1\not\in A$ ($n$ étant le plus grand élément de $A$), on
  en déduit que $x<n+1$.\\
  Si $x<0$, alors $-x\geq 0$. Il existe donc $n\in\N$ tel que
  $n\leq -x<n+1$, donc $-\p{n+1}<x\leq -\p{n+1}+1$. Si $x=-n$, alors $-n$
  convient. Sinon, $-\p{n+1}$ convient.
\end{itemize}
\end{preuve}

\begin{remarques}
  \remarque Si $a\in\Z$ et $b\in\Ns$, $\ent{a/b}$ est le quotient de la division euclidienne de $a$ par $b$.
  \remarque Soit $a>0$ et $x\in\R$. Alors, il existe un unique $n\in\Z$
    tel que $na\leq x < (n+1)a$.
  \remarque  On définit de même la partie entière supérieure de $x\in\R$,
    notée $\lceil x\rceil$, comme l'unique $n\in\Z$ tel que $n-1<x\leq n$. Si $x$
    est entier, alors $\lfloor x\rfloor=\lceil x\rceil=x$. Sinon,
    $\lceil x\rceil=\lfloor x\rfloor+1$.
\end{remarques}

\begin{exos}
\exo Calculer $\ent{x}+\ent{-x}$ pour tout $x\in\R$.
\begin{sol}
Si $x$ est un entier, il existe $n\in \Z$ tel que $x=n=\ent{x}$. Alors $-x=-n=\ent{-x}$ donc la somme cherchée fait $0$.
Sinon, $\ent{x}<x<\ent{x}+1$ donc $-\ent{x}-1<-x<-\ent{x}$ et $-x$ ce qu'on peut réécrire $-\ent{x}-1\leq -x<\p{-\ent{x}-1}+1$ d'où $\ent{-x}=-\ent{x}-1$ donc la somme recherchée vaut $-1$.
\end{sol}
\exo Montrer que la partie entière est une fonction croissante.
\begin{sol}
Soit $x\leq y$. Raisonnons par l'absurde et supposons que $\ent{y}<\ent{x}$ alors $\ent{y}+1\leq \ent{x}$ car ce sont des entiers, d'où $y<\ent{y}+1\leq \ent{x}\leq x$. CONTRADICTION.
\end{sol}
\exo Soit $\alpha\in\interfo{0}{1}$. Montrer qu'il existe un unique
  $n\in\Ns$ tel que
  \[\frac{n-1}{n}\leq \alpha<\frac{n}{n+1}.\]
  \begin{sol}
$$\frac{n-1}{n}\leq \alpha \Longleftrightarrow n(\alpha-1)\geq -1 \Longleftrightarrow n\leq \frac{1}{1-\alpha} $$
et $$\alpha<\frac{n}{n+1} \Longleftrightarrow 1-\alpha > \frac{1}{n+1} \Longleftrightarrow \frac{1}{1-\alpha}<n+1.$$
Ainsi,

$$\frac{n-1}{n}\leq \alpha<\frac{n}{n+1} \Longleftrightarrow n\leq \frac{1}{1-\alpha}<n+1\Longleftrightarrow n=\ent{\frac{1}{1-\alpha}}.$$
  \end{sol}
\end{exos}



\begin{definition}[utile=-3]
Soit $a\in\R$ et $\epsilon>0$. On appelle \emph{valeur approchée} de $a$ à la précision
$\epsilon$ tout réel $b$ tel que $\abs{a-b}\leq\epsilon$.
Si $b\leq a$ (respectivement $b\geq a$), on dit que $b$ est une valeur approchée
de $a$ par \emph{défaut} (respectivement, par \emph{excès}).
\end{definition}

\begin{remarques}
\remarque On note $\Q$ l'ensemble des nombres rationnels et $\R\setminus\Q$ l'ensemble
  des nombres irrationnels. $\Q$ est stable par les opérations usuelles d'addition, de
	soustraction, de multiplication et de division.
\end{remarques}

\begin{definition}[utile=-3]
On dit qu'un réel $a$ est \emph{décimal} lorsqu'il existe $m\in\Z$ et $n\in\N$ tels
que
\[a=m\cdot 10^{-n}.\]
\end{definition}

\begin{remarques}
\remarque Un nombre décimal est rationnel. Cependant $1/3$ est rationnel, mais
  n'est pas décimal.
\remarque L'ensemble $\mathcal{D}$ des nombres décimaux est stable par les
  opérations d'addition, de soustraction, de multiplication, mais pas par
  division.
\end{remarques}

\begin{proposition}
Soit $a\in\R$ et $n\in\N$. Alors, $d=\ent{10^n a}\cdot 10^{-n}\in\mathcal{D}$
est une approximation par défaut de $a$ à la précision $10^{-n}$.
\end{proposition}

\begin{preuve}
$\dfrac{\ent{10^na}}{10^n}\leq a< \dfrac{\ent{10^na}}{10^n}+\dfrac{1}{10^n}$.
\end{preuve}
% \begin{proposition}
% Quels que soient $a\in\R$ et $\epsilon>0$, il existe $r\in\Q$ tel que~:
% \[\abs{a-r}\leq\epsilon\]
%  On dit que $\Q$ est dense dans $\R$.
% \end{proposition}


% \begin{remarques}
% \remarque Autrement dit, si $a\in \R$ et $\epsilon>0$, il existe $r\in\Q$ tel
%   que $\abs{a-r}\leq\epsilon$. La pratique est plus compliquée~:
%   \begin{itemize}
%   \item Parler du développement en fractions continues. Faire un test sur $\pi$,
%     sur le nombre de jours dans l'années.
%   \item Utiliser $\pi=4 \arctan 1$ pour approximer $\pi$.
%   \end{itemize}
% \remarque Soit $a,b\in\R$ tels que $a<b$. Alors il existe $r\in\Q$ tel que
%   $a\leq r\leq b$.\\
%   On peut même avoir l'inégalité stricte si on le souhaite. Il suffit
%   de travailler avec $a'=(2a+b)/3$ et $b'=(a+2b)/3$.
% \remarque De même $\R\setminus\Q$ est dense dans $\R$.
% \end{remarques}

\subsection{Intervalle}


\begin{definition}[utile=-3]
On appelle \emph{droite numérique achevée} et on note $\Rbar$ l'ensemble $\R$ auquel
on adjoint deux éléments notés $-\infty$ et $+\infty$. On munit $\Rbar$ d'une
relation d'ordre totale en prolongeant la relation d'ordre naturelle sur $\R$
et en posant
\[\forall x\in\R \qsep -\infty<x<+\infty.\]
\end{definition}

\begin{remarqueUnique}
\remarque On prolonge aussi de manière naturelle l'addition et la
  multiplication sans toutefois définir
  $\p{+\infty}-\p{+\infty}$ et $0\times\p{\pm \infty}$.
\end{remarqueUnique}

\begin{definition}
On dit qu'une partie de $\R$ est un \emph{intervalle} lorsqu'elle est de la forme
\[\emptyset, \qquad \R, \qquad \interf{a}{b}, \qquad \intero{a}{b}, \qquad 
  \interfo{a}{b}, \qquad \interof{a}{b},\]
\[\interfo{a}{+\infty}, \qquad \intero{a}{+\infty}, \qquad \interof{-\infty}{b},
  \qquad \intero{-\infty}{b}\]
où $a,b\in\R$.
\end{definition}

\begin{remarques}
\remarque En particulier, pour $a=b$, $\interf{a}{b}=\ens{a}$ est un intervalle.
  On dit qu'un intervalle est \emph{non trivial} lorsqu'il contient au moins 2 points.
\remarque Si $I$ est un intervalle non vide, il existe un unique couple $(a,b)\in\Rbar^2$ tel
  que $I=\interf{a}{b}$, $I=\interof{a}{b}$, $I=\interfo{a}{b}$ ou
  $I=\intero{a}{b}$. On dit que $a$ et $b$ sont les
  \emph{extrémités} de $I$. L'intervalle $I$ est dit \emph{ouvert} lorsqu'il ne contient pas ses extrémités
  c'est-à-dire lorsqu'il est vide, ou qu'il est de la forme $\intero{a}{b}$ où $a,b\in\Rbar$.
\remarque Dans ce cours, une partie de $\R$ notée $I$ ou $J$ sera implicitement un intervalle.
\end{remarques}

\section{Fonction réelle d'une variable réelle}
\subsection{Définition}
\begin{definition}[utile=-3]
On appelle \emph{fonction réelle} toute fonction définie sur une partie $\dom$ de $\R$,
à valeurs dans $\R$.
\end{definition}

\begin{remarques}
\remarque Il sera essentiel de ne pas confondre une fonction avec son
  expression.
  Par exemple parler de la fonction $\sin x$ est une erreur grave. On
  parlera plutôt de la fonction définie sur $\R$ qui au réel $x$ associe le
  réel $\sin x$.
\remarque Par abus de langage, il est courant que les énoncés demandent à
  l'élève de donner le domaine de définition d'une fonction donnée par une
  expression (par exemple $\sqrt{x})$. Dans ce cas, il faut donner l'ensemble
  $\dom$ des $x$ pour lesquels cette expression à un sens (ici, $\RP$). La
  fonction $f$ sera alors la fonction de
  $\dom$ dans $\R$, qui à $x$ associe cette expression en $x$.
\end{remarques}

\begin{exoUnique}
\exo Déterminer le domaine de définition de la fonction d'expression
  \[f(x)\defeq\ln\p{x+\sqrt{x^2-1}}.\]
\end{exoUnique}
\begin{sol}
$x\notin]-1;1[$ pour la racine puis si $x>0$, ok. Pour $x\leq -1$, $x+\sqrt{x^2-1}\geq 0 \Longleftrightarrow x^2-1\geq x^2$. Donc $D=[1;+\infty[$.
\end{sol}

\begin{definition}[utile=-3]
Soit $f$ et $g$ deux fonctions définies sur $\dom$.
\begin{itemize}
\item Pour tout $\lambda,\mu\in\R$, on définit la fonction $\lambda f+\mu g$
  par
  \[\forall x\in\dom \qsep \p{\lambda f+\mu g}(x)\defeq
    \lambda f(x)+\mu g(x).\]
\item On définit la fonction $fg$ par
  \[\forall x\in\dom \qsep (fg)(x)\defeq f(x)g(x).\]
\item Si $f$ ne s'annule en aucun point de $\dom$, on définit $1/f$ par
  \[\forall x\in\dom \qsep \p{\frac{1}{f}}(x)\defeq\frac{1}{f(x)}.\]
\end{itemize}
\end{definition}



\subsection{Symétries}
\begin{definition}[utile=-3]
  Soit $f$ une fonction définie sur un domaine $\dom$ \emph{symétrique par rapport
  à 0}, c'est-à-dire tel que
\[\forall x\in\mathcal{D} \qsep -x\in\mathcal{D}.\]
  On dit que
  \begin{itemize}
  \item $f$ est \emph{paire} lorsque
    $$\forall x\in\dom \qsep f(-x)=f(x).$$
  \item $f$ est \emph{impaire} lorsque
    $$\forall x\in\dom \qsep f(-x)=-f(x).$$
  \end{itemize}
\end{definition}

\begin{remarques}
\remarque Si $f$ est paire, la droite $\p{Oy}$ est un axe de symétrie du graphe de
  $f$.
  \begin{center}
\begin{pdfpic}
  \readdata{\listeP}{graph/graphe_dom_paire.txt}
  \psset{xunit=0.5cm,yunit=1.5cm}
  \begin{pspicture}(-11,-0.4)(11,1.2)
  \psaxes[labels=none,ticks=none]{->}(0,0)(-11,-0.4)(11,1.2)
  \dataplot[plotstyle=curve,linewidth=2pt]{\listeP}
  \uput[r](11,0){$x$}
  \uput[r](0,1.2){$y$}
  \uput[ur](0,0){$O$}
  \end{pspicture}
\end{pdfpic}
  \end{center}
\remarque Si $f$ est impaire, $O$ est un centre de symétrie du graphe de $f$.
  \begin{center}
\begin{pdfpic}
  \readdata{\listeP}{graph/graphe_dom_impaire.txt}
  \psset{xunit=0.5cm,yunit=0.9cm}
  \begin{pspicture}(-11,-1.8)(11,1.8)
  \psaxes[labels=none,ticks=none]{->}(0,0)(-11,-1.8)(11,1.8)
  \dataplot[plotstyle=curve,linewidth=2pt]{\listeP}
  \uput[r](11,0){$x$}
  \uput[r](0,1.8){$y$}
  \uput[dr](0,0){$O$}
  \end{pspicture}
\end{pdfpic}
  \end{center}
\remarque Si $f$ est paire ou impaire, pour étudier $f$,  il suffit d'étudier sa restriction à $\mathcal{D}\cap\RP$.
\end{remarques}

\begin{exoUnique}
\exo Montrer que la fonction d'expression
  \[\ln\p{x+\sqrt{x^2+1}}\]
  est impaire.
\end{exoUnique}

\begin{sol} Déjà, elle est bien définie sur $\R$. De plus $\ln\p{-x+\sqrt{x^2+1}}+\ln\p{x+\sqrt{x^2+1}}=\ln(1)=0$ donc $f(-x)=-f(x) \forall x \in \R$.
\end{sol}

\begin{definition}[utile=-3]
Soit $T\in\R$ et $f$ une fonction dont le domaine de définition vérifie
\[\forall x\in\mathcal{D} \qsep x+T\in\mathcal{D} \et x-T\in\mathcal{D}\]
On dit que $f$ est \emph{$T$-périodique}, ou que $T$ est une \emph{période} de $f$, lorsque
\[\forall x\in\mathcal{D} \qsep f\p{x+T}=f(x).\]
Lorsque $f$ admet une période non nulle, on dit que $f$ est \emph{périodique}.
\end{definition}

\begin{remarques}
\remarque Si $f$ est $T$-périodique, alors
  \[\forall x\in\dom\qsep \forall k\in\Z\qsep f(x+kT)=f(x).\]
\remarque Si $f$ est $T$-périodique, la translation de vecteur $T\ve{e_1}$ laisse
  stable le graphe de $f$.
  \begin{center}
\begin{pdfpic}
  \readdata{\listeP}{graph/graphe_dom_periodique.txt}
  \psset{xunit=0.7cm,yunit=1.0cm}
  \begin{pspicture}(-8.5,-1.8)(8.5,1.8)
  \psaxes[labels=none,ticks=none]{->}(0,0)(-8.5,-1.8)(8.5,1.8)
  \dataplot[plotstyle=curve,linewidth=2pt]{\listeP}
  \uput[r](8.5,0){$x$}
  \uput[r](0,1.8){$y$}
  \uput[dr](0,0){$O$}
  \psline{<->}(-3.1415,0.2)(3.1415,0.2)
  \uput[u](-1.5,0.2){$T$}
  \end{pspicture}
\end{pdfpic}
  \end{center}
  Pour étudier $f$, il suffit d'étudier sa restriction à
  $\mathcal{D}\cap[a,a+T]$ pour un certain $a\in\R$.
\remarque S'il existe $T\in\R$ tel que $f\p{T-x}=f(x)$, la droite d'équation $x=T/2$
  est un axe de symétrie du graphe de $f$.
  \begin{preuve}
  $f(T/2+x)=f(T-(T/2+x))=f(T/2-x).$
  \end{preuve}
  \begin{center}
\begin{pdfpic}
  \readdata{\listeP}{graph/graphe_dom_sym.txt}
  \psset{xunit=0.8cm,yunit=1.3cm}
  \begin{pspicture}(-7,-1)(7,1)
  \psaxes[labels=none,ticks=none]{->}(0,0)(-7,-1)(7,1)
  \dataplot[plotstyle=curve,linewidth=2pt]{\listeP}
  \uput[r](7,0){$x$}
  \uput[l](0,1){$y$}
  \uput[dr](0,0){$O$}
  \psline[linestyle=dashed](0.785,-1)(0.785,1)
  \uput[r](0.785,1){$x=T/2$}
  \end{pspicture}
\end{pdfpic}
  \end{center}
\end{remarques}

\begin{exos}
\exo La fonction
  \[f(x)\defeq \sin(2x)+\cos\p{\frac{x}{3}}\]
  est-elle périodique~?
\begin{sol}
Elle est bien $6\pi$-périodique.
\end{sol}
\exo Montrer que le graphe de la fonction
  \[f(x)\defeq\ln\p{x^2+x+1}\]
  admet un axe de symétrie.
  \begin{sol}
$f(-1-x)=f(x)$ donc $x=-1/2$ est axe de symétrie.
\end{sol}
\exo Tracer le graphe d'une fonction quelconque $f$, puis celui des fonctions
 \[x\mapsto f(x)+a, \qquad x\mapsto f(x+a),  \qquad x\mapsto f(a-x), \qquad x\mapsto f(a x), \qquad x\mapsto a f(x).\]
\end{exos}


\begin{proposition}
Soit $A$ et $B$ deux parties de $\R$ et $f$ une bijection de $A$ dans
$B$. Alors le graphe de $f^{-1}$ est le symétrique du graphe de $f$ par
rapport à la première bissectrice des axes $[Ox)$ et $[Oy)$.
\end{proposition}
\begin{preuve}
$\begin{cases}x\in A\\ y=f(x)\end{cases}\Longleftrightarrow \begin{cases}y\in B\\ x=f^{-1}(y)\end{cases}$.
Donc l'ensemble des $(y,f^{-1}(y))$ est l'ensemble des $(f(x),x)$.
Or dans un repère orthonormé, le point de coordonnées $(b,a)$ est le symétrique par rapport à la première bissectrice des axes de ce repère, du point de coordonnées $(a,b)$, d'où le résultat.
\end{preuve}

\subsection{Monotonie}

\begin{definition}[utile=-3]
  Soit $f$ une fonction définie sur $\dom$. On dit que
  \begin{itemize}
  \item $f$ est \emph{croissante} lorsque
    $$\forall x,y\in\dom \qsep x\leq y \implique f(x)\leq f(y).$$
  \item $f$ est \emph{décroissante} lorsque
    $$\forall x,y\in\dom \qsep x\leq y \implique f(x)\geq f(y).$$
  \item $f$ est \emph{strictement croissante} lorsque
    $$\forall x,y\in\dom \qsep x< y \implique f(x)< f(y).$$
  \item $f$ est \emph{strictement décroissante} lorsque
    $$\forall x,y\in\dom \qsep x< y \implique f(x)> f(y).$$
  \end{itemize}
\end{definition}

\begin{remarques}
\remarque Les fonctions constantes sont les seules fonctions qui sont à la fois croissantes et décroissantes.
\remarque Une fonction peut n'être ni croissante, ni décroissante. C'est le cas de la fonction $x\mapsto x^2$ sur $\R$.
% \remarque Une fonction strictement croissante est croissante. La réciproque est fausse, puisqu'une fonction constante est croissante mais pas strictement croissante.
\remarque Si $f$ est strictement monotone, elle est injective.
\remarque Attention, il est possible que $f$ soit croissante sans que
  \[\forall x,y\in\dom \qsep f(x)\leq f(y) \implique x\leq y.\]
  C'est notamment le cas des fonctions constantes qui sont croissantes mais
  pour lesquelles $f(x)\leq f(y)$ quelle que soit la position de $x$ par rapport à $y$.
  Cependant, si $f$ est strictement croissante, par contraposée, on a bien
  \[\forall x,y\in\dom \qsep f(x)\leq f(y) \implique x\leq y.\]
% Si $f$ est strictement croissante alors
%   \[\forall x,y\in\dom \qsep f(x)<f(y) \ssi x< y.\]
%   De même 
%   \[\forall x,y\in\dom \qsep f(x)\leq f(y) \ssi x\leq y.\]
%   Attention, si $f$ est seulement croissante et si $f(x)\leq f(y)$, on ne peut pas en conclure que $x\leq y$. 
\remarque D'après la remarque précédente, si $f$ est strictement croissante et $x_0\in\dom$ est un zéro de $f$, pour placer $x\in\dom$ par rapport à $x_0$, il suffit de déterminer le signe de $f(x)$.
% Par exemple, supposons que $f$ est strictement croissante sur $\R$ et que $x_0\in\R$ est tel que $f(x_0)=0$. Alors, étant donné $x\in\R$, si $f(x)\geq 0$ alors $x\geq x_0$. Si par contre $f(x)\leq 0$, alors $x\leq x_0$.
%   Attention, si on a seulement croissance de $f$, il n'est pas possible de placer $x$ par rapport à $x_0$ à partir du signe de $f(x)$.
  % La réciproque est fausse comme le montre l'exemple de la fonction
  % \[\dspappli{f}{\R}{\R}{x}{\begin{cases}\frac{1}{x}&\text{si $x\neq 0$}\\ 0&\text{si $x=0$.}\end{cases}}\]
  % Cependant, nous verrons que si $f$ est définie sur un intervalle et injective, alors elle est strictement monotone.
  \remarque
Les effets des opérations usuelles sur les propriétés de monotonie sont résumés
dans les tableaux ci-dessous.
\begin{itemize}
\item \emph{Combinaison linéaire positive}
  \begin{center}
  \begin{tabular}{|c|c|c|}
  \hline
  \backslashbox{f}{g} & croissante & décroissante\\
  \hline
  croissante   & croissante & $\times$\\
  \hline
  décroissante & $\times$   & décroissante\\
  \hline
  \end{tabular}
  \end{center}
\item \emph{Produit de fonctions positives}
  \begin{center}
  \begin{tabular}{|c|c|c|}
  \hline
  \backslashbox{f}{g} & croissante & décroissante\\
  \hline
  croissante   & croissante & $\times$\\
  \hline
  décroissante & $\times$   & décroissante\\
  \hline
  \end{tabular}
  \end{center}
\item \emph{Inverse d'une fonction strictement positive ou strictement négative}
  \begin{center}
  \begin{tabular}{|c|c|c|}
  \hline
  f        & croissante   & décroissante\\
  \hline
  1/f      & décroissante & croissante\\
  \hline
  \end{tabular}
  \end{center}
\item \emph{Composition}
  \begin{center}
  \begin{tabular}{|c|c|c|}
  \hline
  \backslashbox{f}{g} & croissante   & décroissante\\
  \hline
  croissante   & croissante   & décroissante\\
  \hline
  décroissante & décroissante & croissante\\
  \hline
  \end{tabular}
  \end{center}
\end{itemize}
\vspace{1ex}
  Lorsque c'est possible, il est souvent bien plus judicieux de
  déterminer la monotonie d'une fonction à partir de ces règles plutôt qu'à
  partir de l'étude du signe de la dérivée. En effet, cette méthode est bien
  plus rapide et source de beaucoup moins d'erreurs.
\end{remarques}


% \begin{definition}[utile=-3]
%   Soit $f$ et $g$ deux fonctions définies respectivement sur $\dom$ et $\dom'$.
%   On suppose que pour tout $x\in\dom$, $f(x)\in\dom'$. On définit alors la
%   fonction $g\circ f$ par
%   $$\forall x\in\dom \qsep \p{g\circ f}(x)\defeq g\p{f(x)}.$$
% \end{definition}

% \subsection{Opérations usuelles, symétries et monotonie}

% Les effets des opérations usuelles sur les propriétés de symétries sont
% résumés dans les tableaux ci-dessous.
% \begin{itemize}
% \item \emph{Combinaison linéaire}
%   \begin{center}
%   \begin{tabular}{|c|c|c|}
%   \hline
%   \backslashbox{f}{g} & paire    & impaire\\
%   \hline
%   paire   & paire    & $\times$\\
%   \hline
%   impaire & $\times$ & impaire\\
%   \hline
%   \end{tabular}
%   \end{center}
% \item \emph{Produit}
%   \begin{center}
%   \begin{tabular}{|c|c|c|}
%   \hline
%   \backslashbox{f}{g} & paire    & impaire\\
%   \hline
%   paire   & paire    & impaire\\
%   \hline
%   impaire & impaire  & paire\\
%   \hline
%   \end{tabular}
%   \end{center}
% \item \emph{Inverse}
%   \begin{center}
%   \begin{tabular}{|c|c|c|}
%   \hline
%   f        & paire  & impaire\\
%   \hline
%   1/f      & paire  & impaire\\
%   \hline
%   \end{tabular}
%   \end{center}
% \item \emph{Composition $g\circ f$}
%   \begin{center}
%   \begin{tabular}{|c|c|c|}
%   \hline
%   \backslashbox{g}{f}  & paire  & impaire\\
%   \hline
%   paire    & paire  & paire\\
%   \hline
%   impaire  & paire  & impaire\\
%   \hline
%   $\times$ & paire  & $\times$\\
%   \hline
%   \end{tabular}
%   \end{center}
% \end{itemize}

\begin{exos}
  \exo Montrer que la fonction
  \[\dspappli{f}{\Rs}{\R}{x}{\frac{1}{x}}\]
  n'est ni croissante, ni décroissante.
\exo Déterminer la monotonie des fonctions d'expressions
  \[\frac{1}{\e^x+\sqrt{1+x}}, \qquad \sqrt{x+1}-\sqrt{x}.\]
  \begin{sol}
Les fonctions $x\mapsto \e^{x}$ et $x \mapsto \sqrt{1+x}$ sont croissantes donc $ x\mapsto \e^x+\sqrt{1+x}$ est croissante et strictement positive, puis  par passage à l'inverse $ x \mapsto \frac{1}{\e^x+\sqrt{1+x}}$ est décroissante.\\
Soit $f$ la fonction $ x \mapsto \sqrt{x+1}-\sqrt{x}$. Écrit sous cette forme, il n'est pas immédiat de savoir si $f$ est monotone. Cependant, en multipliant $f(x)$ par son expression conjuguée, on obtient
  $$\forall x\in\RP \qsep f(x)=\frac{
    \p{\sqrt{x+1}-\sqrt{x}}\p{\sqrt{x+1}+\sqrt{x}}}{\sqrt{x+1}+\sqrt{x}}
    =\frac{1}{\sqrt{x+1}+\sqrt{x}}$$
Par somme puis inverse de fonctions croissantes positives, $f$ est donc décroissante.
  \end{sol}
\end{exos}

\subsection{Fonction majorée, minorée, bornée}

% \begin{definition}[utile=-3]
% On dit qu'une fonction réelle $f$ est~: 
% \begin{itemize}
% \item croissante lorsque
%   \[\forall x_1,x_2\in \mathcal{D} \qsep x_1 \leq x_2 \implique
%     f\p{x_1}\leq f\p{x_2}.\]
% \item décroissante lorsque
%   \[\forall x_1,x_2\in \mathcal{D} \qsep x_1 \leq x_2 \implique
%     f\p{x_1}\geq f\p{x_2}.\]
% \item monotone lorsqu'elle est croissante ou décroissante.
% \item strictement croissante lorsque
%   \[\forall x_1,x_2\in \mathcal{D} \qsep x_1 < x_2 \implique
%     f\p{x_1} < f\p{x_2}.\]
% \item strictement décroissante lorsque
%   \[\forall x_1,x_2\in \mathcal{D} \qsep x_1 < x_2 \implique
%     f\p{x_2} > f\p{x_1}.\]
% \item strictement monotone lorsqu'elle est strictement croissante ou
%   strictement décroissante.
% \end{itemize}
% \end{definition}

% \begin{remarques}
% \remarque Une fonction réelle est strictement monotone si et seulement si elle
%   est monotone et injective.
% \remarque Si $f$ est une fonction réelle strictement croissante, alors
%   \[\forall x_1,x_2\in \mathcal{D} \qsep f\p{x_1} \leq f\p{x_2}
%     \implique x_1 \leq x_2.\]
%   % Cette propriété caractérise d'ailleurs les fonctions strictement
%   % croissantes.
% \end{remarques}

% \begin{proposition}[utile=-3]
% 
% \begin{itemize}
% \item La somme de deux fonctions croissantes (resp. décroissantes) est
%   croissante (resp. décroissante).
% \item Le produit de deux fonctions croissantes (resp. décroissantes) positives est
%   croissant (resp. décroissant).
% \item L'inverse d'une fonction croissante (resp. décroissante) de signe constant
%   est décroissante (resp. croissante).
% \item La composée de deux fonctions monotone est monotone; elle est croissante
%   si les fonctions ont même sens de variations et décroissante si leurs sens
%   de variation sont opposés.
% \item La bijection réciproque d'une fonction strictement croissante (resp.
%   décroissante) est strictement croissante (resp. décroissante).
% \end{itemize}
% \end{proposition}

% \begin{remarqueUnique}
% \remarque Une fonction $f$ est décroissante si et seulement si $-f$ est croissante.
%   En utilisant la proposition précédente, on en déduit par exemple que le produit
%   d'une fonction croissante positive et d'une fonction décroissante négative est
%   décroissante négative.
% \end{remarqueUnique}

% \begin{exoUnique}
% \exo Étudier la monotonie de la fonction
%   $x\mapsto \sin\p{\p{\e^{-x}-1}\pi/2}$ sur $\RP$.
% % \exo $x\mapsto \sqrt{x+1}-\sqrt{x}$ est décroissante sur $\RP$.
% % \exo On ne peut pas obtenir la monotonie de
% %   $x\mapsto x \ln\p{1-1/x}$ sur $\intero{1}{+\infty}$.
% %   C'est le produit de deux fonction croissantes mais on ne peut rien en dire même si
% %   la première est positive. On devra donc dériver la fonction.
% \end{exoUnique}

\begin{definition}[utile=-3]
On dit qu'une fonction réelle $f$ est
\begin{itemize}
\item \emph{majorée} lorsque
  \[\exists M\in\R \qsep \forall x\in\mathcal{D} \qsep
    f(x)\leq M.\]
\item \emph{minorée} lorsque
  \[\exists m\in\R \qsep \forall x\in\mathcal{D} \qsep
    f(x)\geq m.\]
\end{itemize}
\end{definition}

\begin{exoUnique}
\exo Montrer que la fonction d'expression $x \e^{-x}$ est majorée sur $\R$.
% \exo La fonction d'expression $x+1/x$ est minorée par 2 sur $\RPs$.
\end{exoUnique}

\begin{sol}
Dérivée, tableau de variations puis majorée par $f(1)=\dfrac{1}{e}$.
\end{sol}

%% Pour la seconde, penser à l'inégalité arithmetico-géométrique

\begin{definition}[utile=-3]
On dit qu'une fonction réelle ou complexe $f$ est \emph{bornée} lorsque
\[\exists M\in\RP \qsep \forall x\in\mathcal{D} \qsep
  \abs{f(x)}\leq M.\]
\end{definition}

\begin{exoUnique}
\exo Montrer que la fonction d'expression $\frac{x}{1+x^2}$ est bornée par
  $1/2$ sur $\R$.
\end{exoUnique}

\begin{sol}
Déjà, $\frac{x}{1+x^2} \leq \frac{1}{2} \Longleftrightarrow (x-1)^2\geq 0$, ce qui est vrai.
Donc pour $x\geq 0$, on a $0\leq \abs{\frac{x}{1+x^2}} \leq 1/2$. Et si $x<0$, $0\leq \abs{\frac{x}{1+x^2}}=\frac{-x}{1+(-x)^2}\leq 1/2$.
\end{sol}

\begin{proposition}
Une fonction réelle est bornée si et seulement si elle est minorée et majorée.
\end{proposition}

\begin{preuve}
Pour le sens droite-gauche, on la borne par $\abs{m}+\abs{M}$.
\end{preuve}

\begin{definition}[utile=-3]
Soit $f$ et $g$ deux fonctions de domaine $\mathcal{D}$.
On dit que $f$ est inférieure à $g$ et on note $f\leq g$ lorsque
\[\forall x\in\mathcal{D} \qsep f(x)\leq g(x).\]
\end{definition}

\begin{remarques}
\remarque La relation $\leq$ est une relation d'ordre sur
  $\mathcal{F}\p{\mathcal{D},\R}$. Elle n'est pas totale.
\remarque La négation de $f\leq g$ s'écrit
  \[\exists x\in\mathcal{D} \qsep f(x)> g(x).\]
\end{remarques}


\section{Fonction continue, fonction dérivable}

\subsection{Limite}

Dans ce chapitre, on ne définira pas précisément la notion de limite. On se basera sur la
définition intuitive suivante.

\begin{definition}[utile=-3]
Étant donné une fonction $f$ et $a,l\in\Rb$, on dit que $f(x)$ tend vers
$l$ lorsque $x$ tend vers $a$, lorsque, quitte à rendre $x$ proche de $a$, on
peut rendre $f(x)$ aussi proche que l'on souhaite de $l$. Dans ce cas, on
note
$$f(x)\tendvers{x}{a}l.$$
\end{definition}

\begin{proposition}[utile=-3]
Soit $a\in\Rbar$ et $f, g$ deux fonctions telles que
$f(x)$ et $g(x)$ tendent respectivement vers $l_f$ et $l_g\in\R$ lorsque $x$
tend vers $a$. Alors
\begin{itemize}
\item Si $\lambda$ et $\mu$ sont deux réels
  \[\lambda f(x)+\mu g(x)\tendvers{x}{a} \lambda l_f+\mu l_g.\]
\item On a
  \[f(x)g(x) \tendvers{x}{a} l_f l_g.\]
\item Si $l_f\not=0$
  \[\frac{1}{f(x)}\tendvers{x}{a}\frac{1}{l_f}.\]
\item Plus généralement, si $l_g\not=0$
  \[\frac{f(x)}{g(x)}\tendvers{x}{a}\frac{l_f}{l_g}.\]
\end{itemize}
\end{proposition}

\begin{proposition}[utile=-3]
Soit $f$ et $g$ deux fonctions. On suppose que $f(x)$ tend vers $l_f\in\Rbar$
lorsque $x$ tend vers $a\in\Rbar$ et que $g(x)$ tend vers $l_g\in\Rbar$ lorsque $x$
tend vers $l_f$. Alors $g(f(x))$ tend vers $l_g$ lorsque $x$ tend vers $a$. 
\end{proposition}

\begin{remarqueUnique}
\remarque
De nombreuses autres règles existent mélangeant limites finies et
infinies. Elles sont résumées dans les tableaux ci-dessous où la présence
d'une croix représente une forme indéterminée.

\begin{itemize}
\item \emph{Somme}\\  
  Si $f$ et $g$ sont deux fonctions admettant respectivement pour limites
  $l_f$ et $l_g\in\Rbar$, alors $f+g$
  \begin{center}
  \begin{tabular}{|c|c|c|c|}
  \hline
  \backslashbox{$l_f$}{$l_g$}& $-\infty$ & $l_g\in\R$     & $+\infty$\\
  \hline
  $-\infty$ & $-\infty$ & $-\infty$ & $\times$\\
  \hline
  $l_f\in\R$     & $-\infty$ & $l_f+l_g$ & $+\infty$\\
  \hline
  $+\infty$ & $\times$  & $+\infty$ & $+\infty$\\
  \hline
  \end{tabular}
  \end{center}
\item \emph{Opposé}\\
  Si $f$ est une fonction admettant pour limite $l\in\Rbar$, alors $-f$
  \begin{center}
  \begin{tabular}{|c|c|c|c|}
  \hline
  $l$ & $-\infty$ & $l\in\R$ & $+\infty$\\
  \hline
      & $+\infty$ & $-l$ & $-\infty$\\
  \hline
  \end{tabular}
  \end{center}
\item \emph{Multiplication par un scalaire}\\
  Si $f$ est une fonction admettant pour limite $l\in\Rbar$ et $\lambda\in\R$,
  alors $\lambda f$
  \begin{center}
  \begin{tabular}{|c|c|c|c|}
  \hline
  \backslashbox{$\lambda$}{$l$} & $-\infty$ & $l\in\R$    & $+\infty$\\
  \hline
  $\lambda<0$ & $+\infty$ & $\lambda l$ & $-\infty$\\
  \hline
  $\lambda>0$ & $-\infty$  & $\lambda l$ & $+\infty$\\
  \hline
  \end{tabular}
  \end{center}
\item \emph{Produit}\\
  Si $f$ et $g$ sont deux fonctions admettant respectivement pour limites
  $l_f$ et $l_g\in\Rbar$, alors $fg$
  \begin{center}
  \begin{tabular}{|c|c|c|c|c|c|}
  \hline
  \backslashbox{$l_f$}{$l_g$}
      & $-\infty$ & $l_g<0$ & $0$ & $l_g>0$ & $+\infty$\\
  \hline
  $-\infty$ & $+\infty$ & $+\infty$ & $\times$ & $-\infty$ & $-\infty$\\
  \hline
  $l_f<0$ & $+\infty$ & $l_f l_g$ & $0$ & $l_f l_g$ & $-\infty$\\
  \hline
  $l_f=0$ & $\times$ & $0$ & $0$ & $0$ & $\times$\\
  \hline
  $l_f>0$ & $-\infty$ & $l_f l_g$ & $0$ & $l_f l_g$ & $+\infty$\\
  \hline
  $+\infty$ & $-\infty$ & $-\infty$ & $\times$ & $+\infty$ & $+\infty$\\
  \hline
  \end{tabular}
  \end{center}
\item \emph{Inverse}\\
  Si $f$ est une fonction admettant pour limite $l$, alors $1/f$
  \begin{center}
  \begin{tabular}{|c|c|c|c|c|c|c|c|}
  \hline
  $l$ & $-\infty$ & $l<0$ & $0^{-}$ & $0$ & $0^{+}$ & $l>0$ & $+\infty$\\
  \hline
      & $0$ & $1/l$ & $-\infty$ & $\times$ & $+\infty$ & $1/l$ & $0$\\
  \hline
  \end{tabular}
  \end{center}
\item \emph{Exponentiation}\\
  Si $f$ et $g$ sont deux fonctions admettant respectivement pour limites
  $l_f$ et $l_g$, alors $f^g$
  \begin{center}
  \begin{tabular}{|c|c|c|c|c|c|}
  \hline
  \backslashbox{$l_f$}{$l_g$}
            & $-\infty$ & $l_g<0$ & $0$ & $l_g>0$ & $+\infty$\\
  \hline
  $0$ & $+\infty$ & $+\infty$ & $\times$ & $0$ & $0$\\
  \hline
  $0<l_f<1$ & $+\infty$ & $l_f^{l_g}$ & $1$ & $l_f^{l_g}$ & $0$\\
  \hline
  $1$ & $\times$ & $1$ & $1$ & $1$ & $\times$\\
  \hline
  $1<l_f$ & $0$ & $l_f^{l_g}$ & $1$ & $l_f^{l_g}$ & $+\infty$\\
  \hline
  $+\infty$ & $0$ & $0$ & $\times$ & $+\infty$ & $+\infty$\\
  \hline
  \end{tabular}
  \end{center}
  
\end{itemize}
\end{remarqueUnique}

\begin{exoUnique}
\exo Déterminer les limites en 0 de
   $x^x$ et $x^{\frac{1}{\ln x}}$;
   en déduire que \og $0^0$ \fg est une forme indéterminée. De même, déterminer la limite en
   $+\infty$ de $(1+1/x)^x$; en déduire que \og $1^{+\infty}$\fg est une forme indéterminée.
  \begin{sol}
  En effet
  $$x^x=\e^{x\ln x}\tendvers{x}{0} 1\quad\text{car}\quad
        x\ln x\tendvers{x}{0} 0$$
  $$\text{et} \quad x^{\frac{1}{\ln x}}=\e\tendvers{x}{0} \e.$$  
  De même $1^{+\infty}$ est une forme indéterminée. En effet
  $$\p{1+\frac{1}{x}}^x=\e^{x\ln\p{1+\frac{1}{x}}}=\e^{\frac
                        {\ln\p{1+\frac{1}{x}}}{\frac{1}{x}}}
                        \tendvers{x}{+\infty} \e$$
  $$\text{et} \quad 1^x=1\tendvers{x}{+\infty} 1.$$
  \end{sol}
\end{exoUnique}

\subsection{Continuité}

\begin{definition}[utile=-3]
On dit qu'une fonction $f:\mathcal{D}\to\R$ est \emph{continue} en $x_0\in\mathcal{D}$ lorsque
\[f(x)\tendvers{x}{x_0} f\p{x_0}.\]
On dit que $f$ est continue lorsque, quel que soit $x_0\in\mathcal{D}$, $f$ est continue en
$x_0$.
\end{definition}

\begin{proposition}[utile=-3, nom={Théorèmes usuels}]
Soit $f,g:\dom\to\R$ deux fonctions continues. Alors
\begin{itemize}
\item Quels que soient $\lambda,\mu\in\R$, la fonction $\lambda f+\mu g$ est continue.
\item La fonction $fg$ est continue.
\item Si $g$ ne s'annule pas, $f/g$ est continue.
\end{itemize}
\end{proposition}

\begin{proposition}[utile=-3, nom={Théorèmes usuels}]
La composée de deux fonctions continues est continue.
\end{proposition}

\begin{theoreme}[nom={Théorème de la bijection}]
\begin{itemize}
\item Soit $f:[a,b]\to\R$ (où $a,b\in\R$ et $a\leq b$) une fonction continue, strictement croissante. Alors \[f([a,b])=[f(a),f(b)].\]
De plus $f$ réalise une bijection de $[a,b]$ sur
$[f(a),f(b)]$. Autrement dit, pour tout $y\in[f(a),f(b)]$, il existe un unique 
$x\in[a,b]$ tel que $y=f(x)$.
\item Soit $f:]a,b[\to\R$ (où $a,b\in\Rbar$ et $a<b$) une fonction continue, strictement croissante.
  On pose
  \[l_a\defeq\lim_{x\to a} f(x) \et l_b\defeq\lim_{x\to b} f(x).\]
Alors
\[f\p{\intero{a}{b}}=\intero{l_a}{l_b}.\]
De plus $f$ réalise une bijection de $]a,b[$ sur $\intero{l_a}{l_b}$. Autrement dit,
pour tout $y\in\intero{l_a}{l_b}$, il existe un unique $x\in]a,b[$ tel que
$y=f(x)$.
\end{itemize}
\end{theoreme}

\begin{remarques}
% \remarque Dans le cas où le domaine de définition  de $f$ est ouvert, l'existence des
%   limites $l_a$ et $l_b$ est garanti par le théorème de la limite monotone.
\remarque Ce théorème reste valide dans de nombreuses autres situations, par exemple
  lorsque $f$ est strictement décroissante et
  que son domaine de définition est un intervalle semi-ouvert. Par exemple, si $a\in\R$
  et $f:[a,+\infty[\to\R$ est une fonction continue, strictement décroissante, en
  posant
  \[l\defeq\lim_{x\to +\infty} f(x),\]
  alors $f([a,+\infty[)=]l,f(a)]$ et  $f$ réalise une bijection de $[a,+\infty[$ sur $]l,f(a)]$.
\remarque Ce théorème permet
  de calculer $f(A)$ lorsque $A$ est une réunion d'intervalles
  $A=I_1\cup\cdots\cup I_n$ sur lesquels $f$ est continue et strictement monotone.
  Il suffit pour cela de remarquer que
  \[f(A)=f(I_1)\cup\cdots \cup f(I_n).\]
\end{remarques}

\begin{proposition}
Soit $f:I\to\R$ une fonction continue et strictement monotone sur l'intervalle $I$.
D'après le théorème de la bijection, elle réalise une bijection de $I$ sur
l'intervalle $J\defeq f(I)$. De plus
\begin{itemize}
\item $f^{-1}$ est strictement monotone, de même sens de variation que $f$. 
\item $f^{-1}$ est continue.
% \item Si $f$ est dérivable sur $I$, on pose
%   \[A\defeq\enstq{x\in I}{f'(x)\neq 0}.\]
%   Alors $f^{-1}$ est dérivable en tout point de $f(A)$.
\end{itemize}
\end{proposition}

\begin{sol}
Ajout Victor : $\forall x \in A$, $f^{-1}(f(x))=x$ donc $f'(x)(f^{-1})'(f(x))=1$. Ainsi, $\forall u \in f(A)$, en appliquant cela en $x=f^{-1}(u)$ on obtient :
\[(f^{-1})'(u)=\frac{1}{f'(f^{-1}(u))}.\]
\end{sol}

\subsection{Dérivabilité}

\begin{definition}[utile=-3]
Soit $f:\dom\to\R$ et $x_0\in\dom$. On dit que $f$ est \emph{dérivable} en
$x_0$ lorsque
\[\frac{f(x)-f(x_0)}{x-x_0}\]
admet une limite finie lorsque $x$ tend vers $x_0$. Dans ce cas, on note $f'(x_0)$
cette limite que l'on appelle \emph{nombre dérivé} de $f$ en $x_0$.
On dit que $f$ est dérivable lorsqu'elle est dérivable en tout point de $\dom$.
\end{definition}

\begin{remarques}
\remarque Si $f$ est dérivable en $x_0$, la droite d'équation
  $y=f(x_0)+f'(x_0)(x-x_0)$
  est tangente au graphe de $f$ en $x_0$.
\remarque Lorsque
  \[\frac{f(x)-f(x_0)}{x-x_0}\tendvers{x}{x_0}\pm\infty,\]
  le graphe de $f$ admet une tangente verticale en $x_0$. Une telle fonction n'est pas dérivable en $x_0$.
\remarque On dit qu'une fonction $f$ est dérivable à gauche en $x_0$ lorsque
  l'expression
  \[\frac{f(x)-f(x_0)}{x-x_0}\]
  admet une limite finie lorsque $x$ tend vers $x_0$ par la gauche. Si tel est le
  cas, cette limite est notée $f_g'\p{x_0}$. On définit de même la notion de
  dérivabilité à droite. Une fonction est dérivable en $x_0$ si et seulement si
  elle est dérivable à gauche et à droite en $x_0$ et que
  $f_g'\p{x_0}=f_d'\p{x_0}$.
\remarque Si $f$ et $g$ sont des fonctions telles qu'en $x_0$, $f(x_0)=g(x_0)$, on ne peut rien en conclure sur
  $f'(x_0)$ et $g'(x_0)$. En particulier, il est absurde de dire que parce que
  $f(x_0)=0$, on peut en déduire que $f'(x_0)=0$.
  On dira qu'on peut dériver des identités, mais pas des égalités.
  % Par exemple, si $f$ et $g$ sont définies sur $\R$ par
  % \[\forall x\in\R\qsep f(x)\defeq 0 \et g(x)\defeq x,\]
  % alors $f(0)=g(0)$ mais $f'(0)\neq g'(0)$. Cependant, si
  % \[\forall x\in\R \qsep f(x)=g(x)\]
  % et que $f$ est dérivable, on peut en déduire que $g$ est dérivable et
  % \[\forall x\in\R \qsep f'(x)=g'(x).\]
\end{remarques}

\begin{proposition}[utile=-3]
Soit $f:\dom\to\R$ et $x_0\in\dom$. Si $f$ est dérivable en $x_0$,
alors elle est continue en $x_0$.
\end{proposition}

\begin{remarqueUnique}
\remarque La réciproque de cette proposition est fausse comme le montre
  l'exemple de la fonction $x\mapsto\abs{x}$ qui est continue en 0 mais
  qui n'est pas dérivable en 0.
\end{remarqueUnique}

\begin{exoUnique}
\exo Donner une condition nécessaire et suffisante sur $a$ et $b\in\R$
  pour que la fonction $f$ définie sur $\R$ par
  \[\forall x\in\R \qsep f(x)=
    \begin{cases}
    ax+b & \text{si $x<0$}\\
    \e^x & \text{si $x\geq 0$}.
    \end{cases}\]
  soit dérivable en 0.
  \begin{sol}
On cherche d'abord à quelles conditions sur $a$ et $b$ la fonction $f$ est continue en 0. On a
\[f(x)=ax+b\tendversgp{x}{0}b\]
et
\[f(x)=\e^x\tendversd{x}{0}1\]
donc $f$ est continue en 0 si et seulement si $b=1$. On suppose qu'on est dans ce cas. Alors $f$ est dérivable à gauche en 0 et $f_g'(0)=a$. De plus, $f$ est dérivable à droite en 0 et $f_d'(0)=1$. Donc $f$ est dérivable en 0 si et seulement si $a=1$. En conclusion, $f$ est dérivable en 0 si et seulement si $a=1$ et $b=1$. 
  \end{sol}
\end{exoUnique}

\begin{definition}[utile=-3]
Soit $f:\dom\to\R$ une fonction. On note $\dom'$ l'ensemble des $x_0\in\dom$ en
lesquels $f$ est dérivable. On définit la \emph{fonction dérivée} de $f$, notée $f'$
par
\[\dspappli{f'}{\dom'}{\R}{x}{f'(x).}\]
\end{definition}

\begin{remarqueUnique}
\remarque Les fonctions usuelles sont dérivables en tout point de leur ensemble de
  définition, excepté la fonction $x\mapsto\abs{x}$ qui n'est pas dérivable en 0 et les
  fonctions $x\mapsto\sqrt[n]{x}$ qui ne sont pas dérivables en 0 pour $n\geq 2$.
  \begin{center}
  \begin{tabular}{|c|c|c|c|}
  \hline
  $\dom$ & $f(x)$ & $\dom_{f'}$ & $f'(x)$\\
  \hline
  \hline\rule{0pt}{10pt} 
  $\R$ & $x^n \quad (n\in\N)$ & $\R$ & $
    \begin{cases}nx^{n-1}& \text{si $n\geq 1$}\\0&\text{si $n=0$}\end{cases}$\\
  \hline\rule{0pt}{10pt} 
  $\Rs$ & $x^n \quad (n\in\Z)$ & $\Rs$ & $nx^{n-1}$\\
  \hline\rule{0pt}{10pt} 
  $\RPs$ & $x^\alpha \quad (\alpha\in\R)$ & $\RPs$ & $\alpha x^{\alpha-1}$\\
  \hline\rule[-5pt]{0pt}{15pt} 
  $\RP$ & $\sqrt[n]{x}=x^{\frac{1}{n}} \quad (n\in\Ns)$ & $\RPs$ &
    $\frac{1}{n} x^{\frac{1}{n}-1}$\\
  \hline
  \hline
  $\R$ & $\e^x$ & $\R$ & $\e^x$\\
  \hline\rule[-5pt]{0pt}{15pt} 
  $\RPs$ & $\ln x$ & $\RPs$ & $\frac{1}{x}$\\
  \hline\rule[-5pt]{0pt}{15pt} 
  $\Rs$ & $\ln \abs{x}$ & $\Rs$ & $\frac{1}{x}$\\
  \hline
  \hline
  $\R$ & $\cos x$ & $\R$ & $-\sin x$\\
  \hline
  $\R$ & $\sin x$ & $\R$ & $\cos x$\\
  \hline\rule[-5pt]{0pt}{15pt} 
  $\R\setminus\p{\frac{\pi}{2}+\pi\Z}$ & $\tan x$ &
    $\R\setminus\p{\frac{\pi}{2}+\pi\Z}$ & $1+\tan^2 x=\frac{1}{\cos^2 x}$\\
  \hline\rule[-5pt]{0pt}{15pt} 
  $\R\setminus\pi\Z$ & $\cotan x$ & $\R\setminus\pi\Z$ & $-\p{1+\cotan^2 x}=
  -\frac{1}{\sin^2 x}$\\
  \hline
  \end{tabular}
  \end{center}
\end{remarqueUnique}


\begin{proposition}[utile=-3, nom={Théorèmes usuels}]
Soit $f:\dom\to\R$ et $g:\dom\to\R$ deux fonctions dérivables. Alors
\begin{itemize}
\item Quels que soient $\lambda, \mu\in\R$, la fonction
  $\lambda f+\mu g$ est dérivable et
  \[\forall x\in\dom\qsep (\lambda f+\mu g)'(x)=\lambda f'(x)+\mu g'(x)\]
\item La fonction $fg$ est dérivable et
  \[\forall x\in\dom\qsep (fg)'(x)=f'(x)g(x)+f(x)g'(x).\]
\item Si $f$ ne s'annule pas sur $\dom$, $1/f$ est dérivable et
  \[\forall x\in\dom\qsep \p{\frac{1}{f}}'(x)=-\frac{f'(x)}{f^2(x)}.\]
\item Plus généralement, si $g$ ne s'annule pas sur $\dom$, $f/g$ est dérivable
  et
  \[\forall x\in\dom\qsep \p{\frac{f}{g}}'(x)=\frac{f'(x)g(x)-f(x)g'(x)}{g^2(x)}.\]
\end{itemize}
\end{proposition}

\begin{proposition}[utile=-3, nom={Théorèmes usuels}]
Soit $f:\dom_f\to\R$ et $g:\dom_g\to\R$ deux fonctions telles que  $f(\dom_f)\subset\dom_g$. Si 
$f$ et $g$ sont dérivables, alors $g\circ f$ est dérivable et
\[\forall x\in\dom_f\qsep \p{g\circ f}'(x)=f'(x)g'\p{f(x)}.\]
\end{proposition}

\begin{remarqueUnique}
\remarque En particulier, si $f:\dom\to\R$ est une fonction dérivable et $n\in\N$, la fonction $g$ définie sur $\mathcal{D}$ par
\[\forall x\in\mathcal{D}\qsep g(x)\defeq f(x)^n\]
est dérivable et
\[\forall x\in\mathcal{D}\qsep g'(x)=n f'(x) f(x)^{n-1}.\]
% Plus généralement, si $f$ est une fonction strictement positive et $\alpha\in\R$, la fonction $h$ définie sur $\mathcal{D}$ par
% \[\forall x\in\mathcal{D}\qsep h(x)\defeq f(x)^\alpha\]
% est dérivable sur $\mathcal{D}$ et
% \[\forall x\in\mathcal{D}\qsep h'(x)=\alpha f'(x) f(x)^{\alpha-1}.\]
\remarque Si $f(x)\defeq a(x)/b(x)^\alpha$, il est bon d'écrire $f(x)$ sous la forme $f(x)=a(x)b(x)^{-\alpha}$ avant de dériver $f$.
\remarque Attention, ce n'est pas parce que les théorèmes usuels ne peuvent pas s'appliquer en un point qu'on peut en conclure que la fonction n'y est pas dérivable.
\end{remarqueUnique}


\begin{exos}
\exo Montrer que la dérivée d'une fonction paire (resp. impaire, $T$-périodique) est impaire (resp. paire, $T$-périodique).
\exo Étudier la dérivabilité et calculer les dérivées des fonctions
  d'expression
  \[\frac{x+1}{\sqrt{x^2+1}}, \qquad \p{x^3+2x+1}\e^{x^2}, \qquad \ln\sqrt{\frac{1+x}{1-x}}.\]
  % \[\e^{\sin x}, \qquad x^x.\]
  % \begin{sol}
  % Soit $f$ la fonction définie sur $\R$ par
  % $$\forall x\in\R \qsep f(x)\defeq\e^{\sin x}.$$
  % Alors, d'après les théorèmes usuels, $f$ est dérivable sur $\R$ et
  % $$\forall x\in\R \qsep f'(x)=\cos x\,\e^{\sin x}.$$
  % Soit $f$ la fonction définie sur $\RPs$ par
  % $$\forall x>0 \qsep f(x)\defeq x^x=\e^{x\ln x}.$$
  % Alors, d'après les théorèmes usuels, $f$ est dérivable sur $\RPs$ et
  % $$\forall x>0\qsep
  %   f'(x)=\p{\ln x +1}\e^{x\ln x}=\p{\ln x +1}x^x.$$
  % \end{sol}
\exo Étudier la dérivablité et calculer la dérivée de la fonction
  définie sur $\interf{0}{\pi/2}$ par
  \[\forall x\in\interf{0}{\frac{\pi}{2}} \qsep 
    f(x)\defeq\sqrt{1-\cos x}.\]
  \begin{sol}
  Une erreur courante est de penser que si les théorèmes usuels ne
  s'appliquent pas en un point, alors la fonction n'est pas dérivable en ce
  point. Cela est faux, comme le montre l'exemple de la fonction ci-dessous
  \[\dspappli{f}{\interf{0}{\frac{\pi}{2}}}{\R}{x}{\sqrt{1-\cos x}}\]
  D'après les théorèmes usuels, $f$ est continue sur son ensemble de définition
  et dérivable en tout point où $1-\cos x\not=0$, c'est-à-dire sur
  $\interof{0}{\frac{\pi}{2}}$, et
  $$\forall x\in\interof{0}{\frac{\pi}{2}} \qsep f'(x)=
    \frac{\sin x}{2\sqrt{1-\cos x}}.$$
  Donc les théorèmes usuels ne permettent pas de démontrer la dérivabilité de
  $f$ en $0$. Cependant
  \begin{equation*}
  \begin{split}
  \frac{f(0+h)-f(0)}{h} &=\frac{\sqrt{1-\cos h}}{h} =
                          \frac{\sqrt{\cos 0-\cos h}}{h}\\
                        &=\frac{\sqrt{-2\sin\frac{h}{2}\sin\frac{-h}{2}}}{h}=
                          \sqrt{2}\frac{\sin\frac{h}{2}}{h}\\
                        &=\frac{1}{\sqrt{2}}\frac{\sin\frac{h}{2}}{\frac{h}{2}}
                          \tendvers{h}{0} \frac{1}{\sqrt{2}}
  \end{split}
  \end{equation*}
  ce qui montre que $f$ est dérivable en $0$ et que
  $$f'(0)=\frac{1}{\sqrt{2}}.$$    
  \end{sol}
\end{exos}

\begin{proposition}
Soit $f:I\to\R$ une fonction définie sur un intervalle $I$, dérivable et strictement monotone.
Elle réalise donc une bijection de $I$ sur l'intervalle $J\defeq f(I)$. On pose
\[A\defeq\enstq{x\in I}{f'(x)= 0}.\]
Alors $f^{-1}$ est dérivable en tout point de $J\setminus f(A)$ et
\[\forall y\in J\setminus f(A)\qsep (f^{-1})'(y)=\frac{1}{f'(f^{-1}(y))}.\]
\end{proposition}
\begin{sol}
Ajout Victor : $\forall x \in A$, $f^{-1}(f(x))=x$ donc $f'(x)(f^{-1})'(f(x))=1$. Ainsi, $\forall u \in f(A)$, en appliquant cela en $x=f^{-1}(u)$ on obtient :
\[(f^{-1})'(u)=\frac{1}{f'(f^{-1}(u))}.\]
\end{sol}

\subsection{Dérivées successives}


\begin{definition}[utile=-3]
Si $f:\dom\to\R$ est une fonction, on définit par récurrence
la \emph{dérivée $n$-ième} de $f$ de la manière suivante
\begin{itemize}
\item On pose $f^{(0)}\defeq f$.
\item Si $n\in\N$, on définit $f^{(n+1)}$ comme étant la dérivée de $f^{(n)}$.
\end{itemize}
Si $x_0\in\mathcal{D}$, on dit que $f$ est dérivable $n$ fois en $x_0$
lorsque $f^{(n)}$ est définie en $x_0$. On dit que $f$ est dérivable $n$ fois lorsqu'elle
est dérivable $n$ fois en tout point de son domaine de définition.
\end{definition}
 
\begin{proposition}[utile=-3]
Soit $f:\dom\to\R$ et $g:\dom\to\R$ deux fonctions dérivables $n$ fois.
\begin{itemize}
\item Soit $\lambda,\mu\in\R$. Alors $\lambda f+\mu g$ est dérivable $n$
  fois et
  \[\forall x\in\mathcal{D} \qsep
    \p{\lambda f+\mu g}^{(n)}(x)=\lambda f^{(n)}(x)+\mu g^{(n)}(x).\]
\item $fg$ est dérivable $n$ fois.
\item Si $g$ ne s'annule pas, alors $f/g$ est dérivable $n$ fois.
\end{itemize}
\end{proposition}

\begin{preuve}
\begin{itemize}
\item Se montre aisément par récurrence.
\item Démonstration par récurrence similaire au binôme de Newton.
\item Soit $\mathcal{D}$ une partie de $\R$. On montre $\mathcal{H}_n$ : "Si $f$ et $g$ sont dérivables $n$ fois sur $\mathcal{D}$ et $g$ ne s'annule pas sur $\mathcal{D}$, alors $f/g$ est dérivable $n$ fois sur $\mathcal{D}$".
Pour passer à l'hérédité, on applique $\mathcal{H}_n$ à $f'g-g'f$ et $g^2$, leur rapport est dérivable $n$ fois donc $(f/g)'$ est dérivable $n$ fois.
\end{itemize}
\end{preuve}

\begin{proposition}[utile=-3]
Soit $f:\dom_f\to\R$ et $g:\dom_g\to\R$ deux fonctions telles que
$f(\dom_f)\subset\dom_g$. Si $f$ et $g$ sont dérivables $n$ fois, alors $g\circ f$ est
dérivable $n$ fois.  
\end{proposition}

\begin{definition}
On dit qu'une fonction $f:\dom\to\R$ est de classe $\classec{1}$ lorsqu'elle est
dérivable et que sa dérivée est continue.
\end{definition}

\begin{proposition}[utile=-3]
Soit $f:\dom\to\R$ et $g:\dom\to\R$ deux fonctions de classe $\classec{1}$.
\begin{itemize}
\item Soit $\lambda,\mu\in\R$. Alors $\lambda f+\mu g$ est de classe $\classec{1}$.
\item $fg$ est de classe $\classec{1}$.
\item Si $g$ ne s'annule pas, alors $f/g$ est de classe $\classec{1}$.
\end{itemize}
\end{proposition}

\begin{proposition}[utile=-3]
La composée de deux fonctions de classe $\classec{1}$ est de classe $\classec{1}$.
\end{proposition}

\subsection{Dérivation et monotonie}

\begin{proposition}[utile=3]
Soit $f$ une fonction réelle, dérivable sur un intervalle $I$. Alors
\begin{itemize}
\item $f$ est croissante si et seulement si
  \[\forall x\in I \qsep f'(x)\geq 0.\]
\item $f$ est décroissante si et seulement si
  \[\forall x\in I \qsep f'(x)\leq 0.\]
\end{itemize}
\end{proposition}

\begin{remarqueUnique}
\remarque Cette proposition est fausse lorsque le domaine de définition de $f$
  n'est pas un intervalle. Par exemple la fonction
  \[\dspappli{f}{\Rs}{\R}{x}{1/x}\]
  n'est pas décroissante sur $\Rs$ bien qu'elle soit dérivable et que sa dérivée soit
  négative. Cependant ses restrictions
  aux intervalles $\RMs$ et $\RPs$ sont toutes les deux décroissantes.
\end{remarqueUnique}

\begin{exoUnique}
\exo Montrer que
  \[\forall x\in\interf{0}{1}\qsep x\leq \p{1+\frac{x}{2}}\ln(1+x).\]
%  Montrer que
%   \[\forall x\in\interf{0}{\pi/2} \qsep \frac{2}{\pi}x\leq \sin x\leq x, \qquad
%     \forall x\in\intero{0}{1} \qsep x^x (1-x)^{1-x}\geq\frac{1}{2}\]
%   \begin{sol}
% Soit $f$ la fonction définie sur $\interf{0}{\pi/2}$ par
% \[\forall x\in\interf{0}{\frac{\pi}{2}}\qsep f(x)\defeq \frac{2}{\pi}x - \sin x.\]
% D'après les théorèmes usuels, $f$ est deux fois dérivable et
% \begin{eqnarray*}
% \forall x\in\interf{0}{\frac{\pi}{2}}\qsep f'(x)&=&\frac{2}{\pi}-\cos x\\
% f''(x)&=&\sin x \geq 0
% \end{eqnarray*}
% Comme $f''(x)$ ne s'annule qu'en 0, on en déduit que $f'$ est strictement croissante sur $\interf{0}{\pi/2}$. Comme $f'(0)<0$ et $f'(\pi/2)>0$, en utilisant le théorème des valeurs intermédiaires et en utilisant la stricte monotonie de $f'$, il existe $\alpha\in\interf{0}{\pi/2}$ tel que $f'(\alpha)=0$. De plus $f'(x)$ est négative sur $[0,\alpha]$ et positive sur $[\alpha,\pi/2]$. Donc $f$ est décroissante sur $[0,\alpha]$ et croissante sur $[\alpha,\pi/2]$. Comme $f(0)=f(\pi/2)=0$, on en déduit que $f(x)$ est négative sur $[0,\pi/2]$. Donc
% \[\forall x\in\interf{0}{\frac{\pi}{2}}\qsep \frac{2}{\pi}x\leq \sin x.\]
% De même, on définit la fonction $g$ sur $\interf{0}{\pi/2}$ par
% \[\forall x\in\interf{0}{\frac{\pi}{2}}\qsep g(x)\defeq\sin x -x.\]
% D'après les théorèmes usuels, $g$ est dérivable et
% \[\forall x\in\interf{0}{\frac{\pi}{2}}\qsep g'(x)=\cos x-1\leq 0.\]
% Donc $g$ est décroissante. Or $g(0)=0$, donc
% \[\forall x\in\interf{0}{\frac{\pi}{2}}\qsep g(x)\leq 0.\]
% On a donc prouvé que
% \[\forall x\in\interf{0}{\pi/2} \qsep  \sin x \leq x.\]
% Pour la seconde inégalité, on définit la fonction $f$ sur $\intero{0}{1}$ par
% \[\forall x\in\intero{0}{1} \quad f(x)\defeq x^x (1-x)^{1-x}.\]
% D'après les théorèmes usuels, $f$ est dérivable et
% \[\forall x\in\intero{0}{1}\qsep f'(x)=x^{x}(1-x)^{1-x}\ln\left( \frac{x}{1-x}\right).\]
% Donc
% \begin{eqnarray*}
% \forall x\in\intero{0}{1}\qsep f'(x)=0
% &\ssi& \ln\left( \frac{x}{1-x}\right)=0\\
% &\ssi& \frac{x}{1-x}= 1\\
% &\ssi& x=\frac{1}{2}\\
% f'(x)\geq 0
% &\ssi& \ln\left( \frac{x}{1-x}\right)\geq 0\\
% &\ssi& \frac{x}{1-x}\geq 1\\
% &\ssi& x\geq \frac{1}{2}
% \end{eqnarray*}
% Donc $f$ est décroissante sur $\interof{0}{1/2}$ et croissante sur $\interfo{1/2}{1}$. Or $f(1/2)=1/2$, donc
% \[\forall x\in\intero{0}{1}\qsep x^x (1-x)^{1-x}\geq\frac{1}{2}.\]
%   \end{sol}
\end{exoUnique}

\begin{proposition}[utile=3]
Soit $f$ une fonction réelle, dérivable sur un intervalle $I$. Alors $f$ est
constante si et seulement si
\[\forall x\in I \qsep f'(x)=0.\]
\end{proposition}

\begin{remarqueUnique}
\remarque Cette proposition est fausse lorsque le domaine de 
  $f$ n'est pas un intervalle. 
\end{remarqueUnique}

\begin{proposition}[utile=2]
Soit $f$ une fonction réelle, dérivable sur un intervalle $I$. Si
\begin{itemize}
\item $\forall x\in I \qsep f'(x)\geq 0$
%%\item Il n'existe aucun intervalle non trivial sur lequel $f'$ est identiquement
%%  nulle
\item Le nombre de points de $I$ où $f'$ s'annule est fini.
\end{itemize}
alors $f$ est strictement croissante.
\end{proposition}

\begin{remarques}
\remarque La fonction $x\mapsto x^3$ est strictement croissante sur $\R$ bien
  qu'elle soit dérivable et que sa dérivée s'annule en 0.
\remarque Une fonction croissante qui n'est pas strictement croissante est constante sur un intervalle non trivial. Ces fonctions sont donc rares. Cependant, il est toujours plus délicat de montrer qu'une fonction est strictement croissante que croissante. Lorsqu'on a besoin de la stricte monotonie, il convient donc d'être particulièrement attentif. Inversement, il est inutile de prouver la stricte monotonie si seule la monotonie nous est utile.
\end{remarques}

\begin{exoUnique}
\exo Combien de racines réelles possède le polynôme $P(x)\defeq x^3-3x-1$~?
\end{exoUnique}


\subsection{Dérivation des fonctions à valeurs dans $\C$}

\begin{definition}[utile=-3]
Soit $f:\dom\to\C$ et $x_0\in\mathcal{D}$. On dit que $f$ est dérivable en $x_0$ lorsque $\Re(f)$ et $\Im(f)$ le sont. Si c'est le cas, on définit le nombre dérivé de $f$ en $x_0$ par
\[f'(x_0)\defeq \Re(f)'(x_0)+\ii \Im(f)'(x_0).\]
On dit que $f$ est dérivable lorsque $f$ est dérivable en tout point de $\dom$.
\end{definition}

\begin{proposition}[utile=-3, nom={Théorèmes usuels}]
Soit $f:\dom\to\C$ et $g:\dom\to\C$ deux fonctions dérivables. Alors
\begin{itemize}
\item Quels que soient $\lambda,\mu\in\C$, la fonction
  $\lambda f+\mu g$ est dérivable et
  \[\forall x\in\dom\qsep (\lambda f+\mu g)'(x)=\lambda f'(x)+\mu g'(x).\]
\item La fonction $fg$ est dérivable et
  \[\forall x\in\dom\qsep (fg)'(x)=f'(x)g(x)+f(x)g'(x).\]
\item Si $f$ ne s'annule pas sur $\dom$, $1/f$ est dérivable et
  \[\forall x\in\dom\qsep \p{\frac{1}{f}}'(x)=-\frac{f'(x)}{f^2(x)}.\]
\item Plus généralement, si $g$ ne s'annule pas sur $\dom$, $f/g$ est dérivable
  et
  \[\forall x\in\dom\qsep \p{\frac{f}{g}}'(x)=\frac{f'(x)g(x)-f(x)g'(x)}{g^2(x)}.\]
\end{itemize}
\end{proposition}

\begin{proposition}[utile=-3]
Soit $f:\dom\to\C$ une fonction dérivable. Alors la fonction $g$ définie par
\[\forall x\in\mathcal{D}\qsep g(x)\defeq \e^{f(x)}\]
est dérivable et
\[\forall x\in\dom\qsep g'(x)=f'(x)\e^{f(x)}.\]
\end{proposition}


\begin{proposition}[utile=3]
  Soit $f$ une fonction complexe, dérivable sur un intervalle $I$. Alors $f$ est
  constante si et seulement si
  \[\forall x\in I \qsep f'(x)=0.\]
  \end{proposition}

\begin{remarqueUnique}
\remarque On dit que'une fonction $f:\mathcal{D}\to\C$ est de classe $\classec{1}$ lorsque ses
  parties réelles et imaginaires le sont. Comme pour les fonctions à valeurs réelles, on
  montre qu'une combinaison linéaire, un produit, un quotient ainsi que l'exponentielle de fonctions
  de classe $\classec{1}$ sont de classe $\classec{1}$.
\end{remarqueUnique}

\section{Intégration, primitive}

Dans la suite de ce chapitre, $\K$ désignera l'un des corps $\R$ ou $\C$.

% \subsection{Définition, opérations usuelles}

% \begin{definition}[utile=-3]
% Soit $f$ une fonction continue sur un intervalle $I$ et $a,b\in I$. On
% définit l'intégrale
% \[\integ{a}{b}{f(x)}{x}\]
% comme l'aire algébrique comprise entre le graphe de $f$ et l'axe $(Ox)$
% comptée positivement si $a\leq b$ et négativement dans le cas contraire.
% \end{definition}

% \begin{proposition}[utile=-3]
% Soit $f$ et $g$ deux fonctions continues sur un intervalle $I$, $a,b\in I$
% et $\lambda,\mu\in\R$. Alors
% \[\integ{a}{b}{\p{\lambda f(x)+\mu g(x)}}{x}=
%   \lambda\integ{a}{b}{f(x)}{x}+\mu\integ{a}{b}{g(x)}{x}.\]
% \end{proposition}

% \begin{proposition}[utile=-3]
% Soit $f$ une fonction continue sur un intervalle $I$ et $a,b,c\in I$. Alors
% \[\integ{a}{c}{f(x)}{x}=\integ{a}{b}{f(x)}{x}+\integ{b}{c}{f(x)}{x}.\]
% \end{proposition}

% \subsection{Inégalités}

% \begin{proposition}[utile=-3]
% Soit $f$ et $g$ deux fonctions continues sur un intervalle $I$ et $a,b\in I$ tels
% que $a\leq b$. On suppose que
% \[\forall x\in\interf{a}{b} \qsep f(x)\leq g(x).\]
% Alors
% \[\integ{a}{b}{f(x)}{x}\leq \integ{a}{b}{g(x)}{x}.\]
% \end{proposition}

% \begin{exoUnique}
% \exo Montrer que, pour tout $x\in\R$, $0\leq 1-\cos x\leq x^2/2$.
%   En déduire la limite à droite en 0 de
%   \[\integ{x}{3x}{\frac{\cos t}{t}}{t}.\]
%   \begin{sol}
%   L'inégalité se montre en étudiant le signe des fonctions $f$ et $g$ définies sur $\R$ par
% \[\forall x\in\R\qsep f(x)=1-\cos x \quad\text{et}\quad g(x)=1-\cos x-\frac{x^{2}}{2}.\]
% Donc
% \[\forall t>0\qsep \frac{1}{t}+\frac{t}{2}\leq \frac{\cos t}{t}\leq \frac{1}{t}.\]
% Donc, pour tout $x>0$, puisque $x\leq 3x$
% \[\ln 3 -8\cdot\frac{x^{2}}{4}\leq \integ{x}{3x}{\frac{\cos t}{t}}{t} \leq \ln 3\]
% D'après le théorème des gendarmes, on en déduit que
% \[\integ{x}{3x}{\frac{\cos t}{t}}{t}\tendversdp{x}{0} \ln 3.\]
%   \end{sol}
% \end{exoUnique}

\subsection{Primitive}
\begin{definition}[utile=-3]
Soit $f:\dom\to\K$. On appelle
\emph{primitive} de $f$ toute fonction dérivable $F:\dom\to\K$ telle que
\[\forall x\in\dom \qsep F'(x)=f(x).\]
\end{definition}

\begin{proposition}[utile=-3]
Soit $f:I\to\K$ une fonction  admettant une primitive $F$ sur un
intervalle $I$. Alors les primitives de $f$ sont les fonctions $F_C:I\to\K$ définies
sur $I$ par
\[\forall x\in I \qsep F_C(x)\defeq F(x)+C\]
où $C\in\K$.
\end{proposition}

\begin{remarqueUnique}
\remarque Si la fonction d'expression $F(x)$ est une primitive de la fonction d'expression $f(x)$, on écrira
  \[\prim{f(x)}{x}=F(x).\]
  Il faudra rester vigilant avec cette notation. Par exemple
  \[\prim{1}{x}=x \et \prim{1}{x}=x+1\]
  mais $x\neq x+1$. On ne l'utilisera donc que pour calculer des primitives et on s'abstiendra de toute lecture autre que de la gauche vers la droite. On s'abstiendra aussi de l'utiliser avec des inégalités. 
\end{remarqueUnique}

\subsection{Intégration et régularité}

\begin{proposition}[utile=-3]
Soit $f:I\to\K$ une fonction continue sur un intervalle $I$ et $x_0\in I$. On
définit sur $I$ la fonction $F$ par
\[\forall x\in I \qsep F(x)\defeq\integ{x_0}{x}{f(t)}{t}.\]
Alors $F$ est de classe $\classec{1}$ et
  \[\forall x\in I \qsep F'(x)=f(x).\]
En particulier, $F$ est une primitive de $f$.
\end{proposition}


\begin{proposition}[utile=-3]
  Soit $f:I\to\K$ une fonction continue sur un intervalle $I$. Alors $f$ admet une
  primitive. Plus précisément, pour tout $x_0\in I$, il existe une unique
  primitive $F$ de $f$ s'annulant en $x_0$. De plus
  $$\forall x\in I \qsep F(x)=\integ{x_0}{x}{f(t)}{t}.$$
\end{proposition}

\begin{theoreme}[utile=3, nom=Théorème fondamental de l'analyse]
  Soit $f:I\to\K$ une fonction continue sur un intervalle $I$ et $a,b\in I$. Alors, si
  $F$ est une primitive de $f$
  \[\integ{a}{b}{f(x)}{x}=F(b)-F(a).\]
\end{theoreme}

\subsection{Intégration et inégalité}

\begin{proposition}[utile=3]
Soit $f:I\to\R$ et $g:I\to\R$ deux fonctions réelles continues et $a,b\in I$. On suppose que
\[a\leq b \quad\text{et}\quad \cro{\forall x\in\interf{a}{b} \qsep f(x)\leq g(x)}.\]
Alors
\[\integ{a}{b}{f(x)}{x}\leq\integ{a}{b}{g(x)}{x}.\]
\end{proposition}

\begin{preuve}
Il suffit d'appliquer la positivité et la linéarité.
\end{preuve}


\subsection{Intégration par parties, changement de variable}

\begin{proposition}[utile=3, nom=Intégration par parties]
Soit $f:I\to\K$ une fonction de classe $\classec{1}$ et $g:I\to\K$ une fonction
continue sur un intervalle $I$. Soit $G$ une primitive de $g$. Alors, si $a,b\in I$
\[\integppdi{a}{b}{f(x)}{g(x)}{x}=\evaldiff{f(x)G(x)}{a}{b}-
  \integ{a}{b}{f'(x)G(x)}{x}.\]
\end{proposition}

\begin{exoUnique}
\exo Pour tout $n\in\N$, on définit $I_n$ par
  \[I_n\defeq\integ{0}{1}{t^n\sqrt{1-t}}{t}\]
  Calculer $I_0$ et trouver une relation de récurrence entre $I_n$ et $I_{n+1}$.
  \begin{sol}
  Il faut intégrer $\sqrt{1-t}$. On trouve
  \[I_{n+1}=\frac{2(n+1)}{2n+5}I_n\]
  \end{sol}
\end{exoUnique}

\begin{remarqueUnique}
\remarque Si $f$ est dérivable et que $G$ est une primitive de $g$, alors
  \[\primppdi{f(x)}{g(x)}{x}=f(x)G(x)-\prim{f'(x)G(x)}{x}.\]
\end{remarqueUnique}

\begin{proposition}[utile=3, nom=Changement de variables]
  Soit $f:I\to\K$ une fonction continue sur un intervalle $I$. Soit $J$ un intervalle,
  $\bar{x}:J\to I$ une fonction de classe $\classec{1}$ et
  $a_x,b_x \in I$ et $a_t,b_t \in J$ tels que
  $$a_x=\bar{x}\p{a_t} \quad \text{et} \quad b_x=\bar{x}\p{b_t}.$$
  Alors
  \[\integ{a_t}{b_t}{f(\bar{x}(t))\frac{d\bar{x}}{dt}(t)}{t}=\integ{a_x}{b_x}{f(x)}{x}.\]
\end{proposition}

\begin{exoUnique}
\exemple Calculer
  \[\integ{0}{\pi}{\ln\p{1+\cos^2 x}\sin\p{2x}}{x}\]
  \begin{sol}
  C'est 0.
  \end{sol}
% \exemple Montrer que
%   \[\integ{\frac{1}{2}}{2}{\frac{1-x^2}{(1+x^2)\sqrt{1+x^4}}}{x}=0\]
\end{exoUnique}

% \begin{remarqueUnique}
% \remarque Si $\bar{x}$ (que l'on confondra ici avec $x$) est dérivable et que sa dérivée ne s'annule pas, alors, en notant
%   \[\frac{{\rm d}x}{{\rm d}t}=g(x), \quad\text{on a}\quad {\rm d}t=\frac{{\rm d}x}{g(x)}.\]
%   On en déduit que
%   \[\prim{f(\bar{x}(t))}{t}
%   = \prim{\frac{f(x)}{g(x)}}{x}.\]
%   Une fois le calcul d'intégrale terminé, on veillera bien à remplacer $x$ par son expression en $t$. Au cours du calcul, il est important de ne jamais mélanger les $x$ et les $t$ dans la même expression.
% \end{remarqueUnique}

% \begin{proposition}
%   Soit $f$ une fonction continue sur le segment $\interf{-a}{a}$. Alors~:
%   \begin{itemize}
%   \item Si $f$ est paire
%     $$\integ{-a}{0}{f(x)}{x}=\integ{0}{a}{f(x)}{x}.$$
%   \item Si $f$ est impaire
%     $$\integ{-a}{0}{f(x)}{x}=-\integ{0}{a}{f(x)}{x}.$$
%     En particulier
%     $$\integ{-a}{a}{f(x)}{x}=0.$$
%   \end{itemize}
% \end{proposition}

% \begin{proposition}
%   Soit $f$ une fonction $T$-périodique, continue sur $\R$. Alors quel que soit
%   $a\in\R$
%   $$\integ{a}{a+T}{f(x)}{x}=\integ{0}{T}{f(x)}{x}.$$
% \end{proposition}

\subsection{Calcul de primitive}

Étant donnée une fonction $f$ définie sur un intervalle à partir d'une
expression en les fonctions usuelles, on cherche une primitive $F$ de $f$.
Puisque $f$ est une expression en les fonctions usuelles, elle est en
particulier continue, donc admet une primitive. Le problème du calcul de
primitive est d'expliciter une telle fonction.\\

Il est d'abord essentiel de connaitre par cœur les primitives des fonctions usuelles.

\begin{center}
\begin{tabular}{|c|c|c|c|}
\hline
$\dom$ & $f(x)$ & $F(x)$\\
\hline
$\R$ & $x^n \quad (n\in\N)$ & $\frac{x^{n+1}}{n+1}$\\
\hline
$\Rs$ & $x^n \quad (n\in\Z\setminus\ens{-1})$ & $\frac{x^{n+1}}{n+1}$\\
\hline
$\Rs$ & $\frac{1}{x}$ & $\ln\abs{x}$\\
\hline
\hline
$\R$ & $\e^x$ & $\e^x$\\
\hline
$\R$ & $\cos x$ & $\sin x$\\
\hline
$\R$ & $\sin x$ & $-\cos x$\\
\hline
% \hline
% $\R$ & $\ch x$ & $\sh x$\\
% \hline
% $\R$ & $\sh x$ & $\ch x$\\
% \hline
% \hline
% $\R$ & $\frac{1}{1+x^2}$ & $\arctan x$\\
% \hline
% $\intero{-1}{1}$ & $\frac{1}{\sqrt{1-x^2}}$ & $\arcsin x$\\
% \hline
\end{tabular}
\end{center}

Ensuite, il existe de nombreuses techniques à connaitre pour calculer certaines primitives.

\begin{itemize}
\item \emph{\bf Polynômes}\\
  Le calcul d'une primitive d'une fonction polynôme se fait de manière
  immédiate
  $$\prim{\p{a_0+a_1 x+a_2 x^2+\dots+a_n x^n}}{x}=
    a_0 x+\frac{a_1}{2}x^2+\frac{a_2}{3}x^3+\dots+\frac{a_n}{n+1}x^{n+1}.$$
\item \emph{\bf Polynômes-exponentielle}\\
  Le calcul d'une primitive d'une fonction polynôme-exponentielle,
  c'est-à-dire de
  $$\prim{\p{a_0+a_1 x+a_2 x^2+\dots+a_n x^n}\e^{cx}}{x}$$
  se fait facilement par récurrence en effectuant une intégration par
  parties
  $$\primppdi{\p{a_0+a_1 x+a_2 x^2+\dots+a_n x^n}}{\e^{cx}}{x}.$$
  De cette manière, on abaisse le degré du polynôme. Il suffit de réitérer
  le procédé jusqu'à faire disparaitre le terme polynomial.
  \begin{exoUnique}
  \exo Calculer
    \[\prim{(2x+3)\e^x}{x}.\]
    \begin{sol}
  Par exemple~:
  \begin{equation*}
  \begin{split}
  \prim{(2x+3)\e^x}{x}&=\primppdi{\p{2x+3}}{\e^x}{x}\\
                     &=\p{2x+3}\e^x-\prim{2\e^x}{x}\\
                     &=\p{2x+1}\e^x
  \end{split}
  \end{equation*}
    \end{sol}
  \end{exoUnique}
\item \emph{\bf Polynômes-sinus/cosinus}\\
  On calcule de même toute primitive du produit d'une fonction polynôme
  et d'une fonction sinus ou cosinus.
  \begin{exoUnique}
  \exo Calculer
    \[\prim{x\cos x}{x}.\]
    \begin{sol}
  \begin{equation*}
  \begin{split}
  \prim{x\cos x}{x}&=\primppdi{x}{\cos x}{x}\\
                   &=x\sin x-\prim{\sin x}{x}\\
                   &=x\sin x+\cos x    
  \end{split}
  \end{equation*}
    \end{sol}
  \end{exoUnique}
\item \emph{\bf Exponentielle-sinus/cosinus}
  Pour calculer des primitives de la forme
  \[\prim{\e^{a x} \cos( bx)}{x} \quad\text{ou}\quad \prim{\e^{a x} \sin(bx)}{x},\]
  on passe par l'exponentielle complexe. On fait de même si un polynôme est en facteur d'une telle expression.
  \begin{exoUnique}
  \exo Calculer
  \[\prim{\e^{2x} \sin(3x)}{x}.\]  
  \end{exoUnique}
\item \emph{\bf Polynôme-logarithme}\\
  Le calcul d'une primitive d'une fonction polynôme-logarithme, c'est-à-dire
  de
  $$\prim{\p{a_0+a_1 x+a_2 x^2+\dots+a_n x^n}\ln x}{x}$$
  se fait facilement par intégration par parties
  $$\primppid{\p{a_0+a_1 x+a_2 x^2+\dots+a_n x^n}}{\ln x}{x}.$$
  % Cette technique est utile pour calculer une primitive de toute expression faisant intervenir une fonction donc la dérivée est plus simple que la fonction.
    \begin{exoUnique}
  \exo Calculer
  \[\prim{\ln x}{x}\]% \et \prim{\arctan x}{x}.\]
    \begin{sol}
  Par exemple~:
  \begin{equation*}
  \begin{split}
  \prim{\ln x}{x}&=\primppid{1}{\ln x}{x}\\
                 &=x\ln x-\prim{1}{x}\\
                 &=x\ln x-x
  \end{split}
  \end{equation*}
  % \begin{equation*}
  % \begin{split}
  % \prim{\arctan x}{x}&=\primppid{1}{\arctan x}{x}\\
  %                &=x\arctan x-\prim{\frac{x}{1+x^2}}{x}\\
  %                &=x\arctan x-\frac{1}{2}\ln(1+x^2)
  % \end{split}
  % \end{equation*}
    \end{sol}
  \end{exoUnique}
\item \emph{\bf Polynômes en $\sin$ et $\cos$}\\
  Pour le calcul de primitives de polynômes en $\sin$ et $\cos$, c'est-à-dire
  de~:
  $$\prim{\sin^n x \cos^m x}{x} \quad \text{où $n,m\in\N$}$$
  on peut, lorsque $n$ ou $m$ est impair effectuer un changement de variable
  pour se ramener à un calcul de primitive de polynôme.
  \begin{itemize}
  \item Si $m$ est impair, soit $m'\in\N$ tel que $m=2m'+1$. On effectue alors
    le changement de variable $t=\sin x$.
    \begin{equation*}
    \begin{split}
    \prim{\sin^n x\cos^m x}{x}&=\prim{\sin^n x \cos^{2m'+1} x}{x}\\
                              &=\prim{\sin^n x \p{1-\sin^2 x}^{m'}\cos x}{x}\\
                              &=\prim{t^n\p{1-t^2}^{m'}}{t}.
    \end{split}
    \end{equation*}
  \item Si $n$ est impair, on effectue le changement de variable $t=\cos x$.
  \item Si $n$ et $m$ sont pairs, on effectue une linéarisation de
    l'expression. 
  \end{itemize}
\begin{exoUnique}
\exo Calculer
  \[\prim{\sin^2 x\cos^3 x}{x}, \qquad \prim{\cos^2 x\sin^5 x}{x}, \qquad
    \prim{\cos^2 x\sin^2 x}{x}.\]
    \begin{sol}
  Par exemple~:
    \begin{equation*}
    \begin{split}
    \prim{\sin^2 x\cos^3 x}{x}&=\prim{\sin^2 x \p{1-\sin^2 x}\cos x}{x}\\
                              &=\prim{t^2\p{1-t^2}}{t}\\
                              &=\frac{1}{3}t^3-\frac{1}{5}t^5\\
                              &=\frac{1}{3}\sin^3 x-\frac{1}{5}\sin^5 x\\
    \end{split}
    \end{equation*}    
     Par exemple~:
    \begin{eqnarray*}
    \prim{\cos^2 x\sin^5 x}{x}&=&\prim{\cos^2 x\p{1-\cos^2 x}^2\sin x}{x}\\
      &=&-\prim{t^2\p{1-t^2}^2}{t}\\
      &=&-\prim{t^2-2t^4+t^6}{t}\\
      &=&-\frac{1}{3}t^3+\frac{2}{5}t^5-\frac{1}{7}t^7\\
      &=&-\frac{1}{3}\cos^3 x+\frac{2}{5}\cos^5 x-\frac{1}{7}\cos^7 x
    \end{eqnarray*}
Par exemple~:
    \begin{eqnarray*}
    \cos^2 x\sin^2 x&=&\p{\frac{1}{2}\sin\p{2x}}^2\\
     &=&\frac{1}{4}\frac{1-\cos\p{4x}}{2}\\
     &=&\frac{1}{8}-\frac{1}{8}\cos\p{4x}
    \end{eqnarray*}
    Donc~:
    $$\prim{\cos^2 x\sin^2 x}{x}=\frac{1}{8}x-\frac{1}{32}\sin\p{4x}$$
    \end{sol}
\end{exoUnique}
% \item \emph{\bf Fractions rationnelles}\\
% Le calcul de primitives de fractions rationnelles se fait après une décomposition en éléments simples sur $\R$.
% \begin{itemize}
% \item La partie entière s'intègre facilement comme un polynôme.
% \item Les éléments simples de première espèce s'intègrent en
%   \[\priminv{x-\alpha}{x}=\ln\abs{x-\alpha}.\]% \et \priminv{\p{x-\alpha}^n}{x}=\frac{-1}{n-1}\cdot\frac{1}{\p{x-\alpha}^{n-1}} \quad\text{(où $n\geq 2$)}.\]
% \item Les éléments simples de seconde espèce s'intègrent de la manière suivante. On commence par s'occuper du terme en $x$ du numérateur en faisant apparaître une fraction rationelle de la forme $P'/P$.
%   \begin{eqnarray*}
%   \prim{\frac{bx+c}{x^2+\beta x+\gamma}}{x}
%   &=& \frac{b}{2}\prim{\frac{2x+\beta}{x^2+\beta x+\gamma}}{x} + \p{c-\frac{b\beta}{2}}\priminv{x^2+\beta x+\gamma}{x}\\
%   &=& \frac{b}{2}\ln\p{x^2+\beta x+\gamma}+\frac{2c-b\beta}{2}\priminv{x^2+\beta x+\gamma}{x}
%   \end{eqnarray*}
%   Il suffit ensuite de mettre le trinôme (qui rappelons-le n'a pas de racine réelle) sous forme canonique
%   \[x^2+\beta x+\gamma = \p{x+\frac{\beta}{2}}^2+\underbrace{\frac{4\gamma - \beta^2}{4}}_{\defeq \alpha^2>0}=\alpha^2\cro{\p{\frac{2x+\beta}{2\alpha}}^2+1}\]
%   puis de poser $u\defeq (2x+\beta)/(2\alpha)$.
%   \begin{eqnarray*}
%   \priminv{x^2+\beta x+\gamma}{x}
%     &=& \frac{1}{\alpha^2}\priminv{1+\p{\frac{2x+\beta}{2\alpha}}^2}{x}\\
%     &=& \frac{1}{\alpha}\priminv{1+u^2}{u} = \frac{1}{\alpha}\arctan u\\
%     &=& \frac{1}{\alpha}\arctan\frac{2x+\beta}{2\alpha}.
%   \end{eqnarray*}
%   En conclusion
%   \[\prim{\frac{bx+c}{x^2+\beta x+\gamma}}{x}=\frac{b}{2}\ln\p{x^2+\beta x+\gamma}+\frac{2c-b\beta}{2\alpha}\arctan\frac{2x+\beta}{2\alpha}.\]
% \end{itemize}
% \begin{exoUnique}
% \exo Calculer
%   \[\priminv{x^2-1}{x}, \qquad \prim{\frac{2x+1}{x^2+x-2}}{x},  \qquad
%     \priminv{x^3-1}{x}.\]
%   \begin{sol}
%   On trouve~:
%   \begin{itemize}
%   \item Premier~:
%      \[\frac{1}{2}\ln\abs{\frac{1+x}{1-x}}\]
%   \item Réfléchir deux secondes...
%   \item Troisième~:
%      La décomposition en éléments simples donne
%      \[\frac{1}{X^3-1}=\frac{1}{3}\cdot\frac{1}{X-1}-\frac{1}{3}\cdot
%        \frac{X+2}{X^2+X+1}\]
%      Et le calcul de primitive
%      \[\frac{1}{3}\ln\abs{x-1}-\frac{1}{6}\ln\p{x^2+x+1}-\frac{1}{\sqrt{3}}
%        \arctan\p{\frac{2x+1}{\sqrt{3}}}\]
%   \end{itemize}
%   \end{sol}
% \end{exoUnique}
%   \item \emph{\bf Fractions rationnelles en $\e^x$}\\
%     Pour les fractions rationnelles en $\e^x$, un changement de variable $u=\e^x$ permet de se ramener à un calcul de primitive de fraction rationnelle.
%     \begin{exoUnique}
% \exo Calculer
%   \[\priminv{\e^{2x}+1}{x}\]
%   \begin{sol}
%   On trouve
%   \[\frac{1}{X\p{X^2+1}}=\frac{1}{X}-\frac{X}{X^2+1}\]
%   donc
%   \[x-\frac{1}{2}\ln\p{e^{2x}+1}\]
%   \end{sol}
% % \[\priminv{x^2+x+1}{x},\quad
% %   \priminv{x^2-3x+2}{x}.\]
% % \[\prim{\frac{2x+5}{(x-1)^2(x+2)}}{x},\quad
% %   \prim{\frac{x+1}{x(x^2-x+1)}}{x}\]
%     \end{exoUnique}
%   \item \emph{\bf Fractions rationnelles en $\cos x$ et $\sin x$}\\
% Le calcul de primitive de fraction rationelle en $\cos x$ et $\sin x$ s'effectue en suivant les \emph{règles de \textsc{Bioche}}. Supposons que $G$ soit une fraction rationnelle en $\cos x$, $\sin x$.
% \begin{itemize}
% \item Si $G\p{-x}=-G(x)$, il existe une fraction rationnelle $F$ telle que
%   \[G(x)=F\p{\cos x}\sin x.\]
%   On est donc ramené, après le changement de variable $u\defeq\cos x$, à un calcul
%   de primitive d'une fraction rationnelle.
% \item Si $G\p{\pi-x}=-G(x)$, il existe une fraction rationnelle $F$ telle que
%   \[G(x)=F\p{\sin x}\cos x.\]
%   On est donc ramené, après le changement de variable $u\defeq\sin x$, à un calcul
%   de primitive d'une fraction rationnelle.
% \item Si $G\p{\pi+x}=G(x)$, il existe une fraction rationnelle $F$ telle que
%   \[G(x)=F\p{\tan x}\p{1+\tan^2 x}.\]
%   On est donc ramené, après le changement de variable $u\defeq\tan x$, à un calcul
%   de primitive d'une fraction rationnelle.
% \item Sinon, on effectue le changement de variable $u\defeq\tan\p{x/2}$. En
%   remarquant que
%   \[\cos x=\frac{1-u^2}{1+u^2}, \qquad \sin x=\frac{2u}{1+u^2} \et {\rm d}u=\frac{1}{2}(1+u^2){\rm d}t,\]
%   on est ramené à un calcul de primitive d'une fraction rationnelle. Il faut être très soigneux, car ce calcul de primitive se fait sur un intervalle ne contenant pas les $x$ tels que $x\equiv \pi\ [2\pi]$.
% \end{itemize}
% \begin{exoUnique}
% \exo Calculer
%   \[\priminv{\cos x\cos\p{2x}}{x}, \qquad \priminv{\sin x+\sin\p{2x}}{x},
%     \qquad \integinv{0}{2\pi}{2+\cos x}{x}.\]
%   \begin{sol}
%   On trouve~:
%   \begin{itemize}
%   \item Premier~: Changement de variable $u=\sin x$.
%     Décomposition en éléments simples
%     \[\frac{1}{\p{1-2X^2}\p{1-X^2}}=\frac{1}{2}\cro{\frac{1}{X-1}-\frac{1}{X+1}}+
%       \frac{\sqrt{2}}{2}\cro{\frac{1}{X+\frac{\sqrt{2}}{2}}-
%       \frac{1}{X-\frac{\sqrt{2}}{2}}}\]
%     Puis regrouper les pôles opposés. On trouve
%     \[\sqrt{2}\argth\p{\sqrt{2}\sin x}-\argth\p{\sin x}\]
%   \item Second~: Changement de variable $u=\cos x$. Décomposition en éléments
%     simples
%     \[\frac{1}{\p{X^2-1}\p{1+2X}}=\frac{1}{6}\cdot\frac{1}{X-1}+
%       \frac{1}{2}\cdot\frac{1}{X+1}-\frac{4}{3}\cdot\frac{1}{1+2X}\]
%     On trouve donc
%     \[\frac{1}{6}\ln\p{1-\cos x}+\frac{1}{2}\ln\p{1+\cos x}-\frac{2}{3}
%       \ln\abs{1+2\cos x}\]
%   \item Troisième~: Les règles de \textsc{Bioche} ne marchent pas. Il faut
%     faire le changement de variable~: $u=\tan(t/2)$. On trouve
%     \[\frac{2\sqrt{3}}{3}\arctan\p{\frac{\sqrt{3}\tan\p{x/2}}{3}}\]
%     et l'intégrale vaut
%     \[\frac{2\sqrt{3}}{3}\pi\]
%   \end{itemize}
%   \end{sol}
% \end{exoUnique}
%   \item \emph{\bf Fractions rationnelles en $\ch x$ et $\sh x$}\\
% Le calcul de primitive de fraction rationelle en $\ch x$ et $\sh x$ s'effectuent de même. Supposons que $G$ soit une fraction rationnelle en $\ch x$, $\sh x$.
% \begin{itemize}
% \item Si lorsqu'on remplace $\ch$ par $\cos$ et $\sh$ par $\sin$ les règles de
%   \textsc{Bioche} préconisent le changement de variable $u\defeq\cos x$, on effectue le
%   changement de variable $u\defeq\ch x$. 
% \item Si lorsqu'on remplace $\ch$ par $\cos$ et $\sh$ par $\sin$ les règles de
%   \textsc{Bioche} préconisent le changement de variable $u\defeq\sin x$, on effectue le
%   changement de variable $u\defeq\sh x$. 
% \item Si lorsqu'on remplace $\ch$ par $\cos$ et $\sh$ par $\sin$ les règles de
%   \textsc{Bioche} préconisent le changement de variable $u\defeq\tan x$, on effectue le
%   changement de variable $u\defeq\th x$. 
% \end{itemize}
% \begin{exoUnique}
% \exo Calculer
%   \[\prim{\frac{\ch x}{\sh x+\ch x}}{x} \qquad\text{et}\qquad
%     \priminv{\sh^2 x+2}{x}.\]
%   \begin{sol}
%   On trouve
%   \begin{itemize}
%   \item Tout transformer avec des exponentielles. On trouve
%     \[\prim{\frac{\ch x}{\sh x+\ch x}}{x}=\frac{1}{2}x
%       -\frac{1}{4}e^{-2x}\]
% %   Effectuer le changement de variable $u=\th x$. On a~:
% %   \[\frac{1}{\p{1-X}\p{1+X}^2}=\frac{1}{4}\cdot\frac{1}{1+X}+
% %     \frac{1}{2}\cdot\frac{1}{\p{1+X}^2}+\frac{1}{4}\cdot\frac{1}{1-X}\]
% %   On trouve
% %   \[\prim{\frac{\ch x}{\sh x+\ch x}}{x}=\frac{1}{2}x
% %     -\frac{1}{2}\cdot\frac{1}{1+\th x}\]
%   \item Effectuer le changement de variable $t=\th x$. On trouve
%     \[\priminv{\sh^2 x+2}{x}=\frac{\sqrt{2}}{2}\argth\p{\frac{\th x}{\sqrt{2}}}\]
%   \end{itemize}
%   \end{sol}
% \end{exoUnique}
%   \item \emph{\bf Fractions rationnelles en $x$ et $\sqrt[n]{(ax+b)/(cx+d)}$}\\
%   Si $f$ s'écrit
%   \[f(x)\defeq F\p{x,\sqrt[n]{\frac{ax+b}{cx+d}}}\]
%   où $F$ est une fraction rationnelle, le calcul d'une
%   primitive de $f$ se fait en effectuant le changement de variable
%   \[u\defeq\sqrt[n]{\frac{ax+b}{cx+d}}.\]
%   On a alors
%   \[x=\frac{du^n-b}{a-cu^n} \et \text{d}x=G(u)\,\text{d}u\]
%   où $G$ est une fraction rationnelle. On est donc ramené au calcul d'une
%   primitive d'une fraction rationnelle.
% % \item On souhaite désormais calculer une primitive d'une fonction $f$ de la forme
% %   \[f(x)=F\p{x,\sqrt{ax^2+bx+c}}\]
% %   où $F$ est une fraction rationnelle à deux variables.
% %   \begin{itemize}
% %   \item Si $aX^2+bX+c$ admet une racine double réelle $\alpha$ (cas où $\Delta=0$), on
% %     a $\sqrt{ax^2+bx+c}=\sqrt{a}\abs{x-\alpha}$. L'expression est donc
% %     une fraction rationnelle sur chaque intervalle où $x-\alpha$ est de signe
% %     constant.
% %   \item Si $aX^2+bX+c$ n'admet pas de racine réelle (cas où $\Delta<0$), après
% %     mise sous forme canonique de $aX^2+bX+c$ et changement de variable, on
% %     est ramené au calcul d'une primitive d'une fraction rationnelle
% %     en $u$ et $\sqrt{1+u^2}$. Il suffit alors d'effectuer le changement de variable
% %     $u=\sh v$ pour se ramener au calcul d'une primitive d'une fraction
% %     rationnelle en $\ch v$, $\sh v$.
% %   \item Si $aX^2+bX+c$ admet deux racines réelles (cas où $\Delta>0$), après
% %     mise sous forme canonique de $aX^2+bX+c$ et changement de variable, on
% %     est ramené au calcul d'une primitive d'une fraction rationnelle
% %     en $u$ et $\sqrt{1-u^2}$ ou en $u$ et $\sqrt{u^2-1}$. Dans le premier cas, on
% %     effectue le changement de variable $u=\cos v$ alors que dans le second
% %     cas on effectue le changement de variable $u=\pm\ch v$.
% %   \end{itemize}
% % \end{itemize}
% \begin{exoUnique}
% \exo Calculer
%   \[\priminv{x\sqrt{x+1}}{x}.\]% \qquad \prim{\arctan\sqrt{1+x^2}}{x}\]
%   \begin{sol}
%   \begin{itemize}
%   \item On pose $u=\sqrt{1+x}$. On trouve alors
%     \[\priminv{x\sqrt{x+1}}{x}=\ln\p{\frac{\sqrt{1+x}-1}{\sqrt{1+x}+1}}\]
%   \item On trouve
%     \[x\arctan\sqrt{1+x^2}-\argsh x+
%       \sqrt{2}\argth\p{\frac{x}{\sqrt{2\p{1+x^2}}}}\]
%   \end{itemize}
%   \end{sol}
% \end{exoUnique}
\end{itemize}


% Pour clore cette revue des techniques classiques de calcul de primitive, il est important de noter que si d'un point de vue théorique, une fonction continue admet toujours une
% primitive, il n'est pas toujours possible de l'exprimer à l'aide
% fonctions usuelles. Par exemple, avant que le logarithme néperien ne fasse
% son apparition (on ne considérait alors que les fonctions
% rationnelles), les primitives de la fonction $x\mapsto 1/x$ ne pouvait pas
% s'exprimer rationnellement en $x$. L'ajout des fonctions $\ln$ et $\exp$
% dans le jeu des fonctions usuelles n'a pas résolu ce problème. Par exemple,
% on peut démontrer que les primitive de la fonction gaussienne définie sur
% $\R$ par $f(x)\defeq \e^{-x^2}$ ne peuvent pas s'exprimer à l'aide des fonctions usuelles.
%END_BOOK

\end{document}




\section{Exercices}
\setcounter{numeroexercice}{1}
\documentclass{magnolia}

\magtex{tex_driver={pdftex}}
\magfiche{document_nom={Exercices de révision d'analyse},
          auteur_nom={François Fayard},
          auteur_mail={fayard.prof@gmail.com}}
\magexos{exos_matiere={maths},
         exos_niveau={mpsi},
         exos_chapitre_numero={3},
         exos_theme={Compléments d'analyse}}
\magmisenpage{misenpage_presentation={tikzvelvia},
         misenpage_format={a4},
         misenpage_nbcolonnes={1},
         misenpage_preuve={non},
         misenpage_sol={non}}
\maglieudiff{}
\magprocess

\begin{document}

%BEGIN_BOOK

\magsection{Le corps ordonné $\R$}

\magsubsection{La relation d'ordre sur $\R$}


\exercice{nom={Inégalités}}
\begin{questions}
\question Montrer que pour $a,b,c\in\RP$
  \[\p{a^2+1}\p{b^2+1}\p{c^2+1}\geq 8abc.\]
\question Montrer que pour $a,b\in\R$ tels que $0<a\leq b$, on a
  \[\frac{1}{8}\cdot\frac{\p{b-a}^2}{b}\leq\frac{a+b}{2}-\sqrt{ab}
    \leq\frac{1}{8}\cdot\frac{\p{b-a}^2}{a}.\]
\question Montrer que si $a,b,x,y\in\RPs$ sont tels que $a+b=1$, alors
  \[\frac{a}{x}+\frac{b}{y}\geq\frac{1}{ax+by}.\]
\end{questions}
\begin{sol}
$\quad$
\begin{questions}
\question Remarquer que $0\geq 2a\leq a^2+1$. Écrire les mêmes inégalités avec
  $b$ et $c$ et faire le produit des inégalités obtenues.
\question Écrire
  \[\frac{a+b}{2}-\sqrt{ab}=\frac{(\sqrt{a}-\sqrt{b})^2}{2}=
    \frac{(a-b)^2}{2(\sqrt{a}+\sqrt{b})^2}\]
\question Raisonner par équivalence. On aboutit à
  \[ab(x-y)^2\geq 0\]
  qui est bien sur vrai.
\end{questions}
\end{sol}

\exercice{nom={Puissances}}
\begin{questions}
\question Soit $n\in\Ns$ et $(x,y)\in\RP^2$. Montrer que
  \[(x+y)^{\frac{1}{n}}\leq x^{\frac{1}{n}}+y^{\frac{1}{n}}.\]
\question Soit $a\in\RP$ et $n\in\N$. Montrer que $(1+a)^n\geq 1+na$.
\question Soit $a$ et $b$ deux nombres réels tels que $0\leq a\leq b$ et  
  $n\in\Ns$. Montrer que
  \[n(b-a)a^{n-1}\leq b^n-a^n\leq n(b-a)b^{n-1}.\]
\question Soit $a,b,c\in\interf{0}{1}$. Montrer qu'au moins un des trois
  nombres réels
  \[a(1-b), \quad b(1-c), \quad c(1-a)\]
  est inférieur à $\frac{1}{4}$.
\end{questions}
\begin{sol}
$\quad$
\begin{questions}
\question Raisonner par équivalence.
\question Binôme de Newton ou récurrence sur $n$.
\question Factorisation de $a^n-b^n$.
\question Utilisons le fait que $x(1-x)\leq\frac{1}{4}$ pour tout
  $x\in\interf{0}{1}$.
  \begin{itemize}
  \item Si l'un est nul, par exemple $a$, alors $a(1-b)=0\leq 1/4$.
  \item Sinon, prenons le plus grand des $3$. Si c'est par exemple $a$, alors $c(1-a)=\frac{c}{a}a(1-a)\leq 1/4 c/a \leq 1/4$.
  \end{itemize}
\end{questions}
\end{sol}

\exercice{nom={Système non linéaire}}
Soit $x_1,\ldots,x_n$ des nombres réels tels que
\[\sum_{i=1}^nx_i=n \et \sum_{i=1}^nx_i^2=n.\]
Montrer que $x_i=1$ pour tout $i\in\intere{1}{n}$.

\begin{sol}
$$\sum_{i=1}^n(x_i-1)^2=\sum_{i=1}^nx_i^2-2\sum_{i=1}^nx_i+\sum_{i=1}^n 1=0$$
\end{sol}

\exercice{nom={Système non linéaire}}
On suppose que $\p{x,y,z,t}\in\Rs^4$ vérifie le système
\[\syslin{x          &+y          &          +z&=&t\cr
          \frac{1}{x}&+\frac{1}{y}&+\frac{1}{z}&=&\frac{1}{t}}\]
Établir que $\p{x+y}\p{y+z}\p{z+x}=0$, et en déduire la forme générale des
solutions du système ci-dessus.

\begin{sol}
D'une part, $$(x+y)(y+z)(z+x)=2xyz+x^2(y+z)+y^2(x+z)+z^2(x+y).$$
D'autre part, d'après la deuxième équation $$t(yz+xz+xy)=xyz$$ mais en en remplaçant $t$ par $x+y+z$ d'après la première équation, on obtient bien la nullité de $(x+y)(y+z)(z+x)$.
Donc par équation produit nulle. Deux d'entre eux sont opposés et l'autre est égale à $t$. Réciproquement, cela fonctionne.
\end{sol}



% \exercice{nom={Racine carrée}}
% \begin{questions}
% \question Pour $a\in\interfo{1}{+\infty}$, simplifier 
%   $\sqrt{a+2\sqrt{a-1}}+\sqrt{a-2\sqrt{a-1}}$.
% \question Résoudre l'équation 
%   $\sqrt{x+3-4\sqrt{x-1}}+\sqrt{x+8-6\sqrt{x-1}}=1$ d'inconnue $x\in\R$.
% \end{questions}
% \begin{sol}
% \begin{questions}
% \question En élevant au carré, on trouve que ça vaut $2$ si $a\in [1;2]$ et $2\sqrt{a-1}$ sinon.
% \question 
% \end{questions}
% \end{sol}

\magsubsection{Valeur absolue}



\magsubsection{Racine}
\magsubsection{Partie entière, approximation}

\exercice{nom={Rationnels et irrationnels}}
\begin{questions}
\question Soit $a$ et $b$ deux éléments de $\R\setminus\Q$. Peut-on affirmer
  que $a+b$ (respectivement $a\times b$) appartient à $\R\setminus\Q$~?
  Et si $a\in\Q$ et $b\in\R\setminus\Q$~?
\question Montrer, en raisonnant par l'absurde, que $\sqrt{6}-\sqrt{2}-\sqrt{3}$ et
  $\sqrt{2}+\sqrt{3}+\sqrt{5}$ sont irrationnels.
\end{questions}

\exercice{nom={Intersection d'une famille infinie}}
Déterminer 
\[\cap_{n\in\Ns}\interf{-\frac{1}{n}}{\frac{1}{n}} \et
  \cap_{n\in\Ns}\intero{-\frac{1}{n}}{\frac{1}{n}}.\]

\exercice{nom={Autour de la partie entière}}
Soit $x$ et $y$ deux nombres réels et $n\in\Ns$.
\begin{questions}
\question Montrer que
  \[\ent{x+y}\geq \ent{x}+ \ent{y}.\]
  Y a-t-il des cas d'égalité~? D'inégalité stricte~?
\question Montrer que $\ent{\frac{\ent{nx}}{n}}=\ent{x}$.
\end{questions}
\begin{sol}
\begin{questions}
\question On a $$\ent{x}\leq x \qquad \text{et} \qquad \ent{y}\leq y $$ donc $$\ent{x}+\ent{y}\leq x+y$$ et $\ent{x}+\ent{y}$ est un entier donc $\ent{x+y}\geq \ent{x}+\ent{y}$.
$x=1=y$ donne une égalité tandis que $x=y=1,5$ donne une inégalité stricte.
\question On a $\ent{x}\leq x <\ent{x}+1$ donc $n\ent{x}\leq nx <n(\ent{x}+1)$. $n\ent{x}$ étant un entier, on en déduit : $n\ent{x}\leq \ent{nx}(\leq nx) <n(\ent{x}+1)$ d'où en divisant par $n$ :
$$\ent{x}\leq \frac{\ent{nx}}{n} <\ent{x}+1.$$ L'unique entier vérifiant cette dernière inégalité étant $\ent{\frac{\ent{nx}}{n}}$, on en déduit :
$$\ent{\frac{\ent{nx}}{n}}=\ent{x}.$$
\end{questions}
\end{sol}

\exercice{nom={Calcul de somme}}
\begin{questions}
\question Montrer que
  \[\forall x\in\R \qsep \ent{x}+\ent{-x}=
    \begin{cases}
    0  & \text{si $x\in\Z$}\\
    -1 & \text{sinon.}
    \end{cases}\]
\question En déduire que si $p,q\in\Ns$ sont premiers entre eux
  \[\sum_{k=1}^{q-1} \ent{k\cdot \frac{p}{q}}=\frac{\p{p-1}\p{q-1}}{2}\]
\end{questions}
\begin{sol}
\begin{enumerate}
  \item Soit $x \in \R$. 
    \begin{itemize}
      \item Si $x \in \Z$, alors $\ent{x} = x$ et, comme $-x \in \Z$, $\ent{-x} = -x$, donc $\ent{x} + \ent{-x} = 0$. 
      \item Si $x \notin \Z$, alors on a $x-1 < \ent x < x$ et $-x-1 < \ent{-x} < -x$. En sommant cela, on obtient $-2 < \ent x + \ent{-x} < 0$. 
	Comme le seul entier dans $]-2,0[$ est $-1$, on a bien $ \ent x + \ent{-x}$.
    \end{itemize}
    Ainsi, \fbox{$\displaystyle\forall x\in\R \quad \ent{x}+\ent{-x}=
    \begin{cases}
    0  & \text{ si }x\in\Z\\
    -1 & \text{ sinon}\end{cases}.$}
  \item On a, par le changement d'indices $\ell = q-k$ : 
    \begin{align*}
      \displaystyle\sum_{k=1}^{q-1} \ent{k \times \frac{p}{q}} &= \dfrac{1}{2} \displaystyle\sum_{k=1}^{q-1} \ent{k \times \frac{p}{q}} + \dfrac{1}{2}\displaystyle\sum_{\ell=1}^{q-1} \ent{(q-\ell) \times \frac{p}{q}} \\
      \displaystyle\sum_{k=1}^{q-1} \ent{k \times \frac{p}{q}} &= \dfrac{1}{2} \displaystyle\sum_{k=1}^{q-1} \ent{k \times \frac{p}{q}} + \dfrac{1}{2}\displaystyle\sum_{k=1}^{q-1} \ent{p - k\times \frac{p}{q}} \\
                                                       &= \dfrac{1}{2} \displaystyle\sum_{k=1}^{q-1} p + \ent{k \times \frac{p}{q}} + \ent{-k \times \frac{p}{q}}.
    \end{align*}
    Or, si $k \in \Z$, si $\dfrac{kp}{q}$ est entier, alors $q \mid kp$ et comme $p$ et $q$ sont premiers entre eux, par le lemme de Gauss, $q\mid p$.
    Il n'y a aucun multiple de $q$ dans $\llbracket 1,q-1 \rrbracket$, donc si $k\in \llbracket 1,q-1 \rrbracket$, alors $\dfrac{kp}{q} \notin \Z$. 
    Ainsi, par ce qui précède, 
    \begin{equation*}
      \boxed{\displaystyle\sum_{k=1}^{q-1} \ent{k \times \frac{p}{q}} = \dfrac{1}{2} \displaystyle\sum_{k=1}^{q-1} (p-1) = \dfrac{(p-1)(q-1)}{2}.}
    \end{equation*}
\end{enumerate}
\end{sol}

\magsubsection{Intervalle}
\magsection{Fonction réelle d’une variable réelle}

\magsubsection{Définition}

\exercice{nom={Bijection}}
Soit $f$ la fonction définie sur $\R$ par
\[\forall x\in\R\qsep f(x)\defeq\frac{x}{\sqrt{x^2+1}}.\]
\begin{questions}
\question Montrer que $f$ est impaire puis qu'elle réalise une bijection de $\R$ sur un intervalle
  $J$ à préciser. 
\question Déterminer $f^{-1}$, puis tracer les graphes de $f$ et $f^{-1}$.
\end{questions}
\begin{sol}
D'après les théorèmes usuels, $f$ est dérivable sur $\R$ et
\begin{eqnarray*}
\forall x \in \R\qsep f'(x)
&=&\p{1+x^2}^{-1/2}+x(2x)\p{\frac{-1}{2}}\p{1+x^2}^{-3/2}\\
&=&\frac{1}{(x^2+1)^{3/2}}>0
\end{eqnarray*}
donc $f$ est strictement croissante sur $\R$. Puisque
\[\forall x\geq 0\qsep f(x)=\frac{x}{\sqrt{x^2+1}}=\frac{1}{\sqrt{1+\frac{1}{x^2}}}\tendvers{x}{+\infty} 1\]
et que $f$ est impaire, on en déduit que
\[\forall x\in\R\qsep -1<f(x)<1.\]
Montrons que $f$ réalise une bijection de $\R$ dans $]-1,1[$. Puisque $f$ est strictement croissante, elle est injective. Montrons qu'elle réalise une surjection de $\R$ dans $]-1,1[$. Soit $y\in]-1,1[$. Puisque $f$ est continue d'après les théorèmes usuels, et que 
\[f(x)\tendvers{x}{-\infty}-1 \et f(x)\tendvers{x}{+\infty}1\]
on en déduit, d'après le théorème des valeurs intermédiaires, qu'il existe $x\in\R$ tel que $f(x)=y$. Donc $f$ réalise une bijection de $\R$ dans $]-1,1[$. On note désormais $f$ la fonction
\[\dspappli{f}{\R}{]-1,1[}{x}{\frac{x}{\sqrt{x^2+1}}}\]
Calculons sa bijection réciproque. Soit $y\in[0,1[$. Alors
\begin{eqnarray*}
\forall x\in\RP\qsep f(x)=y
&\ssi& \frac{x}{\sqrt{x^2+1}}=y\\
&\ssi& y\sqrt{x^2+1}=x\\
&\ssi& y^2(x^2+1)=x^2 \quad\text{car $y\geq 0$ et $x\geq 0$}\\
&\ssi& x^2=\frac{y^2}{1-y^2}\\
&\ssi& x=\frac{y}{\sqrt{1-y^2}} \quad\text{car $x\geq 0$, $y\geq 0$ et $1-y^2\geq 0$.}
\end{eqnarray*}
Donc
\[\forall y\in[0,1[\qsep f^{-1}(y)=\frac{y}{\sqrt{1-y^2}}.\]
Puisque $f$ est impaire, on en déduit que $f^{-1}$ est impaire donc
\[\forall y\in]-1,1[\qsep f^{-1}(y)=\frac{y}{\sqrt{1-y^2}}.\]
\end{sol}

\magsubsection{Symétries}

\exercice{nom={Symétries de la bijection réciproque}}
Soit $\app{f}{\R}{\R}$ impaire et bijective. Montrer que $f^{-1}$ est impaire.

\exercice{nom={Une fonction périodique étrange}}
Soit $f$ la fonction définie sur $\R$ par
\[\forall x\in\R \qsep
  f(x)\defeq
  \begin{cases}
  1 & \text{si $x\in\Q$}\\
  0 & \text{si $x\in\R\setminus\Q$.}  
  \end{cases}\]
Montrer que l'ensemble des périodes de $f$ est $\Q$. On a donc trouvé une
fonction périodique qui n'admet pas de plus petite période strictement positive.

\magsubsection{Monotonie}


\exercice{nom={Monotonie}}
Soit $f:\R\mapsto\R$ une fonction telle que $f\circ f$ est croissante et
$f\circ f\circ f$ est strictement décroissante. Montrer que $f$ est
strictement décroissante.




\exercice{nom={Monotonie et théorèmes usuels}}
Donner la monotonie (si possible sans dériver) des fonctions d'expressions
\[\textbf{a.}\ \e^{-1/x^2}, \qquad \textbf{b.}\ \e^{1/x^3}, \qquad \textbf{c.}\ x \ln\p{\cos x}
  \ \text{sur $\interfo{0}{\frac{\pi}{2}}$},\]
\[\textbf{d.}\ x\ln\p{1-\frac{1}{x}} \text{sur $\intero{1}{+\infty}$},\qquad
\textbf{e.}\ \sqrt[3]{x+1}-\sqrt[3]{x} \text{ sur $\RPs$},\]
\[\textbf{f.}\ \sin\p{\p{\e^{-x}-1}\pi/2} \text{sur $\RP$.}\]
\begin{sol}
On ne peut pas obtenir la monotonie de
$x\mapsto x \ln\p{1-1/x}$ sur $\intero{1}{+\infty}$.
C'est le produit de deux fonctions croissantes mais on ne peut rien en dire même si la première est positive. On devra donc dériver la fonction.

$a^3-b^3=(a-b)(a^2+ab+b^2)$ donne (ce qui est intéressant sur $\RPs$:
$$\sqrt[3]{x+1}-\sqrt[3]{x}=\frac{\sqrt[3]{x+1}^3-\sqrt[3]{x}^3}{\sqrt[3]{x+1}^2+\sqrt[3]{x+1}\sqrt[3]{x}+\sqrt[3]{x}^2}=\frac{1}{\sqrt[3]{x+1}^2+\sqrt[3]{x+1}\sqrt[3]{x}+\sqrt[3]{x}^2}.$$

Sur $\interfo{0}{\frac{\pi}{2}}$, $\ln(cos)$ est décroissant et négatif donc $-\ln(cos)$ est croissant et positif donc $x\mapsto -x\ln(cos(x))$ est croissant, d'où $x\mapsto x\ln(cos(x))$ croissant.

\end{sol}

\magsubsection{Fonction majorée, minorée, bornée}




% \exercice{nom={Calcul de limites}}
% Déterminer les limites, si elles existent, en $1$ des fonctions
% d'expressions
% \[\frac{\sqrt{x}-\sqrt{x-1}-1}{\sqrt{x^2-1}}, \qquad
%   \frac{x^2-\sqrt{x}}{\sqrt{x}-1}.\]
% Déterminer les limites, si elles existent, des fonctions
% d'expressions
% \[\frac{\sin x-\cos x}{x-\frac{\pi}{4}} \quad \text{en $\frac{\pi}{4}$}, \qquad
%   \frac{\sin (3x)}{1-2\cos x} \quad \text{en $\frac{\pi}{3}$}.\]

% \exercice{nom={Asymptotes}}
% Étudier les variations de
% \[\dspappli{f}{\R}{\R}{x}{\frac{2x^3}{1+x^2}}\]
% ses asymptotes et la position de la courbe de $f$ par rapport à celles-ci. En déduire un graphe précis de $f$.

% \begin{sol}
% $f$ est dérivable sur $\R$ et on a:
% $$\forall x \in \R, f'(x)=\dfrac{2x^4+6x^2}{(1+x^2)^2}\geq 0.$$
% Donc $f$ est croissante sur $\R$.


% Tout d'abord, on constate que $\dfrac{f(x)}{x}\tendvers{x}{\pm\infty}{2}$. On forme donc $$f(x)-2x=\dfrac{-2x}{x^2}\tendvers{x}{\pm\infty}{0},$$. Ainsi, la droite d'équation $y=2x$ est asymptote à $f$ en plus et moins l'infini. De plus, $f(x)-2x<0$ sur $]0;+\infty[$ et $f(x)-2x>0$ sur $]-\infty;0[$ donc la courbe est au-dessus de son asymptote en $-\infty$ et au-dessous en $+\infty$.

% \end{sol}


% \magsection{Continuité}

% \exercice{nom={Résolution d'une équation}}
% Résoudre, l'équation
% \[2x \ln(x)+3(x-1)=0.\]

\magsection{Fonction continue, fonction dérivable}



\magsubsection{Continuité}
\magsubsection{Dérivabilité}




\exercice{nom={Bijection}}
Soit $f$ la fonction définie sur $\RPs$ par
\[\forall x\in\RPs\qsep f(x)\defeq\frac{\e^x}{\e^x-1}.\]
\begin{questions}
\question Montrer que $f$ réalise une bijection de $\RPs$ sur un intervalle $J$ à préciser.
\enonce Dans la suite, on note $g:J\to\RPs$ la bijection réciproque de la corestriction de $f$
  à $J$.
\question Discuter de la monotonie de $g$, de sa continuité et sa dérivabilité. Expliciter
  la dérivée de $g$ sur $J$.
\question Expliciter $g(y)$ en résolvant directement l'équation $f(x)=y$ et retrouver les
  propriétés établies à la question précédente.
\end{questions}

\magsubsection{Dérivation et monotonie}

\exercice{nom={Études de variations}}
Étudier les variations des fonctions suivantes
\[x\mapsto\frac{\ln x}{x} \et x\mapsto\sin x-x+\frac{x^3}{6}.\]

\magsubsection{Dérivation des fonctions à valeurs dans $\C$}

% \exercice{nom={Fonction définie par morceaux}}
% On considère une fonction $f$ dérivable sur le segment $\interf{0}{1}$ avec
% $f\p{0}=f\p{1}$. La fonction $g$ définie par
% \[g(x)\defeq
%   \begin{cases}
%   f\p{2x} & \text{si $0\leq x\leq\frac{1}{2}$}\\
%   f\p{2x-1} & \text{si $\frac{1}{2}<x\leq 1$}
%   \end{cases}\]
% est-elle continue~? dérivable~? Si non, quelles hypothèses faut-il ajouter pour
% que $g$ soit dérivable sur $\interf{0}{1}$~?

% \exercice{nom={Dérivation et symétries}}
% Soit $f$ une fonction dérivable sur $\R$.
% \begin{questions}
% \question On suppose que $f$ est paire. Que peut-on dire de $f'$~?
% \question Même question lorsque $f$ est impaire ou périodique.
% \question Réciproquement, on suppose que $f'$ est impaire
%   (resp. paire, périodique). Que peut-on dire de $f$~?
% \end{questions}






% \exercice{nom={Calculs de dérivées $n$-ièmes}}
% Calculer les dérivées $n$-ièmes des fonctions d'expressions
% \[x\mapsto x^k \quad \text{où $k\in\N$}, \qquad x\mapsto\sin x, \qquad
%   x\mapsto\cos x,\]
% \[x\mapsto\frac{1}{x^k} \quad \text{où $k\in\N$}, \qquad x\mapsto x^\alpha
%   \quad \text{où $\alpha\in\R$}.\]

% \exercice{nom={Une fonction de classe $\classecinf$ définie par morceaux}}
% Soit $f$ la fonction définie sur $\R$ par
% \[\forall x\in\R \qsep f(x)\defeq
%   \begin{cases}
%   \e^{-\frac{1}{x}} & \text{si $x>0$}\\
%   0 & \text{si $x\leq 0$}
%   \end{cases}\]
% \begin{questions}
% \question Montrer que $f$ est dérivable sur $\R$ et calculer sa dérivée.
% \question Montrer que $f$ est de classe $\classec{1}$ sur $\R$.
% \question Montrer que $f$ est de classe $\classec{\infty}$ sur $\Rs$ et que pour tout $n\in\N$, il existe un polynôme $P_n\in\polyR$ tel que
%   $$\forall x\in\Rs \qsep f^{(n)}(x)=
%     \begin{cases}
%     P_n\p{\frac{1}{x}}\e^{-\frac{1}{x}} & \text{si $x>0$}\\
%     0 & \text{si $x<0$.}
%     \end{cases}
%   $$
% % \question En déduire la limite, pour tout entier $n\in\N$, de $f^{(n)}(x)$ lorsque
% %   $x$ tend vers $0$, tout en étant différent de $0$.
% \question À l'aide d'une récurrence, montrer que $f$ est de classe $\classec{\infty}$ sur $\R$.
% \end{questions}
% \begin{sol}
% \begin{questions}
% \question Puisque
% \[\forall x>0\qsep f(x)=\e^{-1/x}\]
% d'après les théorèmes usuels, $f$ est dérivable sur $\RPs$ et
% \[\forall x>0\qsep f'(x)=\frac{1}{x^2}\e^{-1/x}.\]
% Comme de plus
% \[\forall x\leq 0\qsep f(x)=0\]
% on en déduit que $f$ est dérivable sur $\RMs$ et que
% \[\forall x<0\qsep f'(x)=0.\]
% De plus, $f$ est dérivable à gauche en 0 et $f_g'(0)=0$. Montrons que $f$ est dérivable à droite en 0 et que $f_d'(0)=0$. On a, pour $x>0$
% \begin{eqnarray*}
% \frac{f(x)-f(0)}{x}&=&\frac{\e^{-1/x}}{x} \qquad\text{on pose $u\defeq\frac{1}{x}\tendversdp{x}{0}+\infty$.}\\
% &=& u\e^{-u} \tendvers{u}{+\infty} 0
% \end{eqnarray*}
% Donc $f$ est dérivable à droite en 0 et $f_d'(0)=0$. Donc $f$ est dérivable en 0 et $f'(0)=0$. En conclusion
% \begin{conclusion}
% $f$ est dérivable sur $\R$ et
% \[\forall x\in\R\qsep f'(x)=\begin{cases}
% \frac{1}{x^2}\e^{-1/x}&\text{si $x>0$}\\
% 0 &\text{si $x\leq 0$}
% \end{cases}\]
% \end{conclusion}
% \question Montrons que $f$ est de classe $\classec{1}$ sur $\R$. On sait déjà que $f$ est dérivable sur $\R$. Il reste à prouver que sa dérivée est continue. Puisque
% \[\forall x>0\qsep f'(x)=\frac{1}{x^2}\e^{-1/x}\]
% on en déduit, d'après les théorèmes usuels que $f'$ est continue sur $\RPs$. De plus, pour tout $x>0$
% \begin{eqnarray*}
% f'(x)
% &=& \frac{1}{x^2}\e^{-1/x} \qquad\text{on pose $u\defeq\frac{1}{x}\tendversdp{x}{0}+\infty$.}\\
% &=& u^2 \e^{-u} \tendvers{u}{+\infty} 0
% \end{eqnarray*}
% donc
% \[f'(x)\tendversdp{x}{0}0=f'(0).\]
% On en déduit que $f'$ est continue à droite en 0. De plus
% \[\forall x\leq 0\qsep f'(x)=0\]
% donc $f'$ est continue sur $\RMs$ et elle est continue à gauche en 0. En conclusion $f'$ est continue sur $\R$.
% \begin{conclusion}
% $f$ est de classe $\classec{1}$ sur $\R$.
% \end{conclusion}
% \question Puisque
% \[\forall x<0\qsep f(x)=0\]
% on en déduit que $f$ est de classe $\classec{\infty}$ sur $\RMs$ et
% \[\forall n\in\N\qsep \forall x<0\qsep f^{(n)}(x)=0.\]
% De même, puisque
% \[\forall x>0\qsep f(x)=\e^{-1/x}\]
% d'après les théorèmes usuels, $f$ est de classe $\classec{\infty}$ sur $\RPs$. Pour tout $n\in\N$, on pose
% \[\mathcal{H}_n\defeq \text{\og Il existe $P_n\in\polyR$ tel que~:}
%   \ \forall x>0\qsep f^{(n)}(x)=P_n\p{\frac{1}{x}}\e^{-1/x}. \text{\fg}\]
% \begin{itemize}
% \item $\mathcal{H}_0$ est vraie~: En effet, il suffit de poser $P_0\defeq 1$.
% \item $\mathcal{H}_n\implique\mathcal{H}_{n+1}$~: Soit $n\in\N$. On suppose que $\mathcal{H}_n$ est vraie. Montrons que $\mathcal{H}_{n+1}$ est vraie. Puisque $\mathcal{H}_n$ est vraie, il existe $P_n\in\polyR$ tel que
% \[\forall x>0\qsep f^{(n)}(x)=P_n\p{\frac{1}{x}}\e^{-1/x}\]
% donc d'après les théorèmes usuels
% \begin{eqnarray*}
% \forall x>0\qsep f^{(n+1)}(x)
% &=& -\frac{1}{x^2}P_n'\p{\frac{1}{x}}\e^{-1/x} +\frac{1}{x^2}P_n\p{\frac{1}{x}}\e^{-1/x}\\
% &=& P_{n+1}\p{\frac{1}{x}}\e^{-1/x} \quad\text{en posant $P_{n+1}\defeq X^2(P_n-P_n')$.}
% \end{eqnarray*}
% Donc $\mathcal{H}_{n+1}$ est vraie.
% \end{itemize}
% Par récurrence sur $n$, on en déduit que $\mathcal{H}_n$ est vraie pour tout $n\in\N$. En conclusion
% \begin{conclusion}
% $f$ est de classe $\classec{\infty}$ sur $\Rs$ et, pour tout $n\in\N$, il existe $P_n\in\polyR$ tel que
% \[\forall x\in\Rs \qsep f^{(n)}(x)=
%     \begin{cases}
%     P_n\p{\frac{1}{x}}\e^{-\frac{1}{x}} & \text{si $x>0$}\\
%     0 & \text{si $x<0$.}
%     \end{cases}\]
% \end{conclusion}
% \question Commençons par prouver que pour tout $P\in\polyR$
%   \[P\p{\frac{1}{x}}\e^{-1/x}\tendversdp{x}{0}0.\]
%   Soit $P\in\polyR$. Alors, il existe $a_0,\ldots,a_n\in\R$ tels que $P=a_0+a_1X+\cdots+a_n X^n$. Soit $k\in\intere{0}{n}$. Alors
% \begin{eqnarray*}
% \p{\frac{1}{x}}^k\e^{-1/x}&=&\frac{1}{x^k}\e^{-1/x} \qquad\text{on pose $u\defeq\frac{1}{x}\tendversdp{x}{0}+\infty$.}\\
% &=& u^k\e^{-u} \tendvers{u}{+\infty} 0
% \end{eqnarray*}
% donc
% \[\p{\frac{1}{x}}^k\e^{-1/x}\tendversdp{x}{0}0.\]
% On en déduit que
% \[P\p{\frac{1}{x}}\e^{-1/x}=\sum_{k=0}^n \frac{a_k}{x^k}\e^{-1/x}\tendversdp{x}{0}0.\]
% Pour tout $n\in\N$, on pose désormais 
% \[\mathcal{H}_n\defeq \text{\og $f$ est dérivable $n$ fois en 0 et $f^{(n)}(0)=0$. \fg}\]
% \begin{itemize}
% \item $\mathcal{H}_0$ est vraie~: On a bien $f(0)=0$.
% \item $\mathcal{H}_n\implique\mathcal{H}_{n+1}$~: Soit $n\in\N$. On suppose que $\mathcal{H}_n$ est vraie. Montrons que $\mathcal{H}_{n+1}$ est vraie. Puisque $\mathcal{H}_n$ est vraie, $f$ est dérivable $n$ fois en 0 et $f^{(n)}(0)=0$. D'après la question précédente, on en déduit que 
% \[\forall x\in\R\qsep f^{(n)}(x)=
%   \begin{cases}
%     P_n\p{\frac{1}{x}}\e^{-\frac{1}{x}} & \text{si $x>0$}\\
%     0 & \text{si $x\leq 0$.}
%     \end{cases}\]
% On en déduit que $f^{(n)}$ est dérivable à gauche en 0 et que sa dérivée à gauche est nulle. De plus, pour tout $x>0$
% \[\frac{f^{(n)}(x)-f^{(n)}(0)}{x}=\frac{f^{(n)}(x)}{x}=\frac{1}{x}P_n\p{\frac{1}{x}}\e^{-1/x}=Q\p{\frac{1}{x}}\e^{-1/x}\]
% avec $Q\defeq XP_n$. On en déduit que
% \[\frac{f^{(n)}(x)-f^{(n)}(0)}{x}\tendversdp{x}{0}0.\]
% Donc $f^{(n)}$ est dérivable à droite en 0 et sa dérivée est nulle. On en deduit que $f$ est dérivable $n+1$ fois en 0 et que $f^{(n+1)}(0)=0$. Donc $\mathcal{H}_{n+1}$ est vraie.
% \end{itemize}
% Par récurrence sur $n$, on en déduit que pour tout $n\in\N$, $f$ est dérivable $n$ fois en 0. Donc, pour tout $n\in\N$, $f$ est dérivable $n$ fois sur $\R$. En conclusion
% \begin{conclusion}
% $f$ est de classe $\classec{\infty}$ sur $\R$.
% \end{conclusion}
%  % Soit $n\in\N$. Montrons que
%  %  \[f^{(n)}(x)\tendversp{x}{0}0.\]
%  %  Puisque $f^{(n)}(x)=0$ pour $x<0$, on en déduit que la limite pointée à gauche en 0 de $f^{(n)}(x)$ est 0. Montrons que la limite pointée de $f^{(n)}(x)$ à droite en 0 est aussi 0.
% \end{questions}
% \end{sol}

\magsection{Intégration, primitive}


\magsubsection{Primitive}

\exercice{nom={Bijection}}


Soit $f$ la fonction définie sur $\intero{-1}{1}$ par
\[\forall x\in\intero{-1}{1}\qsep f(x)\defeq\frac{1}{\sqrt{1-x^2}}.\]
\emph{Dans cet exercice, il est interdit d'utiliser la fonction $\arcsin$.}
\begin{questions}
\question Montrer que $f$ admet une unique primitive $F$ s'annulant en 0.
\question Montrer que $F(x)$ admet une limite $l\in\Rbar$ lorsque $x$ tend vers 1.
\question Montrer que $F$ est impaire.
\enonce On définit la fonction $\phi$ sur $\intero{-\pi/2}{\pi/2}$ par
  \[\forall x\in\intero{-\pi/2}{\pi/2}\qsep \phi(x)\defeq F(\sin x).\]
\question En dérivant $\phi$, montrer que
	\[\forall x\in\intero{-\pi/2}{\pi/2}\qsep F(\sin x)=x.\]
\question En déduire la valeur de $l$.
\end{questions}

\magsubsection{Intégration et régularité}



\magsubsection{Intégration et inégalité}


\exercice{nom={Étude d'une fonction définie par une intégrale}}
Soit $g$ la fonction d'expression
\[g(x)\defeq\integ{0}{\frac{\pi}{2}}{\frac{\e^t}{1+x\sin t}}{t}.\]
\begin{questions}
\question Montrer que $g$ est définie sur $\intero{-1}{+\infty}$.
\question Montrer que $g$ est décroissante.
% \question Dans cette question, on admet que $f$ est bornée et on se donne un réel
%   $a>-1$. Montrer qu'il existe $M\in\RP$ tel que~:
%   \[\forall x,y\in\interfo{a}{+\infty} \quad \abs{g(x)-g(y)}
%     \leq M\abs{x-y}\]
% \question En déduire que $g$ est continue sur $\intero{-1}{+\infty}$.
\end{questions}

\exercice{nom={Sommes de Riemann de fonctions monotones}}
Soit $f$ une fonction continue et croissante sur $\interf{0}{1}$. On définit
la suite $(u_n)$ par
$$\forall n\in\Ns \qsep u_n\defeq\frac{1}{n}\sum_{k=1}^n f\p{\frac{k}{n}}.$$
\begin{questions}
\question Montrer que pour tout $n\geq 1$, et pour tout $k\in\intere{0}{n-1}$, on
  a
  \[\frac{1}{n}f\p{\frac{k}{n}}\leq\integ{\frac{k}{n}}{\frac{k+1}{n}}{f(t)}{t}
    \leq \frac{1}{n}f\p{\frac{k+1}{n}}.\]
\question En déduire que pour tout entier $n\geq 1$
  \[u_n-\frac{1}{n}\p{f(1)-f(0)} \leq \integ{0}{1}{f(t)}{t}\leq u_n.\]
\question En déduire que la suite $(u_n)$ converge vers
  $\integ{0}{1}{f(t)}{t}$. En donner une interprétation géométrique.
\question Soit $\alpha\in\RPs$. On considère la suite $(v_n)$ définie par
  \[\forall n\in\N \qsep v_n\defeq\sum_{k=0}^n k^\alpha.\]
  Montrer que
  \[v_n\equiS \frac{n^{\alpha+1}}{\alpha+1}\]
  c'est-à-dire que la suite de terme général
  \[v_n\frac{\alpha+1}{n^{\alpha+1}}\]
  converge vers $1$.
\end{questions}
\begin{sol}
\begin{questions}
\question Soit $n\in\Ns$ et $k\in\intere{0}{n-1}$. Alors
  \[0\leq\frac{k}{n}\leq\frac{k+1}{n}\leq 1.\]
  Par croissance de $f$, on en déduit que
  \[\forall t\in\interf{\frac{k}{n}}{\frac{k+1}{n}}\qsep f\p{\frac{k}{n}}\leq f(t)\leq f\p{\frac{k+1}{n}}.\]
  Puisque $k/n\leq (k+1)/n$, on peut intégrer cette inégalité et on obtient
  \[\integ{\frac{k}{n}}{\frac{k+1}{n}}{f\p{\frac{k}{n}}}{t}\leq
    \integ{\frac{k}{n}}{\frac{k+1}{n}}{f(t)}{t}\leq
    \integ{\frac{k}{n}}{\frac{k+1}{n}}{f\p{\frac{k+1}{n}}}{t}\]
  donc
  \[\frac{1}{n}f\p{\frac{k}{n}}\leq\integ{\frac{k}{n}}{\frac{k+1}{n}}{f(t)}{t}
    \leq \frac{1}{n}f\p{\frac{k+1}{n}}.\]
\question Soit $n\in\Ns$. Alors
  \[\forall k\in\intere{0}{n-1}\qsep \frac{1}{n}f\p{\frac{k}{n}}\leq
    \integ{\frac{k}{n}}{\frac{k+1}{n}}{f(t)}{t}
    \leq \frac{1}{n}f\p{\frac{k+1}{n}}.\]
  En sommant cette inégalité pour $k$ allant de 0 à $n-1$, on obtient
  \[\sum_{k=0}^{n-1} \frac{1}{n}f\p{\frac{k}{n}}\leq
    \sum_{k=0}^{n-1} \integ{\frac{k}{n}}{\frac{k+1}{n}}{f(t)}{t}
    \leq \sum_{k=0}^{n-1} \frac{1}{n}f\p{\frac{k+1}{n}}\]
  donc d'après Chasles
  \[\frac{1}{n}\sum_{k=0}^{n-1} f\p{\frac{k}{n}}\leq
    \integ{0}{1}{f(t)}{t}
    \leq \frac{1}{n} \sum_{k=0}^{n-1} f\p{\frac{k+1}{n}}.\]
  En effectuant un changement de variable $k\to k-1$ dans la somme de droite, on obtient
    \[\frac{1}{n}\sum_{k=0}^{n-1} f\p{\frac{k}{n}}\leq
    \integ{0}{1}{f(t)}{t}
    \leq \frac{1}{n}\sum_{k=1}^{n} f\p{\frac{k}{n}}\]
  puis
    \[\frac{1}{n}\cro{\sum_{k=1}^{n} f\p{\frac{k}{n}}-\cro{f(1)-f(0)}}\leq
    \integ{0}{1}{f(t)}{t}
    \leq \frac{1}{n}\sum_{k=1}^{n} f\p{\frac{k}{n}}.\]
  En conclusion
  \[\forall n\in\Ns\qsep u_n-\frac{1}{n}\cro{f(1)-f(0)} \leq \integ{0}{1}{f(t)}{t}\leq u_n.\]
\question Soit $n\in\Ns$. D'après la question précédente
  \[u_n\leq \integ{0}{1}{f(t)}{t}+\frac{1}{n}\cro{f(1)-f(0)}\]
  et
  \[\integ{0}{1}{f(t)}{t}\leq u_n.\]
  Donc
  \[\forall n\in\Ns\qsep \integ{0}{1}{f(t)}{t}\leq u_n \leq\integ{0}{1}{f(t)}{t}+\frac{1}{n}\cro{f(1)-f(0)}.\]
  Or
  \[\integ{0}{1}{f(t)}{t}+\frac{1}{n}\cro{f(1)-f(0)}\tendvers{n}{+\infty}\integ{0}{1}{f(t)}{t}\]
  donc d'après le théorème des gendarmes
  \[u_n\tendvers{n}{+\infty}\integ{0}{1}{f(t)}{t}.\]
\question Soit $\alpha\in\RPs$. On définit la fonction $f$ sur $\interf{0}{1}$ par
  \[\forall t\in\interf{0}{1}\qsep f(t)\defeq t^\alpha.\]
  Alors $f$ est continue et croissante sur $[0,1]$ donc
  \[u_n \tendvers{n}{+\infty} \integ{0}{1}{t^\alpha}{t}=\evaldiff{\frac{1}{\alpha+1}t^{\alpha+1}}{0}{1}=\frac{1}{\alpha+1}.\]
  On définit la suite $(w_n)$ par
  \[\forall n\in\Ns\qsep w_n\defeq \frac{\alpha+1}{n^{\alpha+1}}\sum_{k=0}^n k^\alpha.\]
  Alors
  \begin{eqnarray*}
\forall n\in\Ns\qsep w_n
&=& \frac{\alpha+1}{n}\sum_{k=0}^n \p{\frac{k}{n}}^\alpha\\
&=& \frac{\alpha+1}{n}\sum_{k=0}^n f\p{\frac{k}{n}}\\
&=& \frac{\alpha+1}{n}\sum_{k=1}^n f\p{\frac{k}{n}}\\
&=& (\alpha+1)u_n \tendvers{n}{+\infty} 1.
  \end{eqnarray*}
  Donc
\[\sum_{k=0}^n k^\alpha\equi{n}{+\infty}\frac{n^{\alpha+1}}{\alpha+1}.\]
\end{questions}
\end{sol}

\magsubsection{Intégration par parties, changement de variable}

\exercice{nom={Intégrales de Wallis}}
Pour tout entier $n\in\N$, on définit $I_n$ et $J_n$ par
\[I_n\defeq\integ{0}{\frac{\pi}{2}}{\sin^n t}{t} \qquad
  J_n\defeq\integ{0}{\frac{\pi}{2}}{\cos^n t}{t}\]
\begin{questions}
\question Montrer que pour tout entier $n\in\N$, $I_n=J_n$.
\question Montrer que les suites $(I_n)$ et $(J_n)$ sont positives et
  décroissantes. Calculer $I_0$ et $I_1$.
\question Montrer que pour tout entier $n\in\N$
  \[I_{n+2}=\frac{n+1}{n+2}I_n.\]
\question En déduire que pour tout entier $n\in\N$
  \[I_n I_{n+1}=\frac{\pi}{2(n+1)}.\]
\question En déduire que
  \[I_n\equiS\sqrt{\frac{\pi}{2n}}\]
  c'est-à-dire que la suite de terme général $I_n\sqrt{\frac{2n}{\pi}}$
  converge vers $1$.
\end{questions}
\begin{sol}
\begin{questions}
\question On effectue le changement de variable $t=\dfrac{\pi}{2}-u$ ($dt=-du$). Cela donne :
$$I_n=-\integ{\frac{\pi}{2}}{0}{\sin^n \p{\dfrac{\pi}{2}-u}}{u}=\integ{0}{\frac{\pi}{2}}{\cos^n (u)}{u}=J_n.$$
\question La suite $(I_n)$ est positive car son intégrande l'est. De plus,
$$I_{n+1}-I_n=\integ{0}{\frac{\pi}{2}}{\sin^{n+1} t-\sin^n t}{t}=\integ{0}{\frac{\pi}{2}}{\sin^{n} t(\sin t-1)}{t}\leq 0$$ donc la suite est décroissante. De même pour $J_n$ car $J_n=I_n$.
 On a $I_0=\dfrac{\pi}{2}$ et $I_1=1$.
\question Soit $n\in\N$. On effectue une IPP :
\begin{eqnarray*}
I_{n+2}&=&\integ{0}{\frac{\pi}{2}}{\sin^{n+1} (t)\sin (t)}{t}=\left[-\sin^{n+1} (t) \cos (t)\right]_0^{\frac{\pi}{2}}-\integ{0}{\frac{\pi}{2}}{(n+1)(-\cos (t) \sin^n (t))(-\cos(t))}{t}\\
&=&0-(n+1)\integ{0}{\frac{\pi}{2}}{(1-\sin^2(t))\sin^n(t)}{t}\\
&=&(n+1)(I_n-I_{n+2})
\end{eqnarray*}
D'où le résultat souhaité.
\question Montrons par récurrence pour tout entier $n\in\N$, $H_n$ : "$I_n I_{n+1}=\dfrac{\pi}{2(n+1)}.$"
\begin{itemize}
\item [$\bullet$] $H_0$ est vraie d'après le calcul de $I_0$ et $I_1$ fait précédemment.
\item [$\bullet$] Soit $n\in \N$. Supposons $H_n$ vraie. Alors $$I_{n+1}I_{n+2}=I_{n+1}I_n\frac{n+1}{n+2}\underset{H.R}{=}\frac{\pi}{2(n+1)}\frac{n+1}{n+2}=\frac{\pi}{2(n+2)}.$$ $H_{n+1}$ est donc vraie et on a donc bien démontré le résultat par récurrence.
\end{itemize}
\question Posons $u_n=I_n\sqrt{\frac{2n}{\pi}}$. On a
$$\frac{n+1}{n+2}=\frac{I_{n+2}}{I_n}\leq \frac{I_{n+1}}{I_n}\leq 1 $$ donc $\frac{I_{n+1}}{I_n} \tendvers{n}{+\infty}{1}$ d'après le théorème des gendarmes.
Ainsi, $$u_n^2=\frac{2}{\pi}n I_n^2=\frac{2}{\pi}\times\frac{n}{n+1}\times (n+1)I_n I_{n+1}\times \frac{I_n}{I_{n+1}}\tendvers{n}{+\infty}{\frac{2}{\pi}\frac{\pi}{2}}=1.$$
On peut alors en déduire que $u_n \tendvers{n}{+\infty}{1}.$
\end{questions}
\end{sol}

\exercice{nom={Intégrales}}
Pour tout $n\in\N$, on pose
\[I_n\defeq\integ{0}{1}{(1-x^2)^n}{x}.\]
\begin{questions}
\item Établir une relation de récurrence entre $I_{n}$ et $I_{n+1}$.
\item Calculer $I_{n}$.
\item En déduire \[\sum\limits_{k=0}^{n} \frac{(-1)^k}{2k+1}\binom{n}{k}.\]
\end{questions}

\begin{sol}
\begin{questions}
\item 
\begin{eqnarray*}
I_{n+1}&=&\primppdi{(1-x^2)^{n+1}}{1}{x}\\
&=&\left[x(1-x^2)^{n+1}\right]_0^1-\integ{0}{1}{-(n+1)2x^2(1-x^2)^n}{x}\\
&=&2(n+1)\left[\integ{0}{1}{(x^2-1)(1-x^2)^n}{x}+\integ{0}{1}{(1-x^2)^n}{x}\right]\\
&=&2(n+1)(I_n-I_{n+1}) .
\end{eqnarray*}
Ainsi, $(2n+3)I_{n+1}=2(n+1)I_n$, d'où :
$$\boxed{I_{n+1}=\frac{2n+2}{2n+3}I_n.}$$
\item On a $$\displaystyle I_n=\prod_{k=0}^{n-1}\frac{I_{k+1}}{I_k}\times I_0=\prod_{k=0}^{n-1}\frac{2k+2}{2k+3}\times 1=\frac{2\times 4\times \ldots \times 2n }{3\times 5\times \ldots \times (2n+1)}=\frac{2^{2n}(n!)^2}{(2n+1)!}.$$
\item 
\begin{eqnarray*}
I_{n}&=&\integ{0}{1}{\sum_{k=0}^n\binom{n}{k}(-1)^kx^{2k}}{x}\\
&=&\sum_{k=0}^n\binom{n}{k}(-1)^k\integ{0}{1}{x^{2k}}{x}\\
&=&\sum_{k=0}^n\binom{n}{k}(-1)^k\left[\frac{x^{2k+1}}{2k+1}\right]_0^1\\
&=&\sum_{k=0}^{n} \frac{(-1)^k}{2k+1}\binom{n}{k} .
\end{eqnarray*}
\end{questions}
\end{sol}

\magsubsection{Calcul de primitive}

%%\pagebreak

% \exercice{nom={Inégalités}}
% \begin{questions}
% \question Soit $x,y\in\R$ tels que $0<x\leq y$. Montrer que~:
%   \[\frac{y-x}{2\sqrt{y}}\leq\sqrt{y}-\sqrt{x}\leq\frac{y-x}{2\sqrt{x}}\]
% \question Montrer que pour tout $x,y\in\R$~:
%   \[\abs{\sqrt[3]{2+\sin y}-\sqrt[3]{2+\sin x}}\leq\frac{1}{3}\abs{y-x}\]
% \end{questions}

\exercice{nom={Calcul de primitives}}
Donner le domaine de définition et calculer les primitives suivantes
\[\textbf{a.}\ \prim{\p{x^2+x+1}\e^x}{x}, \qquad \textbf{b.}\ \prim{\p{x^2-1}\cos x}{x}, \qquad
\textbf{c.}\ \prim{x^3\ln x}{x},\]
\[\textbf{d.}\ \prim{\sin^2 x\cos^3 x}{x}, \qquad \textbf{e.}\ \prim{\sin x\cos^2 x}{x}, \qquad
\textbf{f.}\ \prim{\sin^2 x\cos^2 x}{x},\]
\[\textbf{g.}\ \prim{\frac{x}{x^2+1}}{x}, \qquad \textbf{h.}\ \prim{\frac{1}{x\ln x}}{x}, \qquad
\textbf{i.}\ \prim{\ln^n x}{x} \quad (\text{pour $n\in\N$}).\]


\begin{sol}
\begin{questions}
\question On se place sur $\R$.
\begin{eqnarray*}
\primppdi{(x^2+x+1)}{e^x}{x}&=&(x^2+x+1)e^x-\primppdi{(2x+1)}{e^x}{x}\\
&=&(x^2+x+1)e^x-(2x+1)e^x+2\prim{e^x}{x}\\
&=&(x^2-x+2)e^x .
\end{eqnarray*}
\question On se place sur $\R$.
\begin{eqnarray*}
\primppdi{(x^2-1)}{\cos(x)}{x}&=&(x^2-1)\sin(x)-\primppdi{2x}{\sin(x)}{x}\\
&=&(x^2-1)\sin(x)+2x\cos(x)-\prim{2\cos(x)}{x}\\
&=&(x^2-3)\sin(x)+2x\cos(x).
\end{eqnarray*}
\question On se place sur $\RPs$.
\begin{eqnarray*}
\primppid{x^3}{\ln(x)}{x}&=&\frac{x^4}{4}\ln(x)-\prim{\frac{x^4}{4}\frac{1}{x}}{x}\\
&=&\frac{x^4}{4}\ln(x)-\frac{1}{16}x^4.
\end{eqnarray*}

\question On se place sur $\R$.
\begin{equation*}
    \begin{split}
    \prim{\sin^2 x\cos^3 x}{x}&=\prim{\sin^2 x \p{1-\sin^2 x}\cos x}{x}\\
                              &=\prim{t^2\p{1-t^2}}{t}\\
                              &=\frac{1}{3}t^3-\frac{1}{5}t^5\\
                              &=\frac{1}{3}\sin^3 x-\frac{1}{5}\sin^5 x\\
    \end{split}
    \end{equation*}    
\question On se place sur $\R$.

$$\prim{\sin(x)\cos^2(x)}{x}=-\frac{1}{3}\cos^3(x).$$
\question On se place sur $\R$.
    \begin{eqnarray*}
    \cos^2 x\sin^2 x&=&\p{\frac{1}{2}\sin\p{2x}}^2\\
     &=&\frac{1}{4}\frac{1-\cos\p{4x}}{2}\\
     &=&\frac{1}{8}-\frac{1}{8}\cos\p{4x}
    \end{eqnarray*}
    Donc~:
    $$\prim{\cos^2 x\sin^2 x}{x}=\frac{1}{8}x-\frac{1}{32}\sin\p{4x}$$
\question On se place sur $\R$.
$$\prim{\frac{x}{x^2+1}}{x}=\frac{1}{2}\prim{\frac{2x}{x^2+1}}{x}=\frac{1}{2}\ln(x^2+1).$$

\question On se place sur $I=\intero{0}{1}$ ou sur $I=\intero{1}{+\infty}$.
$$\prim{\frac{1}{x\ln(x)}}{x}=\ln|\ln x|.$$

\question Sur $\RPs$, on a
\[I_n(x)=x\sum_{k=0}^n (-1)^k \frac{n!}{(n-k)!}\ln^{n-k} x\]   
\end{questions}
\end{sol}

% \exercice{nom={Calcul d'une valeur approchée de $\sin 1$}}
% \begin{questions}
% \question Montrer que pour tout entier $n$ et tout réel $x$~:
%   $$\sin x=\sum_{k=0}^n \frac{(-1)^k}{(2k+1)!}x^{2k+1}+
%     \integ{0}{x}{\frac{\p{x-t}^{2n+2}}{\p{2n+2}!}\sin^{(2n+3)}(t)}{t}$$
% \question En déduire que~:
%   $$\forall x\in\R \quad \abs{\sin x-\sum_{k=0}^n
%     \frac{(-1)^k}{(2k+1)!}x^{2k+1}} \leq \frac{\abs{x}^{2n+3}}{\p{2n+3}!}$$
% \question En déduire une méthode pour calculer $\sin 1$ avec 100 chiffres exacts 
%   après la virgule avec un ordinateur capable de calculer avec une précision
%   illimitée.
% \end{questions}
% \exercice{nom={Primitives de fractions rationnelles}}
% \begin{questions}
% \question Calculer sur un intervalle $I$ que l'on précisera
%   \[\priminv{1-x^4}{x}, \qquad \prim{\frac{x}{1+x^4}}{x}, \qquad \prim{\frac{x^2}{x^3-1}}{x}.\]
% \question Calculer sur un intervalle $I$ que l'on précisera
%   \[\priminv{1+x+x^2}{x}, \qquad \prim{\frac{2-5x}{1+x^2}}{x},\]
%   \[\prim{\frac{3x+2}{2x^2-4x+3}}{x}, \qquad
%     \prim{\frac{x+3}{x^2-2x+5}}{x}.\]
% % \question Calculer
% %   \[\integ{0}{1}{\frac{x^3+x}{(x^2+x+1)^2}}{x}.\]
% \end{questions}
% \begin{sol}
% $\quad$
% \begin{questions}
% \question On décompose en éléments simples :
% $$\frac{1}{1-X^4}=\frac{1}{2}\frac{1}{1+X^2}+\frac{1}{4}\frac{1}{1-X}+\frac{1}{4}\frac{1}{1+X},$$ d'où :
%   \[\priminv{1-x^4}{x}=\frac{1}{2}\arctan x+\frac{1}{4}\ln\p{\frac{1+x}{1-x}} \]
%   \[\prim{\frac{x}{1+x^4}}{x}=\frac{1}{2}\arctan\p{x^2}\]
%   \[\prim{\frac{x^2}{x^3-1}}{x}=\frac{1}{3}\ln\abs{x^3-1}\]
% \question Dans cette question, on a affaire à des éléments simples de seconde espèce :
% \begin{eqnarray*}
% \prim{\frac{1}{1+x+x^2}}{x}&=&\prim{\frac{1}{\p{x+\frac{1}{2}}^2+\frac{3}{4}}}{x}\\
% &=&\prim{\frac{1}{\frac{3}{4}\p{(\frac{2}{\sqrt{3}}(x+\frac{1}{2}))^2+1}}}{x}\\
% &\underset{u=\frac{2}{\sqrt{3}}(x+\frac{1}{2})}{=}&\frac{4}{3}\frac{\sqrt{3}}{2}\prim{\frac{1}{u^2+1}}{u}\\
% &=&\boxed{\frac{2}{\sqrt{3}}\arctan(\frac{2}{\sqrt{3}}(x+\frac{1}{2})} .
% \end{eqnarray*}

% \begin{eqnarray*}
% \prim{\frac{3x+2}{2x^2-4x+3}}{x}&=&\frac{3}{4}\prim{\frac{4x-4}{2x^2-4x+3}}{x}+\prim{\frac{5}{2x^2-4x+3}}{x}\\
% &=&\frac{3}{4}\ln(2x^2-4x+3)+\frac{5}{2}\prim{\frac{1}{(x-1)^2+\frac{1}{2}}}{x}\\
% &=&\frac{3}{4}\ln(2x^2-4x+3)+\frac{5}{2}\prim{\frac{1}{\frac{1}{2}\p{(\sqrt{2}(x-1))^2+1}}}{x}\\
% &\underset{u=\sqrt{2}(x-1)}{=}&\frac{3}{4}\ln(2x^2-4x+3)+\frac{5}{\sqrt{2}}\prim{\frac{1}{u^2+1}}{u}\\
% &=&\boxed{\frac{3}{4}\ln(2x^2-4x+3)+\frac{5}{\sqrt{2}}\arctan(\sqrt{2}(x-1))} .
% \end{eqnarray*}
% % \question Le mieux est de faire cet exercice en cherchant une primitive sous
% %   une forme donnée.
% %   \[\prim{\frac{x^3+x}{(x^2+x+1)^2}}{x}=\frac{1}{2} \ln\p{1+x+x^2}
% %     -\frac{7\sqrt{3}}{9}
% %     \arctan\p{\frac{2 x+1}{\sqrt{3}}}+\frac{x-1}{3 \p{1+x+x^2}}\]
% %   On trouve
% %   \[\frac{1}{54}\p{18+27\ln 3-7\sqrt{3}\pi}\]
% \end{questions}
% \end{sol}

% \exercice{nom={Primitives de fractions rationnelles}}
% On pose, pour tout $n \in \N$
% \[I_n\defeq\priminv{(1+x^2)^n}{x}.\]
% À l'aide d'une intégration par parties, obtenir une relation de récurrence entre $I_{n+1}$ et $I_n$. En déduire $I_1,I_2$ et $I_3$.

% \begin{sol}
% Par une IPP,
% $$\priminv{(1+x^2)^n}{x}=\frac{x}{(1+x^2)^n}+2n\prim{\frac{x^2}{(1+x^2)^{n+1}}}{x}=\frac{x}{(1+x^2)^n}+2n\p{\prim{\frac{1+x^2}{(1+x^2)^{n+1}}}{x}-\prim{\frac{1}{(1+x^2)^{n+1}}}{x}}.$$
%  D'où :
%  $$I_{n+1}=\frac{x}{2n(1+x^2)^n}+\frac{2n-1}{2n}I_n.$$
%  On trouve alors :
%  $$I_1=\arctan(x)$$
%  $$I_2=\frac{x}{2(1+x^2)}+\frac{1}{2}\arctan(x)$$
%  $$I_3=\frac{x}{4(1+x^2)^2}+\frac{3}{8}\frac{x}{1+x^2}+\frac{3}{8}\arctan(x)$$
% \end{sol}


% \exercice{nom={Calcul de primitives}}
% Calculer les primitives suivantes en précisant à chaque fois le ou les intervalles où ceci est légitime.
% \[\prim{\exp(-\sqrt{x})}{x}, \qquad \prim{\cos(\ln (x))}{x}, \qquad \prim{\ln(1+x^2)}{x},\]
% \[\prim{\arctan(x)}{x}, \qquad \prim{x\sin^2 (x)}{x}\]
% %, \qquad \prim{\frac{x^2}{(x-1)^2 (x^2 + 4)}}{x},\]
% \[\prim{\cos(2x) \cos^2(x)}{x}, \qquad \priminv{\sin(x)+\tan(x)}{x} \text{ sur $\intero{0}{\pi/2}$ en posant $u=\cos(t)$.}\]

% \begin{sol}
% \begin{questions}
% \question On se place sur $\RPs$. On pose $u=\sqrt{x}$ donc $du=\dfrac{1}{2\sqrt{x}}dx=\dfrac{1}{2u}du$ d'où : 
% \begin{eqnarray*}
% \prim{\exp(-\sqrt{x})}{x}&=&\primppid{e^{-u}}{2u}{u}\\
% &=&-2ue^{-u}-\prim{-2e^{-u}}{u}\\
% &=&-2ue^{-u}-2e^{-u}\\
% &=&-2(u+1)e^{-u}\\
% &=&\boxed{-2(\sqrt{x}+1)e^{-\sqrt{x}}} .
% \end{eqnarray*}
% \question On se place sur $\RPs$. On pose $u=\ln{x}$ donc $du=\dfrac{1}{x}dx=\dfrac{1}{e^u}du$ d'où : 
% \begin{eqnarray*}
% \prim{\cos(\ln (x))}{x}&=&\prim{\cos(u)e^u}{u}\\
% &=&\Re\p{\prim{e^{u+iu}}{u}}\\
% &=&\Re\p{\frac{1}{1+i}e^{u+iu}}\\
% &=&\Re\p{\frac{1-i}{2}e^{u}e^{iu}}\\
% &=&\p{\frac{1}{2}\cos(u)+\frac{1}{2}\sin(u)}e^u\\
% &=&\boxed{x\p{\frac{1}{2}\cos(\ln(x))+\frac{1}{2}\sin(\ln(x))}} .
% \end{eqnarray*}

% \question On se place sur $\R$.

% \begin{eqnarray*}
% \prim{\ln(1+x^2)}{x}&=&\primppid{1}{\ln(1+x^2)}{x}\\
% &=&x\ln(1+x^2)-\prim{\frac{2x^2}{1+x^2}}{x}\\
% &=&x\ln(1+x^2)-2\prim{\frac{1+x^2-1}{1+x^2}}{x}\\
% &=&x\ln(1+x^2)-2\prim{1}{x}+2\prim{\frac{1}{1+x^2}}{x}\\
% &=& \boxed{x\ln(1+x^2)-2x+2\arctan(x)} .
% \end{eqnarray*}
% \question \begin{equation*}
%   \begin{split}
%   \prim{\arctan x}{x}&=\primppid{1}{\arctan x}{x}\\
%                  &=x\arctan x-\prim{\frac{x}{1+x^2}}{x}\\
%                  &=\boxed{x\arctan x-\frac{1}{2}\ln(1+x^2)} .
%   \end{split}
%   \end{equation*}
% \question
% \question
% \question
% \end{questions}
% \end{sol}


% \exercice{nom={Primitives de fractions rationnelles en $\cos,\sin$}}
% Calculer sur un intervalle $I$ que l'on précisera les primitives des
% fonctions dont les expressions sont
% \[\frac{1}{\cos x\sin^3 x}, \qquad \frac{1}{\sin x+\cos x}.\]
% \begin{sol}
% $\quad$
% \begin{questions}
%   \question On doit décomposer en éléments simples~:
%   \[\frac{1}{X \left(X^2-1 \right)^2}=
%     \frac{1}{4 (X-1)^2}-\frac{1}{2 (X-1)}+\frac{1}{X}-\frac{1}{4 \
%     (X+1)^2}-\frac{1}{2 (X+1)}\]
%   Puis~:
%   \[\priminv{\cos x\sin^3 x}{x}=\ln\abs{\tan x}-\frac{1}{2\sin^2 x}\]
% % \question
% %   \[\priminv{1+2\cos x}{x}=
% %     \frac{2\sqrt{3}}{3}\arctan\p{\frac{\sqrt{3}\tan\p{x/2}}{3}}\]
% \question
%   \[\priminv{\sin x+\cos x}{x}=\sqrt{2}
%     \argth\p{\frac{\tan\p{\frac{x}{2}}-1}{\sqrt{2}}}\]
% \end{questions}
% \end{sol}


% \exercice{nom={Primitive de fractions rationnelles en $\sh$}}
% Exprimer $\th(x), \ch(x)$ et $\sh(x)$ en fonction de \[t=\th\p{\frac{x}{2}}.\]
% Calculer \[\priminv{\sh(x)}{x}\] en posant $t\defeq\th(x/2)$.

% \exercice{nom={Avec des radicaux}}
% Calculer, sur un intervalle $I$ que l'on précisera
% \[\prim{\frac{1}{2}x\sqrt{\frac{x-1}{x+1}}}{x}.\]
%, \qquad \frac{x}{\sqrt{-x^2+x+2}}.\]
%END_BOOK

\end{document}


\chapter{Compléments d'algèbre}
\setcounter{numeroexercicecours}{1}
\setcounter{numeroexercicecours}{1}
\documentclass{magnoliaold}

\magtex{tex_driver={pdftex},
        tex_packages={tabularx,epigraph,xypic},
        tex_pstricks={pstricks}}
\magfiche{document_nom={Compléments d'algèbre},
          auteur_nom={François Fayard},
          auteur_mail={fayard.prof@gmail.com}}
\magcours{cours_matiere={maths},
          cours_niveau={mpsi},
          cours_chapitre_numero={4},
          cours_chapitre={Compléments d'algèbre}}
\magmisenpage{misenpage_presentation={tikzvelvia},
          misenpage_format={a4},
          misenpage_nbcolonnes={1},
          misenpage_preuve={non},
          misenpage_sol={non}}
% \magmisenpage{}
\maglieudiff{}
\magprocess

\begin{document}

%BEGIN_BOOK
\setlength\epigraphwidth{.8\textwidth}
\epigraph{\og The closer one looks, the more subtle and remarkable Gaussian elimination appears.\fg}{--- \textsc{Nick Trefethen (1955--)}}
% \bigskip
% \hfill\includegraphics[width=0.5\textwidth]{../../Commun/Images/maths-cours-three.jpg}
\magtoc


% \section{Polynômes et fractions rationnelles}

\bigskip

\section{Somme et produit, fonction polynôme}

\subsection{Somme}

\begin{definition}[utile=-3]
Soit $m,n\in\Z$ tels que $m \leq n$ et $u_{m},u_{m+1},\ldots,u_{n-1},u_n\in\C$. On
définit
\[\sum_{k=m}^n u_k\defeq u_{m}+u_{m+1}+\cdots+u_{n-1}+u_n.\]
\end{definition}

\begin{remarqueUnique}
\remarque Lorsque $n=m-1$, la convention est de poser $\sum_{k=m}^n u_k\defeq 0$. Cette
  convention permet d'écrire
  \[\forall n\geq m\qsep \sum_{k=m}^n u_k = u_n + \sum_{k=m}^{n-1} u_k.\]
  % Plus généralement, lorsque $n<m$, on utilise souvent la convention
  % $\sum_{k=m}^n u_k\defeq 0$.
\remarque Si $m,n\in\Z$ sont tels que $n\geq m-1$, alors $\card(\intere{m}{n})=n-m+1$. En
  particulier, quel que soit $a\in\C$
  \[\sum_{k=m}^{n} a =\p{n-m+1}a.\]
\end{remarqueUnique}


\begin{exoUnique}
% \exo Soit $m,n\in\Z$ tels que $m\leq n$ et $a\in\C$. Calculer
%   \[\sum_{k=m}^n a.\]
\exo Écrire avec le symbole $\sum$ les sommes suivantes, sachant que chacune d'elle est composée de $n+1$ termes.
  \[-a_0+a_1-a_2 +a_3+\cdots, \qquad a_1+a_4+a_7+\cdots\]
  \[a_0+2a_1+3a_2+4a_3+\cdots, \qquad a_0-2a_1+4a_2-8a_3+\cdots.\]
\end{exoUnique}

\begin{proposition}
Soit $m,n\in\Z$, $\lambda,\mu\in\C$ et $(u_k)_{k\in\Z}, (v_k)_{k\in\Z}$ deux suites d'éléments de $\C$. Alors
\[\sum_{k=m}^n \p{\lambda u_k+\mu v_k}=\lambda\sum_{k=m}^n u_k + \mu\sum_{k=m}^n v_k.\]
\end{proposition}

\begin{proposition}[utile=-3]
\begin{itemize}
\item Soit $m,n\in\Z$ et $p\in\Z$. Alors
  \[\sum_{k=m}^n u_k=\sum_{k=m-p}^{n-p} u_{k+p}.\]
\item Soit $n\in\N$. Alors
  \[\sum_{k=0}^n u_k=\sum_{k=0}^n u_{n-k}.\]
% \item Soit $n\in\N$. Alors~:
%   \[\sum_{k=0}^n u_k=\sum_{k=0}^{\ent{\frac{n}{2}}} u_{2k}+
%     \sum_{k=0}^{\ent{\frac{n-1}{2}}} u_{2k+1}\]
\end{itemize}
\end{proposition}

\begin{remarques}
\remarque En pratique, lorsque l'on souhaite faire la première transformation,
  on dit qu'on effectue le changement de variable $k\to k+p$.
  \[\sum_{k=m}^n u_k=\text{\og} \sum_{k=m}^{k=n} u_k=\sum_{k+p=m}^{k+p=n} u_{k+p}=
    \sum_{k=m-p}^{k=n-p} u_{k+p}\text{\fg}=\sum_{k=m-p}^{n-p} u_{k+p}.\]
  Le seconde transformation se fait de manière similaire, en utilisant cette fois la
  convention que si les bornes ne sont pas \og dans le bon sens \fg, on les échange; on
  dit dans ce cas qu'on fait le changement de variable $k\to n-k$.
  \[\sum_{k=0}^n u_k=\text{\og} \sum_{k=0}^{k=n} u_k=\sum_{n-k=0}^{n-k=n} u_{n-k}=
  \sum_{k=n}^{k=0} u_{n-k}=\sum_{k=0}^{k=n} u_{n-k}\text{\fg}=\sum_{k=0}^{n} u_{n-k}.\]

\remarque Soit $(u_k)_{k\in\Z}$ une suite. Si $m,n\in\Z$ sont tels que $n\geq m-1$
  \[\sum_{k=m}^n \p{u_{k+1} - u_k} = u_{n+1}-u_m.\]
  On dit qu'une telle somme est \emph{télescopique}.
\end{remarques}

\begin{exos}
\exo
  \begin{questions}
  \question Montrer qu'il existe $a,b\in\R$ tels que
    \[\forall k\in\Ns\qsep \frac{1}{k(k+1)}=\frac{a}{k}+\frac{b}{k+1}.\]
  \question Pour tout $n\in\N$, en déduire une expression simple de
    \[\sum_{k=1}^n \frac{1}{k(k+1)}.\]
  \end{questions}
% \exo Calculer, pour tout $n\in\Ns$
%   \[\sum_{k=1}^n \frac{1}{k(k+1)(k+2)}.\]
%   \begin{sol}
% \[\sum_{k=1}^n \frac{1}{k(k+1)(k+2)}=\frac{n(n+3)}{4(n+1)(n+2)}\]
%   \end{sol}
% \exo Calculer explicitement la suite $(u_n)$ définie par
%   \[u_0\defeq 1 \et \forall n\in\N\qsep u_{n+1}\defeq u_n+2n+3.\]
%   \begin{sol}
% \[u_n=(n+1)^2\]
%   \end{sol}
\exo Pour tout $n\in\N$, en effectuant le changement de variable $k\to k+1$, calculer
  \[\sum_{k=0}^n k 2^k.\]
 \begin{sol}
\[\sum_{k=1}^n k 2^k=2^{n+1}(n-1)+2\]
 \end{sol}
\exo Soit $n\in\Ns$. En considérant les racines $(2n+1)$-ièmes de l'unité, montrer
  que
  \[\sum_{k=1}^n \cos\p{\frac{2k\pi}{2n+1}}=-\frac{1}{2}.\]
% \exo Soit $n\in\N$. Calculer
%   \[\sum_{k=1}^n \p{\frac 1k - \frac 1{n+1-k}}.\]
%   \begin{sol}
% Un changement d'indice montre que cette somme est nulle.
%   \end{sol}
\end{exos}

  
\begin{definition}
Soir $r\in\C$. Une suite $(u_n)$ est dite \emph{en progression arithmétique de raison $r$} lorsque
\[\forall n \in \N\qsep  u_{n+1}= u_n+r.\]
On a alors, pour tout $n \in \N$, $u_n = u_0 + nr$.
\end{definition}

\begin{proposition}
Soit $(u_n)$ une suite en progression arithmétique et $m,n\in\Z$ tels que $n\geq m-1$. Alors
\begin{eqnarray*}
\sum_{k=m}^n u_k
&=& \frac{u_m + u_n}{2}\cdot\p{n-m+1}\\
&=& \frac{\text{premier terme}+\text{dernier terme}}{2}\cdot
    \p{\text{nombre de termes}}.
\end{eqnarray*}
\end{proposition}

\begin{preuve}
On fait la méthode de Gauss en sommant deux fois et regroupant les termes extrémaux, la somme étant constante...
\end{preuve}

\begin{proposition}
  Soit $n\in\N$. Alors
  \[\sum_{k=0}^n k=\frac{n(n+1)}{2} \et \sum_{k=0}^n k^2=\frac{n(n+1)(2n+1)}{6}.\]
  \end{proposition}


% \begin{exoUnique}
% \exo Montrer que pour tout $n\in\N$
%   \[\sum_{k=0}^n k^2 = \frac{n(n+1)(2n+1)}{6}.\]
% \end{exoUnique}
% \begin{proposition}
% Soit $a,b\in\Z$ tels que $b\geq a-1$. Alors
% \[\sum_{k=a}^b k=\frac{(a+b)(b-a+1)}{2}.\]
% \end{proposition}

\begin{definition}
Soit $q\in\C$. Une suite $(u_n)$ est dite \emph{en progression géométrique de raison $q$} lorsque
\[\forall n \in \N\qsep  u_{n+1}= q u_n.\]
On a alors, pour tout $n \in \N$, $u_n = q^n u_0$.
\end{definition}

\begin{proposition}
Soit $(u_n)$ une suite en progression géométrique dont la raison $q\in\C$ est différente de 1 et $m,n\in\Z$ tels que $m\leq n$. Alors
\begin{eqnarray*}
\sum_{k=m}^n u_k
&=& \frac{u_m - u_{n+1}}{1-q}\\
&=& \frac{\text{premier terme}-\text{terme suivant}}{1-q}.
\end{eqnarray*}
\end{proposition}

\begin{exos}
\exo Montrer que la suite de terme général
  \[\sum_{k=0}^n \frac{1}{2^k}\]
  est convergente.
\exo Calculer, pour tout $x\in\R\setminus\ens{1}$
  \[\sum_{k=0}^n k x^k.\]
  \begin{sol}
\[\sum_{k=0}^n k x^k=\frac{x}{(1-x)^2}\cro{1+x^n\p{nx-(n+1)}}\]
  \end{sol}
\end{exos}

\begin{proposition}
Soit $n\in\N$. Alors
  \[\sum_{k=0}^n u_k=\sum_{k=0}^{\ent{\frac{n}{2}}} u_{2k}+
    \sum_{k=0}^{\ent{\frac{n-1}{2}}} u_{2k+1}\]
\end{proposition}

% \begin{remarqueUnique}
% \remarque Lorsque $x\in\R$, $\ent{x}$ est appelée \emph{partie entière} de $x$. C'est l'unique $n\in\Z$ tel que
%   \[n\leq x< n+1.\]
%   Par exemple $\ent{\pi}=3$ et $\ent{-\pi}=-4$.
% \end{remarqueUnique}

\begin{exoUnique}
\exo Calculer la somme
  \[\sum_{k=1}^{2n} (-1)^k k.\]
  \begin{sol}
\[\sum_{k=1}^{2n} (-1)^k k=n\]
  \end{sol}
\end{exoUnique}


\subsection{Produit}

\begin{definition}[utile=-3]
Soit $m,n\in\Z$ tels que $m\leq n$ et $u_{m},u_{m+1},\ldots,u_{n-1},u_n\in\C$. On
définit
\[\prod_{k=m}^n u_k\defeq u_{m}\cdot u_{m+1}\cdots u_{n-1}\cdot u_n.\]
\end{definition}

\begin{remarques}
\remarque Lorsque $n=m-1$, la convention est de poser $\prod_{k=m}^n u_k\defeq 1$. Cette
  convention permet d'écrire
  \[\forall n\geq m\qsep \prod_{k=m}^n u_k = u_n  \prod_{k=m}^{n-1} u_k.\]
  % Plus généralement, lorsque $n<m$, on utilise souvent la convention
  % $\prod_{k=m}^n u_k\defeq 1$.
\remarque Si $m,n\in\Z$ sont tels que $n\geq m-1$. Alors, quel que soit $a\in\C$
  \[\prod_{k=m}^{n} a =a^{n-m+1}.\]
\remarque Pour tout $n\in\N$
   \[n!=\prod_{k=1}^n k.\]
\end{remarques}

\begin{exoUnique}
% \exo Soit $m,n\in\Z$ tels que $m\leq n$ et $a\in\C$. Calculer
%   \[\prod_{k=m}^n a.\]
\exo Exprimer, à l'aide de factorielles, les produits
  \[\prod_{k=1}^n (2k) \et \prod_{k=0}^n (2k+1).\]
\end{exoUnique}

\begin{proposition}
Soit $m,n\in\Z$ et $(u_k)_{k\in\Z}, (v_k)_{k\in\Z}$ deux suites d'éléments de $\C$. Alors
\[\prod_{k=m}^n u_k v_k=\p{\prod_{k=m}^n u_k} \p{\prod_{k=m}^n v_k}.\]
\end{proposition}

\subsection{Somme et produit doubles}

On parle de somme double lorsqu'il y a deux indices. Pour sommer les éléments
$u_{i,j}$ d'un tableau à $n$ lignes et $m$ colonnes, on peut procéder d'au
moins deux manières~: une sommation en lignes ou en colonnes. Évidemment, le résultat est le même. 

\begin{proposition}
Soit $m_1,n_1,m_2,n_2\in\Z$ et $(u_{i,j})$ une famille d'éléments de $\C$. Alors
\[\sum_{i=m_1}^{n_1} \sum_{j=m_2}^{n_2} u_{i, j} =
\sum_{j=m_2}^{n_2} \sum_{i=m_1}^{n_1} u_{i, j}.\]
\end{proposition}

\begin{remarqueUnique}
\remarque Cette somme est parfois notée
\[\sum_{\substack{m_1 \leq i \leq n_1 \\ m_2 \leq j \leq n_2}} u_{i, j} \ou 
  \sum_{(i,j) \in \intere{m_1}{n_1} \times \intere{m_2}{n_2}} u_{i, j}.\]
\end{remarqueUnique}

\begin{exos}
\exo Calculer
  \[\sum_{\substack{0 \leq i \leq n\\0 \leq j \leq n}} ij \et
    \sum_{\substack{0 \leq i \leq n\\0 \leq j \leq n}} (i+j).\]
  \begin{sol}
\[\sum_{0 \leq i,j \leq n} ij=\frac{1}{4}\cro{n(n+1)}^2\]
\[\sum_{0 \leq i,j \leq n} (i+j)=n(n+1)^2\]
  \end{sol}
\exo Calculer
  \[\sum_{\substack{0 \leq i \leq n\\0 \leq j \leq n}} \binom{i}{j}.\]
\exo Calculer
  \[\sum_{k=1}^n k 2^k\]
  en remarquant astucieusement que $k=\sum_{i=1}^k 1$.
  \begin{sol}
  \[\sum_{k=1}^n k 2^k=2^{n+1}(n-1)+2\]
  \end{sol}
\end{exos}

\begin{proposition}
Soit $m_1,n_1,m_2,n_2\in\Z$ et $(u_k)_{k\in\Z}, (v_k)_{k\in\Z}$ deux suites d'éléments de $\C$. Alors
\[\p{\sum_{i=m_1}^{n_1} u_i}\p{\sum_{j=m_2}^{n_2} v_j}=\sum_{i=m_1}^{n_1} \sum_{j=m_2}^{n_2} u_i v_j.\]
\end{proposition}


\begin{exos}
\exo Calculer
  \[\sum_{\substack{1 \leq i \leq n\\1 \leq j \leq n}}\abs{i-j}.\]
  \begin{sol}
  On peut l'écrire comme $\sum_{i=1}^n\p{\sum_{j=1}^{i-1}(i-j)+\sum_{j=1}^n(j-i)}$ puis réindexation dans les deux sommes, etc...
  
  Autre méthode : On met les valeurs de $|i-j|$ dans un tableau à double entrée et on remarque que chaque sous-diagonale à $(n-i)$ fois la valeur $i$. Donc $S=2\sum_{i=1}^{n-1}i(n-i)$. Dans tous les cas on trouve :
\[\sum_{1\leq i, j \leq n}\abs{i-j}=\frac{1}{3}n(n^2-1)\]
  \end{sol}
\exo Calculer
  \[\sum_{k = 1}^n \sum_{s = 0}^{n-k} \dfrac{k}{s + k}.\]
  \begin{sol}
  
  
\[\sum_{k=1}^n \sum_{s=0}^{n-k} \frac{k}{s+k}=\sum_{i=1}^n\sum_{s+k=i}\frac{k}{s+k}=\sum_{i=1}^n\sum_{k=0}^{i}\frac{k}{i}=\frac{1}{4}n(n+3)\]
Si on fait un tableau, on a sommé sur les antidiagonales.
  \end{sol}
\end{exos}



\subsection{Fonction polynôme}

Dans la suite de ce chapitre, $\K$ désigne $\Q$, $\R$ ou $\C$.

\begin{definition}
On appelle \emph{fonction polynôme} à coefficients dans $\K$ toute fonction
$P:\K\to\K$ telle qu'il existe $a_0,\ldots,a_n\in\K$ tels que
\[\forall z\in\K\qsep P(z)\defeq a_0 + a_1 z + \cdots +  a_{n-1} z^{n-1} + a_n z^n.\]
L'ensemble des fonctions polynôme à coefficients dans $\K$ est noté $\mathcal{P}(\K,\K)$.
\end{definition}

\begin{remarqueUnique}
\remarque On dit qu'une fonction polynôme $P\in\mathcal{P}(\K,\K)$ est de \emph{degré}
  $n\in\N$ lorsqu'il existe $a_0,\ldots,a_n\in\K$ tels que $a_n\neq 0$ et
  \[\forall z\in\K\qsep P(z)\defeq a_0 + a_1 z + \cdots +  a_{n-1} z^{n-1} + a_n z^n.\]
\end{remarqueUnique}


\begin{exoUnique}
\exo
  \begin{questions}
  \question Montrer que pour tout $n\in\intere{0}{3}$, il existe une fonction
    polynôme $P_n$ de degré $n$ telle que
    \[\forall x\in\R\qsep \cos(nx)=P_n(\cos x).\]
  \question Généraliser ce résultat pour $n\in\N$.
  \end{questions}
% Soit $n\in\N$. Montrer que $\cos\p{n\theta}$ peut
%   s'exprimer comme un polynôme en $\cos\theta$.
  \begin{sol}
  On a
    \begin{eqnarray*}
    \cos 5\theta+\ii\sin 5\theta 
    &=& \e^{\ii 5\theta}\\
    &=& \p{\e^{\ii\theta}}^5\\
    &=& \p{\cos\theta+\ii\sin\theta}^5\\
    &=& \cos^5\theta+5\ii\cos^4\theta\sin\theta-10\cos^3\theta\sin^2\theta\\
    & & -10\ii\cos^2\theta\sin^3\theta+5\cos\theta\sin^4\theta+\ii\sin^5\theta
    \end{eqnarray*}
    En identifiant les parties réelles, il vient
    \begin{eqnarray*}
    \cos 5\theta
    &=& \cos^5\theta-10\cos^3\theta\p{1-\cos^2\theta}
        +5\cos\theta\p{1-\cos^2\theta}^2\\
    &=& 16\cos^5\theta-20\cos^3\theta+5\cos\theta.
    \end{eqnarray*}
  Plus généralement, soit $n\in\N$. Alors, pour tout $\theta\in\R$
  \begin{eqnarray*}
  \cos(n\theta)+\ii\sin(n\theta)
  &=& \e^{\ii n\theta} = \p{\e^{\ii\theta}}^n\\
  &=& (\cos\theta+\ii\sin\theta)^n = \sum_{k=0}^n \binom{n}{k}\cos^{n-k}\theta\p{\ii\sin\theta}^k\\
  &=& \sum_{k=0}^n \ii^k\binom{n}{k}\cos^{n-k}\theta\p{\sin\theta}^k
  \end{eqnarray*} 
  Or $\ii^k$ est imaginaire pur lorsque $k$ est impair et réel lorsque $k$ est pair. Donc, on prenant la partie réelle de l'identité précédente, on a déduit que
  \begin{eqnarray*}
  \cos(n\theta)
  &=& \sum_{\substack{k=0\\ k\equiv 0\ [2]}}^n \ii^k\binom{n}{k}\cos^{n-k}\theta\p{\sin\theta}^k\\
  &=& \sum_{k=0}^{\floor{\frac{n}{2}}} \ii^{2k} \binom{n}{2k}\cos^{n-2k}\theta\p{\sin\theta}^{2k}\\
  &=& \sum_{k=0}^{\floor{\frac{n}{2}}} \p{-1}^k \binom{n}{2k}\cos^{n-2k}\theta\p{\sin^2\theta}^{k}\\
  &=& \sum_{k=0}^{\floor{\frac{n}{2}}} \p{-1}^k \binom{n}{2k}\cos^{n-2k}\theta\p{1-\cos^2\theta}^{k}
  \end{eqnarray*} 
  Donc $\cos(n\theta)$ s'exprime comme un polynôme en $\cos\theta$. Ces polynômes sont appelés polynômes de Tchebycheff.
  \end{sol}
\end{exoUnique}

\begin{definition}
  Si $P\in\mathcal{P}(\K,\K)$, on appelle \emph{racine} de $P$ tout élément $\alpha\in\K$ tel que $P(\alpha)=0$.
  \end{definition}
  
  \begin{remarques}
  \remarque Le calcul
  des racines des fonctions polynôme de degré 2 se fait en utilisant le discriminant.
  \remarque Il n'y a pas de méthode
  systématique pour trouver les racines des fonctions polynôme de degré supérieur. En effet, on peut montrer qu'il
  n'existe pas de formule générale permettant de calculer les racines des fonctions polynôme de degré 3 ou plus avec
  des radicaux réels. Et même si on s'autorise les racines $n$-ièmes de nombres complexes, il n'existe pas
  de formule générale permettant de déterminer les racines de fonctions polynôme de degré 5 ou plus. Cependant, il
  existe différentes techniques qui sont efficaces pour certaines fonctions polynôme.
  \end{remarques}
  
  \begin{proposition}
  Soit $P\in\mathcal{P}(\K,\K)$ et $\alpha\in\K$. Alors $\alpha$ est une racine de $P$ si et seulement si il existe
  $Q\in\mathcal{P}(\K,\K)$ tel que
  \[\forall z\in\K\qsep P(z)=(z-\alpha)Q(z).\]
  \end{proposition}
  
  \begin{remarques}
  \remarque La factorisation
  effective se fait par division euclidienne. Par exemple, si $P(z)\defeq z^3+3z^2+3z+2$, on
  remarque que $P(-2)=0$ donc $P(z)$ se factorise par $z+2$ et la division euclidienne s'effectue de la manière
  suivante.
  \begin{equation*}
  \renewcommand{\arraystretch}{1.2}
  \renewcommand{\arraycolsep}{2pt}
  \begin{array}{rrrr|l}
  z^3 & +\,3z^2 & +\,3z & +\,2 & z + 2 \\
  \cline{5-5}
  z^3 & +\,2z^2 &   &  & z^2+z+1 \\
  \cline{1-4}
      & z^2 &  +\,3z & +\,2 &     \\
      & z^2 & +\,2z &  &   \\
      \cline{2-4}
      &     & z & +\,2 &   \\
      &     & z & +\,2 &    \\
                \cline{3-4}
      &     &   &0&   
  \end{array}
  \end{equation*}
  
  \noindent donc $z^3+3z^2+3z+2=(z+2)(z^2+z+1)$. Puisque les racines de $Q(z)\defeq z^2+z+1$ sont $\jj$ et
  $\jj^2$, on en déduit que les racines de $P$ sont $-2,\jj$ et $\jj^2$.
  \remarque
  Si $P$ est une fonction polynôme à coefficients entiers, il existe une technique efficace pour déterminer rapidement
  ses racines rationnelles. Soit $a_0,\ldots,a_n\in\Z$ tels que $P(z)\defeq a_n z^n+\cdots+a_1 z+a_0$.
  Si $r$ est une racine rationnelle de $P$, il existe $p\in\Z$ et $q\in\Ns$ premiers entre eux tels que
  $r=p/q$. Puisque $P(r)=0$, on en déduit que
  \[a_n\p{\frac{p}{q}}^n+\cdots+a_1\cdot\frac{p}{q}+a_0=0.\]
  En multipliant par $q^n$, on obtient
  \[a_n p^n + a_{n-1}p^{n-1}q+\cdots+a_1 p q^{n-1}+a_0 q^n=0.\]
  On en déduit que
  \[a_n p^n = -q\p{a_{n-1}p^{n-1}+\cdots+a_1 p q^{n-2}+a_0 q^{n-1}}\]
  et donc que $q$ divise $a_n p^n$. Or $q$ et $p$ sont premiers entre eux donc $q$ et $p^n$ sont premiers entre eux. D'après le lemme de \nom{Gauss}, on en déduit que $q$ divise $a_n$. De même, on montre que $p$ divise $a_0$. Comme il existe un nombre fini de diviseurs d'un entier non nul, les racines rationnelles sont donc à chercher parmi un nombre fini d'éléments. Par exemple, si $P(z)\defeq 3z^3+5z^2+5z+2$, et si $p/q$ est une racine rationnelle mise sous forme irréductible de $P$, alors $p$ divise 2 et $q$ divise 3. Donc $p\in\ens{-2,-1,1,2}$ et $q\in\ens{1, 3}$. Les racines rationnelles éventuelles de $P$ sont donc parmi $\ens{-2,-1,1,2,-2/3,-1/3,1/3,2/3}$. Si on teste tous ces rationnels, on se rend compte que $-2/3$ est une racine de $P$. $P(z)$ se factorise donc par $3z+2$ et une division euclidienne nous donne $P(z)=(3z+2)(z^2+z+1)$. Les racines de $P$ sont donc $-2/3,\jj$ et $\jj^2$.
  \remarque D'autres techniques permettent de trouver les racines d'une fonction polynôme de degré $n\geq 3$. Par exemple, pour certaines fonctions polynôme, ramener la recherche de leurs racines à la recherche des racines $n$-ièmes d'un nombre complexe.
  % \remarque Enfin, si $P$ possède certaines symétries, un changement de variable peut permettre de faire baisser le degré. Par exemple, si $P= a_n X^n+\cdots+a_1 X+a_0$ et que pour tout $k\in\intere{0}{n}$, $a_{n-k}=a_k$ (polynôme symétrique de première espèce), le changement de variable $u=x+1/x$ permet de diviser par 2 le degré du polynôme.
  \end{remarques}
  
  
  % \begin{exoUnique}
  % \exo Résoudre l'inéquation
  %   \[4x+2\leq\sqrt{7x^3+15x^2+11x+3}.\]
  % \end{exoUnique}
  
  % Dans la suite, nous aurons besoin de la version générale de la \emph{division euclidienne}. Elle permet d'effectuer la division de tout polynôme par un polynôme non nul.
  
  % \begin{proposition}
  % Soit $A,B\in\polyK$ avec $B\neq 0$. Alors, il existe un unique couple
  % $\p{Q,R}\in\polyK^2$ tel que
  % \[A=QB+R \et \deg R<\deg B\]
  % $Q$ est appelé \emph{quotient} de la \emph{division euclidienne} de $A$ par $B$, $R$ son \emph{reste}.
  % \end{proposition} 
  
  % \noindent Effectuons par exemple la division euclidienne de $7X^5+4X^4+2X^3-X+5$ par $X^2+2$. Il est important de bien aligner les puissances les unes au dessus des autres et donc de laisser de la place pour les puissances lacunaires (ici, le $X^2$).
  % \begin{equation*}
  % \renewcommand{\arraystretch}{1.2}
  % \renewcommand{\arraycolsep}{2pt}
  % \begin{array}{rrrrrr|l}
  % 7X^5 & +\,4X^4 & +\,2X^3 &  & -\,X & +\, 5& X^2 + 2 \\
  % \cline{7-7}
  % 7X^5 &  & +\,14X^3 &   &  & & 7X^3+4X^2-12X-8 \\
  % \cline{1-6}
  %   & 4X^4 & -\,12X^3 &  & -\,X & +\,5 & \\
  %   & 4X^4 &          & +\,8X^2 & &  &   \\
  % \cline{2-6}
  % & & -\,12X^3 & -\,8X^2 & -\,X & +\, 5 & \\
  % & & -\,12X^3 &         & -24\, X & & \\
  % \cline{3-6}
  % & & & -\,8X^2 & +\,23X & +\,5 & \\
  % & & & -\,8X^2 &        & -\,16 & \\
  % \cline{4-6}
  % & & & & 23X & +\, 21 & 
  % \end{array}
  % \end{equation*}
  
  % \noindent On a donc $7X^5+4X^4+2X^3-X+5=(7X^3+4X^2-12X-8)(X^2+2)+23X+21$. Autrement dit, le quotient de la division euclidienne de $7X^5+4X^4+2X^3-X+5$ par $X^2+2$ est $7X^3+4X^2-12X-8$ et son reste est $23X+21$.\\
  
  % \begin{remarqueUnique}
  % \remarque Si $x$ est une racine de $B$ et $A=QB+R$ est la division euclidienne de $A$ par $B$, alors $A(x)=R(x)$. Cette remarque permet de simplifier le calcul de $A(x)$ lorsqu'on connait un polynôme annulateur de $x$.
  % \end{remarqueUnique}
  
  \begin{proposition}
  Une fonction polynôme $P\in\mathcal{P}(\K,\K)$ de degré $n\in\N$ admet au plus $n$ racines.
  \end{proposition}
  
  \begin{remarques}
  \remarque Soit $P\in\mathcal{P}(\K,\K)$ une fonction polynôme et
   $a_0,\ldots,a_n\in\K$ tels que
    \[\forall z\in\K\qsep P(z)=a_0+a_1 z+\cdots+a_{n-1}z^{n-1}+a_n z^n.\]
    Si $P$ admet au moins $n+1$ racines, alors $a_0=a_1=\cdots=a_n=0$.
  \remarque Soit $a_0,\ldots,a_n,b_0,\ldots,b_n\in\K$ et $R$ une partie de
   $\K$ telle que
    \[\forall z\in R\qsep 
    a_n z^n+a_{n-1}z^{n-1}+\cdots+a_1 z+a_0=b_n z^n+b_{n-1}z^{n-1}+\cdots+b_1 z+b_0.\]
    Si $R$ possède au moins $n+1$ éléments (en particulier si $R$ est infini), alors $a_n=b_n, \ldots, a_0=b_0$.
\remarque Si $P$ est une fonction polynôme non nulle, il
  existe un unique $n\in\N$ et une unique famille $(a_0,\ldots,a_n)\in\K^{n+1}$
  tels que $a_n\neq 0$ et
  \[\forall z\in\K\qsep P(z)=a_0 + a_1 z + \cdots +  a_{n-1} z^{n-1} + a_n z^n.\]
  On dit que $n$ est le \emph{degré} de $P$ et que $a_0,\ldots,a_n$ sont ses
  \emph{coefficients}.
  \end{remarques}

\section{Trigonométrie}

\subsection{Égalité modulaire}

\begin{definition}
Soit $m\in\RPs$ et $a,b\in\R$. On dit que $a$ est \emph{congru à $b$ modulo $m$} et on note
\[a\equiv b\ [m]\]
lorsqu'il existe $k\in\Z$ tel que $a=b+km$.
\end{definition}

\begin{exoUnique}
\exo Soit $a,b\in\R$. Quel est le lien logique entre
  \[\text{\og}a\equiv b\ [2\pi]\text{\fg} \et \text{\og}a\equiv b\ [\pi]\text{\fg} \text{ ?}\]
\end{exoUnique}

\begin{proposition}
\begin{itemize}
\item 
Soit $m\in\RPs$ et $a_1,a_2,b_1,b_2\in\R$ tels que
  \[a_1\equiv b_1\ [m] \et a_2\equiv b_2\ [m].\]
Alors, quels que soient $k_1,k_2\in\Z$,
  \[k_1 a_1+k_2 a_2\equiv k_1 b_1+k_2 b_2\ [m].\]
\item Soit $m\in\RPs$ et $a,b\in\R$ tels que
  \[a\equiv b\ [m].\]
  Alors, si $c\in\RPs$
  \[ac\equiv bc\ [m c].\]
\end{itemize}
\end{proposition}

\begin{remarques}
\remarque On en déduit qu'on peut raisonner avec les \og $\equiv$ \fg de la même manière qu'avec \og $=$ \fg pour résoudre les équations.
  \begin{itemize}
  \item On peut passer une expression d'un côté à l'autre du \og $\equiv$ \fg en changeant son signe.
  \item On peut multiplier les deux côtés du signe \og $\equiv$ \fg par un même coefficient $c\in\RPs$. Il suffit juste de multiplier le modulo par $c$.
  \end{itemize}
  Ces deux transformations permettent de raisonner par équivalence.
\remarque Si $a\in\R$ et $m\in\RPs$, alors l'ensemble des $x\in\R$ tels que $x\equiv a\ [m]$ est noté
  \[a+m\Z\defeq\ensim{a+km}{k\in\Z}.\]
\end{remarques}

\subsection{Formules de trigonométrie}

\begin{center}
\begin{pdfpic}\pspicture*(-1.2,-1.2)(9.2,5.2)
\psset{unit=4cm}
\psline{->}(0,0)(2,0)
\psline{->}(0,0)(0,1.3)
\psline{-}(0,0)(2,1.1547)
\psline[linestyle=dashed](0,0.5)(0.8660,0.5)
\psline[linestyle=dashed](0.8660,0.5)(0.8660,0)
\psline[linestyle=dashed](0,1)(2,1)
\psline[linestyle=dashed](1,0)(1,1)
\psline{->}(0,-0.05)(0.8660,-0.05)
\psline{->}(-0.05,0)(-0.05,0.5)
\psline{->}(0,1.05)(1.7320,1.05)
\psline{->}(1.05,0)(1.05,0.5773)
\psarc{->}(0,0){1.05}{0}{30}
\psarc(0,0){1}{-17.45}{107.45}
\uput[d](0.4330,-0.05){$\cos x$}
\uput[l](-0.05,0.25){$\sin x$}
\uput[r](1.05,0.2886){$\tan x$}
\uput[u](0.8660,1.05){${\rm cotan}\, x$}
\uput[u](0.9093,0.525){$x$}
\endpspicture
\end{pdfpic}
\end{center}

\begin{definition}[utile=-3]
On définit le \emph{sinus}, le \emph{cosinus}, la \emph{tangente} et la \emph{cotangente} d'un angle $x$ exprimé en radians sur le cercle trigonométrique de rayon 1 comme ci-dessus.
En particulier $\tan x$ n'est défini que pour
$x\in\R\setminus\p{\frac{\pi}{2}+\pi\Z}$, $\cotan x$ n'est défini que pour
$x\in\R\setminus \pi\Z$ et
\[\tan x=\frac{\sin x}{\cos x} \et \cotan x=\frac{\cos x}{\sin x}.\]
\end{definition}
\begin{preuve}
Le calcul de $\tan$ et $\cotan$ en fonction de $\sin$ et $\cos$ se fait par
le théorème de Thalès.
\end{preuve}

\begin{remarqueUnique}
\remarque On rappelle les principales valeurs remarquables.
\begin{center}
\begin{pdfpic}
\begin{pspicture}(-1.2,-1.2)(5.2,5.2)
\psset{unit=4.5cm}
\psline{->}(0,0)(1.3,0)
\psline{->}(0,0)(0,1.3)
\psline[linestyle=dashed](0,0.5)(0.8660,0.5)
\psline[linestyle=dashed](0.8660,0.5)(0.8660,0)
\psline[linestyle=dashed](0.5,0.8660)(0.5,0)
\psline[linestyle=dashed](0.5,0.8660)(0,0.8660)
\psline[linestyle=dashed](0.7071,0.7071)(0.7071,0)
\psline[linestyle=dashed](0.7071,0.7071)(0,0.7071)
\psline(0,0)(0.7071,0.7071)
\psline(0,0)(0.5,0.8660)
\psline(0,0)(0.8660,0.5)
\psarc(0,0){1}{-17.45}{107.45}
\psarc{->}(0,0){1.05}{0}{30}
\psarc{->}(0,0){1.05}{0}{45}
\psarc{->}(0,0){1.05}{0}{60}
\uput[r](0.9093,0.525){$\dsp\frac{\pi}{6}$}
\uput[ur](0.7424,0.7424){$\dsp\frac{\pi}{4}$}
\uput[u](0.525,0.9093){$\dsp\frac{\pi}{3}$}
\uput[d](0.5,0){$\frac{1}{2}$}
\uput[d](0.7071,0){$\frac{\sqrt{2}}{2}$}
\uput[d](0.8660,0){$\frac{\sqrt{3}}{2}$}
\uput[l](0,0.5){$\frac{1}{2}$}
\uput[l](0,0.7071){$\frac{\sqrt{2}}{2}$}
\uput[l](0,0.8660){$\frac{\sqrt{3}}{2}$}
\end{pspicture}
\end{pdfpic}
\begin{tabular}{|>{\hfill}p{1.5cm}<{\hfill\rule[-10pt]{0pt}{27pt}}||*{5}{>{\hfill}p{1.5cm}<{\hfill\rule[-12pt]{0pt}{30pt}}|}}
\hline
$x$&$0$&$\dsp\frac{\pi}{6}$&$\dsp\frac{\pi}{4}$&$\dsp\frac{\pi}{3}$&
  $\dsp\frac{\pi}{2}$\\
\hline\hline\rule[-10pt]{0pt}{27pt}
$\sin x$&$0$&$\dsp\frac{1}{2}$&$\dsp\frac{\sqrt{2}}{2}$&$\dsp\frac{\sqrt{3}}{2}$&
  $1$\\
\hline\rule[-10pt]{0pt}{27pt}
$\cos x$&$1$&$\dsp\frac{\sqrt{3}}{2}$&$\dsp\frac{\sqrt{2}}{2}$&$\dsp\frac{1}{2}$&
  $0$\\
\hline\rule[-10pt]{0pt}{27pt}
$\tan x$&$0$&$\dsp\frac{1}{\sqrt{3}}$&$1$&$\sqrt{3}$&indéfini\\
\hline\rule[-10pt]{0pt}{27pt}
$\cotan x$&indéfini&$\sqrt{3}$&$1$&$\dsp\frac{1}{\sqrt{3}}$&$0$\\
\hline
\end{tabular}
\end{center}
\remarque Si $x\in\R\setminus\frac{\pi}{2}\Z$, alors
  \[\cotan x=\frac{1}{\tan x}.\]
  % Remarquons cependant que $\cotan$ est définie en $\pi/2$ alors que $\tan$ ne
  % l'est pas.
\end{remarqueUnique}


\begin{proposition}[utile=-3]
D'après Pythagore, on a
\[\cos^2 x+\sin^2 x=1, \qquad 1+\tan^2 x=\frac{1}{\cos^2 x},\]
\[1+\cotan^2 x=\frac{1}{\sin^2 x}.\]
\end{proposition}

\begin{proposition}[utile=-3, nom={Symétries}]
\begin{align*}
\cos\p{-x}&=\cos x   & \cos\p{\pi+x}&=-\cos x & \cos\p{\pi-x}&=-\cos x\\
\sin\p{-x}&=-\sin x  & \sin\p{\pi+x}&=-\sin x & \sin\p{\pi-x}&=\sin x\\
\tan\p{-x}&=-\tan x  & \tan\p{\pi+x}&=\tan x  & \tan\p{\pi-x}&=-\tan x
\end{align*}
\begin{align*}
\cos\p{\frac{\pi}{2}+x}&=-\sin x &\cos\p{\frac{\pi}{2}-x}&=\sin x\\
\sin\p{\frac{\pi}{2}+x}&=\cos x  &\sin\p{\frac{\pi}{2}-x}&=\cos x\\
\tan\p{\frac{\pi}{2}+x}&=-\cotan x &\tan\p{\frac{\pi}{2}-x}&=\cotan x
\end{align*}
\end{proposition}

\begin{center}
\begin{pdfpic}
\begin{pspicture}(-3.25,-3.25)(3.25,3.25)
\psset{unit=2.5cm}
\pscircle(0,0){1}
\psline{->}(-1.3,0)(1.3,0)
\psline{->}(0,-1.3)(0,1.3)
\psarc{->}(0,0){1.07}{0}{30}
\psarc{->}(0,0){1.07}{0}{150}
\psarc{->}(0,0){1.07}{0}{210}
\psarc{<-}(0,0){1.07}{-30}{0}
\psarc{->}(0,0){1.07}{0}{60}
\psarc{->}(0,0){1.07}{0}{120}
\psline[linestyle=dashed](0.8660,0.5)(-0.8660,0.5)
\psline[linestyle=dashed](0.8660,0.5)(0.8660,-0.5)
\psline[linestyle=dashed](0.5,0.8660)(0.5,0)
\psline[linestyle=dashed](0.5,0.8660)(-0.5,0.8660)
\psline[linestyle=dashed](-0.5,0.8660)(-0.5,0)
\psline[linestyle=dashed](-0.8660,0.5)(-0.8660,-0.5)
\psline[linestyle=dashed](-0.8660,-0.5)(0.8660,-0.5)
\psline(-0.8660,0.5)(0.8660,-0.5)
\psline(0.5,0.8660)(0,0)
\psline(-0.5,0.8660)(0,0)
\psline(-0.8660,-0.5)(0.8660,0.5)
\uput[dl](-0.9266,-0.535){$\pi+x$}
\uput[dr](0.9266,-0.535){$-x$}
\uput[ul](-0.9266,0.535){$\pi-x$}
\uput[ur](0.535,0.9266){$\dsp\frac{\pi}{2}-x$}
\uput[ul](-0.535,0.9266){$\dsp\frac{\pi}{2}+x$}
\uput[ur](0.9266,0.535){$x$}
\end{pspicture}
\end{pdfpic}
\end{center}

\begin{remarques}
\remarque Il est important de retrouver rapidement ces formules en dessinant le cercle trigonométrique et un \og petit \fg angle $x$ vérifiant $0<x<\pi/4$.
\remarque Pour tout $n\in\Z$ et $x\in\R$
  \[\cos\p{x+n\pi}=(-1)^n \cos x \qquad\et\qquad \sin\p{x+n\pi}=(-1)^n \sin x.\]
  En particulier $\cos(n\pi)=(-1)^n$ et $\sin(n\pi)=0$. 
\remarque Quels que soient $x,y\in\R$
  \begin{eqnarray*}
  \cos x = \cos y &\ssi& \cro{x\equiv y \ [2\pi] \qquad\text{ou}\qquad
    x\equiv -y \ [2\pi]}\\
  \sin x = \sin y &\ssi& \cro{x\equiv y \ [2\pi] \qquad\text{ou}\qquad
    x\equiv \pi-y \ [2\pi]}
  \end{eqnarray*}
\remarque Quels que soient $x,y\in\R\setminus\p{\frac{\pi}{2}+\pi\Z}$
  \[\tan x=\tan y \quad\ssi\quad x\equiv y\ [\pi].\]
\end{remarques}

\begin{exos}
\exo Calculer
  \[\cos\p{\frac{7\pi}{6}}, \qquad \sin\p{\frac{2\pi}{3}}, \qquad
    \tan\p{-\frac{3\pi}{4}}.\]
  \begin{sol}
  On a
\[\cos\left( \frac{7\pi}{6} \right)=\cos\left( \pi+\frac{\pi}{6} \right)=-\cos\left( \frac{\pi}{6} \right)=-\frac{\sqrt{3}}{2}\]

\[\sin\left( \frac{2\pi}{3} \right)=\sin\left( \pi-\frac{\pi}{3} \right)=\sin\left( \frac{\pi}{3} \right)=\frac{\sqrt{3}}{2}\]
\[\tan\left( -\frac{3\pi}{4}\right) =\tan\left(\frac{\pi}{4}-\pi\right) =\tan\left(\frac{\pi}{4}\right) =1.\]
  \end{sol}
\exo Résoudre l'équation $\sin x=\cos x$.
\end{exos}

\begin{proposition}[nom={Addition des arcs}]
\begin{align*}
\cos(a+b)&=\cos a\cos b-\sin a\sin b &\sin(a+b)&=\sin a\cos b+\cos a\sin b\\
\cos(a-b)&=\cos a\cos b+\sin a\sin b &\sin(a-b)&=\sin a\cos b-\cos a\sin b
\end{align*}
\begin{align*}
\tan(a+b)&=\frac{\tan a+\tan b}{1-\tan a\tan b}\\
\tan(a-b)&=\frac{\tan a-\tan b}{1+\tan a\tan b}
\end{align*}
\end{proposition}
\begin{preuve}
On montre la formule de $\cos\p{a-b}$ avec le produit scalaire. Les autres en
découlent. Pour $\cos\p{a+b}$, il faut changer $b$ en $-b$. Pour
$\sin\p{a+b}$, écrire $\cos\p{\pi/2-\p{a+b}}=\cos\p{\p{\pi/2-a}-b}$.
Pour les formules avec les tangentes, il faut $a+b\neq \pi/2\ [\pi]$ et
de même pour $a$ et $b$. 
\end{preuve}

\begin{remarqueUnique}
\remarque Si $a,b\in\R$ ne sont pas tous les deux nuls, on pourra factoriser
  $a\cos x+b\sin x$ de la manière suivante.
  \[a\cos x+b\sin x=\sqrt{a^2+b^2}\p{\frac{a}{\sqrt{a^2+b^2}}\cos x+
                      \frac{b}{\sqrt{a^2+b^2}}\sin x}\]
  Puisque
  \[\p{\frac{a}{\sqrt{a^2+b^2}}}^2+\p{\frac{b}{\sqrt{a^2+b^2}}}^2=1,\]
  il existe $\theta_0\in\R$ tel que $\cos\theta_0=a/\sqrt{a^2+b^2}$ et
  $\sin\theta_0=b/\sqrt{a^2+b^2}$. On a alors
  \begin{eqnarray*}
  a\cos x+b\sin x &=& \sqrt{a^2+b^2}\p{\cos\theta_0\cos x+\sin\theta_0\sin x}\\
                  &=& \sqrt{a^2+b^2}\cos\p{x-\theta_0}.
  \end{eqnarray*}
\end{remarqueUnique}

\begin{exoUnique}
\exo Résoudre l'équation $\sqrt{3}\cos x+\sin x=1$.
\begin{sol}
Pour tout $x\in\R$
\begin{eqnarray*}
\sqrt{3}\cos x+\sin x
&=& 2\left(\frac{\sqrt{3}}{2}\cos x+\frac{1}{2}\sin x \right)\\
&=& 2\left( \cos \frac{\pi}{6}\cos x+\sin \frac{\pi}{6}\sin x \right)\\
&=& 2\cos\left( x-\frac{\pi}{6}\right).
\end{eqnarray*}
\end{sol}
\end{exoUnique}

\begin{proposition}[utile=-3, nom={Angle double}]
\begin{align*}
\cos(2x)&=\cos^2 x-\sin^2 x\\
        &=2\cos^2 x-1\\
        &=1-2\sin^2 x\\
\sin(2x)&=2\cos x\sin x\\
\tan(2x)&=\frac{2\tan x}{1-\tan^2 x}
\end{align*}
\begin{align*}
\cos^2 x &=\frac{1+\cos(2x)}{2}\\
\sin^2 x &=\frac{1-\cos(2x)}{2}
\end{align*}
\end{proposition}

\begin{exos}
\exo Exprimer $\cos(\pi/8)$ à l'aide de radicaux.
\exo Soit $a\in\R$. Pour tout $n\in\N$, on pose
  \[p_n\defeq\prod_{k=1}^n \cos\p{\frac{a}{2^k}}.\]
  Simplifier $p_n \sin\p{a/2^n}$ puis en déduire la limite de la suite $\p{p_n}$.
  \begin{sol}
Montrons par récurrence sur $n$ que
  \[\forall n\in\Ns \qsep p_n\sin\frac{a}{2^n}=\frac{1}{2^n}\sin a.\]
  En effet, en posant
  \begin{center}
  $\mathcal{H}_n \defeq$ \flqq\ $p_n\sin\frac{a}{2^n}=\frac{1}{2^n}\sin a$ \frqq
  \end{center}
  \begin{itemize}
  \item $\mathcal{H}_1$ est vraie. En effet~:
    \[p_1\sin\frac{a}{2}=\cos\frac{a}{2}\sin\frac{a}{2}=\frac{1}{2}\sin a.\]
  \item $\mathcal{H}_n \implique \mathcal{H}_{n+1}$. En effet, soit $n\geq 1$.
    On suppose que $\mathcal{H}_n$ est vraie. Montrons que $\mathcal{H}_{n+1}$
    est vraie.
    \begin{eqnarray*}
    p_{n+1}\sin\frac{a}{2^{n+1}}
    &=& \sin\frac{a}{2^{n+1}} \prod_{p=1}^{n+1} \cos\frac{a}{2^p}\\
    &=& \sin\frac{a}{2^{n+1}} \cos\frac{a}{2^{n+1}}
        \prod_{p=1}^n \cos\frac{a}{2^p}\\
    &=& \frac{1}{2} \sin\frac{a}{2^n} \prod_{p=1}^n \cos\frac{a}{2^p}\\
    &=& \frac{1}{2} p_n \sin\frac{a}{2^n}\\
    &=& \frac{1}{2}\cdot \frac{1}{2^n} \sin a
        \quad\text{par hypothèse de récurrence}\\
    &=& \frac{1}{2^{n+1}} \sin a
    \end{eqnarray*}
    Donc $\mathcal{H}_{n+1}$ est vraie.
  \end{itemize}
  Par récurrence, on en déduit que $\mathcal{H}_{n}$ est vraie pour tout
  $n\in\Ns$, donc
  \[\forall n\in\Ns \quad p_n\sin\frac{a}{2^n}=\frac{1}{2^n}\sin a.\]
  Puisque $a$ est non nul, et que
  \[\frac{a}{2^n}\tendvers{n}{+\infty}0,\]
  on en déduit qu'il existe $N\in\Ns$ tel que, pour tout $n\geq N$, $a/2^n\in\intero{-\pi}{\pi}\setminus\ens{0}$ . Soit $n\geq N$. Alors $\sin\p{a/2^n}\neq 0$, donc
  \[p_n=\frac{\sin a}{2^n\sin\frac{a}{2^n}}.\]
  Or, puisque $a/2^n \tendvers{n}{+\infty} 0$ et $\sin(x)/x \tendvers{x}{0} 1$,
  on a
  \[2^n\sin\frac{a}{2^n}=a\frac{\sin\frac{a}{2^n}}{\frac{a}{2^n}}
    \tendvers{n}{+\infty} a.\]
  En conclusion
  \[p_n\tendvers{n}{+\infty} \frac{\sin a}{a}.\]
  \end{sol}
\end{exos}

\begin{proposition}[utile=-3, nom={Linéarisation}]
\begin{align*}
\cos a\cos b&=\frac{1}{2}\cro{\cos\p{a+b}+\cos\p{a-b}}\\
\sin a\sin b&=\frac{1}{2}\cro{\cos\p{a-b}-\cos\p{a+b}}\\
\cos a\sin b&=\frac{1}{2}\cro{\sin\p{a+b}-\sin\p{a-b}}
\end{align*}
\end{proposition}

\begin{exoUnique}
\exo Linéariser $\cos^3 x$, $\cos x\sin^2 x$, puis $\sin^4 x$.
  \begin{sol}
  Pour tout $x\in\R$
  \begin{eqnarray*}
  \cos^{3}x &=& \cos^{2}x \cos x = \frac{1}{2}(\cos (2x) +1)\cos x\\
            &=& \frac{1}{2}(\cos (2x)\cos x +\cos x)=\frac{1}{4}(\cos (3x)+\cos x +2\cos x)\\
            &=& \frac{1}{4}\p{\cos\p{3x}+3\cos x}\\
  \cos x\sin^2 x &=& \frac{1}{2}(\cos x(1-\cos (2x))=\frac{1}{4}(2\cos x -( \cos (3x)+\cos x))\\
                 &=& \frac{1}{4}(\cos x - \cos (3x))\\
  \sin^4 x &=& (1-\cos^{2}x)^{2}=1-2\cos^{2}x+\cos^{4}x\\
           &=& 1-2\times\frac{1}{2}(\cos (2x)-1)+\frac{1}{4}(\cos (2x)-1)^{2}\\
           &=& \frac{1}{8}\p{\cos\p{4x}-4\cos\p{2x}+3}
  \end{eqnarray*}
  \end{sol}
\end{exoUnique}

\begin{proposition}[utile=-3, nom={Factorisation}]
\begin{align*}
\cos p+\cos q&=2\cos\frac{p+q}{2}\cos\frac{p-q}{2} &
\sin p+\sin q&=2\sin\frac{p+q}{2}\cos\frac{p-q}{2}\\
\cos p-\cos q&=-2\sin\frac{p+q}{2}\sin\frac{p-q}{2} &
\sin p-\sin q&=2\cos\frac{p+q}{2}\sin\frac{p-q}{2}
\end{align*}
\begin{align*}
\tan p+\tan q&=\frac{\sin(p+q)}{\cos p\cos q}\\
\tan p-\tan q&=\frac{\sin(p-q)}{\cos p\cos q}
\end{align*}
\end{proposition}
\begin{preuve}
Poser $p=a+b$ et $q=a-b$.
\end{preuve}

% \begin{remarqueUnique}
% \remarque On pourra retenir que si $f$ est la fonction $\sin$ ou $\cos$, on a
%   \[f(p)-f(q)=2f'\p{\frac{p+q}{2}}\sin\p{\frac{p-q}{2}} \et
%     f(p)+f(q)=2f\p{\frac{p+q}{2}}\cos\p{\frac{p-q}{2}}.\]
% \end{remarqueUnique}

\begin{exoUnique}
\exo En multipliant par $\sin(x/2)$, calculer
  \[A\defeq\sum_{k=0}^n \cos\p{kx} \et B\defeq\sum_{k=0}^n \sin^2\p{kx}.\]
  \begin{sol}
  Soit $n\in\N$ et $x\in\R$. On multiplie $A$ par $\sin\p{x/2}$.
  \begin{eqnarray*}
  \sin\left( \frac{x}{2}\right)\sum_{k=0}^n \cos(kx)
  &=& \sum_{k=0}^n \sin\left( \frac{x}{2}\right)\cos (kx)\\
  &=& \frac{1}{2}\sum_{k=0}^n \left(\sin \left( \frac{x}{2}+kx\right) +\sin \left( \frac{x}{2}-kx\right) \right)\\
  &=& \frac{1}{2}\left( \sum_{k=0}^n \sin \left(kx+ \frac{x}{2}\right) -\sum_{k=0}^n \sin \left(kx- \frac{x}{2}\right)    \right)\\
  &=& \frac{1}{2}\left( \sum_{k=1}^{n+1} \sin \left(kx- \frac{x}{2}\right) -\sum_{k=0}^n \sin \left(kx- \frac{x}{2}\right)    \right)\\
  &=& \frac{1}{2}\left(\sin \left( \frac{(n+1)x}{2}\right) +\sin\left(  \frac{x}{2} \right) \right)\\
  &=& \sin \left(\frac{(n+1)x}{2}\right) \cos\left( \frac{nx}{2}\right)
  \end{eqnarray*}
  Or $\sin(x/2)$ est nul si et seulement si $x\equiv 0\ [2\pi]$. Dans ce cas, $\cos(kx)=1$, donc 
  $A=n+1$. En conclusion
  \[\sum_{k=0}^n \cos\p{kx}=
    \begin{cases}
    \frac{\sin\p{\frac{(n+1)x}{2}}\cos\p{\frac{nx}{2}}}{\sin\p{\frac{x}{2}}}
    & \text{si $x\not\equiv 0\ [2\pi]$}\\
    n+1 & \text{sinon.}
    \end{cases}\]
  Pour le calcul de $B$, on a
  \[\sum_{k=0}^n \sin^2\p{kx}=\frac{1}{2 } \sum_{k=0}^{n}\left( 1-\cos (2kx)\right)\]
  donc en utilisant le résultat précédent
  \[\sum_{k=0}^n \sin^2\p{kx}=
    \begin{cases}
    \frac{1}{2}\cro{\p{n+1}-\frac{\sin((n+1)x)\cos(nx)}{\sin x}} &
      \text{si $x\not\equiv 0\ [\pi]$}\\
    0 & \text{sinon.}
    \end{cases}\]
  \end{sol}
\end{exoUnique}

\begin{proposition}[utile=-3]
Soit $x\in\R$ tel que $x\not\equiv \pi\ [2\pi]$. Alors, en posant $t\defeq\tan\p{x/2}$
\[\cos x=\frac{1-t^2}{1+t^2} \et \sin x=\frac{2t}{1+t^2}.\]
Si de plus, $x\not\equiv\frac{\pi}{2}\ [\pi]$, alors
\[\tan x=\frac{2t}{1-t^2}.\]
\end{proposition}

% \begin{remarqueUnique}
% \remarque Remarquons au passage que, puisque $\cos^2 x+\sin^2 x=1$, on retrouve facilement
%   \[\forall t\in\R \qsep \p{\frac{1-t^2}{1+t^2}}^2+\p{\frac{2t}{1+t^2}}^2=1.\]
%   Autrement dit
%   \[\forall t\in\R \qsep \p{t^2-1}^2+\p{2t}^2=\p{1+t^2}^2.\]
%   Cette relation nous donne, pour $t\in\N$, des triplets $\p{a,b,c}\in\N^3$ non triviaux tels que $a^2+b^2=c^2$. Par exemple, pour $t=2$, on obtient $3^2+4^2=5^2$.
%   \begin{sol}
%   En particulier, pour $t=2$ et $t=4$.
%   \end{sol}
% \end{remarqueUnique}


\section{Récurrence linéaire}

 \subsection{Récurrence linéaire d'ordre 1}


\begin{definition}
Si $(a_n)$ et $(b_n)$ sont deux suites à valeurs dans $\K$, l'équation
\[\forall n\in\N\qsep u_{n+1} + a_n u_n = b_n\]
dont l'inconnue est la suite $(u_n)$ est appelée \emph{récurrence linéaire d'ordre 1}.
\end{definition}

\begin{proposition}
Soit $a\in\K$. Alors, les solutions de la récurrence linéaire
\[\forall n\in\N\qsep u_{n+1}=a u_n\]
sont les suites $(u_n)$ définies par
\[\forall n\in\N\qsep u_n\defeq \lambda a^n\]
où $\lambda\in\K$.
\end{proposition}

\begin{proposition}
Soit $(a_n)$ et $(b_n)$ deux suites à valeurs dans $\K$. Si $(v_n)$ est une solution
\og particulière \fg de la récurrence linéaire
\[\forall n\in\N\qsep u_{n+1}+a_n u_n = b_n\]
alors, les solutions de cette récurrence linéaire sont les suites $(v_n+u_n)$ où
$(u_n)$ parcourt l'ensemble des solutions de la récurrence linéaire homogène associée
\[\forall n\in\N\qsep u_{n+1}+a_n u_n=0.\]
\end{proposition}

\begin{exos}
  \exo Soit $a\in\R$. Calculer le $n$-ième terme de la suite $\p{u_n}$
    définie par
    \[u_0\defeq a \et \forall n\in\N \qsep u_{n+1}\defeq\frac{3}{2}u_n+5.\]
    \begin{sol}
    On trouve
    \[\forall n\in\N \quad u_n=\p{10+a}\p{\frac{3}{2}}^n-10\]
    \end{sol}
  \exo On considère la récurrence linéaire
    \[(E) \quad \forall n\in\N\qsep u_{n+1}-2u_n=n^2.\]
    \begin{questions}
    \question Déterminer une solution polynomiale de $(E)$.
    \question En déduire toutes les solutions.
    \end{questions}
    \begin{sol}
  \[u_n=-(n^2+2n+3)+\lambda 2^n\]
    \end{sol}
  \end{exos}

% Étant donné une suite $(v_n)$, on cherche à calculer le terme général d'une suite $(u_n)$ vérifiant
% \[(E) \qquad \forall n\in\N\qsep u_{n+1}=a u_n + v_n\]
% où $a\in\Ks$. Plusieurs méthodes sont disponibles.

\vspace{2ex}
\begin{remarques}
\remarque \emph{Méthode de la similitude}~: Si l'on cherche les suites vérifiant la
  récurrence linéaire $u_{n+1}=a u_n+b$, on introduit la fonction $f:\K\to\K$ définie
  par $f(z)\defeq az+b$.
  \begin{itemize}
  \item Si $a=1$, une récurrence immédiate nous montre que
    \[\forall n\in\N\qsep u_n=u_0 + nb.\]
  \item Sinon, $f$ admet un unique point fixe $\omega\in\K$ et, pour tout $z\in\K$, $f(z)=a(z-\omega)+\omega$.
    On a donc
    \[\forall n\in\N\qsep u_{n+1}-\omega=a(u_n-\omega)\]
    et une récurrence immédiate donne
    \[\forall n\in\N\qsep u_n = a^n(u_0-\omega) +\omega.\]
  \end{itemize}
\remarque \emph{Méthode de la sommation télescopique}~: Si l'on cherche les suites
  vérifiant la récurrence linéaire $u_{n+1}= a u_n + b_n$ où $a\in\Ks$, on commence par remarquer que
\[\forall k\in\N\qsep \frac{u_{k+1}}{a^{k+1}} - \frac{u_k}{a^k}=\frac{b_k}{a^{k+1}},\]
puis on somme cette relation pour $k$ allant de 0 à $n-1$. On obtient une somme télescopique, puis
\[\forall n\in\N\qsep \frac{u_n}{a^n}-u_0=\sum_{k=0}^{n-1} \frac{b_k}{a^{k+1}}\]
et donc
\[\forall n\in\N\qsep u_n=u_0 a^n + \sum_{k=0}^{n-1} b_k a^{n-(k+1)}.\]
% \item \emph{Méthode de superposition}~: On commence par chercher une solution particulière $(w_n)$ de $(E)$.
%   Par exemple, si $v_n$ est un polynôme en $n$, il est naturel de chercher une solution polynomiale en $n$.
%   Une fois qu'on a trouvé une telle solution, il suffit de remarquer que
%   \begin{eqnarray*}
%   \forall n\in\N \qsep u_{n+1}=a u_n+v_n
%   &\ssi& \forall n\in\N \qsep u_{n+1}-a u_n=v_n\\
%   &\ssi& \forall n\in\N \qsep u_{n+1}-a u_n=w_{n+1}-aw_n\\
%   &\ssi& \forall n\in\N \qsep u_{n+1}-w_{n+1}=a(u_n-w_n)\\
%   &\ssi& \exists\lambda\in\K\qsep \forall n\in\N \qsep u_n-w_n=\lambda a^n\\
%   &\ssi& \exists\lambda\in\K\qsep \forall n\in\N \qsep u_n=\lambda a^n+w_n
%   \end{eqnarray*}
%   Les solutions de $(E)$ sont donc les suites de la forme $u_n\defeq \lambda a^n+w_n$ où $\lambda\in\K$.
%   Remarquons que si l'on considère la relation de récurrence linéaire \emph{homogène} associée à $(E)$,
%   obtenue en remplaçant $v_n$ par 0
%   \[(E_H) \qquad \forall n\in\N\qsep u_{n+1}=a u_n,\]
%   alors, les solutions de $(E)$ sont obtenues en ajoutant à une solution particulière de $(E)$
%   l'ensemble des solutions de la récurrence linéaire homogène associée $(E_H)$. Cette méthode
%   se généralisera aux récurrences linéaires d'ordre quelconque.
\end{remarques}




% On pourra aussi utiliser le \emph{théorème de superposition} qui nous dit que si $(w_n)$ est solution particulière de $(E)$, l'ensemble des solutions s'obtient comme somme de cette solution particulière et de la solution générale de l'équation homogène associée; dans ce cas, c'est la relation de récurrence
% \[(E_H) \qquad \forall n\in\N\qsep u_{n+1}=a u_n.\]
% Les solutions de $(E)$ sont donc les suites de la forme
% \[\forall n\in\N\qsep u_n\defeq \lambda a^n+w_n\]
% où $\lambda\in\C$.



\subsection{Récurrence linéaire d'ordre 2}
\begin{proposition}[utile=-3]
Soit $a,b\in\C$. On souhaite trouver les suites $(u_n)$ vérifiant
\[(E)\quad\forall n\in\N \qsep u_{n+2}=au_{n+1}+b u_n.\]
On résout sur $\C$ l'équation caractéristique $z^2=az+b$.
\begin{itemize}
\item Si cette équation admet deux racines distinctes $r_1$ et $r_2\in\C$, alors
  les solutions de $(E)$ sont les suites $(u_n)$ définies par
  \[\forall n\in\N \qsep u_n\defeq\lambda r_1^n +\mu r_2^n\]
  où $\lambda,\mu\in\C$.
\item Si cette équation admet une racine double $r\in\Cs$, alors 
  les solutions de $(E)$ sont les suites $(u_n)$ définies par
  \[\forall n\in\N \qsep u_n\defeq\p{\lambda +\mu n}r^n\]
  où $\lambda,\mu\in\C$.
\end{itemize}
\end{proposition}

\begin{preuve}
Soit $a,b\in\C$. On souhaite trouver les suite $(u_n)$ vérifiant
\[(E)\quad\forall n\in\N \qsep u_{n+2}=au_{n+1}+b u_n.\]
On commence par chercher les suites solutions de $(E)$ de la forme $(r^n)$. Soit $r\in\C$. On définit la suite $(u_n)$ par
\[\forall n\in\N\qsep u_n\defeq r^n.\]
Alors
\begin{eqnarray*}
\text{$(u_n)$ est solution de $(E)$}
&\ssi& \forall n\in\N\qsep u_{n+2}=au_{n+1}+b u_n\\
&\ssi& \forall n\in\N\qsep r^{n+2}=a r^{n+1}+b r^n\\
&\ssi& \forall n\in\N\qsep r^n\p{r^2-ar-b}=0\\
&\ssi& r^2-ar-b=0
\end{eqnarray*}
Dans la dernière équivalence, le sens $\Longleftarrow$ est évident, et le sens $\implique$ s'obtient en prenant $n=0$. En conclusion $(r^n)$ est solution de $E$ si et seulement si $r$ est solution de l'équation caractéristique $z^2=az+b$.
\begin{itemize}
\item Supposons que cette équation admette deux racines distinctes $r_1, r_2\in\C$. Montrons que les solutions $(E)$ sont les suites $(\lambda r_1^n+\mu r_2^n)$ où $\lambda,\mu\in\C$.
\begin{itemize}
\item Ce sont des solutions de $(E)$. En effet, soit $\lambda,\mu\in\C$ et $(u_n)$ la suite définie par
\[\forall n\in\N\qsep u_n\defeq\lambda r_1^n+\mu r_2^n.\]
Alors
\begin{eqnarray*}
\forall n\in\N\qsep u_{n+2}
&=& \lambda r_1^{n+2} + \mu r_2^{n+2}\\
&=& \lambda \p{a r_1^{n+1}+b r_1^n} + \mu \p{a r_2^{n+1}+b r_2^n}\\
& & \text{car $(r_1^n)$ et $(r_2^n)$ sont solutions de $E$}\\
&=& a\p{\lambda r_1^{n+1}+\mu r_2^{n+1}} + b\p{\lambda r_1^n + \mu r_2^n}\\
&=& a u_{n+1}+b u_n.
\end{eqnarray*}
Donc $(u_n)$ est solution de $(E)$.
\item Réciproquement, soit $(u_n)$ une solution de $(E)$. Montrons qu'il existe $\lambda,\mu\in\C$ tels que
\[\forall n\in\N\qsep u_n=\lambda r_1^n+\mu r_2^n.\]
Soit $\lambda,\mu\in\C$. On définit la suite $(v_n)$ par
\[\forall n\in\N\qsep v_n\defeq\lambda r_1^n+\mu r_2^n.\]
Alors
\begin{eqnarray*}
\syslin{v_0&=&u_0\cr v_1&=&u_1}
&\ssi& \syslin{\lambda &+\mu&=&u_0\cr r_1 \lambda&+r_2 \mu&=&u_1}\\
&\ssi& \syslin{\lambda &+\mu&=&u_0\hfill\cr &(r_2-r_1) \mu&=&u_1-r_1 u_0 \quad L_2\gets L_2-r_1 L_1}\\
&\ssi& \syslin{\lambda&=&\frac{u_1-r_2 u_0}{r_1-r_2}\cr \mu&=&\frac{u_1-r_1 u_0}{r_2-r_1}}\\
&    & \text{car $r_2-r_1\neq 0$.}
\end{eqnarray*}
On pose donc $\lambda\defeq (u_1-r_2 u_0)/(r_1-r_2)$ et $\mu\defeq(u_1-r_1 u_0)/(r_2-r_1)$. Les suites $(u_n)$ et $(v_n)$ sont donc deux solutions de $(E)$ telles que $u_0=v_0$ et $u_1=v_1$. Une récurrence montre facilement que pour tout $n\in\N$, $u_n=v_n$. On en déduit que
\[\forall n\in\N\qsep u_n=\lambda r_1^n+\mu r_2^n.\]
\end{itemize}
On a donc prouvé que les solutions de $(E)$ sont les suites $(\lambda r_1^n+\mu r_2^n)$ où $\lambda,\mu\in\C$.
\item Supposons maintenant que l'équation $z^2-az-b=0$ admette une racine double $r\in\Cs$. Puisque la somme des racines est $a$, on en déduit que $2r=a$. Montrons que les solutions de $(E)$ sont les suites $\p{(\lambda+\mu n)r^n}$ où $\lambda,\mu\in\C$.
\begin{itemize}
\item Ce sont des solutions de $(E)$. On sait déjà que $(r^n)$ est une solution de $(E)$. Montrons que la suite $(u_n)$ définie par
\[\forall n\in\N\qsep u_n=n r^n\]
en est une autre. On a
\begin{eqnarray*}
\forall n\in\N\qsep u_{n+2}-\p{au_{n+1}+bu_n}
&=& (n+2)r^{n+2}-\cro{a(n+1)r^{n+1}+b n r^n}\\
&=& r^n\cro{(n+2)r^2-(n+1)ar-bn}\\
&=& r^n\big[\underbrace{\p{r^2-ar-b}}_{=0}n+r\underbrace{(2r-a)}_{=0}\big]\\
&=& 0
\end{eqnarray*}
donc $(n r^n)$ est solution de $E$. On en déduit facilement que si $\lambda,\mu\in\C$ la suite $(\lambda r^n + \mu n r^n)$ est une solution de $(E)$.
\item Réciproquement, soit $(u_n)$ une solution de $(E)$. Montrons qu'il existe $\lambda,\mu\in\C$ tels que
\[\forall n\in\N\qsep u_n=(\lambda+\mu n)r^n.\]
Soit $\lambda,\mu\in\C$. On définit la suite $(v_n)$ par
\[\forall n\in\N\qsep v_n\defeq(\lambda+\mu n)r^n.\]
Alors
\begin{eqnarray*}
\syslin{v_0&=&u_0\cr v_1&=&u_1}
&\ssi& \syslin{\lambda & &=&u_0\cr r \lambda&+r \mu&=&u_1}\\
&\ssi& \syslin{\lambda & &=&u_0\hfill\cr & r \mu&=&u_1-r u_0 \quad L_2\gets L_2-r L_1}\\
&\ssi& \syslin{\lambda&=&u_0\hfill\cr \mu&=&\frac{u_1-r u_0}{r}}\\
&    & \text{car $r\neq 0$ .}
\end{eqnarray*}
On pose donc $\lambda\defeq u_0$ et $\mu\defeq(u_1-r u_0)/r$. Les suites $(u_n)$ et $(v_n)$ sont donc deux solutions de $(E)$ telles que $u_0=v_0$ et $u_1=v_1$. Une récurrence montre facilement que pour tout $n\in\N$, $u_n=v_n$. On en déduit que
\[\forall n\in\N\qsep u_n=(\lambda +\mu n)r^n.\]
\end{itemize}
On a donc prouvé que les solutions de $(E)$ sont les suites $((\lambda+\mu n)r^n)$ où $\lambda,\mu\in\C$.
\end{itemize}
 % \begin{itemize}
% \item On se place dans le cas où l'équation caractéristique admet deux racines distinctes.
% \begin{itemize}
% \item[$\bullet$] Déjà, étant donné $\lambda,\mu \in\C$, on vérifie que $(v_n)$ définie par $v_n=\lambda r_1^n+\mu r_2^n$ est solution.
% Remarquons de plus qu'on a alors $v_0=\lambda+\mu$ et $v_1=\lambda r_1+\mu r_2$, d'où $\mu=\dfrac{v_1-r_1v_0}{r_2-r_1}$ et $\lambda=\dfrac{v_1-r_2v_0}{r_1-r_2}$.
% \item[$\bullet$] Considérons maintenant $(u_n)$ une solution.

% Définissons $(v_n)$ par $v_n=\lambda r_1^n+\mu r_2^n$ avec $\lambda=\dfrac{u_1-r_2u_0}{r_1-r_2}$ et $\mu=\dfrac{u_1-r_1u_0}{r_2-r_1}$. On a alors $v_0=\lambda+\mu=u_0$ et $v_1=\lambda r_1+\mu r_1=u_1$.  De plus, d'après le premier point, $(v_n)$ est solution.

% Les deux suites sont égales aux deux premiers rangs puis vérifiant la relation de récurrence linéaire d'ordre $2$ sont donc égales à tout rang.

% On a montré que $(u_n)$ était du type annoncé.
% \end{itemize}
% \item On se place dans le cas où l'équation caractéristique admet une racine double. Dans ce cas, $r^2=-b$ et $2r=a$.
% \begin{itemize}
% \item[$\bullet$] Déjà, étant donné $\lambda,\mu \in\C$, on vérifie que $(v_n)$ définie par $v_n=\lambda r^n+\mu nr^n$ est solution.

% En effet :
% $$av_{n+1}+b v_n=2r(\lambda r^{n+1}+\mu (n+1)r^{n+1})-r^2(\lambda r^n+\mu nr^n)=\lambda r^{n+2}+\mu (n+2)r^{n+2}=v_{n+2}.$$

% Remarquons de plus qu'on a alors $v_0=\lambda$ et $v_1=(\lambda+\mu) r$, d'où $\lambda=v_0$ et $\mu=\dfrac{v_1}{r}-v_0$.
% \item[$\bullet$] Considérons maintenant $(u_n)$ une solution.

% Définissons $(v_n)$ par $v_n=(\lambda+\mu n) r^n$ avec $\lambda=u_0$ et $\mu=\dfrac{u_1}{r}-u_0$. On a alors $v_0=\lambda=u_0$ et $v_1=(\lambda +\mu) r=u_1$.  De plus, d'après le premier point, $(v_n)$ est solution.

% Les deux suites sont égales aux deux premiers rangs puis vérifiant la relation de récurrence linéaire d'ordre $2$ sont donc égales à tout rang.

% On a montré que $(u_n)$ était du type annoncé.
% \end{itemize}

% \end{itemize}
\end{preuve}


\begin{exos}
\exo Calculer le $n$-ième terme de la suite de \nom{Fibonacci} définie par
  \[F_0\defeq 1, \qquad F_1\defeq 1 \et \forall n\in\N \qsep F_{n+2}\defeq F_{n+1}+F_n.\]
  % À changer, c'était la réponse avec F_0=0
  % \begin{sol}
  % On trouve
  % \[F_n=\frac{\sqrt{5}}{5}\p{\frac{1+\sqrt{5}}{5}}^n-
  %       \frac{\sqrt{5}}{5}\p{\frac{1-\sqrt{5}}{5}}^n\]
  % \end{sol}
% \exo Soit $\p{u_n}$ la suite définie par
%   \[u_0\defeq 1, \qquad u_1\defeq -1 \et \forall n\in\N \qsep 6u_{n+2}+5u_{n+1}-6u_n=0.\]
%   Montrer que $\p{u_n}$ diverge.
%   \begin{sol}
%   On trouve
%   \[u_n=\frac{3}{13}\p{\frac{2}{3}}^n+\frac{10}{13}\p{-\frac{3}{2}}^n\]
%   \end{sol}
\exo Soit $a,b\in\RPs$. Calculer le $n$-ième terme de la suite $\p{u_n}$ définie par
  \[u_0\defeq a, \quad u_1\defeq b \et \forall n\in\N \qsep u_{n+2}\defeq\frac{1}{u_{n+1}^2u_n}.\]
  \begin{sol}
  On trouve
  \[\forall n\in\N \quad u_n=\p{u_1^n u_0^{n-1}}^{\p{-1}^{n+1}}\]
  \end{sol}
\end{exos}

\begin{remarques}
\remarque La suite de \nom{Fibonacci} est ainsi nommée en hommage à \nom{Leonardo Pisano} (Leonard de Pise, 1170--1240) appelé aussi \nom{Leonardo Fibonacci}, qui avait publié cette suite en 1202. Il avait lu le travail de \nom{Al-Khwarizmi} (780--850), un mathématicien Persan. Le livre de \nom{Fibonacci} contient le problème suivant. Combien de couples de lapins peuvent naître d'un couple de lapin en un an ? Pour résoudre ce problème, on sait que~:
\begin{itemize}
\item Jusqu'au premier mois inclus, il n'y a qu'un couple de lapins.
\item Chaque couple de lapin donne naissance à un couple tous les mois.
\item Chaque jeune couple devient fertile à l'âge d'un mois.
\end{itemize}
Avant le travail de \nom{Fibonacci}, la suite $(F_n)$ a déjà été étudiée par les Indiens qui se demandaient combien de rythmes de $n$ temps il était possible de faire avec des noires et des blanches.
\remarque Tout comme pour les récurrences linéaires d'ordre 1, afin de
  résoudre une récurrence linéaire de la forme
  \[\forall n\in\N\qsep u_{n+2}=a u_{n+1}+b u_n + c_n,\]
  il suffit d'en trouver une solution particulière $(v_n)$ et d'y ajouter les solutions
  de la récurrence linéaire homogène associée $u_{n+2}=a u_{n+1}+b u_n$.
\end{remarques}

\begin{proposition}[utile=-3]
Soit $a,b\in\R$. On souhaite trouver les suites $(u_n)$ vérifiant
\[(E) \quad \forall n\in\N \qsep u_{n+2}=au_{n+1}+b u_n.\]
On résout sur $\C$ l'équation caractéristique $z^2=az+b$.
\begin{itemize}
\item Si cette équation admet deux racines réelles distinctes $r_1$ et $r_2$,
  alors les solutions de $(E)$ sont les suites $(u_n)$ définies par
  \[\forall n\in\N \qsep u_n\defeq\lambda r_1^n +\mu r_2^n\]
  où $\lambda,\mu\in\R$.
\item Si cette équation admet une racine double $r\in\Rs$,
  alors les solutions de $(E)$ sont les suites $(u_n)$ définies par
  \[\forall n\in\N \qsep u_n\defeq\p{\lambda +\mu n}r^n\]
  où $\lambda,\mu\in\R$.
\item Si cette équation admet deux racines complexes conjuguées $r\e^{\ii\omega}$ et
  $r\e^{-\ii\omega}$, 
  alors les solutions de $(E)$ sont les suites $(u_n)$ définies par
  \[\forall n\in\N \qsep u_n\defeq\cro{\lambda \cos\p{\omega n}+\mu \sin\p{\omega n}}
    r^n\]
  où $\lambda,\mu\in\R$.
\end{itemize}
\end{proposition}

\begin{preuve}
Les deux premiers cas sont une conséquence de la proposition précédente. Il reste donc à traiter le cas où $\Delta<0$.

L'EC admet alors deux racines conjuguées $r\e^{i\omega}$ et $r\e^{-i\omega}$. 

\underline{Analyse :} Soit $(u_n)$ une solution réelle de $(E)$. D'après la proposition précédente, on peut fixer $(A,B)\in \C^2$ tels que \[\forall n\in\N, u_n=Ar^n\e^{i\omega n}+Br^n\e^{-i\omega n}=r^n\p{A\e^{i\omega n}+B\e^{-i\omega n}}.\]
Or, $\forall n\in\N, u_n=\Re(u_n)=r^n\Re\p{A\e^{i\omega n}+B\e^{-i\omega n}}$ qui est donc de la forme $u_n)=\cro{c_1\cos\p{\omega n}+c_2\sin\p{\omega n}}r^n$.

\underline{Synthèse :} Réciproquement, soit $(c_1,c_2)\in\R^2$. Définissons $(u_n)$ pour tout $n$ entier naturel par $$u_n=\p{c_1\cos\p{\omega n}+c_2\sin\p{\omega n}}r^n.$$
On applique les formules d'Euler pour s'apercevoir que $(u_n)$ ainsi définie est bien solution de $(E)$ car de la forme des solutions de la proposition précédente.

\end{preuve}

\begin{remarqueUnique}
% \remarque L'ensemble des suites réelles vérifiant cette relation de
%   récurrence est un \Rev de dimension 2.
\remarque Dans le cas où l'équation caractéristique admet deux racines
  complexes conjuguées, les solutions de $(E)$ peuvent s'écrire sous la forme
  \[\forall n\in\N \qsep u_n\defeq\lambda\sin(\omega n-\phi) r^n\]
  où $\lambda,\phi\in\R$.
  Lors de la recherche effective de tels coefficients, 
  quitte à changer $\phi$ en $\phi+\pi$, on impose souvent $\lambda\in\RP$.
\end{remarqueUnique}

\begin{exos}
\exo Soit $\p{u_n}$ la suite définie par
  \[u_0\defeq 1\qsep u_1\defeq 0, \et \forall n\in\N \qsep u_{n+2}\defeq 2u_{n+1}-4u_n.\]
  Déterminer les $n\in\N$ pour lesquels $u_n=0$.
  \begin{sol}
  On trouve
  \[u_n=\p{\cos\p{\frac{n\pi}{3}}-\frac{\sqrt{3}}{3}\sin\p{\frac{n\pi}{3}}}2^n\]
  Donc $u_n=0$ pour $n\equiv 1 \ [3]$
  \end{sol}
\exo Soit $\alpha\in\R$. Calculer le $n$-ième terme de la suite $\p{u_n}$
  définie par
  \[u_0,u_1 \in\R \et \forall n\in\N \qsep u_{n+2}\defeq 2\cos\p{\alpha}u_{n+1}-u_n.\]
  \begin{sol}
  On trouve
  \begin{itemize}
  \item Si $\alpha\equiv 0\ [2\pi]$, $u_n=\alpha n+\beta$.
  \item Si $\alpha\equiv \pi\ [2\pi]$, $u_n=\p{\alpha n+\beta}\p{-1}^n$.
  \item Sinon, $u_n=\lambda\cos\p{\alpha n}+\mu\sin\p{\alpha n}$.
  \end{itemize}
  \end{sol}
% \exo Donner une base de l'ensemble des suites réelles vérifiant la relation
%   \[\forall n\in\N \quad u_{n+3}=2u_{n+2}+u_{n+1}-2u_n\]
%   \begin{sol}
%   Voir que $u_{n+1}-u_n$ vérifie une récurrence linéaire d'ordre 2. On trouve
%   $2^n,1,\p{-1}^n$.
%   \end{sol}
\end{exos}
%% Exemple :
%% 1) Suite de Fibonacci (Knuth, volume  1, page 79)
%%    F_0=1 F_1=1 F_(n+2)=F_(n+1)+F_n
%%    1,1,2,3,5,8,13,21,34...
%%    - Leonardo Pisano (Leonardo de Pise) appelé aussi Leonardo Fibonacci
%%      Publié en 1202. Il avait lu le travail de Al-Khwarizmi, 825 (Persan)
%%      (la mot algèbre vient du titre de son oeuvre majeure : Kitab al-jabr
%%      wa'l-muqabala) (Rules of restoring and equating)
%%      Son livre contient le problème suivant : Combien de couples de lapins
%%      peuvent naitre d'un couple de lapin en un an ?
%%      Pour rédoudre ce problème, on sait que :
%%      - chaque couple donne naissance a un couple chaque mois
%%      - chaque couple devient fertile à l'age d'un mois
%%      - les lapins ne meurent jamais
%%    - Avant le travail de Fibonacci, la suite (F_n) a déjà été étudiée par les
%%      indiens. Ils se demandaient combine de rythmes de n temps il était possible
%%      de faire avec des noires et des blanches. Ce nombre est F_(n+1)
%%   
%%    Calcul du n-ième nombre de Fibonacci:
%%    (1+sqrt(5))/2 et (1-sqrt(5))/2
%%    On trouve F_n=(1/sqrt(5))(((1+sqrt(5))/2)^(n+1)-((1-sqrt(5))/2)^(n+1))
  

\section{Système linéaire}


\subsection{Système linéaire à $q$ équations et $p$ inconnues}

\begin{definition}[utile=-3]
On appelle \emph{système linéaire} à $q$ équations et $p$ inconnues tout système
d'équations du type
\[\syslin{a_{1,1}x_1&+a_{1,2}x_2&+\cdots&+a_{1,p}x_p&=&y_1\cr
                   &          &       &          &\hfill\vdots\hfill&\cr
          a_{q,1}x_1&+a_{q,2}x_2&+\cdots&+a_{q,p}x_p&=&y_q}\]
où $a_{1,1},\ldots,a_{q,p},y_1,\ldots,y_q\in\K$ et $x_1,\ldots,x_p\in\K$ sont
les inconnues. On dit que le système est \emph{compatible} lorsqu'il admet au moins une
solution. On dit qu'il est \emph{incompatible} sinon.
\end{definition}

\begin{remarques}
\remarque Pour des raisons de lisibilité, on veillera à toujours placer les inconnues
  les unes en dessous des autres.
\remarque L'ensemble des solutions est l'ensemble $\mathcal{S}$ des
  $(x_1,\ldots,x_p)\in\K^p$ solution du système.
  % \remarque Un système linéaire est composé de $q$ équations aussi appelées \emph{lignes}.
  % Les $q$-uplets $(a_{1,j},\ldots,a_{q,j})\in\K^q$ pour $j\in\intere{1}{p}$ sont appelées
  % \emph{colonnes}. Le $q$-uplet $(y_1,\ldots,y_q)\in\K^q$ est appelé \emph{second membre}.
  % \remarque On dit qu'un système est \emph{homogène} lorsque son second membre est nul,
  % c'est-à-dire lorsque $(y_1,\ldots,y_q)=(0,\ldots,0)$. Dans ce cas, le $p$-uplet
  % $(x_1,\ldots,x_p)=(0,\ldots,0)$ est une solution triviale du système.
  % \remarque Si $\mathcal{S}$ est un système linéaire, on appelle système linéaire homogène
  % associé, le système obtenu en remplaçant le second membre par le $q$-uplet $(0,\ldots,0)$.
\end{remarques}

% \begin{definition}
% On conserve les mêmes notations que la définition précédente.
% \begin{itemize}
% \item On dit que le système est \emph{homogène} lorsque $(y_1,\ldots,y_q)=(0,\ldots,0)$.
% \item On appelle \emph{système homogène associé}, le système
%   \[\syslin{a_{1,1}x_1&+a_{1,2}x_2&+\cdots&+a_{1,p}x_p&=&0\cr
%                    &          &       &          &\hfill\vdots\hfill&\cr
%           a_{q,1}x_1&+a_{q,2}x_2&+\cdots&+a_{q,p}x_p&=&0}\]
%   obtenu en remplaçant les $y_i$ par $0$.
% \end{itemize}
% \end{definition}

\begin{exoUnique}
\exo Résoudre le système suivant par substitution, puis en utilisant la méthode du pivot
de \nom{Gauss}.
  \[\syslin{x&-2y&=&-4\cr
            3x&+y&=&9.\hfill}\]
  \begin{sol}
  $x=2$ et $y=3$.            
  \end{sol}
\end{exoUnique}

\begin{proposition}[utile=-3]
Les opérations suivantes, appelées opérations élémentaires, transforment un
système linéaire en un système linéaire équivalent.
\begin{itemize}
\item Changer l'ordre des équations.
\item Changer l'ordre des inconnues.
\item Multiplier une équation par $\mu\in\Ks$.
\item Ajouter $\lambda$ fois (avec $\lambda\in\K$) une équation à l'une des
  équations suivantes.
\end{itemize}
\end{proposition}

\begin{remarqueUnique}
\remarque En pratique, afin d'expliciter les opérations que l'on vient d'effectuer, on
  utilisera les notations suivantes.
\begin{itemize}
\item $L_i\leftrightarrow L_j$ signifie qu'on a échangé les lignes $i$ et $j$.
\item $L_i\gets \mu L_i$ signifie qu'on a multiplié la ligne $L_i$ par le coefficient $\mu$ non nul.
\item $L_i\gets L_i + \lambda L_j$ signifie qu'on a ajouté $\lambda$ fois la ligne $L_j$ à la ligne $L_i$.
\end{itemize}
% \remarque Si on multiplie une ligne par $\mu= 0$, on perd l'équation correspondante et
%   donc l'équivalence. De même, les opérations du type $L_i \gets L_j + \mu L_i$ font
%   perdre l'équivalence lorsque $\mu=0$. Afin d'éviter de tels écueils, ces opérations
%   sont interdites lors de la résolution d'un système linéaire.
\end{remarqueUnique}

\begin{exoUnique}
\exo Les opérations élémentaires suivantes conservent-elles l'équivalence~?
  \begin{itemize}
  \item $L_1 \gets L_2$.
  \item $L_1 \gets L_1+L_2+L_3$.
  \item $L_1 \gets 2 L_1 + L_2$.
  \item $L_1 \gets \alpha L_1 + \beta L_2$ où $\alpha,\beta\in\K$.
  \item $L_1 \gets L_1+L_2$ et $L_2 \gets L_1-L_2$.
  \end{itemize}
\end{exoUnique}

\begin{proposition}%[nom={Pivot de \nom{Gauss}}]
L'algorithme du pivot de \nom{Gauss} permet de transformer, quitte à échanger les
variables, un système linéaire à $q$ équations et $p$ inconnues en un système linéaire
équivalent de la forme
\[\syslin{a_{1,1}x_1&+a_{1,2}x_2&+\hfill&\cdots   &&+a_{1,p}x_p&=&y_1\hfill\cr
                       & a_{2,2}x_2&+\hfill&\cdots   &&+a_{2,p}x_p&=&y_2\hfill\cr
                       &           &\ddots&         &      &          &=&
      \ \vdots\hfill\cr
                       &           &      &a_{r,r}x_r&+\cdots&+a_{r,p}x_p&=&y_r\hfill\cr
                       &           &      &          &       &          0&=&y_{r+1}\hfill\cr
                       &           &      &          &       &          \vdots\,&=&\ \vdots\hfill\cr
                       &           &      &          &       &          0&=&y_{q}\hfill}\]
où $a_{1,1},\ldots,a_{r,r}$ sont tous non nuls. On dit d'un tel système
qu'il est \emph{échelonné à pivots diagonaux}.
\begin{itemize}
\item Le système est compatible si et seulement si $(y_{r+1},\ldots,y_q)=(0,\ldots,0)$.
\item Le système admet une unique solution si et seulement si il est compatible et
  $r=p$.
\end{itemize}
\end{proposition}

\begin{exos}
\exo Résoudre les systèmes
  \[\syslin{5x&+y&=&1\hfill\cr
            11x&+2y&=&3,} \qquad
    \syslin{2x&+3y&=&1\hfill\cr
            5x&+2y&=&1.}\]
  \begin{sol}
  On trouve $x=1$ et $y=-4$ pour le premier, $x=1/11$ et $y=3/11$ pour le
  second.           
  \end{sol}
\exo Soit $\alpha\in\R$. Résoudre le système
  \[\syslin{x&+y&-z&=&1\hfill\cr
            2x&+y&+2z&=&2\hfill\cr
            3x&+2y&+z&=&\alpha.}\]
  \begin{sol}
  Si $\alpha\neq 3$,, il n'y a pas de solution. si $\alpha=3$, on trouve
  $x=1-3t$, $y=4t$, $z=t$.          
  \end{sol}
\end{exos}

% \begin{definition}[utile=-3]
% On dit qu'un système linéaire à $q$ équations et $p$ inconnues est triangulaire
% lorsqu'il est de la forme~:
% \[\syslin{a_{1,1}x_1&+a_{1,2}x_2&+\hfill&\cdots   &&+a_{1,p}x_p&=&y_1\cr
%                    & a_{2,2}x_2&+\hfill&\cdots   &&+a_{2,p}x_p&=&y_2\cr
%                    &           &\ddots&         &      &          &=&
%   \ \vdots\hfill\cr
%                    &           &      &a_{q,q}x_q&+\cdots&+a_{q,p}x_p&=&y_q}\]
% où les $a_{1,1},a_{2,2},\ldots,a_{q,q}$ sont tous non nuls.
% \end{definition}



% \begin{proposition}[utile=-3]
% À l'aide d'opérations élémentaires, il est possible de transformer tout
% système linéaire en un système triangulaire équivalent.
% \end{proposition}

\begin{remarques}
\remarque Voici une présentation détaillée de l'algorithme du Pivot de \nom{Gauss}.
  \begin{itemize}
  \item \emph{On transforme le système en un système échelonné}.\\
  On commence par déterminer un coefficient $a_{i,j}$ non nul que l'on appelle \emph{pivot}.
  Très souvent $a_{1,1}$ conviendra, mais il est encore plus pratique pour la suite des
  calculs si ce coefficient est $\pm 1$. En effectuant un échange de lignes et
  d'inconnues, on \og remonte \fg ensuite ce coefficient en haut à gauche du
  système. On se retrouve donc dans le cas où $a_{1,1}\neq 0$.
  On utilise alors $a_{1,1}$ comme
  pivot pour éliminer l'inconnue $x_1$ des $q-1$ dernières lignes du système~:
    \[\syslin{a_{1,1}x_1&+a_{1,2}x_2&+\cdots&+a_{1,p}x_p&=&y_1\cr
              a_{2,1}x_1&+a_{2,2}x_2&+\cdots&+a_{2,p}x_p&=&y_2\cr
                       &          &       &          &\hfill\vdots\hfill&\cr
              a_{q,1}x_1&+a_{q,2}x_2&+\cdots&+a_{q,p}x_p&=&y_q} \quad\ssi\quad
    \syslin{a_{1,1}'x_1&+a_{1,2}'x_2&+\cdots&+a_{1,p}'x_p&=&y_1'\hfill\cr
                       &a_{2,2}'x_2&+\cdots&+a_{1,p}'x_p&=&y_2'\hfill\cr
                       &          &       &          &\hfill\vdots\hfill&\cr
                       &a_{q,2}'x_2&+\cdots&+a_{q,p}'x_p&=&y_q'.\hfill}\]
    Pour cela, il suffit d'effectuer les opérations suivantes~:
    \[L_2\gets L_2-\frac{a_{2,1}}{a_{1,1}}\cdot L_1, \qquad 
      L_3\gets L_3-\frac{a_{3,1}}{a_{1,1}}\cdot L_1, \qquad\ldots\qquad
      L_q\gets L_q-\frac{a_{q,1}}{a_{1,1}}\cdot L_1.\]

    On recommence ensuite le même procédé sur les $q-1$ dernières équations du système,
    en ne touchant plus à la première ligne. On cherche d'abord un coefficient
    $a_{i,j}'$ non nul pour lequel $i\geq 2$ et $j\geq 2$. Si un tel coefficient existe,
    un échange de lignes et d'inconnues permet de se ramener
    au cas où $a_{2,2}'\neq 0$ et de continuer l'algorithme.
    On réitère le procédé jusqu'à ce qu'on ne soit plus capable de trouver de
    pivot. Le système est alors échelonné.
    Au cours du calcul, s'il apparaît l'équation $0=0$, on l'élimine du système. Si au
    contraire il apparaît l'équation $0=b$ avec $b\neq 0$, le système n'admet aucune solution et la
    résolution est terminée.
  \item \emph{On introduit les paramètres $t_k$}.\\
    Dans le cas où le système admet au moins une solution, on aboutit à un système
    de la forme
    \[\syslin{a_{1,1}''x_1&+a_{1,2}''x_2&+\hfill&\cdots   &&+a_{1,p}''x_p&=&y_1''\cr
                       & a_{2,2}''x_2&+\hfill&\cdots   &&+a_{2,p}''x_p&=&y_2''\cr
                       &           &\ddots&         &      &          &=&
      \ \vdots\hfill\cr
                       &           &      &a_{r,r}''x_r&+\cdots&+a_{r,p}''x_p&=&y_r''}\]
    où les $a_{1,1}'',a_{2,2}'',\ldots,a_{r,r}''$ sont tous non nuls.
    Afin de paramétrer l'ensemble des solutions, on remarque que ce dernier système est
    équivalent au système triangulaire
    \[\exists t_{r+1},\ldots,t_p\in\K \qsep
      \syslin{a_{1,1}''x_1&+a_{1,2}''x_2&+\hfill& &\cdots   & &+a_{1,p}''x_p&=&y_1''\hfill\cr
                       & a_{2,2}''x_2&+\hfill&  &\cdots   &&+a_{2,p}''x_p&=&y_2''\hfill\cr
                       &           &\ddots&     &    &      &          &=&
      \ \vdots\hfill\cr
                       &           &      &a_{r,r}''x_r&+\cdots& &+a_{r,p}''x_p&=&y_r''\hfill\cr
                       &           &      &           &x_{r+1}& & &=&t_{r+1}\hfill\cr
                       &           &      &           &  &\ddots&&=&\ \vdots\hfill\cr
                       &           &      &           &  & &x_p&=&t_p.\hfill}\]    
    En effet, si $\p{x_1,\ldots,x_p}$ est solution de ce dernier système, on obtient le
    système précédent en ne gardant que les $r$ premières lignes.
    Réciproquement, si $\p{x_1,\ldots,x_p}$ est solution du système échelonné,
    on obtient ce dernier système en posant $t_{r+1}\defeq x_{r+1},\ldots,t_p\defeq x_p$.
  \item \emph{On résout le système triangulaire}.\\
    Ce dernier système se résout simplement en remontant les calculs de la
    dernière ligne à la première, par substitution. On obtient ainsi le système équivalent
    \[\exists t_{r+1},\ldots,t_p\in\K \qsep
      \syslin{x_1&=&c_1&+d_{1,r+1}t_{r+1}&+\cdots&+d_{1,p}t_p\cr
            \vdots\ &=&\vdots\ & & &\vdots\ \cr
              x_r&=&c_r&+d_{r,r+1}t_{r+1}&+\cdots&+d_{r,p}t_p\cr
              x_{r+1}&=& & t_{r+1}& &\cr
              \vdots\ &=& & &\ddots &\cr
              x_p&=& & & &t_p}\]
    qui est un \emph{paramétrage} de l'ensemble des solutions.\\
    Remarquons que lorsque les calculs
    sont complexes, au lieu de résoudre directement le système triangulaire par
    substitution, on peut aussi effectuer un pivot de \nom{Gauss} \og à l'envers \fg
    en commençant par éliminer les $x_p$ des $p-1$ premières équations avec la dernière ligne, puis
    en éliminant les $x_{p-1}$ des $p-2$ premières équations avec l'avant-dernière ligne,
    etc.
  \end{itemize}    
\remarque Lorsqu'on applique l'algorithme du pivot de \nom{Gauss}, on est libre de
  choisir le pivot que l'on souhaite. La seule contrainte est qu'il soit non nul.
  L'expérience montre cependant que certains choix sont plus judicieux que d'autres, car
  ils conduisent à des calculs plus simples.\\
  Par exemple, un pivot égal à 1 est idéal car les opérations sur les lignes sont réduites à $L_i \gets L_i - a_{i,1} L_1$. On évite ainsi les divisions, ce qui offre de nombreux avantages. Par exemple, lorsque les coefficients du système sont entiers, ils le restent après réduction du système. Dans le même ordre d'idées, avant de choisir le pivot,
  lorsque les coefficients d'une même ligne sont des multiples d'un entier $a\in\Zs$, il est bon de simplifier cette ligne par $a$.\\
  Enfin, lorsque les coefficients du système dépendent d'un paramètre $\alpha$, il est toujours préférable d'utiliser un pivot ne dépendant pas de $\alpha$. Cette stratégie permet d'éviter de discuter les cas où ce terme peut s'annuler, et limite la propagation de ce paramètre à tous les autres coefficients du système.
\remarque On montrera plus tard que l'entier $r$ est indépendant du choix des pivots.
  On l'appelle \emph{rang} du système. Dans le cas où le système est compatible, en posant $d\defeq p-r$, on remarque
  que l'ensemble des solutions est paramétré par les $d$ coefficients $t_{r+1},\ldots,t_p$.
  En physique ou en sciences industrielles, on dit
  que l'ensemble des solutions admet $d$ degrés de liberté. Par exemple,
  si $d=1$, l'ensemble des solutions est une droite affine. Si $d=2$, l'ensemble des
  solutions est un plan affine. Le fait que $p=d+r$ est un résultat que nous
  retrouverons dans le cours d'algèbre linéaire sous le nom de \emph{théorème du rang}.
% \remarque Enfin, il est parfois pratique dans les calculs d'omettre le nom des variables. On utilise ce qu'on appelle une \og matrice augmentée \fg.
% \begin{eqnarray*}
% \forall x,y\in\R\qsep \syslin{x&+2y&=&1\hfill\cr
%             x&+3y&=&4}
% &\ssi& \begin{pmatrix}[cc|c]
%   1 & 2 & 1\\
%   1 & 3 & 4
% \end{pmatrix}\\
% &\ssi& \begin{pmatrix}[cc|c]
%   1 & 2 & 1\\
%   0 & 1 & 3 
% \end{pmatrix}
% \begin{matrix}
% \\
% L_2\gets L_2-L_1
% \end{matrix}\\
% &\ssi& \begin{pmatrix}[cc|c]
%   1 & 0 & -5\\
%   0 & 1 & 3 
% \end{pmatrix}
% \begin{matrix}
% L_1\gets L_1-2L_2\\
% \,
% \end{matrix}\\
% &\ssi& \syslin{x&=&-5\hfill\cr
%         y&=&3.\hfill}
% \end{eqnarray*}
% Cette technique à l'avantage d'écrire le minimum nécessaire et de vous obliger à aligner les coefficients les uns au dessus des autres. Son seul inconvénient est de rendre difficile le changement de l'ordre des variables qui peut parfois être pratique pour trouver un bon pivot.
\end{remarques}

\begin{exoUnique}
\exo Discuter et résoudre les systèmes suivants, selon la valeur de
  $\alpha\in\R$~:
  \[\syslin{-3a&+\alpha b& & &=&0\hfill\cr
            -3a&-b&+2\alpha c& &=&0\hfill\cr
               &-2b&+c&+3\alpha d&=&0\hfill\cr
               &   &-c&+3d&=&0,} \qquad
    \syslin{a&+b&+c&+d&=&3\hfill\cr
            a&+\alpha b&+c&-\alpha d&=&\alpha+2\hfill\cr
            \alpha a&-b&-\alpha c&-\alpha d&=&-1.\hfill}\]
  \begin{sol}%
  Pour la première, si $\alpha\neq 1$, la seule solution est la solution nulle.
  Sinon les solutions sont données par $a=t$, $b=3t$, $c=3t$, $d=t$ avec
  $t\in\R$.\\
  Pour la seconde, si $\alpha\neq 0$ et $\alpha\neq 1$, on a
  \[a=1-\frac{\alpha+1}{2\alpha}t, \quad b=1+\frac{\alpha+1}{\alpha-1}t, \quad
    c=1-\frac{3\alpha^2+1}{2\alpha\p{\alpha-1}}t, \quad d=t\]
  Si $\alpha=0$, les solutions sont
  \[a=2-t, \quad b=1, \quad c=t, \quad d=0\]
  Si $\alpha=1$, les solutions sont
  \[a=1, \quad b=2-t, \quad c=t, \quad d=0\]
  \end{sol}
\end{exoUnique}

\begin{remarques}
\remarque Il est parfois pratique dans les calculs d'omettre le nom des variables. On
  utilise alors ce qu'on appelle une \emph{matrice augmentée}.
\begin{eqnarray*}
\syslin{x&+2y&=&1\hfill\cr
            x&+3y&=&4}
&\ssi& \begin{pmatrix}[cc|c]
  1 & 2 & 1\\
  1 & 3 & 4
\end{pmatrix}\\
&\underset{L_2\gets L_2-L_1}{\ssi}& \begin{pmatrix}[cc|c]
  1 & 2 & 1\\
  0 & 1 & 3 
\end{pmatrix}
\\
&\underset{L_1\gets L_1-2L_2}{\ssi}& \begin{pmatrix}[cc|c]
  1 & 0 & -5\\
  0 & 1 & 3 
\end{pmatrix}
\\
&\ssi& \syslin{x&=&-5\hfill\cr
        y&=&3.\hfill}
\end{eqnarray*}
Cette technique a l'avantage d'écrire le minimum nécessaire et de vous obliger à aligner
les coefficients les uns au-dessus des autres. Son seul inconvénient est de rendre
impossible le changement d'ordre des inconnues.
\remarque L'échange des inconnues étant impossible, il arrive que cette méthode nous 
  empêche de réduire le système à un système échelonné à pivots diagonaux.
  On se contente donc d'un \emph{système échelonné}, c'est-à-dire
  un système où chaque ligne commence par un nombre de zéros strictement supérieur à
  celui de la ligne précédente, comme dans l'exemple suivant
  \[\begin{pmatrix}[cccc|c]
    a_{1,1} & \star &   \star &   \star & \star \\
          0 &     0 & a_{2,3} &   \star & \star \\
          0 &     0 &       0  & a_{3,4} & \star
  \end{pmatrix}\]
  où $a_{1,1}, a_{2,3}$ et $a_{3,4}$ sont des coefficients non nuls qu'on appelle encore
  \emph{pivots}.
  On réintroduit ensuite les inconnues, avant de les réordonner
  pour obtenir un système échelonné à pivots diagonaux et finir la résolution
  du système.
\end{remarques}

\begin{definition}
On considère un système linéaire à $q$ équations et $p$ inconnues.
\[(E) \qquad \syslin{a_{1,1}x_1&+a_{1,2}x_2&+\cdots&+a_{1,p}x_p&=&y_1\cr
                    &          &       &          &\hfill\vdots\hfill&\cr
          a_{q,1}x_1&+a_{q,2}x_2&+\cdots&+a_{q,p}x_p&=&y_q}\]
\begin{itemize}
\item On dit qu'il est \emph{homogène} lorsque $(y_1,\ldots,y_q)=(0,\ldots,0)$.
\item On appelle \emph{système linéaire homogène} associé à $(E)$, le système obtenu en
  remplaçant les $y_i$ par $0$.
\end{itemize}
\end{definition}

\begin{remarqueUnique}
\remarque Le $p$-uplet $(x_1,\ldots,x_p)=(0,\ldots,0)$ est toujours solution d'un système homogène. On dit
  que c'est la \emph{solution triviale}. Il est possible que ce soit la seule solution
  ou qu'il y en ait d'autres.
\end{remarqueUnique}

% \begin{proposition}
% Soit $(E)$ un système linéaire à $n$ équations et $n$ inconnues. On suppose que la solution
% triviale est l'unique solution du système linéaire homogène associé. Alors
% $(E)$ admet une et une seule solution.
% \end{proposition}  

% \begin{remarqueUnique}
% \remarque Cette proposition est utile pour résoudre des systèmes linéaires
%   possèdant des symétries. En effet, l'algorithme du pivot de \nom{Gauss} a
%   l'énorme inconvénient de casser ces symétries. On procède donc en effectuant une 
%   analyse. Cette proposition permet d'éviter la synthèse.
% \end{remarqueUnique}

% \begin{exoUnique}
% \exo Soit $a,b,c\in\C$. Résoudre le système
%     \[\syslin{x&+y&+z&=&a\hfill\cr
%             x&+\jj y&+\jj^2 z&=&b\hfill\cr
%             x&+\jj^2 y&+ \jj z&=&c.}\]
% \end{exoUnique}


\subsection{Interprétation géométrique lorsque $p=2$ ou $p=3$}

Dans cette partie, nous allons donner une interprétation géométrique des résultats précédents aux cas $p=2$ et $p=3$. Commençons par le cas $p=2$. On munit le plan d'un repère orthonormé $\mathcal{R}=(O,\ve{e_1},\ve{e_2})$. On rappelle qu'un point $M$ du plan est déterminé de manière unique par ses coordonnées $(x,y)\in\R^2$ définies par \[\ve{OM}=x\ve{e_1}+y\ve{e_2}.\]

\begin{proposition}
\begin{itemize}
\item Soit $(a,b,c)\in\R^3$ tel que $(a,b)\neq(0,0)$. Alors l'ensemble d'équation $ax+by=c$ est une droite $\mathcal{D}$ de vecteur normal $\ve{u}$ de coordonnées $(a,b)$.
\item Réciproquement, soit $\mathcal{D}$ une droite de vecteur normal $\ve{u}$ de coordonnées $(a,b)$. Alors, il existe $c\in\R$ tel que $ax+by=c$ est une équation de $\mathcal{D}$.
\end{itemize}
\end{proposition}

\begin{preuve}
On fait le lien avec le produit scalaire. 
\end{preuve}

Résoudre un système linéaire à $q$ équations et 2 inconnues $x$ et $y$ dont chaque ligne contient un coefficient non nul revient donc à déterminer l'intersection de $q$ droites du plan. Le cas où $q=2$ est important.
\begin{itemize}
\item Si les droites $\mathcal{D}_1$ et $\mathcal{D}_2$ ne sont pas parallèles, alors elles se coupent en un unique point. Le système admet donc une unique solution.
\item Si les droites $\mathcal{D}_1$ et $\mathcal{D}_2$ sont parallèles, deux cas se présentent.
\begin{itemize}
\item Si elles ne sont pas confondues, elles ne se coupent pas. Le système n'admet donc aucune solution.
\item Si elles sont confondues, le système admet une infinité de solutions. Ce sont les coordonnées des points de cette droite.
\end{itemize}
\end{itemize}

% \begin{proposition}
% Le système linéaire
% \[\syslin{
% a_1 x&+y_1 y&=&c_1\cr
% a_2 x&+y_2 y&=&c_2}\]  
% admet une unique solution si et seulement si son déterminant
% \[\begin{vmatrix}
% a_1 & y_1\\
% a_2 & y_2
% \end{vmatrix}=a_1 y_2 - a_2 y_1\]
% est non nul.
% \end{proposition}

Pour obtenir une interprétation géométrique du cas $p=3$, on munit l'espace d'un repère orthonormé \[\mathcal{R}=(O,\ve{e_1},\ve{e_2},\ve{e_3}).\] On rappelle qu'on point $M$ de l'espace est déterminé de manière unique par ses coordonnées $(x,y,z)\in\R^3$ définies par \[\ve{OM}=x\ve{e_1}+y\ve{e_2}+z\ve{e_3}.\]

\begin{proposition}
\ 
\begin{itemize}
\item Soit $(a,b,c,d)\in\R^4$ tel que $(a,b,c)\neq(0,0,0)$. Alors l'ensemble d'équation $ax+by+cz=d$ est un plan $\mathcal{P}$ de vecteur normal $\ve{u}$ de coordonnées $(a,b,c)$.
\item Réciproquement, soit $\mathcal{P}$ un plan de vecteur normal $\ve{u}$ de coordonnées $(a,b,c)$. Alors, il existe $d\in\R$ tel que $ax+by+cz=d$ est une équation de $\mathcal{P}$.
\end{itemize}
\end{proposition}

Résoudre un système linéaire à $q$ équations et 3 inconnues $x$, $y$ et 
$z$ dont chaque ligne contient un coefficient non nul revient donc à déterminer l'intersection de $q$ plans dans l'espace. Le cas où $q=3$ est important.
\begin{itemize}
\item Si le rang du système est égal à 3, les 3 plans s'intersectent en un unique point.
\item S'il est strictement inférieur à 3 et que le système n'est pas compatible, l'intersection des plans est vide. Sinon~:
\begin{itemize}
\item Soit le rang est égal à 2 et l'intersection des 3 plans est une droite.
\item Soit le rang est égal à 1 et l'intersection des 3 plans est un plan.
\end{itemize} 
\end{itemize}

% \subsection{Résolution numérique d'un système linéaire}

% La méthode du pivot de \nom{Gauss} est utilisée pour résoudre des systèmes de Cramer de manière numérique. Dans ce cas, les coefficients du système sont représentés par des nombres flottants. Nous verrons que la puissance des ordinateurs actuels permettent de résoudre en un temps raisonnable des systèmes possédant jusqu'à $10\,000$ équations et $10\,000$ inconnues. De manière générale, deux questions se posent quant à la résolution d'un système linéaire à $n$ équations et $n$ inconnues.
% \begin{itemize}
% \item Est-ce qu'un ordinateur est capable de résoudre le système dans un temps raisonnable~?
% \item Peut-on faire confiance au résultat trouvé par l'ordinateur~?\end{itemize}
% On peut répondre à la première question par un calcul de complexité. La seconde question nécessite une preuve de stabilité numérique qui est en dehors du cadre des classes préparatoires. Nous allons cependant donner quelques réponses à cette question.\\

% Commençons par le calcul de complexité. Avant toute chose, il est essentiel de savoir si l'ensemble des coefficients du système tient dans la mémoire de l'ordinateur. Pour des raisons de stabilité numérique, on travaille le plus souvent avec des nombres flottants de 64 bits. Chaque nombre flottant occupe donc 8 octets. Le système étant composé de $n^2$ coefficients, il occupe $8 n^2$ octets dans la mémoire de l'ordinateur. Pour $n=10\,000$, cela représente de l'ordre de $0.8$~Go et rentre aisément dans la mémoire vive (RAM) d'un ordinateur personnel. Cependant, pour $n=100\,000$, la taille du système approche 1~To et rentre difficilement dans la mémoire vive des ordinateurs les plus puissants. Pour des raisons de complexité spatiale, il est donc utopique de chercher à résoudre des systèmes ayant plusieurs centaines de milliers d'inconnues.\\
% Intéressons nous maintenant à la complexité temporelle de la méthode du pivot de Gauss, c'est-à-dire au nombre d'opérations élémentaires sur les nombres flottants qu'elle nécessite. Tout d'abord, on cherche à déterminer le nombre d'opérations nécessaires à l'élimination de la première inconnue. Pour des raisons de stabilité numérique, il est préférable de chercher pour pivot, non pas un coefficient quelconque non nul, mais plutôt le coefficient ayant la valeur absolue la plus grande possible. Cette stratégie, appelée méthode du pivot total, nécessite de faire une recherche parmi $n^2$ coefficients. Pour des raisons de performance, on utilise le plus souvent la méthode du pivot partiel qui consiste à chercher pour pivot le coefficient de la première colonne ayant la plus grande valeur absolue. Cette stratégie nécessite de faire une recherche parmi $n$ éléments. Si tous ces coefficients sont nuls, le système n'est pas de \nom{Cramer} et la résolution du système peut s'arrêter. Sinon, nous procédons ensuite à l'élimination. Pour chaque ligne $i$ entre 2 et $n$, il est nécessaire de calculer $\alpha_i\defeq a_{i,1}\cdot(1/a_{1,1})$, remplacer les $n-1$ coefficients $a_{i,j}$ par $a_{i,j} - \alpha_i a_{1,j}$ et remplacer les $y_i$ par $y_i- \alpha_i y_1$. On doit donc effectuer au total une division, $n(n-1)$ additions et $n(n-1)$ multiplications. Bien qu'une division de deux nombres flottants est 10 fois plus lente qu'une addition ou une multiplication, le temps passé à effecter cette division est négligeable devant le temps passé à faire les additions et les multiplications. On ne retiendra donc que ces opérations. Afin de réduire le système à un système échelonné, il est nécessaire d'effectuer cette élimination avec des systèmes de taille $n,n-1,\ldots,2$. Comme
% \[\sum_{k=2}^{n} k \approx \frac{1}{2}n^2 \et \sum_{k=2}^n k(k-1) \approx \frac{1}{3}n^3,\]
% on en déduit que pour réduire un système de Cramer à un système échelonné, la méthode du pivot partiel nécessite de faire des recherches parmi $(1/2)n^2$ éléments et d'effectuer $(1/3)n^3$ additions et $(1/3)n^3$ multiplications. Pour terminer la résolution, les opérations de substitutions nécessitent de l'ordre de $(1/2)n^2$ additions et $(1/2)n^2$ multiplications et sont donc négligeables devant les opérations précédentes. Les temps passé lors de la recherche des pivots est donc négligeable devant le temps passé à effectuer des additions et des multiplications. On retiendra que la résolution du système est limitée en temps par le calcul de $(1/3)n^3$ additions et $(1/3)n^3$ multiplications sur les nombres flottants, soit $(2/3)n^3$ opérations élémentaires. Le processeur actuel d'un ordinateur portable (2 GHz, 8 coeurs, unités vectorielles de 256 bits) peut effectuer de l'ordre de 200 milliards opérations sur les nombres flottants par seconde (200 GFlops). Si $n=10\,000$, il faudra quelques secondes pour résoudre un tel système. Même si on avait assez de mémoire, pour $n=100\,000$, il faudrait près d'une heure pour résoudre un tel système.\\
% En conclusion, les machines actuelles sont limitées à la résolution de systèmes de \nom{Cramer} à quelques dizaines de milliers de variables. Cette limitation est dûe à la fois à la taille de leur mémoire vive (RAM) et à la puissance de leurs processeurs.\\

% Il se pose désormais la question de la confiance que l'on peut accorder aux résultats donnés par l'ordinateur. Comme on le sait, les additions et les multiplications sur les nombres flottants ne sont pas exactes et des erreurs d'arrondis vont se propager lors de la résolution numérique du système linéaire. Remarquons tout d'abord que ces erreurs d'arrondi peuvent rendre non nul des coefficients qui auraient dû l'être. La méthode du pivot de Gauss va donc presque toujours trouver que le système est de rang $n$ même si ce n'est pas le cas. Elle trouvera donc presque toujours une unique solution même si le système n'en n'a pas ou en a une infinité. Et même si le système admet une unique solution, le confiance que l'on peut avoir dans le résultat est relative. Une analyse de stabilité numérique montre que~:
% \begin{itemize}
% \item Si on remplace les inconnues $x_1,\ldots,x_n$ dans le système par les solutions calculées par la méthode du pivot partiel, on obtiendra (presque toujours) un résultat proche des $y_i$.
% \item Cependant, même pour des valeurs raisonnables de $n$ (par exemple $n=20$), il est possible que les valeurs des $x_j$ soient très éloignées des valeurs exactes de la solution. On dit dans ce cas que le système est \emph{mal conditionné}.
% \end{itemize} 
% Autrement dit la solution calculée numériquement par la méthode du pivot de \nom{Gauss} est la solution exacte d'un problème proche de notre problème initial.  Cependant, si le système est mal conditionné, il est possible que cette solution soit très éloignée de la solution exacte.
% Remarquons enfin que si nous n'adoptons pas la stratégie du pivot partiel et que nous choisissons pour pivot le premier coefficient non nul que nous trouvons, aucune garantie ne peut être apportée sur la solution trouvée, même pour $n=2$.
%END_BOOK

\end{document}



\section{Exercices}
\setcounter{numeroexercice}{1}
\documentclass{magnolia}

\magtex{tex_driver={pdftex}}
\magfiche{document_nom={Compléments d'algèbre},
          auteur_nom={François Fayard},
          auteur_mail={fayard.prof@gmail.com}}
\magexos{exos_matiere={maths},
         exos_niveau={mpsi},
         exos_chapitre_numero={4},
         exos_theme={Compléments d'algèbre}}
\magmisenpage{}
\maglieudiff{}
\magprocess

\begin{document}

%BEGIN_BOOK


% \exercice{nom={Décomposition en éléments simples}}
% Décomposer en éléments simples sur $\R$ les fractions
% \[\frac{10X^3}{(X^2+1)(X^2-4)} \et
%   \frac{X^3-1}{(X-1)(X-2)(X-3)}.\]
% % \[\frac{X^2}{(X^2+X+1)^2}, \qquad \frac{(X^2+4)^2}{(X^2+1)(X^2-2)^2}.\]
% \begin{sol}
% $\quad$
% \begin{questions}
% \question 
%   \[\frac{10X^3}{(X^2+1)(X^2-4)}=
%     \frac{4}{X-2}+\frac{4}{X+2}+\frac{2X}{X^2+1}\]
% % \question 
% %   \[\frac{X^4+1}{X^4+X^2+1}=
% %     1-\frac{1}{2}\cdot\frac{X}{X^2-X+1}+\frac{1}{2}\cdot\frac{X}{X^2+X+1}\]
% \question 
%   \[\frac{X^3-1}{(X-1)(X-2)(X-3)}=
%     1-\frac{7}{X-2}+\frac{13}{X-3}\]
% % \question 
% %   \[\frac{X^2}{(X^2+X+1)^2}=
% %     \frac{1}{X^2+X+1}-\frac{X+1}{\p{X^2+X+1}^2}\]\
% % \question 
% %   \begin{eqnarray*}
% %   \frac{(X^2+4)^2}{(X^2+1)(X^2-2)^2}&=&
% %     \frac{1}{X^2+1}+\frac{3\sqrt{2}}{4}\cdot\frac{1}{X+\sqrt{2}}
% %     +\frac{3}{2}\cdot\frac{1}{\p{X+\sqrt{2}}^2}\\
% %    & &
% %     -\frac{3\sqrt{2}}{4}\cdot\frac{1}{X-\sqrt{2}}
% %     +\frac{3}{2}\cdot\frac{1}{\p{X-\sqrt{2}}^2}
% %   \end{eqnarray*}
% \end{questions}
% \end{sol}


\magsection{Somme et produit}
\magsubsection{Somme}

\exercice{nom={Sommes}}
Simplifier les sommes suivantes.
\[\textbf{a.}\ \sum_{k=0}^n k(k-1), \qquad \textbf{b.}\ \sum_{k=1}^n (2k-1), \qquad \textbf{c.}\ \sum_{k=1}^n (-1)^k\]
\[\textbf{d.}\ \sum_{k=0}^n (k+n), \qquad \textbf{e.}\ \sum_{k=1}^{n+1} \frac{2^k}{3^{2k-1}}, \qquad
\textbf{f.}\ \sum_{k=2}^n \ln\p{1-\frac{1}{k^2}}.\]

\exercice{nom={Décomposition en éléments simples}}
\begin{questions}
\question Montrer qu'il existe $a,b\in\R$ tels que
  \[\forall k\in\N\qsep \frac{1}{(3k+1)(3k+4)}=\frac{a}{3k+1}+\frac{b}{3k+4}.\]
\question En déduire
  \[\sum_{k=0}^n \frac{1}{(3k+1)(3k+4)}.\]
\end{questions}

\exercice{nom={Sommes}}
\begin{questions}
\question Montrer que pour tout $n\in\N$
  \[\sum_{k=0}^n k^3=\p{\frac{n(n+1)}{2}}^2.\]
\question Montrer que pour tout $n\in\N$
  \[\sum_{k=0}^n (-1)^k k^2=(-1)^n \frac{n(n+1)}{2}.\]
\end{questions}

\exercice{nom={Récurrence}}
Montrer par récurrence que pour tout $n\in\Ns$
\[\sum_{k=1}^{2n} \frac{(-1)^{k-1}}{k}=\sum_{k=1}^n \frac{1}{n+k}.\]

\exercice{nom={Somme lacunaire de coefficients binomiaux}}
En considérant $(1+1)^n$ et $(1-1)^n$ calculer, pour tout $n\in\Ns$, les sommes
\[A_n\defeq \sum_{\substack{0\leq k\leq n\\\text{$k$ pair}}} \binom{n}{k}
  \qquad\text{et}\qquad
  B_n\defeq \sum_{\substack{0\leq k\leq n\\\text{$k$ impair}}} \binom{n}{k}.\]

\exercice{nom={Coefficients binomiaux}}
Montrer par récurrence que pour tout $p\in\N$ et $n\geq p$
\[\sum_{k=p}^n \binom{k}{p}=\binom{n+1}{p+1}.\]

\exercice{nom={Coefficients binomiaux}}
Simplifier, pour tout $n\in\N$, la somme
\[\sum_{k=0}^n \frac{(-1)^k}{k+1}\binom{n}{k}.\]

\magsubsection{Produit}

\exercice{nom={Produits}}
Simplifier les produits suivants en les exprimant le plus possible à l'aide
de puissances et de factorielles.
\[\textbf{a.}\ \prod_{k=1}^n \sqrt{k(k+1)},\qquad
  \textbf{b.}\ \prod_{k=1}^n (-5)^{k^2-k},\qquad
  \textbf{c.}\ \prod_{k=1}^n \frac{4^k}{k^2}\]
\[\textbf{d.}\ \prod_{k=0}^n (2k+1),\qquad
  \textbf{e.}\ \prod_{k=1}^n (4k^2-1),\qquad
  \textbf{f.}\ \prod_{p=0}^{n-1} \sum_{k=0}^p 2^{p! k}.\]

\exercice{nom={Produit}}
Soit $z\in\C$. Pour tout $n\in\N$, on pose
\[P_n\defeq \prod_{k=0}^n \p{1+z^{2^k}}.\]
Calculer $(1-z)P_n$ et en déduire une expression simple de $P_n$.

\exercice{nom={Majoration}}
\begin{questions}
\question
\begin{questions}
\question Montrer que
  \[\forall k\geq 2\qsep 1+\frac{1}{k^2} \leq \frac{k^2}{(k-1)(k+1)}.\]
\question En déduire que
  \[\forall n\in\Ns\qsep \prod_{k=1}^n \p{1+\frac{1}{k^2}}\leq 4.\]
\end{questions}
\question Montrer que pour tout $n\in\Ns$
  \[\prod_{k=1}^n \p{1+\frac{1}{k^3}}\leq 3-\frac{1}{n}.\]
\end{questions}

\exercice{nom={Limite}}
\begin{questions}
\question Factoriser $k^3-1$ par $k-1$ et $k^3+1$ par $k+1$ pour tout $k\geq 2$.
\question En déduire, sans récurrence, que pour tout $n\geq 2$
  \[\prod_{k=2}^n \frac{k^3-1}{k^3+1}=\frac{2}{3}\cdot\frac{n^2+n+1}{n(n+1)}.\]
\question En déduire la limite de la suite de terme général
  \[u_n\defeq \prod_{k=2}^n \frac{k^3-1}{k^3+1}.\] 
\end{questions}


\magsubsection{Somme et produit doubles}

\exercice{nom={Sommes}}
Simplifier les sommes suivantes.
\[\textbf{a.}\ \sum_{1\leq i\leq j\leq n} i, \qquad
  \textbf{b.}\ \sum_{1\leq i < j\leq n} j, \qquad
  \textbf{c.}\ \sum_{1\leq i,j\leq n} (i+j)\]
\[\textbf{d.}\ \sum_{1\leq i,j\leq n} x^{i+j},\qquad
  \textbf{e.}\ \sum_{1\leq i\leq j\leq n} \frac{i}{j+1},\qquad
  \textbf{f.}\ \sum_{1\leq i\leq j\leq n} (j-i)\]
\[\textbf{g.}\ \sum_{1\leq i\leq j\leq n} \frac{i^2}{j},\qquad
  \textbf{h.}\ \sum_{1\leq i,j\leq n} (i+j)^2, \qquad
  \textbf{i.}\ \sum_{1\leq i,j\leq n} \max(i,j).\]
% \exercice{nom={Une somme de série}}
% Calculer la limite lorsque $n$ tend vers $+\infty$ de
% \[\sum_{k=2}^n \frac{3k^2-1}{(k-1)^2k^2(k+1)^2}.\]
% On pourra calculer $P'/P^2$ avec $P\defeq(X-1)X(X+1)$.
% \begin{sol}
% On pose $P\defeq (X-1)X(X+1)$. Alors, il existe $a,b,c\in\R$ tels que
% \[\frac{1}{P}=\frac{1}{(X-1)X(X+1)}=\frac{a}{X-1}+\frac{b}{X}+\frac{c}{X+1}.\]


% On trouve
% \[\frac{3X^2-1}{(X-1)^2X^2(X+1)^2}=\frac{1}{2}\cdot\frac{1}{(X-1)^2}
%   -\frac{1}{X^2}+\frac{1}{2}\cdot\frac{1}{(X+1)^2}\] en primitivant la DES de $-1/P$.
% (Astuce de Anis Marrakchi : Pour connaître le coefficient au dessus de
% $1/(X+1)$ et $-1/(X-1)$ on regroupe ces éléments simples, on multiplie par
% $X^2$ et on fait tendre $x$ vers $+\infty$ // Autre astuce de Anis Marrakchi :
% on peut remarquer que $F=P'/P^2$).
% On obtient donc une somme téléscopique. La somme de la série est donc
% $\frac{3}{8}$.
% \end{sol}


\exercice{nom={Avec des racines $n$-ièmes}}
Soit $n\geq 2$ et $z\in\C$. On pose $\omega\defeq \e^{\ii\frac{2\pi}{n}}$ et
\[S_n\defeq \sum_{k=0}^{n-1} \p{z+\omega^k}^n.\]
Calculer $S_n$.


\exercice{nom={Somme double}}
Soit $n\in\N$. Calculer
\[\sum_{0 \leq i,j \leq n} 2^{i-j} \et \sum_{\substack{(i,j)\in\intere{0}{n}^2\\0\leq i+j \leq n}} \binom{n}{i+j}.\]
\begin{sol}
\begin{questions}
\question
\[\sum_{i=0}^n \sum_{j=0}^n 2^{i-j}=2^{n+2}-4+2^{-n}.\]
\question
\[\sum_{i=0}^n \sum_{j=0}^n \binom{n}{i+j}=2^{n-1}(n+2).\]
\end{questions}
\end{sol}



\exercice{nom={Somme de {\sc Gauss}}}
Soit $n\in\N$ un entier impair. On pose
\[\omega\defeq\e^{\frac{2\ii\pi}{n}} \quad\text{et}\quad S\defeq \sum_{k=0}^{n-1} \omega^{k^2}.\]
\begin{questions}
\question Écrire $\abs{S}^2$ comme une somme double, puis montrer que
  \[\abs{S}^2=\sum_{k=0}^{n-1}\sum_{p=-k}^{n-k-1} \omega^{2pk+p^2}.\]
\question Soit $k\in\intere{0}{n-1}$.
  \begin{questions}
  \question Montrer que la fonction
  \[\dspappli{\phi}{\Z}{\C}{p}{\omega^{2pk+p^2}}\]
  est $n$-périodique.
  \question En déduire pour tout $k\in\intere{0}{n-1}$ une écriture simplifiée de
  \[\sum_{p=-k}^{n-k-1} \omega^{2pk+p^2}.\]
  \end{questions}
\question Simplifier
  \[\sum_{k=0}^{n-1} \omega^{2pk}\]
  pour tout $p\in\Z$.
\question En déduire que $\abs{S}=\sqrt{n}$.
\end{questions}

\begin{sol}
\begin{questions}
\question 
On calcule :
\begin{eqnarray*}
|S|^2&=& \conj{S}S\\
&=& \p{\sum_{k=0}^{n-1} \omega^{-k^2}}\p{\sum_{s=0}^{n-1} \omega^{s^2}}\\
&=& \sum_{k=0}^{n-1}\p{\omega^{-k^2}\sum_{s=0}^{n-1} \omega^{s^2}}\\
&=& \sum_{k=0}^{n-1}\p{\omega^{-k^2}\sum_{p=-k}^{n-k-1} \omega^{(p+k)^2}}\\
&=& \sum_{k=0}^{n-1}\sum_{p=-k}^{n-k-1} \omega^{2pk+p^2}
\end{eqnarray*}
\question Soit $k\in\intere{0}{n-1}$.
  \begin{questions}
  \question Facile car $\omega^n=1$.
  \question Ainsi, on a :
  \[\sum_{p=-k}^{n-k-1} \omega^{2pk+p^2}=\sum_{p=0}^{n-1} \omega^{2pk+p^2}.\]
  \end{questions}
\question On a pour tout $p\in\Z$ :
  \[\sum_{k=0}^{n-1} \omega^{2pk}=\sum_{k=0}^{n-1}\p{\omega^{2p}}^k.\]
  On a une somme géométrique, reste à savoir si $\omega^{2p}=1$ et c'est le cas si et seulement si $n$ divise $2p$.
  
  $\bullet$ 1er cas : Si $n$ ne divise pas $2p$, la somme géométrique fait $0$.
   $\bullet$ 2ème cas : Si $n$ divise $2p$, la somme fait $n$.
 
\question Déjà, faisons le bilan des deux premières questions. 


\[\abs{S}^2=\sum_{k=0}^{n-1}\sum_{p=-k}^{n-k-1} \omega^{2pk+p^2}=\sum_{k=0}^{n-1}\sum_{p=0}^{n-1} \omega^{2pk+p^2}=\sum_{p=0}^{n-1}\sum_{k=0}^{n-1}\omega^{2pk+p^2}=\sum_{p=0}^{n-1}\omega^{p^2}\sum_{k=0}^{n-1}\omega^{2pk}\]
Comme $n$ est impair, le seul $p$ entre $0$ et $n-1$ tel que $n$ divise $2p$ est $0$. Donc seul le terme en $p=0$ n'est pas nul et vaut $n$. D'où $|S|^2=n$ et on a le résultat souhaité.

\end{questions}
\end{sol}

\magsubsection{Fonction polynôme}

\exercice{nom={Polynôme à coefficients symétriques}}
\begin{questions}
\question Montrer que le changement de variable $u\defeq z+1/z$ simplifie l'équation
  \[z^4-5z^3+6z^2-5z+1=0\]
  en une équation du second degré en $u$.
\question En déduire l'ensemble de ses solutions sur $\C$.
\question Sur le même modèle, résoudre l'équation
  \[z^4+z^3-10z^2-z+1=0.\]
\end{questions}
\begin{sol}
$\quad$
\begin{questions}
\question On trouve $2\pm\sqrt{3}$.
\question Faire le changement de variable $u=x-1/x$. On trouve
  \[\frac{-\p{1+\sqrt{33}}\pm\sqrt{50+2\sqrt{33}}}{4}, \qquad
    \frac{-1+\sqrt{33}\pm\sqrt{50-2\sqrt{33}}}{4}\]
\end{questions}
\end{sol}

\exercice{nom={Simplification de racine}}
On pose
\[a\defeq \sqrt[5]{\frac{5\sqrt{5}+11}{2}}, \quad b\defeq \sqrt[5]{\frac{5\sqrt{5}-11}{2}} \et
  x\defeq a-b.\]
On souhaite montrer que $x=1$.
\begin{questions}
\question Calculer $ab$.
\question En déduire que $x$ est racine de $A(t)\defeq t^5+5t^3+5t-11$.
\question Montrer que $1$ est la seule racine positive de $A$ et conclure.
\end{questions}
\begin{sol}
Écrire $x=a-b$ et remarquer que $ab=1$. Développer $x^5$ et en déduire que
$x^5+5x^3+5x-11=0$. Factoriser ce polynôme par $\p{x-1}$~:
$\p{x-1}\p{x^4+x^3+6x^2+6x+11}$. Comme le second polynôme n'a pas de racine
positive, on en déduit que $x=1$.
\end{sol}

% \exercice{nom={Équations algèbriques}}
% \begin{questions}
% \question Résoudre les équations
%   \[x=\sqrt{2-x}, \qquad \sqrt{x+7}=\sqrt{x+2}+\sqrt{x-1}.\]
% \question Résoudre les équations suivantes sur $\R$
%   \[1+x+x^2+x^3+x^4+x^5=0, \qquad 1-x+x^2-x^3+x^4=0.\]
% \end{questions}
% \begin{sol}
% $\quad$
% \begin{questions}
% \question Seul 1 est racine pour la première. Seul 2 est racine pour la
%   seconde.
% \question Pour la première, seul 1 est solution. Pour la seconde, il n'y a pas
%   de solution.
% \end{questions}
% \end{sol}

\exercice{nom={Inéquation}}
Résoudre l'inéquation
\[4x+2\leq\sqrt{7x^3+15x^2+11x+3}.\]
On commencera bien entendu par donner son domaine de définition.

\begin{sol}
$7x^3+15x^2+11x+3=(x+1)(7x^2+8x+3)$ avec le terme du second degré toujours positif. Donc la racine est définie pour $x\geq -1$.
Si $x\leq -1/2$, $4x+2\leq 0$ donc l'inégalité est vérifiée pour $x\in [-1;-1/2]$.

Si $x\geq 1/2$, on peut élever les deux termes au carré tout en gardant l'équivalence et l'inégalité est équivalente à $(x-1)(7x^2+6x+1)\geq 0$. L'étude du trinôme + tableau de signe montre que $S=[-1;-1/2]\cup[-1/2;\dfrac{-3+\sqrt{2}}{7}]\cup [1;+\infty[.$
\end{sol}

\exercice{nom={Coefficients binomiaux}}
Écrire la somme
\[\sum_{k=0}^n (1+z)^k\]
de deux manières différentes. En déduire
\[\sum_{k=j}^n \binom{k}{j}.\]

\begin{sol}
D'une part,
$$\sum_{k=0}^n (1+X)^k=\sum_{k=0}^n \sum_{j=0}^k \binom{k}{j}X^j=\sum_{j=0}^n\p{\sum_{k=j}^n\binom{k}{j}}X^j.$$
D'autre part, $$\sum_{k=0}^n (1+X)^k=\frac{1-(1+X)^{n+1}}{1-(1+X)}=\frac{\sum_{k=0}^{n+1}\binom{n+1}{k}X^k -1}{X}=\sum_{k=1}^{n+1}\binom{n+1}{k}X^{k-1}=\sum_{j=0}^{n}\binom{n+1}{j+1}X^{j}.$$
Par identification des coefficients, on a pour tout $j\in \llbracket 0,n\rrbracket, \sum_{k=j}^n \binom{k}{j}=\binom{n+1}{j+1}$.
\end{sol}


\exercice{nom={Méthode de {\sc Cardan}}}
Soit $a,b,c\in\C$ et $Q(z)\defeq z^3+az^2+bz+c$. 
\begin{questions}
\question Pour quels $\alpha\in\C$ le polynôme $P(z)\defeq Q(z+\alpha)$ n'a-t-il pas de coefficient en $z^2$~?
\enonce On choisit un tel $\alpha$ et on définit $p,q\in\C$ tels que $P(z)=z^3-3pz+2q$.
\question 
\begin{questions}
\question Soit $u\in\Cs$. Montrer que
  \[u+\frac{p}{u}\]
  est racine de $P$ si et seulement si $w\defeq u^3$ est racine d'un trinôme $R$ que l'on déterminera.
\question On note $w_1$ et $w_2$ les racines complexes de $R$. Soit $u_1$ une racine cubique de $w_1$. Montrer qu'il existe une unique racine cubique $u_2$ de $w_2$ telle que $u_1u_2=p$.
\question Déterminer les racines de $Q$ en fonction de $\alpha$, $u_1$, $u_2$ et $\jj$.
\end{questions} 
\question Déterminer les racines des polynômes
  \[z^3+3z^2+6z+2,\qquad z^3-3z-1.\]
\end{questions}

\textit{C'est pour résoudre de telles équations, pour lesquelles $P$ admet trois racines réelles, mais $R$ n'en n'a pas, que {\sc Rafael Bombelli} (1526--1572) a introduit les nombres complexes.}

\begin{sol}
\begin{questions}
\question On trouve $\alpha=-a/3$.
\question 
\begin{questions}
\question On trouve que $w\defeq u^3$ est racine de $X^2+2qX+p^3$.
\question On a $w_1w_2=p^3$ donc $w_2=p^3/u_1^3$. Ainsi, les racines cubiques de $w_2$ sont $\dfrac{p}{u_1}$, $j\dfrac{p}{u_1}$ et $j^2\dfrac{p}{u_1}$. 
\question 
\begin{eqnarray*}
u\in\ens{u_1,ju_1,j^2u_1,u_2,ju_2,j^2u_2} &\Longleftrightarrow& u^3 \text{ est racine de } R \\
&\Longleftrightarrow& u^+\dfrac{p}{u} \text{ est racine de } P\\
&\Longleftrightarrow& u^+\dfrac{p}{u} \text{ est racine de } Q(X+\alpha)\\
&\Longleftrightarrow& u^+\dfrac{p}{u}+\alpha \text{ est racine de } Q
\end{eqnarray*}
Ainsi, les racines de $Q$ sont $$\mathcal{S}=\ens{u_1+u_2+\alpha; ju_1+j^2u_2+\alpha; j^2u_1+ju_2+\alpha}.$$
\end{questions} 
\question Déterminer les racines des polynômes
  \[X^3+3X^2+6X+2,\qquad X^3-3X-1.\]
\end{questions}
\end{sol}

\exercice{nom={Méthode de {\sc Ferrari}}}
Soit $a,b,c,d\in\C$ et $P(z)\defeq z^4+az^3+bz^2+cz+d$. 
\begin{questions}
\question Montrer qu'il existe un unique $\alpha\in\C$ tel que $Q(z)\defeq P(z+\alpha)$ ne possède pas de terme en $z^3$.
\enonce Pour la suite, $\alpha$ désignera cette valeur. Il existe donc $p,q,r\in\C$ tels que $Q(z)=z^4+p z^2+q z+ r$.
\question Soit $v\in\C$. Montrer que
  \[\forall z\in\C\qsep Q(z)=0 \quad\ssi\quad \p{z^2+\frac{v}{2}}^2=\p{v-p}z^2-qz+\p{\frac{v^2}{4}-r}.\]
  On pose alors $A(z)\defeq \p{v-p}z^2-qz+\p{v^2/4-r}$.
\question Montrer que $A$ admet une racine double si et seulement si $v$ est racine d'un polynôme $B$ de degré 3 que l'on déterminera.
\enonce Pour la suite, on suppose que $P(z)\defeq z^4-2z^3+3z^2+4zQ-10$.
\question Montrer que $B$ admet une racine évidente. En déduire les racines de $P$.
\end{questions}

\textit{On a ainsi prouvé, dans les deux exercices précédents, que l'on pouvait calculer les racines des équations de degré 3 et 4 à l'aide de racines $n$-ièmes de nombres complexes. Ces résultats étaient connus dès le 16e siècle. {\sc Niels Abel} (1802--1829), puis {\sc Évariste Galois} (1811--1832), ont démontré qu'il n'était pas possible de résoudre l'équation générale de degré 5 en utilisant des racines $n$-ièmes de nombres complexes.}

\exercice{nom={Principe du maximum pour les polynômes}}
Soit $P(z)\defeq\sum_{k=0}^n a_k z^k$ une fonction polynôme à coefficients complexes de degré inférieur ou égal à $n$. On se donne un réel positif $M$ tel que
$$\forall z\in\U \qsep \abs{P(z)}\leq M.$$
\begin{questions}
\question On pose $\omega\defeq\exp\p{\ii\frac{2\pi}{n+1}}$. Calculer
  $$\sum_{j=0}^n P\p{w^j}.$$
\question En déduire que $\abs{P(0)}\leq M$.
\end{questions}

\begin{sol}
\begin{questions}
\question 
  $$\sum_{j=0}^n P\p{w^j}=\sum_{j=0}^n\sum_{k=0}^na_k\omega^{jk}=\sum_{k=0}^na_k\sum_{j=0}^n\p{\omega^{k}}^j=(n+1)a_0.$$
\question Ainsi, $|(n+1)a_0|=\abs{\sum_{j=0}^n P\p{w^j}|}\leq (n+1)M$ par I.T permet de conclure.
\end{questions}
\end{sol}

\magsection{Trigonométrie}
\magsubsection{Égalité modulaire}
\magsubsection{Formules de trigonométrie}

% \exercice{nom={Calcul du cosinus et du sinus de $\frac{\pi}{12}$ et
%   $-\frac{\pi}{8}$}}
% \begin{questions}
% \question En utilisant les formules de trigonométrie, calculer
%   \[\sin^2\p{\frac{\pi}{12}}.\]
% \question En déduire le sinus, le cosinus et la tangente de l'angle de mesure
%   $\pi/12$ en utilisant des racines carrées.
% \question En s'inspirant de cette méthode, calculer le sinus, le cosinus et la
%   tangente de l'angle $-\pi/8$.
% \end{questions}
% \begin{sol}
% $\quad$
% \begin{questions}
% \question On a
%   \[\sin^2\p{\frac{\pi}{12}}=\frac{2-\sqrt{3}}{4}\]
% \question On trouve
%   \[\sin\p{\frac{\pi}{12}}=\frac{\sqrt{2-\sqrt{3}}}{2}, \quad \cos\p{\frac{\pi}{12}}=\frac{\sqrt{2+\sqrt{3}}}{2}, \quad
%     \tan\p{\frac{\pi}{12}}=2-\sqrt{3}\]
% \question On trouve
%   \[\sin\p{-\frac{\pi}{8}}=-\frac{\sqrt{2-\sqrt{2}}}{2}, \quad \cos\p{-\frac{\pi}{8}}=\frac{\sqrt{2+\sqrt{2}}}{2}, \quad
%     \tan\p{-\frac{\pi}{8}}=1-\sqrt{2}\]
% \end{questions}
% \end{sol}

% \exercice{nom={Sur les angles d'un triangle}}
% Soit $a,b,c$ trois réels tels que $a+b+c=\pi$. Montrer que
% \[\cos^2 a+\cos^2 b+\cos^2 c+2\cos a \cos b\cos c=1.\]

% \exercice{nom={Équations trigonométriques}}
% Résoudre les équations suivantes sur $\R$
% \[\cos\p{2x}+\cos(x)=-1\]
% \[\sqrt{3}\cos\p{5x}=\cos\p{2x}+\cos\p{12x} \qquad
%   \tan\p{2x}=3\tan{x}\]
% \[\cos^6 x+\sin^6 x=\frac{1}{4} \qquad
%   2\cos\p{\frac{x}{3}}-\sin\p{\frac{x}{2}}=2.\]
% \begin{sol}
% \quad
% \begin{questions}
% \question Exprimer $\cos\p{2x}$ en fonction de $\cos x$. On trouve
%   \[x\equiv\frac{\pi}{2}\ [\pi] \ou x\equiv\frac{2\pi}{3}\ [2\pi] \ou
%     x\equiv -\frac{2\pi}{3}\ [2\pi]\]
% \question Factoriser le second membre. L'équation se factorise alors par
%   $\cos\p{5x}$. On trouve
%   \[x\equiv\frac{\pi}{10}\ \cro{\frac{\pi}{5}} \ou
%     x\equiv\frac{\pi}{42}\ \cro{\frac{2\pi}{7}} \ou
%     x\equiv -\frac{\pi}{42}\ \cro{\frac{2\pi}{7}}\]
% \question Il faut $x\not\equiv\pi/4\ [\pi/2]$. On change tout en $\sin$, $\cos$
%   puis on factoriser par $\sin x$. On obtient $\sin x\p{1-2\cos\p{2x}}=0$.
%   Les solutions sont donc
%   \[x\equiv 0\ [\pi] \ou x\equiv\frac{\pi}{6}\ [\pi] \ou 
%     x\equiv -\frac{\pi}{6}\ [\pi]\]
% \question On pose $u=\cos^2 x$. On obtient un polynôme du second degré en $u$
%   qui a une unique racine $1/2$. On trouve
%   \[x\equiv\frac{\pi}{4}\ \cro{\frac{\pi}{2}}\]
% \question On pose $x=6u$. On factorise $2\p{\cos\p{2u}-1}-\sin\p{3u}$ par
%   $\sin u$ et on obtient $\sin u\p{4\sin^2 u-4\sin u-3}=0$.
%   Une racine est en dehors de $\interf{-1}{1}$. L'autre est $-1/2$. Finalement
%   \[x\equiv 0\ [6\pi] \ou x\equiv-\pi\ [12\pi] \ou x\equiv-5\pi\ [12\pi]\]
% \end{questions}
% \end{sol}

% \exercice{nom={Inéquations trigonométriques}}
% Résoudre les inéquations suivantes sur $\R$
% \[\cos x+\cos\p{x+\pi/3}\geq 0, \qquad 2\cos x+\sin x \leq 2.\]

% \exercice{nom={Calcul de somme}}
% \begin{questions}
% \question Calculer $\tan p-\tan q$.
% \question En déduire la valeur de
%   \[S_n=\sum_{k=1}^n \frac{1}{\cos\p{k\theta}\cos\p{\p{k+1}\theta}}\]
%   où $\theta\in\R$.
% \end{questions}
% \begin{sol}
% Il faut que $\theta\not\equiv\frac{\pi}{2k} \cro{\frac{\pi}{k}}$ pour tout
% $k\in\intere{1}{n+1}$. En suite on trouve
% \[S_n=\begin{cases}
%       \frac{\sin(n\theta)}{\cos((n+1)\theta)\cos(\theta)\sin \theta} &
%         \text{si $\theta\not\equiv 0\ [\pi]$}\\
%       n & \text{si $\theta=2k\pi$}\\
%       -n & \text{si $\theta=(2k+1)\pi$}
%       \end{cases}\]
% \end{sol}

\exercice{nom={Inégalité}}
Montrer que pour tout $x\in \R$ et $n \in \N$
\[|\sin(nx)| \leq n |\sin(x)|.\]
\begin{sol}
Récurrence et I.T
\end{sol}

\exercice{nom={Équations trigonométriques}}
Résoudre les équations suivantes d'inconnue $x$.
\[\textbf{a.}\ \cos(3x)=\sin(x), \qquad \textbf{b.}\ \cos x+\sin x=1+\tan x,\qquad \textbf{c.}\ \sin x+\sin(2x)=0\]
\[\textbf{d.}\ \tan(2x)=3\tan x,\qquad \textbf{e.}\ 2\sin x+\sin(3x)=0,\qquad \textbf{f.}\ 3\tan x=2\cos x\]
\[\textbf{g.}\ \cos x=\sqrt{3}\sin x, \qquad \textbf{h.}\ 2\cos(4x)+\sin x=\sqrt{3}\cos x.\]


\exercice{nom={Équations trigonométriques}}
\begin{questions}
\question Résoudre les équations \[\cos(x)-\cos(2x)=\sin(3x), \qquad \cos(5x)+2\cos (3x)+3\cos(x)=0\]
\question Résoudre l'équation
  \[\cos(x)+\sin(x)=\tan \left( \frac{x}{2} \right)\]
  en posant (soigneusement) $t\defeq\tan(x/2)$.
\end{questions}

\begin{sol}
\begin{questions}
\question On exprime les $\cos(nx)$ comme polynôme en $\cos(x)$. On obtient donc
  \begin{eqnarray*}
\forall x\in\R\qsep & & \cos(5x)+2\cos (3x)+3\cos(x)=0\\
&\ssi& 16\cos^5(x)-12\cos^3(x)+2\cos(x)=0\\
&\ssi& u(16 u^4-12u^2+2)=0 \quad\text{en posant $u=\cos(x)$}\\
&\ssi& u = 0 \ou u^2=\frac{1}{4} \ou u^2=\frac{1}{2}\\
&\ssi& u =0 \ou u =\pm\frac{1}{2} \ou u=\pm\frac{1}{\sqrt{2}}\\
&\ssi& \cos(x)=0 \ou \cos(x)=\pm\cos\frac{\pi}{3} \ou \cos(x)=\pm\cos\frac{\pi}{4}\\
&\ssi& x\equiv\frac{\pi}{2}\ [\pi]\ou\cos(x)\equiv\pm\frac{\pi}{3}\, [2\pi]\ou\cos(x)\equiv\pm\frac{2\pi}{3}\, [2\pi]\ou\\
&    & x\equiv\pm\frac{\pi}{4}\ [2\pi]\ou x\equiv\pm\frac{3\pi}{4}\, [2\pi].
  \end{eqnarray*}
Pour $\cos(x)-\cos(2x)=\sin(3x)$. On utilise $\cos(a)-\cos(b)$ pour la gauche et $\sin(2a)$ pour la droite. On obtient :
$$-2\sin(3x/2)\sin(-x/2)=2\sin(3x/2)\cos(3x/2).$$
On peut transformer $\sin(x/2)=\cos(\pi/2-x/2)$ puis les différents cas mènent à $x\equiv \pi/4 [\pi]$, $x\equiv -\pi/2 [2\pi]$ et $x\equiv 0 [2\pi/3]$.
\question Le domaine de définition de l'équation est $\mathcal{D}=\R\setminus\p{\pi+2\pi\Z}$. Soit $x\in\mathcal{D}$. On pose $t\defeq\tan(x/2)$. Alors
  \begin{eqnarray*}
\cos x+\sin x=\tan\frac{x}{2}
&\ssi& \frac{1-t^2}{1+t^2}+\frac{2t}{1+t^2}=t\\
&\ssi& 1-t^2+2t=t(1+t^2)\\
&\ssi& t^3+t^2-t-1=0\\
&\ssi& (t-1)(t+1)^2=0\\
&\ssi& t=1 \ou t=-1\\
&\ssi& \tan\frac{x}{2}=\tan\frac{\pi}{4} \ou \tan\frac{x}{2}=\tan\p{-\frac{\pi}{4}}\\
&\ssi& \frac{x}{2}\equiv\frac{\pi}{4}\ [\pi] \ou \frac{x}{2}\equiv-\frac{\pi}{4}\ [\pi]\\
&\ssi& x\equiv\frac{\pi}{2}\ [2\pi] \ou x\equiv-\frac{\pi}{2}\ [2\pi]\\
&\ssi& x\equiv\frac{\pi}{2}\ [\pi]
  \end{eqnarray*}
\end{questions}
\end{sol}

\exercice{nom={Calcul de somme}}
Pour tout $n\in\Ns$, calculer
\[\sum_{k=0}^{n-1} 3^k \sin^3 \p{\frac{\theta}{3^{k+1}}}.\]

\begin{sol}
On linéarise $\sin^3(x)=3/4\sin(x)-1/4\sin(3x)$ puis quand on écrit la somme elle devient télescopique. On obtient
\[\frac{1}{4}\cro{\sin(\theta)-3^n\sin\p{\frac{\theta}{3^n}}}.\]
\end{sol}

\exercice{nom={Mon capitaine}}
Pour tout $n\in\N$, calculer
\[\sum_{k=0}^n k\binom{n}{k}\cos (k \theta).\]


\magsection{Récurrence linéaire}
\magsubsection{Récurrence linéaire d'ordre 1}
\exercice{nom={Récurrences d'ordre 1}}
Déterminer une expression explicite des suites définies par
\begin{questions}
\question $u_0\defeq 0$ et $\forall n\in\N\qsep u_{n+1}\defeq 2 u_n+1$.
\question $u_0\defeq 1$ et $\forall n\in\N\qsep u_{n+1}\defeq 3 - \frac{u_n}{2}$.
\question $u_0\defeq 1$ et $\forall n\in\N\qsep u_{n+1}\defeq 2 u_n^2$.
\end{questions}


\magsubsection{Récurrence linéaire d'ordre 2}

\exercice{nom={Récurrences doubles}}
Déterminer les suites définies par
\[\textbf{a.}\ a_0\defeq 0,\quad a_1\defeq -1 \et \forall n\in\N\qsep a_{n + 2} \defeq 5a_{n + 1}-6a_n,\]
\[\textbf{b.}\ b_0\defeq 1,\quad b_1\defeq 1 \et \forall n\in\N\qsep b_{n + 2} \defeq -b_{n + 1}-b_n,\]
\[\textbf{c.}\ c_0\defeq 1,\quad c_1\defeq -1 \et \forall n\in\N\qsep c_{n + 2} \defeq 4c_{n + 1}-4c_n+n.\]

\begin{sol}
Pour le premier, on trouve $a_n=2^n-3^n$. Pour le second, on trouve $b_n=2\sin(\pi/6+n2\pi/3)$.
Pour la troisième, on cherche une solution particulière du type $v_n=n+\alpha$. On trouve $\alpha=2$. On résout l'équation homogène puis on ajoute les $2$ et ensuite avec les valeurs initiales on obtient :
$$c_n=(-1-n)2^n+n+2.$$
\end{sol}

\exercice{nom={Récurrence double}}
Déterminer une expression explicite de la suite $(u_n)$ définie par
\[u_0\defeq 1,\qquad
  u_1\defeq 2\qquad\text{et}\qquad
  \forall n\in\N\qsep u_{n+2}\defeq\frac{u_{n+1}^6}{u_n^5}.\]

\exercice{nom={Récurrence double avec second membre polynomial}}
On dira qu'une suite réelle $(u_n)$ est solution de $(E)$ lorsque
\[\forall n\in\N\qsep u_{n+2}=2u_{n+1}+8u_n+9n^2.\]
\begin{questions}
\question Montrer que $(E)$ possède une solution de la forme $(an^2+bn+c)$ où $a,b,c\in\R$.
\question En déduire une expression explicite de l'unique solution de $(E)$ telle que $u_0=0$ et $u_1=1$.
\end{questions}

\exercice{nom={Équation fonctionnelle}}
Déterminer les fonctions $\app{f}{\RP}{\RP}$ telles que
\[\forall x \geq 0\qsep f(f(x))=6x-f(x).\]

\begin{sol}
On procède par analyse-synthèse. On suppose qu'une telle solution existe.\\

Fixons $x_0\geq 0$ et définissons $(u_n)$ par $u_n=f^n(x_0)$. On démontre (en appliquant l'égalité en $f(x)$) $u_{n+2}+u_{n+1}-6u_n=0$. D'où : $$u_n=\lambda (-3)^n +\mu 2^n \quad \text{ avec en particulier } \lambda=\dfrac{2}{5}x_0-\dfrac{1}{5}f(x_0).$$

On peut écrire $u_n=\displaystyle (-3)^n\p{\lambda+\mu\p{\frac{-2}{3}}^n}$. On peut avec cette forme montrer que $u_n\geq 0 \forall n\in \N$ implique $\lambda=0$. Ainsi $f(x_0)=2x_0$ et cela pour $x_0$ quelconque dans $\RP$.

Synthèse : cela fonctionne.
\end{sol}


\magsection{Système linéaire}
\magsubsection{Système linéaire à $q$ équations et $p$ inconnues}


% \exercice{nom={Résolution de systèmes linéaires}}
% Résoudre les systèmes linéaires suivants~:
% \[\syslin{ x&+3y&=&3\hfill\cr
%           3x&-y &=&1\hfill} \qquad
%   \syslin{3x&-2y&=&1\hfill\cr
%            x&+y &=&2\hfill}\]
% \[\syslin{ 2x&-y &=&2\hfill\cr
%           -4x&+2y&=&0\hfill} \qquad
%   \syslin{ 2x&-y &=&2\hfill\cr
%           -4x&+2y&=&-4\hfill}\]

\exercice{nom={Résolution de systèmes linéaires}}
\begin{questions}
\question Résoudre les systèmes linéaires suivants
  \[\syslin{2x&+y &+z&=&1\hfill\cr
             x&-y &  &=&0\hfill\cr
            3x&+5y&+z&=&3,\hfill} \qquad
    \syslin{2x&-y &+2z&=&2\hfill\cr
             x&+y &+z &=&1\hfill\cr
             x&   &+z &=&1.\hfill}\]
\question Soit $a,b,c,d\in\R$. Résoudre le système linéaire
  \[\syslin{x&+y&+z&+t&=&a\hfill\cr
            x&+2y&+3z&+4t&=&b\hfill\cr
            x&+3y&+6z&+10t&=&c\hfill\cr
            x&+4y&+10z&+20t&=&d.\hfill}\]
\end{questions}
\begin{sol}
\begin{questions}
\question Pour le premier, on trouve $x=y=2/5$ et $z=-1/5$. Pour le deuxième, $y=0$ et $x+z=1$ donc $x=1-t_1$ et $t_1=z$.
\question Le système admet une unique solution
  \[\syslin{x&=&4a&-6b&+4c&-d\cr
            y&=&-6a&+14b&-11c&+3d\cr
            z&=&4a&-11b&+10c&-3d\cr
            t&=&-a&+3b&-3c&+d}\]
\end{questions}
\end{sol}

\exercice{nom={Équation fonctionnelle}}
Déterminer les fonctions $f:\RPs\to\R$ telles que
\[\forall x\in\RPs\qsep f(x)+3f\p{\frac{1}{x}}=x^2.\]

\exercice{nom={Sytème de type \nom{Vandermonde}}}
Soit $a,b,c$ trois réels deux à deux distincts. Résoudre le système
  \[\syslin{x&+ay&+a^2z&=&a^4\hfill\cr
            x&+by&+b^2z&=&b^4\hfill\cr
            x&+cy&+c^2z&=&c^4.}\]
\begin{sol}
Pour le deuxième système, on trouve 
$$z=a^2+b+c^2+ab+ac+bc$$
$$y=-(bc^2+b^2c+a^2b+ab^2+a^2c+ac^2)-2abc$$
$$x=a^2bc+ab^2c+abc^2.$$
\end{sol}



\exercice{nom={Résolution de systèmes linéaires}}
\begin{questions}
\question Soit $m\in\R$. Résoudre le système
  \[\syslin{mx&+ y&+ z&=&1\hfill\cr
             x&+my&+ z&=&m\hfill\cr
             x&+ y&+mz&=&m^2.\hfill}\]
\question Soit $m,a,b,c,\in\R$. On considère le système
  \[\syslin{-mx&+\p{m-1}y&+mz&=&a\hfill\cr
            \p{2m-1}x&+\p{m-1}y&-mz&=&b\hfill\cr
            -2x&-4y&+2mz&=&c.\hfill}\]
  Calculer le rang du système en fonction de $m$. Le résoudre dans ces différents cas.
\end{questions}
\begin{sol}
$\quad$
\begin{questions}
\question Si $m=1$, $x=1-t_1-t_2$, $y=t_1$, $z=t_2$. Si $m=-2$ le système n'a
  pas de solution. Sinon
  \[x=-\frac{m+1}{m+2} \qquad y=\frac{1}{m+2} \quad z=\frac{\p{m+1}^2}{m+2}\]
\question Le système n'admet pas une unique solution si et seulement si
  $m=0,\frac{1}{3},1$.
  \begin{itemize}
  \item Si $m=0$, alors
    \begin{itemize}
    \item Si $a+b-\frac{c}{2}\neq 0$, alors le système n'a pas de solution.
    \item Sinon $x=a-b$, $y=-a$ et $z=t$.
    \end{itemize}
  \item Si $m=\frac{1}{3}$, alors
    \begin{itemize}
    \item Si $4a+2b-c\neq 0$, alors le système n'admet pas de solution.
    \item Sinon, $x=-2t-\frac{3}{2}b-\frac{3}{2}a$, $y=t$,
      $z=\frac{3}{2}a-\frac{3}{2}b$.
    \end{itemize}
  \item Si $m=1$, alors
    \begin{itemize}
    \item Si $a+b\neq 0$, alors le système n'admet pas de solution.
    \item Sinon, $x=t-a$, $y=\frac{1}{2}a-\frac{1}{4}c$, $z=t$.
    \end{itemize}
  \end{itemize}
\end{questions}
\end{sol}

\exercice{nom={Calcul de rang}}
Soit $m\in\R$. Résoudre et déterminer le rang du système linéaire
\[\syslin{ x&+ y& +z& +t&=&3\hfill\cr
           x&+my& +z&-mt&=&m+2\hfill\cr
          mx& -y&-mz& -t&=&-1.\hfill}\]
\begin{sol}
En deux étapes (Pivot en haut à gauche puis échange des colonnes $t$ et $y$ et des lignes $2$ et $3$ puis $L_3\leftarrow L_3-L_2$) on se ramène à un système échelonné :
\[\syslin{ x&+ t& +z& +y&=&3\hfill\cr
           &  -(m+1)t&-2mz&-(m+1)y&=&-1-3m \hfill\cr
          & & 2mz& +2my&=&4m\hfill}\]
\begin{itemize}
\item[$\bullet$] Si $m=0$ :
\[\syslin{ x&+ t& +z& +y&=&3\hfill\cr
           &  t& &+y&=&1 \hfill}\]
           Le rang du système est $2$ et on peut prendre $z$ et $y$ en paramètre...
\item[$\bullet$] Si $m\neq 0$, le système devient 
\[\syslin{ x&+ t& +z& +y&=&3\hfill\cr
           &  -(m+1)t&-2mz&-(m+1)y&=&-1-3m \hfill\cr
          & & z& +y&=&2\hfill}\]
          \begin{itemize}
\item[$\bullet$] Si $m=-1$ : le système devient 
\[\syslin{ x&+ t& +z& +y&=&3\hfill\cr
           &  &z&&=&1 \hfill\cr
          & & z& +y&=&2\hfill}\]
          On a donc $y=1$, $z=1$ et $x+t=1$ donc c'est de rang $3$ et on passe $t$ en paramètre.
\item[$\bullet$] Si $m\neq -1$, le système est de rang $3$, on peut fixer $y$ comme paramètre et exprimer les autres en fonction. 
\end{itemize}
\end{itemize}
\end{sol}

\exercice{nom={Somme lacunaire de coefficients binomiaux}}
On définit $A,B$ et $C$ par
$$A\defeq\sum_{\substack{k=0\\k\egmodulo{0}{3}}}^{n} \binom{n}{k}, \qquad
  B\defeq\sum_{\substack{k=0\\k\egmodulo{1}{3}}}^{n} \binom{n}{k} \qquad\text{et}\qquad
  C\defeq\sum_{\substack{k=0\\k\egmodulo{2}{3}}}^{n} \binom{n}{k}.$$
\begin{questions}
\question Calculer $A+B+C$, $A+\jj B+\jj^2C$ et $A+\jj^2B+\jj C$.
\question En déduire $A$.
\question Sur le même modèle, étant donnés $b\in\Ns$ et $a\in\intere{0}{b-1}$, calculer
  $$A=\sum_{\substack{k=0\\k\egmodulo{a}{b}}}^{n} \binom{n}{k}.$$
\end{questions}
\begin{sol}
$\quad$
\begin{questions}
\question
\question On trouve
  \[A=\frac{2^n+2\cos\frac{n\pi}{3}}{3}\]
\question
  \[A=\frac{2^n+2^{\frac{n+2}{2}}\cos\frac{n\pi}{4}}{4}\]
\end{questions}
\end{sol}



\magsubsection{Interprétation géométrique lorsque $p=2$ ou $p=3$}
%END_BOOK

\end{document}


\end{document}
